\documentclass[10pt,a4paper,oneside]{article}
\usepackage[utf8]{inputenc}
\usepackage[T1]{fontenc}
\usepackage{lmodern}
\usepackage[a4paper,left=18mm,right=17mm,top=21mm,bottom=20mm,headheight=14pt]{geometry}
\usepackage{xcolor}
\definecolor{ink}{HTML}{18364A}
\definecolor{teal}{HTML}{0A6E75}
\definecolor{muted}{HTML}{60737A}
\usepackage[most]{tcolorbox}
\usepackage{hyperref}
\hypersetup{colorlinks=true,linkcolor=teal,urlcolor=teal}
\usepackage{fancyhdr}
\usepackage{enumitem}
\usepackage{microtype}
\usepackage{needspace}
\usepackage{titlesec}
\usepackage{listings}
\usepackage{tabularx}
\usepackage{longtable}
\usepackage{array}
\setlength{\parindent}{0pt}
\setlength{\parskip}{4pt}
\setlist[enumerate]{topsep=3pt,itemsep=1pt,parsep=0pt}
\titleformat{\section}{\LARGE\bfseries\color{ink}}{\thesection}{0.6em}{}
\titleformat{\subsection}{\Large\bfseries\color{teal}}{\thesubsection}{0.6em}{}
\titleformat{\subsubsection}{\large\bfseries\color{ink}}{\thesubsubsection}{0.5em}{}
\pagestyle{fancy}
\fancyhf{}
\lhead{\footnotesize Grafika va veb - Attestatsiya}
\rhead{\footnotesize 2026}
\cfoot{\footnotesize\thepage}
\renewcommand{\headrulewidth}{0.3pt}
\lstset{basicstyle=\ttfamily\footnotesize,breaklines=true,breakatwhitespace=false,
  columns=fullflexible,frame=single,framesep=4pt,rulecolor=\color{muted},
  xleftmargin=5pt,xrightmargin=5pt,aboveskip=6pt,belowskip=6pt,keepspaces=true,
  showstringspaces=false}
\begin{document}
\begin{titlepage}
\vspace*{30mm}
{\color{teal}\rule{\textwidth}{1.8pt}}\par
\vspace{18mm}
{\Large INFORMATIKA ATTESTATSIYASI \ \ 2026}\par
\vspace{7mm}
{\fontsize{31}{36}\selectfont\bfseries\color{ink}GRAFIKA VA VEB}\par
\vspace{5mm}
{\Large Mavzulashtirilgan testlar kitobi}\par
\vspace{19mm}
{\large Kompyuter grafikasi va media \quad|\quad HTML va CSS}\par
\vspace{13mm}
{\large\textbf{502 ta amaldagi savol}}\par
{\large 125 nashr qilingan, 377 tekshiruvdagi}\par
\vspace{15mm}
{\normalsize Y1 - bitta javob \quad Y2 - moslashtirish \quad Y3 - ketma-ketlik}\par
\vfill
{\color{muted}Barcha savollar avval javobsiz berilgan. \\
Javoblar va bazadagi izohlar kitob oxirida.}\par
\vspace{8mm}
{\color{teal}\rule{\textwidth}{1.8pt}}
\end{titlepage}
\section*{Foydalanish tartibi}
\addcontentsline{toc}{section}{Foydalanish tartibi}
Avval savollarni mustaqil bajaring. Javoblar va izohlar kitobning oxirida joylashgan.
Y1 -- bitta to'g'ri javobni tanlang; Y2 -- chap va o'ng ustunlarni moslang; Y3 -- amallar ketma-ketligini yozing.
\par\textbf{Nashr holati:} I qismdagi 125 savol \textit{published}; II qismdagi 377 savol
\textit{review} holatida. Review holati mustaqil ekspert tasdig'ini anglatmaydi.
452 ta \textit{archived} yozuv yangi faol savol sifatida takroran kiritilmadi.
\par\textbf{Rasmli savollar:} 12 ta review savolida tashqi tasvirga havola berilgan.
PDFdan foydalanayotganda rasm havolasini ochish lozim. Internet bo'lmagan sharoitda
bu savollar to'liq vizual dalilga ega emas.
\par\textbf{Manba:} attestatsiya100\% Supabase, public.questions, question\_options, question\_keys, lessons (2026-10-23 holati).
\tableofcontents
\clearpage
\section{I qism. Nashr qilingan testlar}
\subsection{Kompyuter grafikasi va media}
\subsubsection{Kompyuter grafikasi turlari}
\needspace{5\baselineskip}
\hypertarget{q1}{\textbf{\color{ink}1.}}  {\color{muted}\scriptsize Y1 \textbullet\ daraja 1 \textbullet\ DB 031d74be}
\par
"Rastr grafika" atamasiga eng aniq ta'rifni toping.
\begin{enumerate}[label=\Alph*),leftmargin=9mm]
\item tasvirni piksellar to'ri ko'rinishida saqlaydigan grafika turi
\item tasvirni geometrik obyektlar va formulalar orqali ifodalaydigan grafika turi
\item rastr tasvirning rang qiymatiga ega eng kichik elementi
\item rasm yoki ekranning bir dyuymiga to'g'ri keladigan piksel soni
\end{enumerate}
\vspace{2mm}\hrule height 0.2pt \vspace{4mm}
\needspace{5\baselineskip}
\hypertarget{q2}{\textbf{\color{ink}2.}}  {\color{muted}\scriptsize Y1 \textbullet\ daraja 1 \textbullet\ DB 21445d93}
\par
Quyidagi ta'rif qaysi atamaga tegishli? Tasvirni geometrik obyektlar va formulalar orqali ifodalaydigan grafika turi.
\begin{enumerate}[label=\Alph*),leftmargin=9mm]
\item Piksel
\item Vektor grafika
\item PPI
\item Rang chuqurligi
\end{enumerate}
\vspace{2mm}\hrule height 0.2pt \vspace{4mm}
\needspace{5\baselineskip}
\hypertarget{q3}{\textbf{\color{ink}3.}}  {\color{muted}\scriptsize Y1 \textbullet\ daraja 1 \textbullet\ DB 2b21719c}
\par
"Rastr va vektor farqi" atamasiga eng aniq ta'rifni toping.
\begin{enumerate}[label=\Alph*),leftmargin=9mm]
\item kodek oqimni kodlashi, konteyner esa oqimlarni bir faylda saqlashini ajratish
\item rastr piksel, vektor geometrik obyektlarga asoslanishini ajratish
\item validatsiya qoidaga moslikni, verifikatsiya aslga moslikni tekshirishini ajratish
\item Internet tarmoqlar infratuzilmasi, WWW esa undagi xizmatlardan biri ekanini ajratish
\end{enumerate}
\vspace{2mm}\hrule height 0.2pt \vspace{4mm}
\needspace{5\baselineskip}
\hypertarget{q4}{\textbf{\color{ink}4.}}  {\color{muted}\scriptsize Y1 \textbullet\ daraja 1 \textbullet\ DB a4a4a3e9}
\par
Qaysi rang modelida rang komponentlari ko'pincha o'n oltilik (hex) ko'rinishda ifodalanadi?
\begin{enumerate}[label=\Alph*),leftmargin=9mm]
\item RGB
\item CMYK
\item HSB
\item WBRK
\end{enumerate}
\vspace{2mm}\hrule height 0.2pt \vspace{4mm}
\needspace{5\baselineskip}
\hypertarget{q5}{\textbf{\color{ink}5.}}  {\color{muted}\scriptsize Y1 \textbullet\ daraja 1 \textbullet\ DB da483ac5}
\par
Quyidagi ta'rif qaysi atamaga tegishli? Tasvirni piksellar to'ri ko'rinishida saqlaydigan grafika turi.
\begin{enumerate}[label=\Alph*),leftmargin=9mm]
\item Rastr grafika
\item Vektor grafika
\item Piksel
\item PPI
\end{enumerate}
\vspace{2mm}\hrule height 0.2pt \vspace{4mm}
\needspace{5\baselineskip}
\hypertarget{q6}{\textbf{\color{ink}6.}}  {\color{muted}\scriptsize Y1 \textbullet\ daraja 1 \textbullet\ DB f9a00874}
\par
"Vektor grafika" atamasiga eng aniq ta'rifni toping.
\begin{enumerate}[label=\Alph*),leftmargin=9mm]
\item rastr tasvirning rang qiymatiga ega eng kichik elementi
\item tasvirni geometrik obyektlar va formulalar orqali ifodalaydigan grafika turi
\item rasm yoki ekranning bir dyuymiga to'g'ri keladigan piksel soni
\item bitta piksel rangini kodlash uchun ajratilgan bitlar soni
\end{enumerate}
\vspace{2mm}\hrule height 0.2pt \vspace{4mm}
\needspace{5\baselineskip}
\hypertarget{q7}{\textbf{\color{ink}7.}}  {\color{muted}\scriptsize Y1 \textbullet\ daraja 2 \textbullet\ DB 2a0289ca}
\par
Logotip o'lchami oshirilganda qirralari silliq qoldi. Bu vaziyatdagi asosiy tushunchani aniqlang.
\begin{enumerate}[label=\Alph*),leftmargin=9mm]
\item Piksel
\item Vektor grafika
\item PPI
\item Rang chuqurligi
\end{enumerate}
\vspace{2mm}\hrule height 0.2pt \vspace{4mm}
\needspace{5\baselineskip}
\hypertarget{q8}{\textbf{\color{ink}8.}}  {\color{muted}\scriptsize Y1 \textbullet\ daraja 2 \textbullet\ DB 99abbf48}
\par
Fotosurat kattalashtirilganda alohida kvadrat nuqtalar ko'rina boshladi. Bu vaziyatdagi asosiy tushunchani aniqlang.
\begin{enumerate}[label=\Alph*),leftmargin=9mm]
\item Rastr grafika
\item Vektor grafika
\item Piksel
\item PPI
\end{enumerate}
\vspace{2mm}\hrule height 0.2pt \vspace{4mm}
\needspace{5\baselineskip}
\hypertarget{q9}{\textbf{\color{ink}9.}}  {\color{muted}\scriptsize Y1 \textbullet\ daraja 2 \textbullet\ DB f4c6b8d0}
\par
Quyidagi ta'rif qaysi atamaga tegishli? Rastr piksel, vektor geometrik obyektlarga asoslanishini ajratish.
\begin{enumerate}[label=\Alph*),leftmargin=9mm]
\item Kodek va konteyner farqi
\item Validatsiya va verifikatsiya farqi
\item Internet va WWW farqi
\item Rastr va vektor farqi
\end{enumerate}
\vspace{2mm}\hrule height 0.2pt \vspace{4mm}
\needspace{5\baselineskip}
\hypertarget{q10}{\textbf{\color{ink}10.}}  {\color{muted}\scriptsize Y1 \textbullet\ daraja 3 \textbullet\ DB 0cab7503}
\par
Fotosurat piksellandi, logotip esa kattalashganda silliq qoldi. Bu vaziyatdagi asosiy tushunchani aniqlang.
\begin{enumerate}[label=\Alph*),leftmargin=9mm]
\item Kodek va konteyner farqi
\item Rastr va vektor farqi
\item Validatsiya va verifikatsiya farqi
\item Internet va WWW farqi
\end{enumerate}
\vspace{2mm}\hrule height 0.2pt \vspace{4mm}
\needspace{5\baselineskip}
\hypertarget{q11}{\textbf{\color{ink}11.}}  {\color{muted}\scriptsize Y1 \textbullet\ daraja 3 \textbullet\ DB 64126691}
\par
Fotosuratdagi har bir mayda rang o'tishini tahrirlash kerak. Qaysi grafika turi tabiiyroq tanlov?
\begin{enumerate}[label=\Alph*),leftmargin=9mm]
\item Rastr grafika
\item Faqat vektor chizma
\item Faqat matn fayli
\item Audio format
\end{enumerate}
\vspace{2mm}\hrule height 0.2pt \vspace{4mm}
\needspace{5\baselineskip}
\hypertarget{q12}{\textbf{\color{ink}12.}}  {\color{muted}\scriptsize Y1 \textbullet\ daraja 3 \textbullet\ DB 95fa440a}
\par
Grafik muharrirda chizish uchun varaqsimon ish maydoni, egri chiziq chizish buyrug'i va palitradan tashqari rangni tanlash oynasi so'raldi. Mos atamalar ketma-ketligini tanlang.
\begin{enumerate}[label=\Alph*),leftmargin=9mm]
\item Canvas - Curve - Edit Colors
\item Menu - Line - Colors
\item Selection - Pencil - Layers
\item List - Bezier - Texture
\end{enumerate}
\vspace{2mm}\hrule height 0.2pt \vspace{4mm}
\needspace{5\baselineskip}
\hypertarget{q13}{\textbf{\color{ink}13.}}  {\color{muted}\scriptsize Y1 \textbullet\ daraja 3 \textbullet\ DB c27b8d76}
\par
Grafik muharrir vazifasi qaysi javobda to'g'ri ko'rsatilgan?
\begin{enumerate}[label=\Alph*),leftmargin=9mm]
\item Grafik tasvirlarni kompyuterdan ko'chirib olishga mo'ljallangan amaliy dastur
\item Grafik tasvirlarni yaratish va qayta ishlashga mo'ljallangan amaliy dastur
\item Matnlarni yaratish va qayta ishlashga mo'ljallangan amaliy dastur
\item Jadvallarni yaratish va qayta ishlashga mo'ljallangan amaliy dastur
\end{enumerate}
\vspace{2mm}\hrule height 0.2pt \vspace{4mm}
\needspace{5\baselineskip}
\hypertarget{q14}{\textbf{\color{ink}14.}}  {\color{muted}\scriptsize Y1 \textbullet\ daraja 3 \textbullet\ DB cf9a3f8c}
\par
Logotipni kattalashtirganda qirralar tiniq qolishi kerak. Qaysi grafika turi afzal?
\begin{enumerate}[label=\Alph*),leftmargin=9mm]
\item Vektor grafika
\item Past aniqlikdagi rastr
\item Faqat BMP rastr
\item Skrinshot
\end{enumerate}
\vspace{2mm}\hrule height 0.2pt \vspace{4mm}
\needspace{5\baselineskip}
\hypertarget{q15}{\textbf{\color{ink}15.}}  {\color{muted}\scriptsize Y1 \textbullet\ daraja 3 \textbullet\ DB e2910000}
\par
Bir logotipni tashqi reklamada 15 baravar kattalashtirish va konturlarini mustaqil tahrirlash talab qilinadi. Manba sifatida qaysi fayl turi eng mos?
\begin{enumerate}[label=\Alph*),leftmargin=9mm]
\item SVG vektor manbasi
\item JPEG rastr nusxasi
\item GIF animatsiya nusxasi
\item BMP rastr nusxasi
\end{enumerate}
\vspace{2mm}\hrule height 0.2pt \vspace{4mm}
\needspace{5\baselineskip}
\hypertarget{q16}{\textbf{\color{ink}16.}}  {\color{muted}\scriptsize Y2 \textbullet\ daraja 3 \textbullet\ DB e2910000}
\par
Dizaynerga kerakli grafik manba yoki eksport turini vazifasi bilan to'g'ri juftlang.
\textbf{\footnotesize Chap ustun:}
\begin{enumerate}[label=\arabic*),leftmargin=9mm]
\item SVG
\item PSD
\item PNG
\item JPEG
\end{enumerate}
\textbf{\footnotesize O'ng ustun:}
\begin{enumerate}[label=\Alph*),leftmargin=9mm]
\item Qatlamlari tahrirlanadigan ishchi loyiha
\item Fotosurat uchun yo'qotishli siqilgan tarqatish nusxasi
\item Masshtabga chidamli geometrik logotip manbasi
\item Alpha shaffofligini saqlagan rastr chiqish
\end{enumerate}
\vspace{2mm}\hrule height 0.2pt \vspace{4mm}
\needspace{5\baselineskip}
\hypertarget{q17}{\textbf{\color{ink}17.}}  {\color{muted}\scriptsize Y1 \textbullet\ daraja 4 \textbullet\ DB e2910000}
\par
Dizayner ekran uchun tayyorlagan RGB tasvirini ofset bosmaxonasiga jo'natadi. Qaysi mulohaza eng asosli?
\begin{enumerate}[label=\Alph*),leftmargin=9mm]
\item SVG fayl bo'lsa rang modeli ahamiyatsiz.
\item RGB barcha bosma bo'yoqlarni bevosita bildiradi.
\item Alpha kanali RGB ni avtomatik CMYKga aylantiradi.
\item CMYK bosma jarayoni uchun ishlatiladi; rang moslashtirish kerak bo'lishi mumkin.
\end{enumerate}
\vspace{2mm}\hrule height 0.2pt \vspace{4mm}
\needspace{5\baselineskip}
\hypertarget{q18}{\textbf{\color{ink}18.}}  {\color{muted}\scriptsize Y3 \textbullet\ daraja 4 \textbullet\ DB e2910000}
\par
Shaffof fonli logotipning vektor manbasini saqlab, vebga tayyor nusxa chiqarish uchun amallarning mantiqiy ketma-ketligini tuzing.
\begin{enumerate}[label=\Alph*),leftmargin=9mm]
\item Ranglar va fon/shaffoflik talabini tahrirlang.
\item Tahrirlanadigan SVG manba nusxasini alohida saqlang.
\item Vektor manbada logotipning geometrik obyektlarini yarating.
\item Vebdagi kerakli ko'rinishni eksport qilib, shaffofligini tekshiring.
\end{enumerate}
\vspace{2mm}\hrule height 0.2pt \vspace{4mm}
\needspace{5\baselineskip}
\hypertarget{q19}{\textbf{\color{ink}19.}}  {\color{muted}\scriptsize Y1 \textbullet\ daraja 5 \textbullet\ DB e2910000}
\par
120x80 pikselga ega, har pikseli 24 bit bilan kodlanadigan rastrda faqat xom rang ma'lumoti saqlansa, hajm necha bayt bo'ladi? (Siqish, sarlavha va xizmat baytlari yo'q.)
\begin{enumerate}[label=\Alph*),leftmargin=9mm]
\item 9 600 bayt
\item 38 400 bayt
\item 28 800 bayt
\item 230 400 bayt
\end{enumerate}
\vspace{2mm}\hrule height 0.2pt \vspace{4mm}
\needspace{5\baselineskip}
\hypertarget{q20}{\textbf{\color{ink}20.}}  {\color{muted}\scriptsize Y1 \textbullet\ daraja 5 \textbullet\ DB e2910000}
\par
1920x1080 rasmning eni va bo'yi alohida 50\% ga kamaytirildi. Yangi rasm necha pikseldan tashkil topadi?
\begin{enumerate}[label=\Alph*),leftmargin=9mm]
\item 1 036 800 piksel
\item 518 400 piksel
\item 259 200 piksel
\item 2 073 600 piksel
\end{enumerate}
\vspace{2mm}\hrule height 0.2pt \vspace{4mm}
\subsubsection{Rastr va vektor tasvirlar ustida amallar}
\needspace{5\baselineskip}
\hypertarget{q21}{\textbf{\color{ink}21.}}  {\color{muted}\scriptsize Y1 \textbullet\ daraja 1 \textbullet\ DB 255a9e14}
\par
"Alfa kanal" atamasiga eng aniq ta'rifni toping.
\begin{enumerate}[label=\Alph*),leftmargin=9mm]
\item tasvirni piksellar to'ri ko'rinishida saqlaydigan grafika turi
\item tasvirni geometrik obyektlar va formulalar orqali ifodalaydigan grafika turi
\item rastr tasvirning rang qiymatiga ega eng kichik elementi
\item pikselning shaffoflik darajasini saqlaydigan qo'shimcha kanal
\end{enumerate}
\vspace{2mm}\hrule height 0.2pt \vspace{4mm}
\needspace{5\baselineskip}
\hypertarget{q22}{\textbf{\color{ink}22.}}  {\color{muted}\scriptsize Y1 \textbullet\ daraja 1 \textbullet\ DB 54099ea3}
\par
"PPI" atamasiga eng aniq ta'rifni toping.
\begin{enumerate}[label=\Alph*),leftmargin=9mm]
\item bitta piksel rangini kodlash uchun ajratilgan bitlar soni
\item qizil, yashil va ko'k yorug'likni qo'shishga asoslangan rang modeli
\item sian, magenta, sariq va qora bo'yoqlarni ayrish tamoyiliga asoslangan model
\item rasm yoki ekranning bir dyuymiga to'g'ri keladigan piksel soni
\end{enumerate}
\vspace{2mm}\hrule height 0.2pt \vspace{4mm}
\needspace{5\baselineskip}
\hypertarget{q23}{\textbf{\color{ink}23.}}  {\color{muted}\scriptsize Y1 \textbullet\ daraja 1 \textbullet\ DB 6a0548fa}
\par
Asl landshaft rasm tahrirdan keyin balandligi o'zgarmagan, eni qisqargan va tasvirdagi obyektlar cho'zilmagan; chap-o'ng qismlar kesilgan. Qanday o'zgartirish bajarilgan?
\begin{enumerate}[label=\Alph*),leftmargin=9mm]
\item Portret nisbatida qirqish (Crop - Aspect Ratio - Portrait)
\item Tomonlar nisbatini saqlab o'lchamni o'zgartirish
\item Tomonlar nisbatini saqlamasdan o'lchamni o'zgartirish
\item Landshaft nisbatida qirqish (Crop - Aspect Ratio - Landscape)
\end{enumerate}
\vspace{2mm}\hrule height 0.2pt \vspace{4mm}
\needspace{5\baselineskip}
\hypertarget{q24}{\textbf{\color{ink}24.}}  {\color{muted}\scriptsize Y1 \textbullet\ daraja 1 \textbullet\ DB 87b11ccd}
\par
Tasvir yoki ekranning kengligi bilan balandligi orasidagi nisbat nima deb ataladi?
\begin{enumerate}[label=\Alph*),leftmargin=9mm]
\item Filtrlash
\item Kadr
\item Tomonlar nisbati (mutanosiblik)
\item Vizuallashtirish
\end{enumerate}
\vspace{2mm}\hrule height 0.2pt \vspace{4mm}
\needspace{5\baselineskip}
\hypertarget{q25}{\textbf{\color{ink}25.}}  {\color{muted}\scriptsize Y1 \textbullet\ daraja 1 \textbullet\ DB a2611dca}
\par
"Piksel" atamasiga eng aniq ta'rifni toping.
\begin{enumerate}[label=\Alph*),leftmargin=9mm]
\item rasm yoki ekranning bir dyuymiga to'g'ri keladigan piksel soni
\item bitta piksel rangini kodlash uchun ajratilgan bitlar soni
\item rastr tasvirning rang qiymatiga ega eng kichik elementi
\item qizil, yashil va ko'k yorug'likni qo'shishga asoslangan rang modeli
\end{enumerate}
\vspace{2mm}\hrule height 0.2pt \vspace{4mm}
\needspace{5\baselineskip}
\hypertarget{q26}{\textbf{\color{ink}26.}}  {\color{muted}\scriptsize Y1 \textbullet\ daraja 1 \textbullet\ DB a99c5e9b}
\par
"Rang chuqurligi" atamasiga eng aniq ta'rifni toping.
\begin{enumerate}[label=\Alph*),leftmargin=9mm]
\item bitta piksel rangini kodlash uchun ajratilgan bitlar soni
\item qizil, yashil va ko'k yorug'likni qo'shishga asoslangan rang modeli
\item sian, magenta, sariq va qora bo'yoqlarni ayrish tamoyiliga asoslangan model
\item pikselning shaffoflik darajasini saqlaydigan qo'shimcha kanal
\end{enumerate}
\vspace{2mm}\hrule height 0.2pt \vspace{4mm}
\needspace{5\baselineskip}
\hypertarget{q27}{\textbf{\color{ink}27.}}  {\color{muted}\scriptsize Y1 \textbullet\ daraja 1 \textbullet\ DB ab346b95}
\par
"CMYK" atamasiga eng aniq ta'rifni toping.
\begin{enumerate}[label=\Alph*),leftmargin=9mm]
\item pikselning shaffoflik darajasini saqlaydigan qo'shimcha kanal
\item tasvirni piksellar to'ri ko'rinishida saqlaydigan grafika turi
\item sian, magenta, sariq va qora bo'yoqlarni ayrish tamoyiliga asoslangan model
\item tasvirni geometrik obyektlar va formulalar orqali ifodalaydigan grafika turi
\end{enumerate}
\vspace{2mm}\hrule height 0.2pt \vspace{4mm}
\needspace{5\baselineskip}
\hypertarget{q28}{\textbf{\color{ink}28.}}  {\color{muted}\scriptsize Y1 \textbullet\ daraja 1 \textbullet\ DB f963fb55}
\par
"RGB" atamasiga eng aniq ta'rifni toping.
\begin{enumerate}[label=\Alph*),leftmargin=9mm]
\item sian, magenta, sariq va qora bo'yoqlarni ayrish tamoyiliga asoslangan model
\item qizil, yashil va ko'k yorug'likni qo'shishga asoslangan rang modeli
\item pikselning shaffoflik darajasini saqlaydigan qo'shimcha kanal
\item tasvirni piksellar to'ri ko'rinishida saqlaydigan grafika turi
\end{enumerate}
\vspace{2mm}\hrule height 0.2pt \vspace{4mm}
\needspace{5\baselineskip}
\hypertarget{q29}{\textbf{\color{ink}29.}}  {\color{muted}\scriptsize Y1 \textbullet\ daraja 2 \textbullet\ DB 0e8cc7d7}
\par
Pencil2D da PIXEL so'zi segmentlar bilan chiziladi: P=5, I=2, X=4, E=5, L=3 bo'lak. Har bo'lak 1 kadr, har tugagan harfdan keyin (oxirgisidan tashqari) 1 bo'sh kadr qo'shiladi. 15 FPS da davomiylik qancha?
\begin{enumerate}[label=\Alph*),leftmargin=9mm]
\item 1,53 s
\item 1,5 s
\item 0,65 s
\item 1,7 s
\end{enumerate}
\vspace{2mm}\hrule height 0.2pt \vspace{4mm}
\needspace{5\baselineskip}
\hypertarget{q30}{\textbf{\color{ink}30.}}  {\color{muted}\scriptsize Y1 \textbullet\ daraja 2 \textbullet\ DB 16f098b3}
\par
Quyidagi ta'rif qaysi atamaga tegishli? Sian, magenta, sariq va qora bo'yoqlarni ayrish tamoyiliga asoslangan model.
\begin{enumerate}[label=\Alph*),leftmargin=9mm]
\item Alfa kanal
\item Rastr grafika
\item CMYK
\item Vektor grafika
\end{enumerate}
\vspace{2mm}\hrule height 0.2pt \vspace{4mm}
\needspace{5\baselineskip}
\hypertarget{q31}{\textbf{\color{ink}31.}}  {\color{muted}\scriptsize Y1 \textbullet\ daraja 2 \textbullet\ DB 4187ac43}
\par
Quyidagi ta'rif qaysi atamaga tegishli? Rasm yoki ekranning bir dyuymiga to'g'ri keladigan piksel soni.
\begin{enumerate}[label=\Alph*),leftmargin=9mm]
\item Rang chuqurligi
\item RGB
\item CMYK
\item PPI
\end{enumerate}
\vspace{2mm}\hrule height 0.2pt \vspace{4mm}
\needspace{5\baselineskip}
\hypertarget{q32}{\textbf{\color{ink}32.}}  {\color{muted}\scriptsize Y1 \textbullet\ daraja 2 \textbullet\ DB 4ffbb820}
\par
Quyidagi ta'rif qaysi atamaga tegishli? Rastr tasvirning rang qiymatiga ega eng kichik elementi.
\begin{enumerate}[label=\Alph*),leftmargin=9mm]
\item PPI
\item Rang chuqurligi
\item Piksel
\item RGB
\end{enumerate}
\vspace{2mm}\hrule height 0.2pt \vspace{4mm}
\needspace{5\baselineskip}
\hypertarget{q33}{\textbf{\color{ink}33.}}  {\color{muted}\scriptsize Y1 \textbullet\ daraja 2 \textbullet\ DB 754ccca4}
\par
Quyidagi ta'rif qaysi atamaga tegishli? Pikselning shaffoflik darajasini saqlaydigan qo'shimcha kanal.
\begin{enumerate}[label=\Alph*),leftmargin=9mm]
\item Rastr grafika
\item Vektor grafika
\item Piksel
\item Alfa kanal
\end{enumerate}
\vspace{2mm}\hrule height 0.2pt \vspace{4mm}
\needspace{5\baselineskip}
\hypertarget{q34}{\textbf{\color{ink}34.}}  {\color{muted}\scriptsize Y1 \textbullet\ daraja 2 \textbullet\ DB 7859c0a9}
\par
Quyidagi ta'rif qaysi atamaga tegishli? Qizil, yashil va ko'k yorug'likni qo'shishga asoslangan rang modeli.
\begin{enumerate}[label=\Alph*),leftmargin=9mm]
\item CMYK
\item RGB
\item Alfa kanal
\item Rastr grafika
\end{enumerate}
\vspace{2mm}\hrule height 0.2pt \vspace{4mm}
\needspace{5\baselineskip}
\hypertarget{q35}{\textbf{\color{ink}35.}}  {\color{muted}\scriptsize Y1 \textbullet\ daraja 2 \textbullet\ DB cce6a8c0}
\par
Pencil2D animatsiyasida TEMPO so'zining tasvirdagi segmentlari T=4, E=5, M=5, P=5, O=6 ta elementdan iborat. Har bir element 2 kadrda chiziladi; har bir tugallangan harfdan keyin, oxirgisidan tashqari, 3 tadan bo'sh kadr qo'yiladi. Ijro tezligi 10 FPS bo'lsa, animatsiya qancha davom etadi?
\begin{enumerate}[label=\Alph*),leftmargin=9mm]
\item 6,2 soniya
\item 6,(2) soniya
\item 5 soniya
\item 3,7 soniya
\end{enumerate}
\vspace{2mm}\hrule height 0.2pt \vspace{4mm}
\needspace{5\baselineskip}
\hypertarget{q36}{\textbf{\color{ink}36.}}  {\color{muted}\scriptsize Y1 \textbullet\ daraja 2 \textbullet\ DB d32e942f}
\par
300 PPI bosma rasmda har dyuymga 300 piksel joylashdi. Bu vaziyatdagi asosiy tushunchani aniqlang.
\begin{enumerate}[label=\Alph*),leftmargin=9mm]
\item Rang chuqurligi
\item RGB
\item CMYK
\item PPI
\end{enumerate}
\vspace{2mm}\hrule height 0.2pt \vspace{4mm}
\needspace{5\baselineskip}
\hypertarget{q37}{\textbf{\color{ink}37.}}  {\color{muted}\scriptsize Y1 \textbullet\ daraja 2 \textbullet\ DB e80a0b68}
\par
Quyidagi ta'rif qaysi atamaga tegishli? Bitta piksel rangini kodlash uchun ajratilgan bitlar soni.
\begin{enumerate}[label=\Alph*),leftmargin=9mm]
\item Rang chuqurligi
\item RGB
\item CMYK
\item Alfa kanal
\end{enumerate}
\vspace{2mm}\hrule height 0.2pt \vspace{4mm}
\needspace{5\baselineskip}
\hypertarget{q38}{\textbf{\color{ink}38.}}  {\color{muted}\scriptsize Y1 \textbullet\ daraja 2 \textbullet\ DB ebe7dc39}
\par
Ekrandagi bitta nuqtaning RGB qiymati o'zgartirildi. Bu vaziyatdagi asosiy tushunchani aniqlang.
\begin{enumerate}[label=\Alph*),leftmargin=9mm]
\item PPI
\item Rang chuqurligi
\item Piksel
\item RGB
\end{enumerate}
\vspace{2mm}\hrule height 0.2pt \vspace{4mm}
\needspace{5\baselineskip}
\hypertarget{q39}{\textbf{\color{ink}39.}}  {\color{muted}\scriptsize Y1 \textbullet\ daraja 3 \textbullet\ DB 3c051833}
\par
Monitor uch rang kanalining intensivligini o'zgartirib rang hosil qildi. Vaziyatni tahlil qilib, unga eng aniq mos tushunchani toping.
\begin{enumerate}[label=\Alph*),leftmargin=9mm]
\item CMYK
\item Alfa kanal
\item Rastr grafika
\item RGB
\end{enumerate}
\vspace{2mm}\hrule height 0.2pt \vspace{4mm}
\needspace{5\baselineskip}
\hypertarget{q40}{\textbf{\color{ink}40.}}  {\color{muted}\scriptsize Y1 \textbullet\ daraja 3 \textbullet\ DB 5f8612ce}
\par
24 bitli tasvirda har piksel uchun uch bayt ajratildi. Vaziyatni tahlil qilib, unga eng aniq mos tushunchani toping.
\begin{enumerate}[label=\Alph*),leftmargin=9mm]
\item RGB
\item CMYK
\item Rang chuqurligi
\item Alfa kanal
\end{enumerate}
\vspace{2mm}\hrule height 0.2pt \vspace{4mm}
\needspace{5\baselineskip}
\hypertarget{q41}{\textbf{\color{ink}41.}}  {\color{muted}\scriptsize Y2 \textbullet\ daraja 3 \textbullet\ DB b572e1f0}
\par
Chapdagi Zoom, Crop, Resize va Canvas size buyruqlarini o'ngdagi tasvirga ta'siri bilan moslashtiring. Piksel soni, ishchi maydon va ekrandagi ko'rinish o'zgarishini farqlang.
\textbf{\footnotesize Chap ustun:}
\begin{enumerate}[label=\arabic*),leftmargin=9mm]
\item Zoom
\item Crop
\item Resize
\item Canvas size
\end{enumerate}
\textbf{\footnotesize O'ng ustun:}
\begin{enumerate}[label=\Alph*),leftmargin=9mm]
\item Rasmning ishchi maydoni chegarasini o'zgartirish
\item Belgilangan kadrlash chegarasidan tashqari qismini olib tashlash
\item Rastrdagi eni va bo'yi pikselda qayta namunalash
\item Manba piksel sonini o'zgartirmay ko'rinishni kattalashtirish
\end{enumerate}
\vspace{2mm}\hrule height 0.2pt \vspace{4mm}
\needspace{5\baselineskip}
\hypertarget{q42}{\textbf{\color{ink}42.}}  {\color{muted}\scriptsize Y1 \textbullet\ daraja 3 \textbullet\ DB d553ba7a}
\par
Rastr rasmni juda kattalashtirganda kvadrat nuqtalar ko'rina boshladi. Buning sababi nima?
\begin{enumerate}[label=\Alph*),leftmargin=9mm]
\item Piksellar kattalashib ko'rinmoqda
\item Vektor formulalari ko'paymoqda
\item Audio kanal kamaymoqda
\item HTML tegi buzildi
\end{enumerate}
\vspace{2mm}\hrule height 0.2pt \vspace{4mm}
\needspace{5\baselineskip}
\hypertarget{q43}{\textbf{\color{ink}43.}}  {\color{muted}\scriptsize Y1 \textbullet\ daraja 3 \textbullet\ DB e2910000}
\par
600x400 pikselli rasmga Zoom=200\% qo'llandi, lekin Resize bajarilmadi. Manba rasmning piksel soni qanday o'zgaradi?
\begin{enumerate}[label=\Alph*),leftmargin=9mm]
\item 240 000 bo'lib qoladi.
\item PPI ikki baravar ortadi.
\item 960 000 taga aylanadi.
\item 120 000 taga kamayadi.
\end{enumerate}
\vspace{2mm}\hrule height 0.2pt \vspace{4mm}
\needspace{5\baselineskip}
\hypertarget{q44}{\textbf{\color{ink}44.}}  {\color{muted}\scriptsize Y1 \textbullet\ daraja 3 \textbullet\ DB eb58fddc}
\par
Rasmning keraksiz chet qismini olib tashlash kerak. Qaysi amal mos?
\begin{enumerate}[label=\Alph*),leftmargin=9mm]
\item Crop (kesish)
\item Rotate faqat
\item Zoom faqat
\item Undo tarixini o'chirish
\end{enumerate}
\vspace{2mm}\hrule height 0.2pt \vspace{4mm}
\needspace{5\baselineskip}
\hypertarget{q45}{\textbf{\color{ink}45.}}  {\color{muted}\scriptsize Y1 \textbullet\ daraja 4 \textbullet\ DB 4ae973a6}
\par
RGB modelida qora rang qanday hosil bo'ladi?
\begin{enumerate}[label=\Alph*),leftmargin=9mm]
\item R=G=B=255
\item faqat K=100\%
\item R=G=B=0
\item alfa=0 bo'lsa har doim
\end{enumerate}
\vspace{2mm}\hrule height 0.2pt \vspace{4mm}
\needspace{5\baselineskip}
\hypertarget{q46}{\textbf{\color{ink}46.}}  {\color{muted}\scriptsize Y1 \textbullet\ daraja 4 \textbullet\ DB ae15590a}
\par
Bosmaxona qog'ozdagi rangni to'rtta bo'yoq kanali bilan tayyorladi. Yaqin tushunchalarni farqlab, aynan qaysi atama vaziyatni eng aniq ifodalashini toping.
\begin{enumerate}[label=\Alph*),leftmargin=9mm]
\item Alfa kanal
\item Rastr grafika
\item Vektor grafika
\item CMYK
\end{enumerate}
\vspace{2mm}\hrule height 0.2pt \vspace{4mm}
\needspace{5\baselineskip}
\hypertarget{q47}{\textbf{\color{ink}47.}}  {\color{muted}\scriptsize Y1 \textbullet\ daraja 4 \textbullet\ DB e2910000}
\par
Rastr rasm eni 900 piksel. Uni 300 PPI aniqlikda bosib chiqarish rejalashtirildi. Bosmadagi eni nechaga teng?
\begin{enumerate}[label=\Alph*),leftmargin=9mm]
\item 9 dyuym
\item 0,3 dyuym
\item 300 dyuym
\item 3 dyuym
\end{enumerate}
\vspace{2mm}\hrule height 0.2pt \vspace{4mm}
\needspace{5\baselineskip}
\hypertarget{q48}{\textbf{\color{ink}48.}}  {\color{muted}\scriptsize Y3 \textbullet\ daraja 4 \textbullet\ DB e2910000}
\par
Rasmning 300x200 piksellik qismidan veb uchun ixcham nusxa yaratish ish tartibini belgilang.
\begin{enumerate}[label=\Alph*),leftmargin=9mm]
\item Mos JPEG/PNG formatida eksport qiling va natijani tekshiring.
\item Asl yuqori sifatli rasmni oching va uning nusxasini saqlang.
\item Crop amali bilan belgilangan maydondan tashqarini olib tashlang.
\item Select orqali kerakli 300x200 hududni belgilang.
\end{enumerate}
\vspace{2mm}\hrule height 0.2pt \vspace{4mm}
\needspace{5\baselineskip}
\hypertarget{q49}{\textbf{\color{ink}49.}}  {\color{muted}\scriptsize Y1 \textbullet\ daraja 5 \textbullet\ DB a34cc4ff}
\par
Shaffof fonli logotip va o'lchamni sifat yo'qotmasdan o'zgartirish uchun eng mos tanlov qaysi?
\begin{enumerate}[label=\Alph*),leftmargin=9mm]
\item BMP rastr
\item SVG vektor
\item JPEG rastr
\item 24 bitli WAV
\end{enumerate}
\vspace{2mm}\hrule height 0.2pt \vspace{4mm}
\needspace{5\baselineskip}
\hypertarget{q50}{\textbf{\color{ink}50.}}  {\color{muted}\scriptsize Y1 \textbullet\ daraja 5 \textbullet\ DB cc8f21ff}
\par
Logotip fon bilan yarim shaffof aralashdi. Yaqin tushunchalarni farqlab, aynan qaysi atama vaziyatni eng aniq ifodalashini toping.
\begin{enumerate}[label=\Alph*),leftmargin=9mm]
\item Alfa kanal
\item Rastr grafika
\item Vektor grafika
\item Piksel
\end{enumerate}
\vspace{2mm}\hrule height 0.2pt \vspace{4mm}
\needspace{5\baselineskip}
\hypertarget{q51}{\textbf{\color{ink}51.}}  {\color{muted}\scriptsize Y1 \textbullet\ daraja 5 \textbullet\ DB e2910000}
\par
800x600 pikselli suratda aynan 600x500 piksellik ichki soha Crop qilindi. Crop'dan so'ng qoladigan piksel sonini toping.
\begin{enumerate}[label=\Alph*),leftmargin=9mm]
\item 480 000
\item 360 000
\item 300 000
\item 120 000
\end{enumerate}
\vspace{2mm}\hrule height 0.2pt \vspace{4mm}
\needspace{5\baselineskip}
\hypertarget{q52}{\textbf{\color{ink}52.}}  {\color{muted}\scriptsize Y1 \textbullet\ daraja 5 \textbullet\ DB e2910000}
\par
Ikki vektor to'rtburchak qisman ustma-ust joylashgan. Intersection amali bajarilsa, qaysi hudud qoladi?
\begin{enumerate}[label=\Alph*),leftmargin=9mm]
\item Ikki shaklning barcha hududlari birlashmasi
\item Faqat ustma-ust tushgan umumiy hudud
\item Birinchi shakldan ikkinchisining ayirmasi
\item Ikkala shakl faqat bir guruhga birlashadi
\end{enumerate}
\vspace{2mm}\hrule height 0.2pt \vspace{4mm}
\subsubsection{MS Paint va Adobe Photoshop: asosiy amallar}
\needspace{5\baselineskip}
\hypertarget{q53}{\textbf{\color{ink}53.}}  {\color{muted}\scriptsize Y1 \textbullet\ daraja 2 \textbullet\ DB 7a30d683}
\par
MS Paintdagi rasm yakunda dastlabki holatiga nisbatan 180 deg  ga aylangan ko'rinishga kelishi kerak. Qaysi ikki buyruq ketma-ketligi aynan shu natijani beradi?
\begin{enumerate}[label=\Alph*),leftmargin=9mm]
\item Rotate 180 deg  =\textgreater{} Flip vertical
\item Flip vertical =\textgreater{} Flip horizontal
\item Rotate left 90 deg  =\textgreater{} Flip vertical
\item Flip horizontal =\textgreater{} Rotate 180 deg 
\end{enumerate}
\vspace{2mm}\hrule height 0.2pt \vspace{4mm}
\needspace{5\baselineskip}
\hypertarget{q54}{\textbf{\color{ink}54.}}  {\color{muted}\scriptsize Y1 \textbullet\ daraja 3 \textbullet\ DB 0692c20e}
\par
Adobe Photoshop'da asl rasm qatlamini saqlab qolgan holda tajriba tariqasida tahrir qilish kerak. Eng xavfsiz birinchi qadam qaysi?
\begin{enumerate}[label=\Alph*),leftmargin=9mm]
\item Asl qatlamni darhol o'chirish
\item Asl qatlamdan nusxa olib, nusxa ustida ishlash
\item Barcha qatlamlarni birlashtirish
\item Rasmni faqat JPEG formatida saqlash
\end{enumerate}
\vspace{2mm}\hrule height 0.2pt \vspace{4mm}
\needspace{5\baselineskip}
\hypertarget{q55}{\textbf{\color{ink}55.}}  {\color{muted}\scriptsize Y1 \textbullet\ daraja 3 \textbullet\ DB 06bcb05b}
\par
Photoshop'da portretning orqa fonini vaqtincha yashirib, keyin uni yana tiklash imkonini saqlash kerak. Asl piksel ma'lumotini o'chirmasdan qaysi amal maqsadga muvofiq?
\begin{enumerate}[label=\Alph*),leftmargin=9mm]
\item Asl qavatni Eraser bilan butunlay o'chirish
\item Layer Mask qo'shib yashiriladigan sohaga qora rang bilan bo'yash
\item Layerlarni birlashtirib Flatten qilish
\item Faylni faqat JPEG sifatida saqlash
\end{enumerate}
\vspace{2mm}\hrule height 0.2pt \vspace{4mm}
\needspace{5\baselineskip}
\hypertarget{q56}{\textbf{\color{ink}56.}}  {\color{muted}\scriptsize Y2 \textbullet\ daraja 3 \textbullet\ DB 13f94a6a}
\par
Photoshop ish muhiti qismlarini vazifasi bilan moslang.
\textbf{\footnotesize Chap ustun:}
\begin{enumerate}[label=\arabic*),leftmargin=9mm]
\item Tools panel
\item Options bar
\item Layers panel
\item Document window
\end{enumerate}
\textbf{\footnotesize O'ng ustun:}
\begin{enumerate}[label=\Alph*),leftmargin=9mm]
\item Joriy tanlangan asbob sozlamalari
\item Asboblar to'plami
\item Rasmning ishchi hujjati ko'rinadigan soha
\item Qatlam ko'rinishi va tartibini boshqarish
\end{enumerate}
\vspace{2mm}\hrule height 0.2pt \vspace{4mm}
\needspace{5\baselineskip}
\hypertarget{q57}{\textbf{\color{ink}57.}}  {\color{muted}\scriptsize Y1 \textbullet\ daraja 3 \textbullet\ DB 2fa097c1}
\par
Photoshop'da qatlam vaqtincha ko'rinmasin, lekin o'chirilmasin. Nima qilish kerak?
\begin{enumerate}[label=\Alph*),leftmargin=9mm]
\item Layer visibility ni o'chirish
\item Layerni Delete qilish
\item Rasmni Crop qilish
\item Faylni yopish
\end{enumerate}
\vspace{2mm}\hrule height 0.2pt \vspace{4mm}
\needspace{5\baselineskip}
\hypertarget{q58}{\textbf{\color{ink}58.}}  {\color{muted}\scriptsize Y1 \textbullet\ daraja 3 \textbullet\ DB 6a7ef315}
\par
Photoshop'da obyektni fon ustida alohida tahrirlash va keyin joyini o'zgartirish kerak. Eng qulay yondashuv qaysi?
\begin{enumerate}[label=\Alph*),leftmargin=9mm]
\item Obyektni alohida qatlamda saqlash
\item Barcha qatlamni darhol birlashtirish
\item Faqat JPEGga eksport qilish
\item Har bir amal oldidan dasturni yopish
\end{enumerate}
\vspace{2mm}\hrule height 0.2pt \vspace{4mm}
\needspace{5\baselineskip}
\hypertarget{q59}{\textbf{\color{ink}59.}}  {\color{muted}\scriptsize Y1 \textbullet\ daraja 3 \textbullet\ DB e2910000}
\par
Photoshop qatlamiga oq fonli niqob (Layer Mask) qo'shildi. Undagi kichik soha qora rang bilan bo'yaldi. Odatdagi maska rejimida nima bo'ladi?
\begin{enumerate}[label=\Alph*),leftmargin=9mm]
\item Qora bo'yalgan soha vaqtincha yashirinadi, asl piksel saqlanishi mumkin.
\item Qora soha fotosuratga qaytarib bo'lmaydigan tarzda qo'shiladi.
\item Qora soha avtomatik oppoq rangli matnga aylanadi.
\item Butun ishchi fayl avtomatik JPEGga aylanadi.
\end{enumerate}
\vspace{2mm}\hrule height 0.2pt \vspace{4mm}
\needspace{5\baselineskip}
\hypertarget{q60}{\textbf{\color{ink}60.}}  {\color{muted}\scriptsize Y2 \textbullet\ daraja 3 \textbullet\ DB e2910000}
\par
Grafik muharrirlardagi ushbu to'rtta asbob va vositani ularning amaliy vazifasi bilan moslashtiring. Har bir vositaga bitta mos ta'rifni tanlang.
\textbf{\footnotesize Chap ustun:}
\begin{enumerate}[label=\arabic*),leftmargin=9mm]
\item Paint Color Picker
\item Photoshop Lasso
\item Layer Mask
\item Clone Stamp
\end{enumerate}
\textbf{\footnotesize O'ng ustun:}
\begin{enumerate}[label=\Alph*),leftmargin=9mm]
\item Erkin shaklli hududni tanlash
\item Pikselni o'chirmasdan ko'rinishni cheklash
\item Rasmdagi mavjud rangni namuna olish
\item Rasmning boshqa sohasidan namuna olib ko'chirish
\end{enumerate}
\vspace{2mm}\hrule height 0.2pt \vspace{4mm}
\needspace{5\baselineskip}
\hypertarget{q61}{\textbf{\color{ink}61.}}  {\color{muted}\scriptsize Y1 \textbullet\ daraja 3 \textbullet\ DB ee80805f}
\par
Photoshop'ning Windows versiyasida rasmning tanlangan obyektini joyidan siljitib, shu bilan birga uning asl nusxasini eski joyida qoldirish kerak. Qaysi usul mos?
\begin{enumerate}[label=\Alph*),leftmargin=9mm]
\item Move Tool bilan Alt tugmasini bosib tanlangan sohani sudrash
\item Crop Tool bilan sohaga chegaralar chizish
\item Ctrl+Z bilan oxirgi amalni bekor qilish
\item Image \textgreater{} Flatten Image ni tanlash
\end{enumerate}
\vspace{2mm}\hrule height 0.2pt \vspace{4mm}
\needspace{5\baselineskip}
\hypertarget{q62}{\textbf{\color{ink}62.}}  {\color{muted}\scriptsize Y3 \textbullet\ daraja 4 \textbullet\ DB 005f436e}
\par
Paintda asl tasvirni saqlagan holda yangi nomda nusxa chiqarishni tartiblang.
\begin{enumerate}[label=\Alph*),leftmargin=9mm]
\item Kerakli belgilash yoki chizish tahrirlarini bajaring
\item Save As buyrug'ida yangi nom va mos formatni belgilang
\item Yangi fayl joylashuvi va tahrirlangan tasvirni tekshiring
\item Asl rasmni Paintda oching
\end{enumerate}
\vspace{2mm}\hrule height 0.2pt \vspace{4mm}
\needspace{5\baselineskip}
\hypertarget{q63}{\textbf{\color{ink}63.}}  {\color{muted}\scriptsize Y1 \textbullet\ daraja 4 \textbullet\ DB 0b8a993a}
\par
Photoshop'da oddiy qatlamga aylantirilgan tasvirni Crop Tool bilan qirqish kerak. Kesilgan chetlarni keyin qayta tiklash uchun Crop sozlamasini qanday qo'yish lozim?
\begin{enumerate}[label=\Alph*),leftmargin=9mm]
\item Delete Cropped Pixels tanlovini yoqish
\item Crop o'rniga Flatten Image qilish
\item Delete Cropped Pixels tanlovini o'chirib qo'yish
\item Qirqishdan so'ng faqat JPEGga eksport qilish
\end{enumerate}
\vspace{2mm}\hrule height 0.2pt \vspace{4mm}
\needspace{5\baselineskip}
\hypertarget{q64}{\textbf{\color{ink}64.}}  {\color{muted}\scriptsize Y2 \textbullet\ daraja 4 \textbullet\ DB 1cf03d19}
\par
Photoshop qatlam boshqaruvini natijasi bilan moslang.
\textbf{\footnotesize Chap ustun:}
\begin{enumerate}[label=\arabic*),leftmargin=9mm]
\item Eye icon
\item Opacity
\item Lock
\item Bring layer forward
\end{enumerate}
\textbf{\footnotesize O'ng ustun:}
\begin{enumerate}[label=\Alph*),leftmargin=9mm]
\item Qatlam tahririni cheklash
\item Qatlamni ustma-ust ko'rinishda yuqoriroq qilish
\item Qatlamning umumiy shaffofligini sozlash
\item Qatlam ko'rinishini yoqish yoki o'chirish
\end{enumerate}
\vspace{2mm}\hrule height 0.2pt \vspace{4mm}
\needspace{5\baselineskip}
\hypertarget{q65}{\textbf{\color{ink}65.}}  {\color{muted}\scriptsize Y3 \textbullet\ daraja 4 \textbullet\ DB 4a1d83cd}
\par
Photoshopda layer mask bilan fonning bir qismini yashirishni tartiblang.
\begin{enumerate}[label=\Alph*),leftmargin=9mm]
\item Natijani tekshiring va kerak bo'lsa oq bilan qayta ko'rsating
\item Kerakli foto qatlamini tanlang
\item Layer mask qo'shing
\item Maskada yashiriladigan joyni qora bilan bo'yang
\end{enumerate}
\vspace{2mm}\hrule height 0.2pt \vspace{4mm}
\needspace{5\baselineskip}
\hypertarget{q66}{\textbf{\color{ink}66.}}  {\color{muted}\scriptsize Y2 \textbullet\ daraja 4 \textbullet\ DB 8ebe0aac}
\par
Photoshop non-destructive ishlov berish vositalarini moslang.
\textbf{\footnotesize Chap ustun:}
\begin{enumerate}[label=\arabic*),leftmargin=9mm]
\item Layer mask
\item Adjustment layer
\item Smart Object
\item Duplicate layer
\end{enumerate}
\textbf{\footnotesize O'ng ustun:}
\begin{enumerate}[label=\Alph*),leftmargin=9mm]
\item Ayrim transformatsiyalarni qayta sozlashga mos obyekt
\item Rang yoki tonni asl pikselni bevosita o'zgartirmay tuzatish
\item Oq-qora niqob orqali ko'rinishni boshqarish
\item Asl qatlam nusxasidan alohida tahrir qilish
\end{enumerate}
\vspace{2mm}\hrule height 0.2pt \vspace{4mm}
\needspace{5\baselineskip}
\hypertarget{q67}{\textbf{\color{ink}67.}}  {\color{muted}\scriptsize Y3 \textbullet\ daraja 4 \textbullet\ DB ab17deb0}
\par
Photoshopda obyektni aniq tanlab alohida qatlamga olish ish tartibini tuzing.
\begin{enumerate}[label=\Alph*),leftmargin=9mm]
\item Tasvirda ajratiladigan obyekt chegarasini aniqlang
\item Tanlangan obyektni nusxalab alohida qatlamga joylashtiring
\item Tanlov chegarasini zaruratga ko'ra tekshirib tuzating
\item Marquee, Lasso yoki Quick Selectiondan mos asbobni tanlang
\end{enumerate}
\vspace{2mm}\hrule height 0.2pt \vspace{4mm}
\needspace{5\baselineskip}
\hypertarget{q68}{\textbf{\color{ink}68.}}  {\color{muted}\scriptsize Y1 \textbullet\ daraja 4 \textbullet\ DB e2910000}
\par
Photoshopda fon va tasvir alohida qatlamlarda turibdi. Matnning obyekt ustida ko'rinishi kerak. Qanday qilish eng to'g'ri?
\begin{enumerate}[label=\Alph*),leftmargin=9mm]
\item Barcha qatlamlarni Lock qilib matnni yuqoriga olib chiqish.
\item Matnni JPG formatda saqlash bilan ustiga chiqadi.
\item Hujjatning Zoom qiymatini kamaytirish.
\item Matn qatlamini obyekt qatlamidan yuqoriga ko'tarish va ko'rinishini yoqish.
\end{enumerate}
\vspace{2mm}\hrule height 0.2pt \vspace{4mm}
\needspace{5\baselineskip}
\hypertarget{q69}{\textbf{\color{ink}69.}}  {\color{muted}\scriptsize Y3 \textbullet\ daraja 4 \textbullet\ DB e2910000}
\par
Photoshopda rasm fonini asl obyektga zarar yetkazmasdan vaqtincha yashirish jarayonini tartiblang.
\begin{enumerate}[label=\Alph*),leftmargin=9mm]
\item Reveal Selection rejimida Layer Mask qo'shing va tanlovdan tashqaridagi fon yashirilganini tekshiring.
\item Rasmni ochib, tahrir qilinadigan qatlamni ajrating.
\item Qatlamli PSD manbani saqlab, tekshirish uchun alohida eksport qiling.
\item Ko'rinishi saqlanadigan asosiy obyektni belgilash asbobi bilan tanlang.
\end{enumerate}
\vspace{2mm}\hrule height 0.2pt \vspace{4mm}
\needspace{5\baselineskip}
\hypertarget{q70}{\textbf{\color{ink}70.}}  {\color{muted}\scriptsize Y3 \textbullet\ daraja 5 \textbullet\ DB 7e4b62cd}
\par
Photoshopda rangini buzmasdan tuzatish va eksport tartibini tuzing.
\begin{enumerate}[label=\Alph*),leftmargin=9mm]
\item Asl foto qatlamini tayyor holda saqlang
\item Adjustment layer yarating
\item Brightness/Contrast yoki Hue/Saturation sozlamalarini o'zgartiring
\item PSD manbani saqlab, kerakli PNG/JPEG eksportini oling
\end{enumerate}
\vspace{2mm}\hrule height 0.2pt \vspace{4mm}
\needspace{5\baselineskip}
\hypertarget{q71}{\textbf{\color{ink}71.}}  {\color{muted}\scriptsize Y1 \textbullet\ daraja 5 \textbullet\ DB e2910000}
\par
MS Paint dasturida Curved line (Egri chiziq) asbobi qo'llanmoqda. To'g'ri ish ketma-ketligini aniqlang.
\begin{enumerate}[label=\Alph*),leftmargin=9mm]
\item Curved line faqat yopiq ko'pburchakni chizadi.
\item Avval qalam bilan 3 ta alohida nuqta chizib, Fill qilinadi.
\item Dastlab to'g'ri segment tortiladi, keyin kerakli joylaridan egiladi.
\item Birinchi egri chiziqni Eraser bilan tayyorlash kerak.
\end{enumerate}
\vspace{2mm}\hrule height 0.2pt \vspace{4mm}
\needspace{5\baselineskip}
\hypertarget{q72}{\textbf{\color{ink}72.}}  {\color{muted}\scriptsize Y1 \textbullet\ daraja 5 \textbullet\ DB e2910000}
\par
Photoshopda qatlamli plakatni ertaga qayta tahrirlash rejalashtirilgan, ayni paytda mijozga bitta tekis rasm jo'natish kerak. Eng ma'qul ikkita saqlash natijasi qaysi?
\begin{enumerate}[label=\Alph*),leftmargin=9mm]
\item Faqat JPEG, chunki barcha qatlamlar saqlanadi.
\item PSD ishchi loyiha va PNG/JPEG eksport nusxasi.
\item HTML va CSS, chunki Photoshop shu formatlarni majburiy saqlaydi.
\item Faqat BMP, chunki maskalar BMPda qatlamlarga aylanadi.
\end{enumerate}
\vspace{2mm}\hrule height 0.2pt \vspace{4mm}
\needspace{5\baselineskip}
\hypertarget{q73}{\textbf{\color{ink}73.}}  {\color{muted}\scriptsize Y2 \textbullet\ daraja 5 \textbullet\ DB f911c8ee}
\par
Tasvirga ta'sir qiladigan Photoshop sozlamalarini moslang.
\textbf{\footnotesize Chap ustun:}
\begin{enumerate}[label=\arabic*),leftmargin=9mm]
\item Brightness
\item Contrast
\item Hue/Saturation
\item Gaussian Blur
\end{enumerate}
\textbf{\footnotesize O'ng ustun:}
\begin{enumerate}[label=\Alph*),leftmargin=9mm]
\item Och va to'q sohalar farqi
\item Konturlarni yumshatish
\item Umumiy yorug'lik darajasi
\item Rang toni va to'yinganligi
\end{enumerate}
\vspace{2mm}\hrule height 0.2pt \vspace{4mm}
\subsubsection{MS Paint va Photoshop - nazariya V2}
\needspace{5\baselineskip}
\hypertarget{q74}{\textbf{\color{ink}74.}}  {\color{muted}\scriptsize Y1 \textbullet\ daraja 1 \textbullet\ DB 97dc1781}
\par
Photoshopda tasvir shaklini geometrik buzish yoki egish uchun qaysi filtrlar guruhi ishlatiladi?
\begin{enumerate}[label=\Alph*),leftmargin=9mm]
\item Blur
\item Sharpen
\item Distort
\item Clone
\end{enumerate}
\vspace{2mm}\hrule height 0.2pt \vspace{4mm}
\needspace{5\baselineskip}
\hypertarget{q75}{\textbf{\color{ink}75.}}  {\color{muted}\scriptsize Y1 \textbullet\ daraja 2 \textbullet\ DB 617050a3}
\par
Paintdagi rasm yakunda chap-o'ng yo'nalishda akslantirilgan, ammo tepasi va pasti o'zgarmagan holatga kelishi kerak. Qaysi ikki buyruq ketma-ketligi shu natijani beradi?
\begin{enumerate}[label=\Alph*),leftmargin=9mm]
\item Rotate 180 deg  =\textgreater{} Flip horizontal
\item Rotate 90 deg  =\textgreater{} Flip horizontal
\item Rotate 180 deg  =\textgreater{} Flip vertical
\item Flip horizontal =\textgreater{} Rotate 180 deg 
\end{enumerate}
\vspace{2mm}\hrule height 0.2pt \vspace{4mm}
\needspace{5\baselineskip}
\hypertarget{q76}{\textbf{\color{ink}76.}}  {\color{muted}\scriptsize Y1 \textbullet\ daraja 2 \textbullet\ DB ebb0a916}
\par
Paintdagi rasm yakunda faqat chap-o'ng yo'nalishda akslantirilgan, tepasi va pasti esa o'zgarmagan holatga kelishi kerak. Qaysi ikki buyruq ketma-ketligi shu natijani beradi?
\begin{enumerate}[label=\Alph*),leftmargin=9mm]
\item Rotate 180 deg  =\textgreater{} Flip horizontal
\item Rotate 90 deg  =\textgreater{} Flip horizontal
\item Rotate 180 deg  =\textgreater{} Flip vertical
\item Flip horizontal =\textgreater{} Rotate 180 deg 
\end{enumerate}
\vspace{2mm}\hrule height 0.2pt \vspace{4mm}
\subsection{HTML va CSS}
\subsubsection{HTML: matn, rasm va ro'yxatlar}
\needspace{5\baselineskip}
\hypertarget{q77}{\textbf{\color{ink}77.}}  {\color{muted}\scriptsize Y1 \textbullet\ daraja 1 \textbullet\ DB 9c293464}
\par
24-bitli RGB rang yozuvida \texttt{\textless{}body bgcolor="\#EE00EE"\textgreater{}} qanday rangni bildiradi?
\begin{enumerate}[label=\Alph*),leftmargin=9mm]
\item Yellow
\item Orange
\item Magenta
\item Green
\end{enumerate}
\vspace{2mm}\hrule height 0.2pt \vspace{4mm}
\needspace{5\baselineskip}
\hypertarget{q78}{\textbf{\color{ink}78.}}  {\color{muted}\scriptsize Y1 \textbullet\ daraja 1 \textbullet\ DB f73078dc}
\par
HTML hujjatida \texttt{\textless{}title\textgreater{}} elementi qayerda yoziladi?
\begin{enumerate}[label=\Alph*),leftmargin=9mm]
\item \texttt{\textless{}body\textgreater{}} ichida
\item \texttt{\textless{}head\textgreater{}} ichida
\item Hujjatning istalgan joyida
\item \texttt{\textless{}p\textgreater{}} ichida
\end{enumerate}
\vspace{2mm}\hrule height 0.2pt \vspace{4mm}
\needspace{5\baselineskip}
\hypertarget{q79}{\textbf{\color{ink}79.}}  {\color{muted}\scriptsize Y1 \textbullet\ daraja 3 \textbullet\ DB 3169514c}
\par
HTML sahifada asosiy sarlavha yaratish uchun qaysi teg mos?
\begin{enumerate}[label=\Alph*),leftmargin=9mm]
\item \textless{}h1\textgreater{}
\item \textless{}img\textgreater{}
\item \textless{}ul\textgreater{}
\item \textless{}br\textgreater{}
\end{enumerate}
\vspace{2mm}\hrule height 0.2pt \vspace{4mm}
\needspace{5\baselineskip}
\hypertarget{q80}{\textbf{\color{ink}80.}}  {\color{muted}\scriptsize Y1 \textbullet\ daraja 3 \textbullet\ DB 31c99ed4}
\par
\texttt{Olma, Nok, Uzum} elementlarini tartibsiz ro'yxat sifatida chiqarish kerak. Qaysi tashqi teg ishlatiladi?
\begin{enumerate}[label=\Alph*),leftmargin=9mm]
\item \textless{}ul\textgreater{}
\item \textless{}ol\textgreater{}
\item \textless{}table\textgreater{}
\item \textless{}form\textgreater{}
\end{enumerate}
\vspace{2mm}\hrule height 0.2pt \vspace{4mm}
\needspace{5\baselineskip}
\hypertarget{q81}{\textbf{\color{ink}81.}}  {\color{muted}\scriptsize Y1 \textbullet\ daraja 3 \textbullet\ DB 3f9897c4}
\par
\texttt{https://example.com} manziliga bosiladigan havola yaratish uchun qaysi HTML mos?
\begin{enumerate}[label=\Alph*),leftmargin=9mm]
\item \textless{}a href="https://example.com"\textgreater{}Sayt\textless{}/a\textgreater{}
\item \textless{}img href="https://example.com"\textgreater{}
\item \textless{}p src="https://example.com"\textgreater{}Sayt\textless{}/p\textgreater{}
\item \textless{}link\textgreater{}Sayt\textless{}/link\textgreater{}
\end{enumerate}
\vspace{2mm}\hrule height 0.2pt \vspace{4mm}
\needspace{5\baselineskip}
\hypertarget{q82}{\textbf{\color{ink}82.}}  {\color{muted}\scriptsize Y3 \textbullet\ daraja 3 \textbullet\ DB 3fd7c6dc}
\par
Oddiy HTML sahifa tuzilmasini yaratish qadamlarini tartiblang.
\begin{enumerate}[label=\Alph*),leftmargin=9mm]
\item body qismiga ko'rinadigan mazmunni joylashtirish
\item html hujjat konteynerini boshlash
\item Teglarni yopib hujjatni brauzerda tekshirish
\item head qismida title kabi metama'lumotni berish
\end{enumerate}
\vspace{2mm}\hrule height 0.2pt \vspace{4mm}
\needspace{5\baselineskip}
\hypertarget{q83}{\textbf{\color{ink}83.}}  {\color{muted}\scriptsize Y1 \textbullet\ daraja 3 \textbullet\ DB 7515c326}
\par
Quyidagi HTML kod qanday ro'yxat iyerarxiyasini hosil qiladi?

\par
\begin{lstlisting}
<ul>
 <li>Aralov
  <ul>
   <li>Frontend
    <ol><li>HTML</li><li>CSS</li><li>Ana</li></ol>
   </li>
   <li>Backend
    <ul><li>Salom</li><li>Maktab</li></ul>
   </li>
  </ul>
 </li>
</ul>
\end{lstlisting}

\begin{enumerate}[label=\Alph*),leftmargin=9mm]
\item Aralov ostida Frontend va Backend; Frontend ostida tartibli HTML, CSS, Ana; Backend ostida markersiz Salom, Maktab
\item Aralov ostida faqat Frontend; HTML, CSS, Ana va Backend bir xil darajada
\item Frontend va Backend raqamli tashqi ro'yxat; Aralov ularning ostida
\item Barcha bandlar bitta tartibli ro'yxatda joylashadi
\end{enumerate}
\vspace{2mm}\hrule height 0.2pt \vspace{4mm}
\needspace{5\baselineskip}
\hypertarget{q84}{\textbf{\color{ink}84.}}  {\color{muted}\scriptsize Y1 \textbullet\ daraja 3 \textbullet\ DB 91c3580d}
\par
Veb-sahifaga \texttt{logo.png} rasmini joylashtirish kerak. Qaysi yozuv mos?
\begin{enumerate}[label=\Alph*),leftmargin=9mm]
\item \textless{}img src="logo.png" alt="Logo"\textgreater{}
\item \textless{}p href="logo.png"\textgreater{}
\item \textless{}ul src="logo.png"\textgreater{}
\item \textless{}h1 image="logo.png"\textgreater{}
\end{enumerate}
\vspace{2mm}\hrule height 0.2pt \vspace{4mm}
\needspace{5\baselineskip}
\hypertarget{q85}{\textbf{\color{ink}85.}}  {\color{muted}\scriptsize Y2 \textbullet\ daraja 3 \textbullet\ DB b1b42ae5}
\par
HTML teglarini ularning vazifalari bilan juftlang.
\textbf{\footnotesize Chap ustun:}
\begin{enumerate}[label=\arabic*),leftmargin=9mm]
\item \textless{}h1\textgreater{}
\item \textless{}a\textgreater{}
\item \textless{}img\textgreater{}
\item \textless{}ul\textgreater{}
\end{enumerate}
\textbf{\footnotesize O'ng ustun:}
\begin{enumerate}[label=\Alph*),leftmargin=9mm]
\item Rasm joylashtirish
\item Asosiy sarlavha
\item Tartibsiz ro'yxat konteyneri
\item Giperhavola
\end{enumerate}
\vspace{2mm}\hrule height 0.2pt \vspace{4mm}
\needspace{5\baselineskip}
\hypertarget{q86}{\textbf{\color{ink}86.}}  {\color{muted}\scriptsize Y1 \textbullet\ daraja 3 \textbullet\ DB bfa456da}
\par
Paragraf ichida yangi qatorga majburiy o'tish kerak. Qaysi teg mos?
\begin{enumerate}[label=\Alph*),leftmargin=9mm]
\item \textless{}br\textgreater{}
\item \textless{}hr\textgreater{} faqat jadval uchun
\item \textless{}img\textgreater{}
\item \textless{}title\textgreater{}
\end{enumerate}
\vspace{2mm}\hrule height 0.2pt \vspace{4mm}
\needspace{5\baselineskip}
\hypertarget{q87}{\textbf{\color{ink}87.}}  {\color{muted}\scriptsize Y3 \textbullet\ daraja 4 \textbullet\ DB e2910000}
\par
Tashqi ro'yxatda "Web" bandining ichida "HTML" va "CSS" nomli ichki tartibli ro'yxat yaratish bosqichlarini to'g'ri tartiblang.
\begin{enumerate}[label=\Alph*),leftmargin=9mm]
\item Tashqi \texttt{\textless{}li\textgreater{}Web} bandini oching.
\item Tashqi \texttt{\textless{}/li\textgreater{}\textless{}/ul\textgreater{}} teglarini to'g'ri yopib natijani sinang.
\item Tashqi \texttt{\textless{}ul\textgreater{}} konteynerini oching.
\item Shu li ichiga \texttt{\textless{}ol\textgreater{}\textless{}li\textgreater{}HTML\textless{}/li\textgreater{}\textless{}li\textgreater{}CSS\textless{}/li\textgreater{}\textless{}/ol\textgreater{}} yozing.
\end{enumerate}
\vspace{2mm}\hrule height 0.2pt \vspace{4mm}
\needspace{5\baselineskip}
\hypertarget{q88}{\textbf{\color{ink}88.}}  {\color{muted}\scriptsize Y1 \textbullet\ daraja 5 \textbullet\ DB e2910000}
\par
Quyidagi kod natijasiga qaysi tavsif MOS KELMAYDI?\par
\begin{lstlisting}
<p style="font-family:serif;color:#008000;text-align:right">Salom</p>
\end{lstlisting}
\par
Brauzer standart CSS ni qo'shimcha o'zgartirmagan deb oling.
\begin{enumerate}[label=\Alph*),leftmargin=9mm]
\item Matn o'ng tomonga tekislanadi.
\item Matn serif shriftlar oilasida ko'rsatiladi.
\item Matn chap tomonga tekislanadi.
\item Matn yashil rangda bo'ladi.
\end{enumerate}
\vspace{2mm}\hrule height 0.2pt \vspace{4mm}
\needspace{5\baselineskip}
\hypertarget{q89}{\textbf{\color{ink}89.}}  {\color{muted}\scriptsize Y1 \textbullet\ daraja 5 \textbullet\ DB e2910000}
\par
Ichma-ich ro'yxatning natijasini aniqlang.\par
\begin{lstlisting}
<ol start="3">
 <li>Veb
  <ul><li>HTML</li><li>CSS</li></ul>
 </li>
 <li>Grafika</li>
</ol>
\end{lstlisting}

\begin{enumerate}[label=\Alph*),leftmargin=9mm]
\item HTML va CSS asosiy 3,4 bandlar; Veb va Grafika markerli quyi ro'yxat.
\item Veb 3, Grafika 4; HTML va CSS Veb bandi ichidagi markerli pastki ro'yxat.
\item Veb 1, Grafika 2; HTML va CSS ham 3,4 raqam oladi.
\item Faqat Grafika ko'rinadi, chunki ul ol ichida yozilgan.
\end{enumerate}
\vspace{2mm}\hrule height 0.2pt \vspace{4mm}
\subsubsection{HTML: jadval va forma elementlari}
\needspace{5\baselineskip}
\hypertarget{q90}{\textbf{\color{ink}90.}}  {\color{muted}\scriptsize Y1 \textbullet\ daraja 1 \textbullet\ DB 86471e56}
\par
HTML jadvalida vertikal joylashgan kataklarni qaysi atribut birlashtiradi?
\begin{enumerate}[label=\Alph*),leftmargin=9mm]
\item \texttt{colspan}
\item \texttt{rowspan}
\item \texttt{unity}
\item \texttt{union}
\end{enumerate}
\vspace{2mm}\hrule height 0.2pt \vspace{4mm}
\needspace{5\baselineskip}
\hypertarget{q91}{\textbf{\color{ink}91.}}  {\color{muted}\scriptsize Y1 \textbullet\ daraja 1 \textbullet\ DB b3f6c20e}
\par
HTML jadvalida vertikal (ustma-ust joylashgan) kataklarni birlashtirish uchun qaysi atribut ishlatiladi?
\begin{enumerate}[label=\Alph*),leftmargin=9mm]
\item \texttt{colspan}
\item \texttt{rowspan}
\item \texttt{unity}
\item \texttt{union}
\end{enumerate}
\vspace{2mm}\hrule height 0.2pt \vspace{4mm}
\needspace{5\baselineskip}
\hypertarget{q92}{\textbf{\color{ink}92.}}  {\color{muted}\scriptsize Y1 \textbullet\ daraja 2 \textbullet\ DB 23e61347}
\par
Quyidagi HTML parcha brauzerda nimani ko'rsatadi?

\par
\begin{lstlisting}
<h1>A) Ism familiya</h1>
Name: <input type="text" value="Otabek">
Family: <input type="text" value="Aralov">
<br><br>
<button>ok</button> <button>cancel</button>
\end{lstlisting}

\begin{enumerate}[label=\Alph*),leftmargin=9mm]
\item \texttt{A) Ism familiya} sarlavhasi, Otabek va Aralov qiymatlari yozilgan ikki matn maydoni hamda ok/cancel tugmalari
\item Sarlavha, ikkita bo'sh parol maydoni va bitta submit tugmasi
\item Otabek va Aralovdan iborat ikkita ochiluvchi ro'yxat
\item Faqat oddiy matn; kiritish maydonlari va tugmalar ko'rinmaydi
\end{enumerate}
\vspace{2mm}\hrule height 0.2pt \vspace{4mm}
\needspace{5\baselineskip}
\hypertarget{q93}{\textbf{\color{ink}93.}}  {\color{muted}\scriptsize Y1 \textbullet\ daraja 2 \textbullet\ DB 82e02b50}
\par
Quyidagi HTML kodida \texttt{rowspan} yoki \texttt{colspan} atributi bo'lmagan \texttt{\textless{}td\textgreater{}} elementlaridagi sonlar yig'indisini barcha \texttt{\textless{}td\textgreater{}} elementlari soniga ko'paytiring.

\par
\begin{lstlisting}
<table>
  <tr><td>160</td><td>80</td><td>-160</td><td rowspan="2">160</td></tr>
  <tr><td colspan="2">80</td><td>85</td></tr>
</table>
\end{lstlisting}

\begin{enumerate}[label=\Alph*),leftmargin=9mm]
\item 660
\item 990
\item 750
\item 800
\end{enumerate}
\vspace{2mm}\hrule height 0.2pt \vspace{4mm}
\needspace{5\baselineskip}
\hypertarget{q94}{\textbf{\color{ink}94.}}  {\color{muted}\scriptsize Y1 \textbullet\ daraja 2 \textbullet\ DB a62d680b}
\par
HTML jadvalida sarlavhalar \texttt{\textless{}th\textgreater{}Mahsulot\textless{}/th\textgreater{}\textless{}th\textgreater{}Narxi\textless{}/th\textgreater{}\textless{}th\textgreater{}Miqdori\textless{}/th\textgreater{}}. Olma-5000 qatoridagi \texttt{10 kg} katagi \texttt{rowspan=2} bilan Banan-7000 qatorini ham qamraydi. Keyingi qatorda \texttt{Yangi mevalar} \texttt{colspan=2}, uchinchi katak \texttt{5 kg}. \texttt{border=1}, \texttt{cellspacing=0}, \texttt{cellpadding=5}. Natijani tavsiflang.
\begin{enumerate}[label=\Alph*),leftmargin=9mm]
\item Sarlavhalar oddiy; rowspan va colspan to'g'ri
\item Sarlavhalar qalin; 10 kg ikki qatorni, Yangi mevalar ikki ustunni egallaydi; kataklar orasida bo'shliq yo'q
\item Barcha kataklar orasida qo'sh chiziqli bo'shliq bor
\item Barcha matnlar qalin ko'rinadi
\end{enumerate}
\vspace{2mm}\hrule height 0.2pt \vspace{4mm}
\needspace{5\baselineskip}
\hypertarget{q95}{\textbf{\color{ink}95.}}  {\color{muted}\scriptsize Y1 \textbullet\ daraja 2 \textbullet\ DB cf6c0312}
\par
HTML jadvalida Facebook katagi \texttt{colspan=2} bilan birlashtirilgan va uning ichidagi nom \texttt{\textless{}ol\textgreater{}} tartibli ro'yxatida; Telegram katagi \texttt{rowspan=2} va \texttt{\textless{}ul\textgreater{}} ro'yxatida; Instagram hamda YouTube oddiy kataklardagi \texttt{\textless{}ol\textgreater{}} ro'yxatlarida berilgan. Bir vaqtning o'zida birlashtirilgan katakda va tartibli ro'yxatda turgan ijtimoiy tarmoq asoschisini toping.
\begin{enumerate}[label=\Alph*),leftmargin=9mm]
\item Aka-uka Durovlar
\item Mark Zuckerberg
\item Steve Chen
\item Kevin Systrom
\end{enumerate}
\vspace{2mm}\hrule height 0.2pt \vspace{4mm}
\needspace{5\baselineskip}
\hypertarget{q96}{\textbf{\color{ink}96.}}  {\color{muted}\scriptsize Y1 \textbullet\ daraja 2 \textbullet\ DB feb86531}
\par
400px kenglikdagi uch ustunli jadvalning uchinchi ustuni 25\% kenglikda. Birinchi qatordagi \texttt{Lenovo} katagi uchinchi ustunda; ikkinchi qatordagi \texttt{Asus} katagi esa dastlabki ikki ustunni \texttt{colspan=2} bilan egallaydi. Lenovo va Asus kataklari kengliklari yig'indisi a, eng qisqa katak kengligi b bo'lsa, a+b ni toping.
\begin{enumerate}[label=\Alph*),leftmargin=9mm]
\item 400px
\item 450px
\item 500px
\item 600px
\end{enumerate}
\vspace{2mm}\hrule height 0.2pt \vspace{4mm}
\needspace{5\baselineskip}
\hypertarget{q97}{\textbf{\color{ink}97.}}  {\color{muted}\scriptsize Y1 \textbullet\ daraja 3 \textbullet\ DB 07f9773b}
\par
Quyidagi HTML jadval kodining natijasini aniqlang:\par
\begin{lstlisting}
<table>
<tr><td>A</td><td rowspan="2">B</td></tr>
<tr><td>C</td></tr>
</table>
\end{lstlisting}
\par
\texttt{B} yozuvi brauzerda qanday ko'rinadi?
\begin{enumerate}[label=\Alph*),leftmargin=9mm]
\item B ikki qatorni qamragan yagona katakda ko'rinadi
\item B ikkinchi qatorda ikki marta takrorlanadi
\item B faqat C o'rnida ko'rinadi
\item B alohida uchinchi qatorda paydo bo'ladi
\end{enumerate}
\vspace{2mm}\hrule height 0.2pt \vspace{4mm}
\needspace{5\baselineskip}
\hypertarget{q98}{\textbf{\color{ink}98.}}  {\color{muted}\scriptsize Y1 \textbullet\ daraja 3 \textbullet\ DB 1af0c63f}
\par
Forma ichida foydalanuvchi ism kiritishi uchun bir qatorli matn maydoni kerak. Qaysi element mos?
\begin{enumerate}[label=\Alph*),leftmargin=9mm]
\item \textless{}input type="text"\textgreater{}
\item \textless{}img type="text"\textgreater{}
\item \textless{}table type="text"\textgreater{}
\item \textless{}a type="text"\textgreater{}
\end{enumerate}
\vspace{2mm}\hrule height 0.2pt \vspace{4mm}
\needspace{5\baselineskip}
\hypertarget{q99}{\textbf{\color{ink}99.}}  {\color{muted}\scriptsize Y1 \textbullet\ daraja 3 \textbullet\ DB 85e12f2d}
\par
HTML jadvalida yangi qator yaratish uchun qaysi teg ishlatiladi?
\begin{enumerate}[label=\Alph*),leftmargin=9mm]
\item \textless{}tr\textgreater{}
\item \textless{}td\textgreater{}
\item \textless{}th\textgreater{}
\item \textless{}input\textgreater{}
\end{enumerate}
\vspace{2mm}\hrule height 0.2pt \vspace{4mm}
\needspace{5\baselineskip}
\hypertarget{q100}{\textbf{\color{ink}100.}}  {\color{muted}\scriptsize Y1 \textbullet\ daraja 3 \textbullet\ DB e2910000}
\par
HTML jadvalida qaysi ko'rinish hosil bo'ladi?\par
\begin{lstlisting}
<table>
  <tr><td rowspan="2">A</td><td>B</td></tr>
  <tr><td>C</td></tr>
</table>
\end{lstlisting}

\begin{enumerate}[label=\Alph*),leftmargin=9mm]
\item A chapdagi ikki qatorni egallaydi; B yuqori o'ngda, C pastki o'ngda.
\item A pastki o'ng katakda; B va C chap ustunda.
\item A birinchi qatorning ikkita ustunini egallaydi; B va C keyingi qatorlarda.
\item B ikkita qatorni, A va C ikkita ustunni egallaydi.
\end{enumerate}
\vspace{2mm}\hrule height 0.2pt \vspace{4mm}
\needspace{5\baselineskip}
\hypertarget{q101}{\textbf{\color{ink}101.}}  {\color{muted}\scriptsize Y2 \textbullet\ daraja 3 \textbullet\ DB e2910000}
\par
HTML parchasini brauzerdagi vazifa yoki ko'rinishi bilan moslang.
\textbf{\footnotesize Chap ustun:}
\begin{enumerate}[label=\arabic*),leftmargin=9mm]
\item \texttt{\textless{}h2\textgreater{}Bo'lim\textless{}/h2\textgreater{}}
\item \texttt{\textless{}ol start="5"\textgreater{}\textless{}li\textgreater{}A\textless{}/li\textgreater{}\textless{}/ol\textgreater{}}
\item \texttt{\textless{}td colspan="2"\textgreater{}X\textless{}/td\textgreater{}}
\item \texttt{\textless{}img src="x.png" alt="Rasm"\textgreater{}}
\end{enumerate}
\textbf{\footnotesize O'ng ustun:}
\begin{enumerate}[label=\Alph*),leftmargin=9mm]
\item Manzilga bog'langan tasvir va muqobil matn
\item Ikkinchi darajali sarlavha
\item Katak ikkita mantiqiy ustunni egallaydi
\item Birinchi ko'rinadigan band raqami 5
\end{enumerate}
\vspace{2mm}\hrule height 0.2pt \vspace{4mm}
\needspace{5\baselineskip}
\hypertarget{q102}{\textbf{\color{ink}102.}}  {\color{muted}\scriptsize Y1 \textbullet\ daraja 3 \textbullet\ DB fcda7730}
\par
Foydalanuvchi faqat bitta variantni tanlashi kerak bo'lgan guruh uchun qaysi forma elementi mos?
\begin{enumerate}[label=\Alph*),leftmargin=9mm]
\item Radio tugmalar
\item Bir nechta mustaqil checkbox
\item Rasm tegi
\item Jadval qatori
\end{enumerate}
\vspace{2mm}\hrule height 0.2pt \vspace{4mm}
\needspace{5\baselineskip}
\hypertarget{q103}{\textbf{\color{ink}103.}}  {\color{muted}\scriptsize Y1 \textbullet\ daraja 4 \textbullet\ DB d78fee6d}
\par
Quyidagi forma yuborilsa, qaysi HTTP usuli, manzil va \texttt{ism} qiymati ishlatiladi?\par
\begin{lstlisting}
<form action="/qabul" method="post">
  <input name="ism" value="Ali">
  <button type="submit">Yuborish</button>
</form>
\end{lstlisting}

\begin{enumerate}[label=\Alph*),leftmargin=9mm]
\item GET so'rovi /qabul?ism=Ali manziliga
\item POST so'rovi /qabul manziliga, tana ma'lumoti ism=Ali
\item POST so'rovi /qabul manziliga, tana ma'lumoti value=ism
\item GET so'rovi /qabul manziliga, tana ma'lumoti ism=Ali
\end{enumerate}
\vspace{2mm}\hrule height 0.2pt \vspace{4mm}
\needspace{5\baselineskip}
\hypertarget{q104}{\textbf{\color{ink}104.}}  {\color{muted}\scriptsize Y1 \textbullet\ daraja 4 \textbullet\ DB e2910000}
\par
HTML forma jo'natilganda foydalanuvchiga ko'rinmaydigan, lekin forma qiymati sifatida yuboriladigan maydon qaysi?\par
\begin{lstlisting}
<form method="get">
<input name="id" value="7">
<input type="hidden" name="level" value="B">
<button type="submit">Jo'natish</button>
</form>
\end{lstlisting}

\begin{enumerate}[label=\Alph*),leftmargin=9mm]
\item level maydoni yuborilmaydi, chunki hidden.
\item Yuborilish uchun hidden maydon required bo'lishi shart.
\item Jo'natilganda faqat button matni yuboriladi.
\item id=7 va level=B yuborilishi mumkin; hidden maxfiy ma'lumot uchun ishonchli himoya emas.
\end{enumerate}
\vspace{2mm}\hrule height 0.2pt \vspace{4mm}
\subsubsection{CSS stillarini HTML elementlariga qo'llash}
\needspace{5\baselineskip}
\hypertarget{q105}{\textbf{\color{ink}105.}}  {\color{muted}\scriptsize Y1 \textbullet\ daraja 2 \textbullet\ DB 043038f5}
\par
Quyidagi CSS va HTML fragmentlarining natijasini aniqlang:

\par
\begin{lstlisting}
body {
  font-family: "Times New Roman";
  color: black;
}
h2 {
  text-align: right;
  color: red;
}
p {
  font-size: 16px;
  margin: 5px;
}
.t1 {
  color: green;
  font-weight: bold;
}
\end{lstlisting}


\par
\begin{lstlisting}
<body>
  <h2>HTML</h2>
  <p>HTML - bu belgilash tili.</p>
  <p class="t1">HTML web sahifa tuzish uchun ishlatiladi.</p>
</body>
\end{lstlisting}


\par
Qaysi tavsif to'g'ri?
\begin{enumerate}[label=\Alph*),leftmargin=9mm]
\item Sarlavha o'ngda qizil; birinchi abzats qora 16px; ikkinchisi yashil va qalin; abzats marginlari 5px
\item Sarlavha o'ngda qizil; ikkala abzats orasida juda katta bo'shliq; ikkinchisi yashil va qalin
\item Sarlavha o'ngda qizil; birinchi abzats qizil, ikkinchisi ko'k va oddiy
\item Sarlavha o'ngda qizil; ikkinchi abzats yashil va qalin, ammo barcha matn markazda
\end{enumerate}
\vspace{2mm}\hrule height 0.2pt \vspace{4mm}
\needspace{5\baselineskip}
\hypertarget{q106}{\textbf{\color{ink}106.}}  {\color{muted}\scriptsize Y1 \textbullet\ daraja 2 \textbullet\ DB ad13504b}
\par
Quyidagi HTML fragmentining natijasini tekshiring:

\par
\begin{lstlisting}
<body>
  <h2 style="color: red; text-align: right; font-size: 40px;">
    Welcome to Web Design
  </h2>
  <img
    src="web.png"
    width="150"
    height="100"
    style="float: right;"
    alt="web image"
  >
  <p>Oddiy abzats matni.</p>
</body>
\end{lstlisting}


\par
Qaysi fikr to'g'ri?
\begin{enumerate}[label=\Alph*),leftmargin=9mm]
\item Rasm brauzer oynasining o'ng tomoniga joylashadi
\item Sarlavha qizil va chapga tekislanadi
\item h2 uchun 40px amal qilmaydi
\item p matni qizil bo'ladi
\end{enumerate}
\vspace{2mm}\hrule height 0.2pt \vspace{4mm}
\needspace{5\baselineskip}
\hypertarget{q107}{\textbf{\color{ink}107.}}  {\color{muted}\scriptsize Y1 \textbullet\ daraja 2 \textbullet\ DB e01e2f7a}
\par
CSSda \texttt{p \{ margin: 15px; \}} berilgan. Ketma-ket joylashgan ikkita blokli \texttt{\textless{}p\textgreater{}} elementi orasidagi vertikal masofa odatiy margin collapsing qoidasiga ko'ra qancha bo'ladi?
\begin{enumerate}[label=\Alph*),leftmargin=9mm]
\item 15px
\item 0px
\item 30px
\item 10px
\end{enumerate}
\vspace{2mm}\hrule height 0.2pt \vspace{4mm}
\needspace{5\baselineskip}
\hypertarget{q108}{\textbf{\color{ink}108.}}  {\color{muted}\scriptsize Y1 \textbullet\ daraja 2 \textbullet\ DB ec6030dc}
\par
CSSda \texttt{p \{ margin: 10px 25px; \}} berilgan. Oddiy hujjat oqimida ketma-ket joylashgan ikkita blokli \texttt{\textless{}p\textgreater{}} elementi orasidagi vertikal masofa margin collapsing qoidasiga ko'ra qancha bo'ladi?
\begin{enumerate}[label=\Alph*),leftmargin=9mm]
\item 10px
\item 20px
\item 25px
\item 35px
\end{enumerate}
\vspace{2mm}\hrule height 0.2pt \vspace{4mm}
\needspace{5\baselineskip}
\hypertarget{q109}{\textbf{\color{ink}109.}}  {\color{muted}\scriptsize Y1 \textbullet\ daraja 2 \textbullet\ DB ffe8defe}
\par
Quyidagi CSS va HTML parchada \texttt{A} matnining yakuniy rangi qaysi?

\par
\begin{lstlisting}
<style>
 table, th { background-color: lightblue; }
 tr { background-color: red; }
 td { color: green; }
</style>
<table><tr><td>A</td><td>2</td></tr></table>
\end{lstlisting}

\begin{enumerate}[label=\Alph*),leftmargin=9mm]
\item yashil
\item och ko'k
\item qizil
\item qora
\end{enumerate}
\vspace{2mm}\hrule height 0.2pt \vspace{4mm}
\needspace{5\baselineskip}
\hypertarget{q110}{\textbf{\color{ink}110.}}  {\color{muted}\scriptsize Y1 \textbullet\ daraja 3 \textbullet\ DB 054472de}
\par
Quyidagi CSS va HTML kodlari berilgan. Boshqa muallif qoidalari va \texttt{!important} yo'q.

\par
\begin{lstlisting}
* {
  color: purple;
  font-weight: bold;
}
body {
  color: orange;
}
ol {
  list-style-type: upper-roman;
  color: orange;
}
li {
  color: brown;
  font-weight: normal;
}
p {
  color: teal;
  font-style: italic;
}
.p2 {
  color: black;
}
#id3 {
  color: darkred;
}
\end{lstlisting}


\par
\begin{lstlisting}
<body>
  <h2>Sarlavha</h2>
  <ol>
    <li class="p2">Birinchi band</li>
    <li class="p2" id="id3">
      <p>Ikkinchi band</p>
    </li>
    <li>Uchinchi band</li>
  </ol>
</body>
\end{lstlisting}


\par
Sarlavha va I-III ro'yxat bandlarining rangi, markerlari hamda shrift ko'rinishini aniqlang.
\begin{enumerate}[label=\Alph*),leftmargin=9mm]
\item Sarlavha purple va qalin; I qora va oddiy; II markeri to'q qizil, ichki p teal, kursiv va qalin; III jigarrang va oddiy.
\item Sarlavha orange va qalin; I qora va oddiy; II markeri to'q qizil, ichki p teal, kursiv va qalin; III jigarrang va oddiy.
\item Sarlavha purple va qalin; I jigarrang va oddiy; II markeri qora, ichki p to'q qizil, kursiv va oddiy; III jigarrang va oddiy.
\item Sarlavha purple va oddiy; I qora va qalin; II markeri to'q qizil, ichki p jigarrang, kursiv va oddiy; III jigarrang va qalin.
\end{enumerate}
\vspace{2mm}\hrule height 0.2pt \vspace{4mm}
\needspace{5\baselineskip}
\hypertarget{q111}{\textbf{\color{ink}111.}}  {\color{muted}\scriptsize Y1 \textbullet\ daraja 3 \textbullet\ DB 0c448427}
\par
Quyidagi ikkita bo'sh \texttt{div} normal blok oqimida bir-birining darhol ketidan keladi; ularda border, padding, flex/grid konteksti yo'q. CSS berilgan:

\par
\begin{lstlisting}
div { margin-top: 25px; margin-bottom: 25px; }
\end{lstlisting}


\par
Ikkala qo'shni vertikal margin orasidagi haqiqiy masofa qancha?
\begin{enumerate}[label=\Alph*),leftmargin=9mm]
\item 0 px
\item 25 px
\item 50 px
\item 75 px
\end{enumerate}
\vspace{2mm}\hrule height 0.2pt \vspace{4mm}
\needspace{5\baselineskip}
\hypertarget{q112}{\textbf{\color{ink}112.}}  {\color{muted}\scriptsize Y3 \textbullet\ daraja 3 \textbullet\ DB 29c2b2da}
\par
CSS qoidasini HTML elementga qo'llash jarayonini tartiblang.
\begin{enumerate}[label=\Alph*),leftmargin=9mm]
\item Kerakli xossa va qiymatlarni qoida ichida yozish
\item Uslub beriladigan element yoki guruhni aniqlash
\item Brauzerda natijani tekshirib zarur bo'lsa tuzatish
\item Mos CSS selektorini tanlash
\end{enumerate}
\vspace{2mm}\hrule height 0.2pt \vspace{4mm}
\needspace{5\baselineskip}
\hypertarget{q113}{\textbf{\color{ink}113.}}  {\color{muted}\scriptsize Y1 \textbullet\ daraja 3 \textbullet\ DB 5f422b90}
\par
Faqat \texttt{id="menu"} bo'lgan elementga CSS qo'llash uchun qaysi selektor kerak?
\begin{enumerate}[label=\Alph*),leftmargin=9mm]
\item \#menu
\item .menu
\item menu*
\item @menu
\end{enumerate}
\vspace{2mm}\hrule height 0.2pt \vspace{4mm}
\needspace{5\baselineskip}
\hypertarget{q114}{\textbf{\color{ink}114.}}  {\color{muted}\scriptsize Y1 \textbullet\ daraja 3 \textbullet\ DB 69e5eb76}
\par
\texttt{class="warning"} berilgan bir nechta elementni tanlash uchun qaysi CSS selektor mos?
\begin{enumerate}[label=\Alph*),leftmargin=9mm]
\item .warning
\item \#warning
\item warning\#
\item @warning
\end{enumerate}
\vspace{2mm}\hrule height 0.2pt \vspace{4mm}
\needspace{5\baselineskip}
\hypertarget{q115}{\textbf{\color{ink}115.}}  {\color{muted}\scriptsize Y1 \textbullet\ daraja 3 \textbullet\ DB 854ef8e9}
\par
Rasmning CSS orqali kengligini 200 piksel qilish uchun qaysi qoida mos?
\begin{enumerate}[label=\Alph*),leftmargin=9mm]
\item img \{ width: 200px; \}
\item img \{ font-size: 200px; \}
\item img \{ color: 200px; \}
\item img \{ href: 200px; \}
\end{enumerate}
\vspace{2mm}\hrule height 0.2pt \vspace{4mm}
\needspace{5\baselineskip}
\hypertarget{q116}{\textbf{\color{ink}116.}}  {\color{muted}\scriptsize Y1 \textbullet\ daraja 3 \textbullet\ DB 94069084}
\par
Barcha \texttt{\textless{}p\textgreater{}} elementlari matnini ko'k qilish uchun qaysi CSS qoida to'g'ri?
\begin{enumerate}[label=\Alph*),leftmargin=9mm]
\item p \{ color: blue; \}
\item p = color(blue)
\item \textless{}p color=blue\textgreater{} CSS
\item color:p:blue
\end{enumerate}
\vspace{2mm}\hrule height 0.2pt \vspace{4mm}
\needspace{5\baselineskip}
\hypertarget{q117}{\textbf{\color{ink}117.}}  {\color{muted}\scriptsize Y1 \textbullet\ daraja 3 \textbullet\ DB b0feef99}
\par
Element fonini sariq qilish uchun qaysi CSS xossa ishlatiladi?
\begin{enumerate}[label=\Alph*),leftmargin=9mm]
\item background-color
\item font-family
\item text-align
\item border-collapse
\end{enumerate}
\vspace{2mm}\hrule height 0.2pt \vspace{4mm}
\needspace{5\baselineskip}
\hypertarget{q118}{\textbf{\color{ink}118.}}  {\color{muted}\scriptsize Y1 \textbullet\ daraja 3 \textbullet\ DB e2910000}
\par
Quyidagi qoida qaysi element(lar)ga tatbiq etiladi?\par
\begin{lstlisting}
<style>.box > p {color:blue;}</style>
<div class="box">
 <p id="one">A</p>
 <section><p id="two">B</p></section>
</div>
\end{lstlisting}

\begin{enumerate}[label=\Alph*),leftmargin=9mm]
\item Faqat A ga
\item A va B ikkalasiga
\item Faqat B ga
\item Hech biriga
\end{enumerate}
\vspace{2mm}\hrule height 0.2pt \vspace{4mm}
\needspace{5\baselineskip}
\hypertarget{q119}{\textbf{\color{ink}119.}}  {\color{muted}\scriptsize Y2 \textbullet\ daraja 3 \textbullet\ DB e2910000}
\par
Quyidagi CSS selektorini aynan qaysi elementlar tanlanishi bilan moslashtiring.
\textbf{\footnotesize Chap ustun:}
\begin{enumerate}[label=\arabic*),leftmargin=9mm]
\item \texttt{.card p}
\item \texttt{.card \textgreater{} p}
\item \texttt{h1, h2}
\item \texttt{p.note}
\end{enumerate}
\textbf{\footnotesize O'ng ustun:}
\begin{enumerate}[label=\Alph*),leftmargin=9mm]
\item Barcha h1 va h2 elementlari
\item card ichidagi barcha avlod p lar
\item note classiga ega p elementlari
\item cardning faqat bevosita farzand p lari
\end{enumerate}
\vspace{2mm}\hrule height 0.2pt \vspace{4mm}
\needspace{5\baselineskip}
\hypertarget{q120}{\textbf{\color{ink}120.}}  {\color{muted}\scriptsize Y2 \textbullet\ daraja 3 \textbullet\ DB fa021355}
\par
CSS xossalarini ularning ta'siri bilan juftlang.
\textbf{\footnotesize Chap ustun:}
\begin{enumerate}[label=\arabic*),leftmargin=9mm]
\item color
\item background-color
\item width
\item text-align
\end{enumerate}
\textbf{\footnotesize O'ng ustun:}
\begin{enumerate}[label=\Alph*),leftmargin=9mm]
\item Element kengligini belgilaydi
\item Matn rangini belgilaydi
\item Matnning gorizontal tekislanishini belgilaydi
\item Element fon rangini belgilaydi
\end{enumerate}
\vspace{2mm}\hrule height 0.2pt \vspace{4mm}
\needspace{5\baselineskip}
\hypertarget{q121}{\textbf{\color{ink}121.}}  {\color{muted}\scriptsize Y1 \textbullet\ daraja 4 \textbullet\ DB e2910000}
\par
CSS box-modelda \texttt{box-sizing:content-box}, \texttt{width:200px}, chap/o'ng \texttt{padding:10px}, har ikki yondan \texttt{border:3px} va \texttt{margin:5px} bo'lsa, jami tashqi kenglik necha piksel?
\begin{enumerate}[label=\Alph*),leftmargin=9mm]
\item 226px
\item 246px
\item 200px
\item 236px
\end{enumerate}
\vspace{2mm}\hrule height 0.2pt \vspace{4mm}
\needspace{5\baselineskip}
\hypertarget{q122}{\textbf{\color{ink}122.}}  {\color{muted}\scriptsize Y3 \textbullet\ daraja 4 \textbullet\ DB e2910000}
\par
Bir elementning matn rangiga ta'sir qiluvchi CSS qoidalari ichidan yakuniy rangni topish ish tartibini belgilang. Oddiy kaskad qoidalarini hisobga oling.
\begin{enumerate}[label=\Alph*),leftmargin=9mm]
\item Teng specificity holatida keyin kelgan yutuvchi qoidani tanlab natijani tekshiring.
\item Kaskadning origin, importance va layer kabi ustuvorlik bosqichlarini solishtiring.
\item Teng sharoitda mos selektorlarning specificity qiymatlarini solishtiring.
\item Elementga tegishli deklaratsiyalarni barcha CSS manbalardan aniqlang.
\end{enumerate}
\vspace{2mm}\hrule height 0.2pt \vspace{4mm}
\needspace{5\baselineskip}
\hypertarget{q123}{\textbf{\color{ink}123.}}  {\color{muted}\scriptsize Y1 \textbullet\ daraja 5 \textbullet\ DB e2910000}
\par
Quyidagi koddagi "Natija" matni qaysi rangda ko'rinadi? Oddiy muallif uslublari; \texttt{!important} va boshqa CSS qoida yo'q.\par
\begin{lstlisting}
<style>
 p {color:blue;}
 .note {color:green;}
 #msg {color:maroon;}
</style>
<p id="msg" class="note">Natija</p>
\end{lstlisting}

\begin{enumerate}[label=\Alph*),leftmargin=9mm]
\item Ko'k
\item Oq
\item To'q qizil (maroon)
\item Yashil
\end{enumerate}
\vspace{2mm}\hrule height 0.2pt \vspace{4mm}
\needspace{5\baselineskip}
\hypertarget{q124}{\textbf{\color{ink}124.}}  {\color{muted}\scriptsize Y1 \textbullet\ daraja 5 \textbullet\ DB fe9c6a0c}
\par
CSS da

\par
\begin{lstlisting}
img {
  width: 240px;
  height: 140px;
  object-fit: cover
}
\end{lstlisting}


\par
qo'llangan. Asl tasvir juda keng (16:9), lekin quti nisbatan baland. \texttt{cover} natijasida nima yuz berishi mumkin?
\begin{enumerate}[label=\Alph*),leftmargin=9mm]
\item Rasm avtomatik SVG formatiga aylanadi.
\item Tasvir qutini to'ldiradi, ortiqcha keng tomoni kesilishi mumkin.
\item Asl tasvir butunlay ko'rinadi va yonlarda doim bo'sh joy qoladi.
\item object-fit tasvirni har doim asl piksellarini o'zgartirib qayta saqlaydi.
\end{enumerate}
\vspace{2mm}\hrule height 0.2pt \vspace{4mm}
\subsubsection{HTML: jadvallar va formalar - V2}
\needspace{5\baselineskip}
\hypertarget{q125}{\textbf{\color{ink}125.}}  {\color{muted}\scriptsize Y1 \textbullet\ daraja 3 \textbullet\ DB 2008dfd5}
\par
360 px kenglikdagi HTML jadvalida: Samsung katagi 25\%, Xiaomi katagi 20\% va Huawei katagi 100 px. Honor katagi qolgan kenglikni egallaydi. Keyingi qatorda Nokia katagi Samsung va Honor ustunlarini \texttt{colspan=2} bilan birlashtiradi; iPhone Huawei ustunida, Xiaomi esa \texttt{rowspan=2}. Kontent uzunligi kenglikka ta'sir qilmasa, eng keng va eng tor kataklardagi yozuvlarni toping.
\begin{enumerate}[label=\Alph*),leftmargin=9mm]
\item Nokia / Xiaomi
\item Honor / iPhone
\item Samsung / Xiaomi
\item Nokia / iPhone
\end{enumerate}
\vspace{2mm}\hrule height 0.2pt \vspace{4mm}
\section{II qism. Tekshiruvdagi testlar}
\subsection{Kompyuter grafikasi va media}
\subsubsection{MS Paint va Adobe Photoshop: asosiy amallar}
\needspace{5\baselineskip}
\hypertarget{q126}{\textbf{\color{ink}126.}}  {\color{muted}\scriptsize Y1 \textbullet\ daraja 3 \textbullet\ DB f068e832}
\par
Photoshop'da \textbf{rasm a} oq fonli bitta portret edi. \textbf{Rasm b}da portretning foni moviy rangga almashtirilgan, har bir portret atrofiga qora ingichka ramka tushirilgan va ayni portretdan to'rtta nusxa 2x2 tartibda hosil qilingan. Qaysi asbob va amallar mos?
\begin{enumerate}[label=\Alph*),leftmargin=9mm]
\item Paint Bucket Tool - rang; Layer Style =\textgreater{} Stroke - ramka; Alt + Move - nusxa ko'chirish
\item Fill Tool - rang; Layer Style =\textgreater{} Line - ramka; Alt + Pointer - nusxa ko'chirish
\item Stamp Tool - rang; Layer Style =\textgreater{} Frame - ramka; Ctrl+C - nusxa ko'chirish
\item Color Picker Tool - rang; Layer Style =\textgreater{} Shadow - ramka; Alt + Crop - nusxa ko'chirish
\end{enumerate}
\begin{quote}\footnotesize\color{muted}\textbf{Rasm:} Photoshopdagi dastlabki va yakuniy portretlar\\
\href{https://github.com/sarvar9417/attestatsiya/blob/main/public/kuzgi-demo/demo_q17_photoshop.svg}{Rasmni onlayn ko'rish}
\end{quote}
\vspace{2mm}\hrule height 0.2pt \vspace{4mm}
\subsubsection{Grafika turlari - manbali nazariya V2}
\needspace{5\baselineskip}
\hypertarget{q127}{\textbf{\color{ink}127.}}  {\color{muted}\scriptsize Y1 \textbullet\ daraja 1 \textbullet\ DB 310ef54d}
\par
Vektorli tasvir asosan qanday ifodalanadi?
\begin{enumerate}[label=\Alph*),leftmargin=9mm]
\item Faqat HTML teglari bilan
\item Faqat audio signallari bilan
\item Faqat piksel matritsasi bilan
\item Geometrik obyektlar va yo'llar bilan
\end{enumerate}
\vspace{2mm}\hrule height 0.2pt \vspace{4mm}
\needspace{5\baselineskip}
\hypertarget{q128}{\textbf{\color{ink}128.}}  {\color{muted}\scriptsize Y1 \textbullet\ daraja 1 \textbullet\ DB 3ac39921}
\par
Photoshopning qatlamli loyiha formati odatda qaysi?
\begin{enumerate}[label=\Alph*),leftmargin=9mm]
\item HTML
\item PSD
\item MP3
\item JPG
\end{enumerate}
\vspace{2mm}\hrule height 0.2pt \vspace{4mm}
\needspace{5\baselineskip}
\hypertarget{q129}{\textbf{\color{ink}129.}}  {\color{muted}\scriptsize Y1 \textbullet\ daraja 1 \textbullet\ DB 49eb8c08}
\par
CMYK modeli ko'proq qaysi muhit bilan bog'liq?
\begin{enumerate}[label=\Alph*),leftmargin=9mm]
\item Poligrafik chop etish
\item 3D animatsiya tezligi
\item Ovoz siqish
\item HTML forma
\end{enumerate}
\vspace{2mm}\hrule height 0.2pt \vspace{4mm}
\needspace{5\baselineskip}
\hypertarget{q130}{\textbf{\color{ink}130.}}  {\color{muted}\scriptsize Y1 \textbullet\ daraja 1 \textbullet\ DB 69564ecc}
\par
3D modeldagi uchinchi koordinata qaysi?
\begin{enumerate}[label=\Alph*),leftmargin=9mm]
\item z
\item g
\item alpha
\item r
\end{enumerate}
\vspace{2mm}\hrule height 0.2pt \vspace{4mm}
\needspace{5\baselineskip}
\hypertarget{q131}{\textbf{\color{ink}131.}}  {\color{muted}\scriptsize Y1 \textbullet\ daraja 1 \textbullet\ DB 733a2958}
\par
Rastrli tasvirning eng kichik tasvir elementi nima?
\begin{enumerate}[label=\Alph*),leftmargin=9mm]
\item Yo'l (path)
\item Piksel
\item Qatlam palitrasi
\item Kontur tuguni
\end{enumerate}
\vspace{2mm}\hrule height 0.2pt \vspace{4mm}
\needspace{5\baselineskip}
\hypertarget{q132}{\textbf{\color{ink}132.}}  {\color{muted}\scriptsize Y1 \textbullet\ daraja 1 \textbullet\ DB 7f126819}
\par
Fraktal shakl oxirida PNGga eksport qilinsa, to'g'ri xulosa qaysi?
\begin{enumerate}[label=\Alph*),leftmargin=9mm]
\item Matematik model asosida yaratilgan natija rastrga aylantirilgan bo'lishi mumkin
\item PNGning barcha piksellari avtomatik vektorga o'tadi
\item Fraktal tushunchasi faqat PNG ekanini isbotlaydi
\item Fraktal rasmni ekranda ko'rsatish imkonsiz
\end{enumerate}
\vspace{2mm}\hrule height 0.2pt \vspace{4mm}
\needspace{5\baselineskip}
\hypertarget{q133}{\textbf{\color{ink}133.}}  {\color{muted}\scriptsize Y1 \textbullet\ daraja 1 \textbullet\ DB 86650176}
\par
PNGning tipik afzalligi nima?
\begin{enumerate}[label=\Alph*),leftmargin=9mm]
\item Faqat 3D model saqlash
\item Yo'qotishsiz saqlash va shaffoflik
\item Faqat shrift tavsifi
\item Faqat yo'qotishli siqish
\end{enumerate}
\vspace{2mm}\hrule height 0.2pt \vspace{4mm}
\needspace{5\baselineskip}
\hypertarget{q134}{\textbf{\color{ink}134.}}  {\color{muted}\scriptsize Y1 \textbullet\ daraja 1 \textbullet\ DB 90f2d1cf}
\par
SVG faylining asosiy tasvir tavsifi qaysi turga mansub?
\begin{enumerate}[label=\Alph*),leftmargin=9mm]
\item Faqat rastrli
\item Vektorli
\item Faqat jadval
\item Faqat audio
\end{enumerate}
\vspace{2mm}\hrule height 0.2pt \vspace{4mm}
\needspace{5\baselineskip}
\hypertarget{q135}{\textbf{\color{ink}135.}}  {\color{muted}\scriptsize Y1 \textbullet\ daraja 1 \textbullet\ DB 9592fe46}
\par
RGB modelidagi R harfi qaysi rangni bildiradi?
\begin{enumerate}[label=\Alph*),leftmargin=9mm]
\item Rasterize
\item Red (qizil)
\item Resolution
\item Rastr
\end{enumerate}
\vspace{2mm}\hrule height 0.2pt \vspace{4mm}
\needspace{5\baselineskip}
\hypertarget{q136}{\textbf{\color{ink}136.}}  {\color{muted}\scriptsize Y1 \textbullet\ daraja 1 \textbullet\ DB a4b658fa}
\par
JPEG odatda qaysi ma'lumotni saqlaydi?
\begin{enumerate}[label=\Alph*),leftmargin=9mm]
\item CSS xossalarni
\item HTML tuzilmani
\item Rastr tasvirni
\item Vektor tugunlari ro'yxatini
\end{enumerate}
\vspace{2mm}\hrule height 0.2pt \vspace{4mm}
\needspace{5\baselineskip}
\hypertarget{q137}{\textbf{\color{ink}137.}}  {\color{muted}\scriptsize Y1 \textbullet\ daraja 1 \textbullet\ DB cce7d23e}
\par
7-sinf darsligida rastr va vektordan tashqari formula hamda o'ziga o'xshashlik asosida shakl hosil qiluvchi tur qaysi?
\begin{enumerate}[label=\Alph*),leftmargin=9mm]
\item Oddiy JPG fayli
\item Fraktal grafika
\item Piksel jadvali
\item To'g'ri chiziq asbobi
\end{enumerate}
\vspace{2mm}\hrule height 0.2pt \vspace{4mm}
\needspace{5\baselineskip}
\hypertarget{q138}{\textbf{\color{ink}138.}}  {\color{muted}\scriptsize Y1 \textbullet\ daraja 1 \textbullet\ DB f68dd4ce}
\par
RGBA dagi A nimani ifodalaydi?
\begin{enumerate}[label=\Alph*),leftmargin=9mm]
\item Aspect ratio qiymatini
\item Anti-aliasing algoritmini
\item Alpha, shaffoflik kanalini
\item Audio kanalni
\end{enumerate}
\vspace{2mm}\hrule height 0.2pt \vspace{4mm}
\needspace{5\baselineskip}
\hypertarget{q139}{\textbf{\color{ink}139.}}  {\color{muted}\scriptsize Y1 \textbullet\ daraja 2 \textbullet\ DB 21e1b761}
\par
HSB (HSV) rang modelidagi Hue komponenti nimani bildiradi?
\begin{enumerate}[label=\Alph*),leftmargin=9mm]
\item Rang tonini
\item Faylning balandligini
\item Raster piksel sonini
\item Shaffoflik kanalini
\end{enumerate}
\vspace{2mm}\hrule height 0.2pt \vspace{4mm}
\needspace{5\baselineskip}
\hypertarget{q140}{\textbf{\color{ink}140.}}  {\color{muted}\scriptsize Y1 \textbullet\ daraja 2 \textbullet\ DB 4ad0a329}
\par
Qaysi juftlik RGB ekrani va CMYK bosmasini to'g'ri farqlaydi?
\begin{enumerate}[label=\Alph*),leftmargin=9mm]
\item Piksel va HTML
\item Nur qo'shilishi va bo'yoq aralashishi
\item Bo'yoq va ovoz
\item Sikl va shart
\end{enumerate}
\vspace{2mm}\hrule height 0.2pt \vspace{4mm}
\needspace{5\baselineskip}
\hypertarget{q141}{\textbf{\color{ink}141.}}  {\color{muted}\scriptsize Y1 \textbullet\ daraja 2 \textbullet\ DB 4b838082}
\par
SVG logotipni JPEGga eksport qilish qanday oqibat beradi?
\begin{enumerate}[label=\Alph*),leftmargin=9mm]
\item Fayl avtomatik 3D bo'ladi
\item JPEG geometriyani saqlaydi
\item CSS ishlamay qolishi majburiy
\item Eksport qilingan fayl rastrga aylanishi mumkin
\end{enumerate}
\vspace{2mm}\hrule height 0.2pt \vspace{4mm}
\needspace{5\baselineskip}
\hypertarget{q142}{\textbf{\color{ink}142.}}  {\color{muted}\scriptsize Y1 \textbullet\ daraja 2 \textbullet\ DB 6ea417a6}
\par
Qaysi tasvirni PSD sifatida saqlash foydali?
\begin{enumerate}[label=\Alph*),leftmargin=9mm]
\item MP3 audiotrek
\item Matnli CSV jadval
\item Faqat web sahifaga bir zumda qo'yiladigan ikon
\item Qatlamlar bilan keyin ham tahrirlanadigan kompozitsiya
\end{enumerate}
\vspace{2mm}\hrule height 0.2pt \vspace{4mm}
\needspace{5\baselineskip}
\hypertarget{q143}{\textbf{\color{ink}143.}}  {\color{muted}\scriptsize Y1 \textbullet\ daraja 2 \textbullet\ DB 71ca28f8}
\par
Qaysi vazifa uchun vektor manba eng mos?
\begin{enumerate}[label=\Alph*),leftmargin=9mm]
\item Skanerlangan fotosuratni aynan qayta saqlash
\item Videoning ovozini kodlash
\item JPEG artefaktini yashirish
\item Har o'lchamga chiziladigan logotip
\end{enumerate}
\vspace{2mm}\hrule height 0.2pt \vspace{4mm}
\needspace{5\baselineskip}
\hypertarget{q144}{\textbf{\color{ink}144.}}  {\color{muted}\scriptsize Y1 \textbullet\ daraja 2 \textbullet\ DB 78c970de}
\par
Rasm kattalashganda piksellanish muammosi ko'proq nimaga xos?
\begin{enumerate}[label=\Alph*),leftmargin=9mm]
\item Matematik formulaga
\item Vektor manbadagi geometriyaga
\item HTML matn tuzilishiga
\item Past o'lchamli rastrga
\end{enumerate}
\vspace{2mm}\hrule height 0.2pt \vspace{4mm}
\needspace{5\baselineskip}
\hypertarget{q145}{\textbf{\color{ink}145.}}  {\color{muted}\scriptsize Y1 \textbullet\ daraja 2 \textbullet\ DB 86d5cf81}
\par
Qaysi rasm formati odatda an'anaviy 256 rangli palitrali animatsiya uchun ishlatiladi?
\begin{enumerate}[label=\Alph*),leftmargin=9mm]
\item BMP
\item PSD
\item JPG
\item GIF
\end{enumerate}
\vspace{2mm}\hrule height 0.2pt \vspace{4mm}
\needspace{5\baselineskip}
\hypertarget{q146}{\textbf{\color{ink}146.}}  {\color{muted}\scriptsize Y1 \textbullet\ daraja 2 \textbullet\ DB 956366e0}
\par
RGB (255,255,255) odatda qaysi rang?
\begin{enumerate}[label=\Alph*),leftmargin=9mm]
\item Yashil
\item Ko'k
\item Oq
\item Qora
\end{enumerate}
\vspace{2mm}\hrule height 0.2pt \vspace{4mm}
\needspace{5\baselineskip}
\hypertarget{q147}{\textbf{\color{ink}147.}}  {\color{muted}\scriptsize Y1 \textbullet\ daraja 2 \textbullet\ DB 9774a88b}
\par
Qaysi juftlikda birinchisi vektor, ikkinchisi rastr?
\begin{enumerate}[label=\Alph*),leftmargin=9mm]
\item GIF va SVG
\item JPEG va SVG
\item SVG va PNG
\item BMP va PNG
\end{enumerate}
\vspace{2mm}\hrule height 0.2pt \vspace{4mm}
\needspace{5\baselineskip}
\hypertarget{q148}{\textbf{\color{ink}148.}}  {\color{muted}\scriptsize Y1 \textbullet\ daraja 2 \textbullet\ DB a0099621}
\par
Qaysi tasdiq to'g'ri?
\begin{enumerate}[label=\Alph*),leftmargin=9mm]
\item PNG doim vektorli
\item SVG faqat foto uchun yaratilgan
\item JPEG odatda alpha saqlaydi
\item PNG rastrli, lekin uning ichida shaffof piksel bo'lishi mumkin
\end{enumerate}
\vspace{2mm}\hrule height 0.2pt \vspace{4mm}
\needspace{5\baselineskip}
\hypertarget{q149}{\textbf{\color{ink}149.}}  {\color{muted}\scriptsize Y1 \textbullet\ daraja 2 \textbullet\ DB db2cf6d6}
\par
RGB (0,0,0) odatda qaysi rang?
\begin{enumerate}[label=\Alph*),leftmargin=9mm]
\item Oq
\item Qora
\item Sariq
\item Magenta
\end{enumerate}
\vspace{2mm}\hrule height 0.2pt \vspace{4mm}
\needspace{5\baselineskip}
\hypertarget{q150}{\textbf{\color{ink}150.}}  {\color{muted}\scriptsize Y1 \textbullet\ daraja 2 \textbullet\ DB eaf591b6}
\par
Vektorda stroke nimani bildiradi?
\begin{enumerate}[label=\Alph*),leftmargin=9mm]
\item Kontur chizig'ini
\item Ichki fon rangini
\item Fayl siqish tezligini
\item Piksel koordinatalarini
\end{enumerate}
\vspace{2mm}\hrule height 0.2pt \vspace{4mm}
\needspace{5\baselineskip}
\hypertarget{q151}{\textbf{\color{ink}151.}}  {\color{muted}\scriptsize Y1 \textbullet\ daraja 3 \textbullet\ DB 04f458cf}
\par
Rasmdagi 5x4 pikselli jadvalda nechta turli rang bor? Har bir rangni turiga ko'ra sanang.
\begin{enumerate}[label=\Alph*),leftmargin=9mm]
\item 5 ta
\item 4 ta
\item 20 ta
\item 3 ta
\end{enumerate}
\begin{quote}\footnotesize\color{muted}\textbf{Rasm:} 5x4 pikselli to'rt xil rangli jadval\\
\href{https://github.com/sarvar9417/attestatsiya/blob/main/public/graphics-web-v2/palette_grid.svg}{Rasmni onlayn ko'rish}
\end{quote}
\vspace{2mm}\hrule height 0.2pt \vspace{4mm}
\needspace{5\baselineskip}
\hypertarget{q152}{\textbf{\color{ink}152.}}  {\color{muted}\scriptsize Y1 \textbullet\ daraja 3 \textbullet\ DB 0d1f992b}
\par
RGB (0,255,0) qiymatiga eng mos rangni toping.
\begin{enumerate}[label=\Alph*),leftmargin=9mm]
\item qizil
\item sariq
\item yashil
\item ko'k
\end{enumerate}
\vspace{2mm}\hrule height 0.2pt \vspace{4mm}
\needspace{5\baselineskip}
\hypertarget{q153}{\textbf{\color{ink}153.}}  {\color{muted}\scriptsize Y3 \textbullet\ daraja 3 \textbullet\ DB 0f5d52c5}
\par
Quyidagi amallarni to'g'ri tartibda bajaring: 1) Talab qilinadigan chiqish o'lchamini aniqlash; 2) Asl vektor loyihani saqlash; 3) SVG/PNG eksport nusxasini olish.
\begin{enumerate}[label=\Alph*),leftmargin=9mm]
\item Talab qilinadigan chiqish o'lchamini aniqlash
\item Asl vektor loyihani saqlash
\item SVG/PNG eksport nusxasini olish
\end{enumerate}
\vspace{2mm}\hrule height 0.2pt \vspace{4mm}
\needspace{5\baselineskip}
\hypertarget{q154}{\textbf{\color{ink}154.}}  {\color{muted}\scriptsize Y1 \textbullet\ daraja 3 \textbullet\ DB 2071a7ad}
\par
RGB(255,255,255) va RGB(0,0,0) juftligini qanday talqin qilish to'g'ri?
\begin{enumerate}[label=\Alph*),leftmargin=9mm]
\item Qizil va yashil
\item Bir xil oq
\item Oq va qora
\item Ikki xil yashil
\end{enumerate}
\vspace{2mm}\hrule height 0.2pt \vspace{4mm}
\needspace{5\baselineskip}
\hypertarget{q155}{\textbf{\color{ink}155.}}  {\color{muted}\scriptsize Y3 \textbullet\ daraja 3 \textbullet\ DB 485676cb}
\par
Fotosuratni JPEG qilib tarqatishdan oldingi asosiy ish tartibini belgilang.
\begin{enumerate}[label=\Alph*),leftmargin=9mm]
\item Foto uchun kerakli kadr va o'lchamni tanlang
\item JPEG eksportida sifat yoki siqish sozlamasini belgilang
\item Eksport tasvirini siqish artefaktlariga tekshiring
\item Asl sifatli nusxani saqlang
\end{enumerate}
\vspace{2mm}\hrule height 0.2pt \vspace{4mm}
\needspace{5\baselineskip}
\hypertarget{q156}{\textbf{\color{ink}156.}}  {\color{muted}\scriptsize Y3 \textbullet\ daraja 3 \textbullet\ DB 48d15bb1}
\par
Masshtablanadigan logotip tayyorlash tartibini tiklang.
\begin{enumerate}[label=\Alph*),leftmargin=9mm]
\item Tahrirlanadigan vektor manba faylini saqlang
\item Vektor muharririda kontur va ranglarni tahrirlang
\item Logotipning chiziq va shakllarini geometrik obyekt sifatida belgilang
\item Foydalanish joyiga mos SVG yoki rastr eksportini oling
\end{enumerate}
\vspace{2mm}\hrule height 0.2pt \vspace{4mm}
\needspace{5\baselineskip}
\hypertarget{q157}{\textbf{\color{ink}157.}}  {\color{muted}\scriptsize Y2 \textbullet\ daraja 3 \textbullet\ DB 4cde7f81}
\par
Rang modelini uni tavsiflovchi parametrlar bilan moslang.
\textbf{\footnotesize Chap ustun:}
\begin{enumerate}[label=\arabic*),leftmargin=9mm]
\item RGB
\item CMYK
\item HSB
\item CIELAB
\end{enumerate}
\textbf{\footnotesize O'ng ustun:}
\begin{enumerate}[label=\Alph*),leftmargin=9mm]
\item Sian, magenta, sariq va qora bosma bo'yoqlar
\item Ton, to'yinganlik va yorqinlik
\item L* yorug'lik, a* va b* rang o'qlari
\item Qizil, yashil va ko'k yorug'lik kanallari
\end{enumerate}
\vspace{2mm}\hrule height 0.2pt \vspace{4mm}
\needspace{5\baselineskip}
\hypertarget{q158}{\textbf{\color{ink}158.}}  {\color{muted}\scriptsize Y2 \textbullet\ daraja 3 \textbullet\ DB 51c0b65a}
\par
Fayl turini asosiy ish vazifasi bilan moslang.
\textbf{\footnotesize Chap ustun:}
\begin{enumerate}[label=\arabic*),leftmargin=9mm]
\item SVG
\item PSD
\item JPEG
\item PNG
\end{enumerate}
\textbf{\footnotesize O'ng ustun:}
\begin{enumerate}[label=\Alph*),leftmargin=9mm]
\item Masshtablanadigan vektor konturlar
\item Fotosuratni odatda yo'qotishli siqish
\item Shaffof fonli rastr tasvirni yo'qotishsiz saqlash
\item Qatlamlarni saqlaydigan Photoshop loyihasi
\end{enumerate}
\vspace{2mm}\hrule height 0.2pt \vspace{4mm}
\needspace{5\baselineskip}
\hypertarget{q159}{\textbf{\color{ink}159.}}  {\color{muted}\scriptsize Y3 \textbullet\ daraja 3 \textbullet\ DB 6439a536}
\par
Shaffof fonli veb-belgi yaratish jarayonini tartiblang.
\begin{enumerate}[label=\Alph*),leftmargin=9mm]
\item Mos o'lchamdagi rasm yoki vektor belgini tayyorlang
\item Alfani qo'llaydigan PNG yoki mos SVG nusxasini chiqaring
\item Alpha sohasi kutilganidek shaffofligini tekshiring
\item Shaffof fon talabi borligini aniqlang
\end{enumerate}
\vspace{2mm}\hrule height 0.2pt \vspace{4mm}
\needspace{5\baselineskip}
\hypertarget{q160}{\textbf{\color{ink}160.}}  {\color{muted}\scriptsize Y1 \textbullet\ daraja 3 \textbullet\ DB 690bca73}
\par
RGB (255,0,0) qiymatiga eng mos rangni toping.
\begin{enumerate}[label=\Alph*),leftmargin=9mm]
\item qizil
\item yashil
\item ko'k
\item sariq
\end{enumerate}
\vspace{2mm}\hrule height 0.2pt \vspace{4mm}
\needspace{5\baselineskip}
\hypertarget{q161}{\textbf{\color{ink}161.}}  {\color{muted}\scriptsize Y2 \textbullet\ daraja 3 \textbullet\ DB 76725731}
\par
Moslashtiring: 1) PNG; 2) SVG; 3) TIFF. Ta'riflar: a) Rastr format; b) Vektor format; c) Skan uchun tasvir formati.
\textbf{\footnotesize Chap ustun:}
\begin{enumerate}[label=\arabic*),leftmargin=9mm]
\item PNG
\item SVG
\item TIFF
\end{enumerate}
\textbf{\footnotesize O'ng ustun:}
\begin{enumerate}[label=\Alph*),leftmargin=9mm]
\item Rastr format
\item Vektor format
\item Skan uchun tasvir formati
\end{enumerate}
\vspace{2mm}\hrule height 0.2pt \vspace{4mm}
\needspace{5\baselineskip}
\hypertarget{q162}{\textbf{\color{ink}162.}}  {\color{muted}\scriptsize Y1 \textbullet\ daraja 3 \textbullet\ DB 8db60ef2}
\par
RGB (255,255,0) qiymatiga eng mos rangni toping.
\begin{enumerate}[label=\Alph*),leftmargin=9mm]
\item qizil
\item sariq
\item yashil
\item ko'k
\end{enumerate}
\vspace{2mm}\hrule height 0.2pt \vspace{4mm}
\needspace{5\baselineskip}
\hypertarget{q163}{\textbf{\color{ink}163.}}  {\color{muted}\scriptsize Y1 \textbullet\ daraja 3 \textbullet\ DB 9d6e5008}
\par
RGB (0,0,255) qiymatiga eng mos rangni toping.
\begin{enumerate}[label=\Alph*),leftmargin=9mm]
\item ko'k
\item sariq
\item yashil
\item qizil
\end{enumerate}
\vspace{2mm}\hrule height 0.2pt \vspace{4mm}
\needspace{5\baselineskip}
\hypertarget{q164}{\textbf{\color{ink}164.}}  {\color{muted}\scriptsize Y1 \textbullet\ daraja 3 \textbullet\ DB a70459a5}
\par
Grafik model sinovi: o'ziga o'xshash shakllarni formula bilan iterativ qurish topshirildi. Qaysi tasvir turini tanlaysiz?
\begin{enumerate}[label=\Alph*),leftmargin=9mm]
\item Rastr fotosurat montaji
\item Oddiy HTML forma
\item Fraktal
\item Faqat RGB kanali
\end{enumerate}
\vspace{2mm}\hrule height 0.2pt \vspace{4mm}
\needspace{5\baselineskip}
\hypertarget{q165}{\textbf{\color{ink}165.}}  {\color{muted}\scriptsize Y1 \textbullet\ daraja 3 \textbullet\ DB b171b634}
\par
logotipni har xil miqyosda bosishga qaysi format eng mos?
\begin{enumerate}[label=\Alph*),leftmargin=9mm]
\item GIF
\item SVG
\item BMP
\item JPEG
\end{enumerate}
\vspace{2mm}\hrule height 0.2pt \vspace{4mm}
\needspace{5\baselineskip}
\hypertarget{q166}{\textbf{\color{ink}166.}}  {\color{muted}\scriptsize Y1 \textbullet\ daraja 3 \textbullet\ DB bc0292f0}
\par
2-bitli palitrada eng ko'pi bilan nechta rang indeksi bo'lishi mumkin?
\begin{enumerate}[label=\Alph*),leftmargin=9mm]
\item 2
\item 5
\item 8
\item 4
\end{enumerate}
\vspace{2mm}\hrule height 0.2pt \vspace{4mm}
\needspace{5\baselineskip}
\hypertarget{q167}{\textbf{\color{ink}167.}}  {\color{muted}\scriptsize Y2 \textbullet\ daraja 3 \textbullet\ DB d6a10687}
\par
Grafika turini uning ichki tasvirlash modeli bilan moslang.
\textbf{\footnotesize Chap ustun:}
\begin{enumerate}[label=\arabic*),leftmargin=9mm]
\item Rastr
\item Vektor
\item Fraktal
\item 3D
\end{enumerate}
\textbf{\footnotesize O'ng ustun:}
\begin{enumerate}[label=\Alph*),leftmargin=9mm]
\item x, y, z fazoviy geometriyasi
\item Rangli piksellar matritsasi
\item Geometrik yo'l va obyektlar
\item Formula va o'ziga o'xshash takrorlar
\end{enumerate}
\vspace{2mm}\hrule height 0.2pt \vspace{4mm}
\needspace{5\baselineskip}
\hypertarget{q168}{\textbf{\color{ink}168.}}  {\color{muted}\scriptsize Y1 \textbullet\ daraja 3 \textbullet\ DB d71fd2ec}
\par
saytda shaffof fondagi rastr ikonka uchun qaysi format eng mos?
\begin{enumerate}[label=\Alph*),leftmargin=9mm]
\item TXT
\item JPEG
\item MP3
\item PNG
\end{enumerate}
\vspace{2mm}\hrule height 0.2pt \vspace{4mm}
\needspace{5\baselineskip}
\hypertarget{q169}{\textbf{\color{ink}169.}}  {\color{muted}\scriptsize Y1 \textbullet\ daraja 4 \textbullet\ DB 0acbe86b}
\par
1280x720 pikselli rastr tasvir nechta pikselga ega?
\begin{enumerate}[label=\Alph*),leftmargin=9mm]
\item 460800
\item 921600
\item 2000
\item 1843200
\end{enumerate}
\vspace{2mm}\hrule height 0.2pt \vspace{4mm}
\needspace{5\baselineskip}
\hypertarget{q170}{\textbf{\color{ink}170.}}  {\color{muted}\scriptsize Y1 \textbullet\ daraja 4 \textbullet\ DB 103ab181}
\par
300x200 pikselli rastr tasvir nechta pikselga ega?
\begin{enumerate}[label=\Alph*),leftmargin=9mm]
\item 500
\item 30000
\item 120000
\item 60000
\end{enumerate}
\vspace{2mm}\hrule height 0.2pt \vspace{4mm}
\needspace{5\baselineskip}
\hypertarget{q171}{\textbf{\color{ink}171.}}  {\color{muted}\scriptsize Y3 \textbullet\ daraja 4 \textbullet\ DB 39c86c4b}
\par
3D sahnani tekis tasvirga chiqarish zanjirini tartiblang.
\begin{enumerate}[label=\Alph*),leftmargin=9mm]
\item Material va rang xossalarini sozlang
\item Rendering orqali ikki o'lchovli kadr oling
\item Obyektning fazoviy shaklini modellashtiring
\item Kamera nuqtasi va yoritishni joylashtiring
\end{enumerate}
\vspace{2mm}\hrule height 0.2pt \vspace{4mm}
\needspace{5\baselineskip}
\hypertarget{q172}{\textbf{\color{ink}172.}}  {\color{muted}\scriptsize Y1 \textbullet\ daraja 4 \textbullet\ DB 47b450da}
\par
HSV/HSBda Saturation 0\% qilib belgilansa, rang odatda qanday tus oladi?
\begin{enumerate}[label=\Alph*),leftmargin=9mm]
\item Faqat qora
\item Mutlaqo shaffof
\item Doim yorqin qizil
\item Kulrangga yaqin, to'yingan bo'lmagan
\end{enumerate}
\vspace{2mm}\hrule height 0.2pt \vspace{4mm}
\needspace{5\baselineskip}
\hypertarget{q173}{\textbf{\color{ink}173.}}  {\color{muted}\scriptsize Y3 \textbullet\ daraja 4 \textbullet\ DB 52e8c4b9}
\par
Qatlamli grafika loyihasidan veb uchun yakuniy rasm olishni tartiblang.
\begin{enumerate}[label=\Alph*),leftmargin=9mm]
\item Tahrirlanadigan PSD loyihani saqlang
\item Veb uchun kerakli piksel o'lchami va shaffoflikni aniqlang
\item PNG/JPEG nusxasini eksport qilib tekshiring
\item Loyihadagi qatlamlar va ularning ko'rinishini tayyorlang
\end{enumerate}
\vspace{2mm}\hrule height 0.2pt \vspace{4mm}
\needspace{5\baselineskip}
\hypertarget{q174}{\textbf{\color{ink}174.}}  {\color{muted}\scriptsize Y1 \textbullet\ daraja 4 \textbullet\ DB 639b795f}
\par
Bir nechta skanlangan tasvir yuqori sifatda arxivlanadi. Darslik formatlari ichida qaysi format ko'proq mos keladi?
\begin{enumerate}[label=\Alph*),leftmargin=9mm]
\item EXE
\item HTML
\item TIFF
\item MP3
\end{enumerate}
\vspace{2mm}\hrule height 0.2pt \vspace{4mm}
\needspace{5\baselineskip}
\hypertarget{q175}{\textbf{\color{ink}175.}}  {\color{muted}\scriptsize Y1 \textbullet\ daraja 4 \textbullet\ DB 640468d1}
\par
1024x768 pikselli rastr tasvir nechta pikselga ega?
\begin{enumerate}[label=\Alph*),leftmargin=9mm]
\item 1572864
\item 1792
\item 393216
\item 786432
\end{enumerate}
\vspace{2mm}\hrule height 0.2pt \vspace{4mm}
\needspace{5\baselineskip}
\hypertarget{q176}{\textbf{\color{ink}176.}}  {\color{muted}\scriptsize Y3 \textbullet\ daraja 4 \textbullet\ DB 76f7f27d}
\par
Bosma banner uchun rastr tasvir o'lchamini tekshirishni tartiblang.
\begin{enumerate}[label=\Alph*),leftmargin=9mm]
\item Bosmada kerakli fizik kenglikni dyuymda belgilang
\item Mavjud rasmning gorizontal piksel sonini aniqlang
\item Piksellarni fizik kenglikka bo'lib haqiqiy PPI ni toping
\item PPI talab bilan solishtirib mos tasvir yoki qayta ishlashni tanlang
\end{enumerate}
\vspace{2mm}\hrule height 0.2pt \vspace{4mm}
\needspace{5\baselineskip}
\hypertarget{q177}{\textbf{\color{ink}177.}}  {\color{muted}\scriptsize Y1 \textbullet\ daraja 4 \textbullet\ DB 79192b54}
\par
Logo SVGda yaratiladi, saytda esa PNG nusxasi ishlatiladi. Qaysi xulosa to'g'ri?
\begin{enumerate}[label=\Alph*),leftmargin=9mm]
\item Asl manba vektor, veb nusxa rastr
\item Ikkalasi vektor
\item PNG manba geometriyani saqlaydi
\item Ikkalasi 3D
\end{enumerate}
\vspace{2mm}\hrule height 0.2pt \vspace{4mm}
\needspace{5\baselineskip}
\hypertarget{q178}{\textbf{\color{ink}178.}}  {\color{muted}\scriptsize Y1 \textbullet\ daraja 4 \textbullet\ DB 920ce394}
\par
8 bitli indekslangan rang xususiyatida maksimal nechta indeks bo'ladi?
\begin{enumerate}[label=\Alph*),leftmargin=9mm]
\item 512
\item 257
\item 256
\item 128
\end{enumerate}
\vspace{2mm}\hrule height 0.2pt \vspace{4mm}
\needspace{5\baselineskip}
\hypertarget{q179}{\textbf{\color{ink}179.}}  {\color{muted}\scriptsize Y1 \textbullet\ daraja 4 \textbullet\ DB c252348e}
\par
640x480 pikselli rastr tasvir nechta pikselga ega?
\begin{enumerate}[label=\Alph*),leftmargin=9mm]
\item 307200
\item 1120
\item 614400
\item 153600
\end{enumerate}
\vspace{2mm}\hrule height 0.2pt \vspace{4mm}
\needspace{5\baselineskip}
\hypertarget{q180}{\textbf{\color{ink}180.}}  {\color{muted}\scriptsize Y2 \textbullet\ daraja 4 \textbullet\ DB c39c727a}
\par
Rang qiymati va uning RGB yoki alpha talqinini moslang.
\textbf{\footnotesize Chap ustun:}
\begin{enumerate}[label=\arabic*),leftmargin=9mm]
\item RGB(0,0,0)
\item RGB(255,255,255)
\item RGB(255,255,0)
\item Alpha=0 (8 bit)
\end{enumerate}
\textbf{\footnotesize O'ng ustun:}
\begin{enumerate}[label=\Alph*),leftmargin=9mm]
\item Qora
\item To'liq shaffof
\item Sariq
\item Oq
\end{enumerate}
\vspace{2mm}\hrule height 0.2pt \vspace{4mm}
\needspace{5\baselineskip}
\hypertarget{q181}{\textbf{\color{ink}181.}}  {\color{muted}\scriptsize Y2 \textbullet\ daraja 4 \textbullet\ DB c94dab82}
\par
Rasm haqidagi texnik atamalarni ma'nolari bilan moslang.
\textbf{\footnotesize Chap ustun:}
\begin{enumerate}[label=\arabic*),leftmargin=9mm]
\item Piksel
\item Alpha
\item PPI
\item Rang chuqurligi
\end{enumerate}
\textbf{\footnotesize O'ng ustun:}
\begin{enumerate}[label=\Alph*),leftmargin=9mm]
\item Piksel shaffofligi kanali
\item Rastr matritsasidagi rang elementi
\item Bitta piksel rangini ifodalovchi bitlar soni
\item Bosma tasvirdagi bir dyuymga to'g'ri keladigan piksel
\end{enumerate}
\vspace{2mm}\hrule height 0.2pt \vspace{4mm}
\needspace{5\baselineskip}
\hypertarget{q182}{\textbf{\color{ink}182.}}  {\color{muted}\scriptsize Y2 \textbullet\ daraja 4 \textbullet\ DB d12407bd}
\par
Moslashtiring: 1) Fraktal; 2) Rastr; 3) Vektor. Ta'riflar: a) Formula bilan o'ziga o'xshashlik; b) Geometrik kontur; c) Piksel matritsasi.
\textbf{\footnotesize Chap ustun:}
\begin{enumerate}[label=\arabic*),leftmargin=9mm]
\item Fraktal
\item Rastr
\item Vektor
\end{enumerate}
\textbf{\footnotesize O'ng ustun:}
\begin{enumerate}[label=\Alph*),leftmargin=9mm]
\item Formula bilan o'ziga o'xshashlik
\item Geometrik kontur
\item Piksel matritsasi
\end{enumerate}
\vspace{2mm}\hrule height 0.2pt \vspace{4mm}
\needspace{5\baselineskip}
\hypertarget{q183}{\textbf{\color{ink}183.}}  {\color{muted}\scriptsize Y1 \textbullet\ daraja 4 \textbullet\ DB d904ad97}
\par
1920x1080 pikselli rastr tasvir nechta pikselga ega?
\begin{enumerate}[label=\Alph*),leftmargin=9mm]
\item 4147200
\item 1036800
\item 2073600
\item 3000
\end{enumerate}
\vspace{2mm}\hrule height 0.2pt \vspace{4mm}
\needspace{5\baselineskip}
\hypertarget{q184}{\textbf{\color{ink}184.}}  {\color{muted}\scriptsize Y3 \textbullet\ daraja 4 \textbullet\ DB de9be24a}
\par
Quyidagi amallarni to'g'ri tartibda bajaring: 1) Vektor modelni chizish; 2) Export nusxasini tekshirish; 3) Eksport formatini tanlash.
\begin{enumerate}[label=\Alph*),leftmargin=9mm]
\item Vektor modelni chizish
\item Export nusxasini tekshirish
\item Eksport formatini tanlash
\end{enumerate}
\vspace{2mm}\hrule height 0.2pt \vspace{4mm}
\needspace{5\baselineskip}
\hypertarget{q185}{\textbf{\color{ink}185.}}  {\color{muted}\scriptsize Y2 \textbullet\ daraja 4 \textbullet\ DB ead25e92}
\par
Rasm formati va uning muhim cheklovi yoki kuchli tomonini moslang.
\textbf{\footnotesize Chap ustun:}
\begin{enumerate}[label=\arabic*),leftmargin=9mm]
\item GIF
\item TIFF
\item BMP
\item SVG
\end{enumerate}
\textbf{\footnotesize O'ng ustun:}
\begin{enumerate}[label=\Alph*),leftmargin=9mm]
\item Ba'zan katta hajmli, siqilmagan rastr
\item Skan va ishchi sifatli rastr saqlash
\item Cheklangan palitra va oddiy animatsiya
\item Konturlarni kattalashtirganda qayta chizish
\end{enumerate}
\vspace{2mm}\hrule height 0.2pt \vspace{4mm}
\needspace{5\baselineskip}
\hypertarget{q186}{\textbf{\color{ink}186.}}  {\color{muted}\scriptsize Y1 \textbullet\ daraja 4 \textbullet\ DB fd935bab}
\par
1200x800 pikselli rastr tasvir nechta pikselga ega?
\begin{enumerate}[label=\Alph*),leftmargin=9mm]
\item 2000
\item 480000
\item 960000
\item 1920000
\end{enumerate}
\vspace{2mm}\hrule height 0.2pt \vspace{4mm}
\needspace{5\baselineskip}
\hypertarget{q187}{\textbf{\color{ink}187.}}  {\color{muted}\scriptsize Y1 \textbullet\ daraja 5 \textbullet\ DB 102eddec}
\par
Kichik matnli logo reklama banneri uchun kattalashtiriladi. Qaysi manba dizayni qulay?
\begin{enumerate}[label=\Alph*),leftmargin=9mm]
\item Vektorli konturli original
\item Past sifatli skrinshotni kesish
\item 48x48 BMPni ekrandan ko'chirib olish
\item 72x72 pikselli JPEGni tinimsiz cho'zish
\end{enumerate}
\vspace{2mm}\hrule height 0.2pt \vspace{4mm}
\needspace{5\baselineskip}
\hypertarget{q188}{\textbf{\color{ink}188.}}  {\color{muted}\scriptsize Y1 \textbullet\ daraja 5 \textbullet\ DB 107bbf36}
\par
3D model JPEGda rasm holida saqlandi. Qaysi ma'lumot yo'qolishi mumkin?
\begin{enumerate}[label=\Alph*),leftmargin=9mm]
\item 3D geometriyani mustaqil tahrirlash imkoniyati
\item JPEG kengaytmasi
\item 2D piksel qiymatlari
\item Fayl nomi
\end{enumerate}
\vspace{2mm}\hrule height 0.2pt \vspace{4mm}
\needspace{5\baselineskip}
\hypertarget{q189}{\textbf{\color{ink}189.}}  {\color{muted}\scriptsize Y3 \textbullet\ daraja 5 \textbullet\ DB 28fe24c1}
\par
Grafika fayli turini aniqlashda to'g'ri tahlil ketma-ketligini tuzing.
\begin{enumerate}[label=\Alph*),leftmargin=9mm]
\item Faylning asl mazmuni va yaratilgan dasturini aniqlang
\item Topilgan modelga asoslanib rastr/vektor xulosasini bering
\item Masshtablash va mustaqil obyekt tahriri mavjudligini sinang
\item Ichki tasvir modeli pikselmi yoki geometrik obyektmi tekshiring
\end{enumerate}
\vspace{2mm}\hrule height 0.2pt \vspace{4mm}
\needspace{5\baselineskip}
\hypertarget{q190}{\textbf{\color{ink}190.}}  {\color{muted}\scriptsize Y2 \textbullet\ daraja 5 \textbullet\ DB 38dc8040}
\par
Dizayn vaziyatini eng ma'qul manba yoki chiqish turi bilan moslang.
\textbf{\footnotesize Chap ustun:}
\begin{enumerate}[label=\arabic*),leftmargin=9mm]
\item Kattalashtiriladigan logotip manbasi
\item Qatlamlari keyin tahrirlanadigan maket
\item Murakkab fotosuratni ixcham tarqatish
\item Shaffof fonli veb-piktogramma
\end{enumerate}
\textbf{\footnotesize O'ng ustun:}
\begin{enumerate}[label=\Alph*),leftmargin=9mm]
\item Vektor/SVG
\item JPEG eksport
\item PSD loyiha
\item PNG eksport
\end{enumerate}
\vspace{2mm}\hrule height 0.2pt \vspace{4mm}
\needspace{5\baselineskip}
\hypertarget{q191}{\textbf{\color{ink}191.}}  {\color{muted}\scriptsize Y1 \textbullet\ daraja 5 \textbullet\ DB 50a505af}
\par
640x480 piksel, 32 bit/piksel RGBA xom piksel massivi necha bayt? (siqish va header yo'q)
\begin{enumerate}[label=\Alph*),leftmargin=9mm]
\item 1228800
\item 921600
\item 2457600
\item 307200
\end{enumerate}
\vspace{2mm}\hrule height 0.2pt \vspace{4mm}
\needspace{5\baselineskip}
\hypertarget{q192}{\textbf{\color{ink}192.}}  {\color{muted}\scriptsize Y1 \textbullet\ daraja 5 \textbullet\ DB 5f3ab62e}
\par
PNGdagi shaffoflik va GIFdagi palitra cheklovlarini bir xil deb hisoblash to'g'rimi?
\begin{enumerate}[label=\Alph*),leftmargin=9mm]
\item Ha, ikkisi faqat 1 xil rangni saqlaydi
\item Yo'q, format imkoniyatlari va shaffoflik modellari farq qiladi
\item Ha, ikkalasi MP4 kodek
\item Ha, har ikkalasi faqat vektor
\end{enumerate}
\vspace{2mm}\hrule height 0.2pt \vspace{4mm}
\needspace{5\baselineskip}
\hypertarget{q193}{\textbf{\color{ink}193.}}  {\color{muted}\scriptsize Y1 \textbullet\ daraja 5 \textbullet\ DB 62c50491}
\par
Rasmdagi ikki fayldan biri fotosurat uchun JPEG, ikkinchisi logotip uchun SVG. Qaysi xulosa asosli?
\begin{enumerate}[label=\Alph*),leftmargin=9mm]
\item Ikkalasi ham faqat audio format
\item JPEG barcha qatlamlarni o'zgartirmasdan saqlaydi
\item Har bir format o'z vazifasiga mos; biri rastr, ikkinchisi vektor
\item SVG faqat piksellar sonini saqlaydi
\end{enumerate}
\vspace{2mm}\hrule height 0.2pt \vspace{4mm}
\needspace{5\baselineskip}
\hypertarget{q194}{\textbf{\color{ink}194.}}  {\color{muted}\scriptsize Y1 \textbullet\ daraja 5 \textbullet\ DB 65efe398}
\par
Fonda shaffof piksel kerak, ammo JPEGga oddiy eksport qilindi. Qaysi xulosa xavfsiz?
\begin{enumerate}[label=\Alph*),leftmargin=9mm]
\item JPEG har doim barcha qatlamni saqlaydi
\item JPEG oddiy alfa shaffoflikni saqlash uchun mos format emas
\item JPEG doim 8-kanalli alfa qo'shadi
\item JPEG vektor konturlarini o'zicha saqlaydi
\end{enumerate}
\vspace{2mm}\hrule height 0.2pt \vspace{4mm}
\needspace{5\baselineskip}
\hypertarget{q195}{\textbf{\color{ink}195.}}  {\color{muted}\scriptsize Y3 \textbullet\ daraja 5 \textbullet\ DB 68360d33}
\par
Fraktal tasvirni yasash va ko'rsatishning mantiqiy tartibini tuzing.
\begin{enumerate}[label=\Alph*),leftmargin=9mm]
\item Hosil bo'lgan natijani ekranda piksellar bilan tasvirlang
\item Qaytariladigan matematik formulani tanlang
\item Iteratsiya chegarasi va hisoblash parametrlarini kiriting
\item Formula bo'yicha nuqtalar yoki sohalarni hisoblang
\end{enumerate}
\vspace{2mm}\hrule height 0.2pt \vspace{4mm}
\needspace{5\baselineskip}
\hypertarget{q196}{\textbf{\color{ink}196.}}  {\color{muted}\scriptsize Y1 \textbullet\ daraja 5 \textbullet\ DB 68bea628}
\par
RGB uchun 0 deg /120 deg /240 deg  HSB rang tonlari darslikda qaysi uch rang atrofi sifatida beriladi?
\begin{enumerate}[label=\Alph*),leftmargin=9mm]
\item Qora, moviy, oq
\item Oq, qora, kulrang
\item Qizil, yashil, ko'k
\item Sariq, kulrang, qora
\end{enumerate}
\vspace{2mm}\hrule height 0.2pt \vspace{4mm}
\needspace{5\baselineskip}
\hypertarget{q197}{\textbf{\color{ink}197.}}  {\color{muted}\scriptsize Y2 \textbullet\ daraja 5 \textbullet\ DB 69f86a02}
\par
Grafika jarayonini foydalanilgan asosiy tushuncha bilan moslang.
\textbf{\footnotesize Chap ustun:}
\begin{enumerate}[label=\arabic*),leftmargin=9mm]
\item 3D modellashtirish
\item 3D rendering
\item Rasterlash
\item Tracing
\end{enumerate}
\textbf{\footnotesize O'ng ustun:}
\begin{enumerate}[label=\Alph*),leftmargin=9mm]
\item Vektor tasvirni piksel to'riga aylantirish
\item Rastr konturidan vektor yo'lni tiklashga urinish
\item Modeldan ko'rinadigan 2D tasvirni hisoblash
\item Fazoviy geometriyani qurish
\end{enumerate}
\vspace{2mm}\hrule height 0.2pt \vspace{4mm}
\needspace{5\baselineskip}
\hypertarget{q198}{\textbf{\color{ink}198.}}  {\color{muted}\scriptsize Y1 \textbullet\ daraja 5 \textbullet\ DB 87c6476c}
\par
LAB modelida L* kanali qanday xususiyat bilan bog'liq?
\begin{enumerate}[label=\Alph*),leftmargin=9mm]
\item Grafik faylning megabayti
\item HTML havolasi
\item Yorug'lik darajasi
\item Faqat alfa shaffoflik
\end{enumerate}
\vspace{2mm}\hrule height 0.2pt \vspace{4mm}
\needspace{5\baselineskip}
\hypertarget{q199}{\textbf{\color{ink}199.}}  {\color{muted}\scriptsize Y1 \textbullet\ daraja 5 \textbullet\ DB 9448c06f}
\par
Kichik shriftli diagramma kengaytiriladi
\begin{enumerate}[label=\Alph*),leftmargin=9mm]
\item Manba SVG yoki yuqori aniqlikdagi rastrdan foydalanish
\item JPEGni GIF deb nomlash
\item Kichik JPEGni 5000 foiz kattalashtirishning o'zi
\item Bir rang kanalini o'chirish
\end{enumerate}
\vspace{2mm}\hrule height 0.2pt \vspace{4mm}
\needspace{5\baselineskip}
\hypertarget{q200}{\textbf{\color{ink}200.}}  {\color{muted}\scriptsize Y1 \textbullet\ daraja 5 \textbullet\ DB a7eafdcf}
\par
3D sahna 2D PNGga eksport qilindi. Keyingi tahrir uchun nimani yo'qotish xavfi bor?
\begin{enumerate}[label=\Alph*),leftmargin=9mm]
\item Fazoviy geometriya va obyektlarning alohida tahrirlash imkonini
\item Piksel ranglarini butunlay
\item Rasmni monitorda ochish imkonini
\item Fayl nomini tanlash imkonini
\end{enumerate}
\vspace{2mm}\hrule height 0.2pt \vspace{4mm}
\needspace{5\baselineskip}
\hypertarget{q201}{\textbf{\color{ink}201.}}  {\color{muted}\scriptsize Y1 \textbullet\ daraja 5 \textbullet\ DB bb2f6225}
\par
HSVda Hue o'zgarmas, Saturation 100\% dan 0\%ga tushsa eng asosli kutilma qaysi?
\begin{enumerate}[label=\Alph*),leftmargin=9mm]
\item Rang to'yinganligi yo'qolib kulrang tomonga o'tadi
\item Yangi piksel avtomatik yaratiladi
\item Tasvir doim shaffof bo'ladi
\item Gorizontal rezolyutsiya ortadi
\end{enumerate}
\vspace{2mm}\hrule height 0.2pt \vspace{4mm}
\needspace{5\baselineskip}
\hypertarget{q202}{\textbf{\color{ink}202.}}  {\color{muted}\scriptsize Y1 \textbullet\ daraja 5 \textbullet\ DB d43e3140}
\par
Rastr fotosuratdan logotip chizmasini vektor sifatida qayta ishlash
\begin{enumerate}[label=\Alph*),leftmargin=9mm]
\item PNGni ZIPga solish
\item Faylni .svg deb qayta nomlash
\item JPEG sifatini kamaytirish
\item Konturlarni tracing orqali ajratib, tugunlarni tekshirish
\end{enumerate}
\vspace{2mm}\hrule height 0.2pt \vspace{4mm}
\subsubsection{Rastr va vektor amallari - nazariya V2}
\needspace{5\baselineskip}
\hypertarget{q203}{\textbf{\color{ink}203.}}  {\color{muted}\scriptsize Y1 \textbullet\ daraja 1 \textbullet\ DB 17207b64}
\par
Indekslangan 8 bitli rangda 1 piksel odatda nimani saqlaydi?
\begin{enumerate}[label=\Alph*),leftmargin=9mm]
\item Palitradagi 256 tagacha indeksdan birini
\item Albatta 3 bayt to'liq RGBni
\item Faqat qora va oq bitta bitni
\item Har doim 32 bitli alpha kanali
\end{enumerate}
\vspace{2mm}\hrule height 0.2pt \vspace{4mm}
\needspace{5\baselineskip}
\hypertarget{q204}{\textbf{\color{ink}204.}}  {\color{muted}\scriptsize Y1 \textbullet\ daraja 1 \textbullet\ DB 41520e5a}
\par
Canvas size odatda nimani o'zgartiradi?
\begin{enumerate}[label=\Alph*),leftmargin=9mm]
\item Tasvir formatidagi JPEG algoritmini
\item Ish maydonining chegaralarini
\item Har bir pikseldagi rang bitini
\item Printer tezligini
\end{enumerate}
\vspace{2mm}\hrule height 0.2pt \vspace{4mm}
\needspace{5\baselineskip}
\hypertarget{q205}{\textbf{\color{ink}205.}}  {\color{muted}\scriptsize Y1 \textbullet\ daraja 1 \textbullet\ DB 7d63a85a}
\par
Crop amali qanday vazifani bajaradi?
\begin{enumerate}[label=\Alph*),leftmargin=9mm]
\item Audio faylni qisqartiradi
\item Rang profilini almashtiradi
\item Tasvirning kerakli qismini qoldirib qolganini kesadi
\item Vektor birlashmasini hisoblaydi
\end{enumerate}
\vspace{2mm}\hrule height 0.2pt \vspace{4mm}
\needspace{5\baselineskip}
\hypertarget{q206}{\textbf{\color{ink}206.}}  {\color{muted}\scriptsize Y1 \textbullet\ daraja 1 \textbullet\ DB 8e2dc024}
\par
Resize amali nimani o'zgartiradi?
\begin{enumerate}[label=\Alph*),leftmargin=9mm]
\item Faqat fayl kengaytmasini
\item Tasvirning ko'rsatilgan o'lchamlarini
\item Faqat qatlam nomini
\item Faqat matn kodlanishini
\end{enumerate}
\vspace{2mm}\hrule height 0.2pt \vspace{4mm}
\needspace{5\baselineskip}
\hypertarget{q207}{\textbf{\color{ink}207.}}  {\color{muted}\scriptsize Y1 \textbullet\ daraja 1 \textbullet\ DB 9c00f222}
\par
32-bit RGBAda bir piksel uchun xom ma'lumot qancha?
\begin{enumerate}[label=\Alph*),leftmargin=9mm]
\item 4 bayt
\item 3 bayt
\item 32 bayt
\item 8 bayt
\end{enumerate}
\vspace{2mm}\hrule height 0.2pt \vspace{4mm}
\needspace{5\baselineskip}
\hypertarget{q208}{\textbf{\color{ink}208.}}  {\color{muted}\scriptsize Y1 \textbullet\ daraja 1 \textbullet\ DB a63a899f}
\par
Clipping mask nimani boshqaradi?
\begin{enumerate}[label=\Alph*),leftmargin=9mm]
\item Rasmning GPS joyini
\item CMYKning kanal nomini
\item Ko'rinadigan hudud chegarasini
\item Faylning MD5 summasini
\end{enumerate}
\vspace{2mm}\hrule height 0.2pt \vspace{4mm}
\needspace{5\baselineskip}
\hypertarget{q209}{\textbf{\color{ink}209.}}  {\color{muted}\scriptsize Y1 \textbullet\ daraja 1 \textbullet\ DB bc912025}
\par
Resize va Zoom orasidagi asosiy farq nimada?
\begin{enumerate}[label=\Alph*),leftmargin=9mm]
\item Ikkalasi fayl o'lchamini har doim ikki baravar qiladi
\item Resize piksel o'lchamini o'zgartirishi mumkin, Zoom faqat ko'rish masshtabini
\item Zoom rasmning piksel sonini oshiradi
\item Resize hech qachon piksel sonini o'zgartirmaydi
\end{enumerate}
\vspace{2mm}\hrule height 0.2pt \vspace{4mm}
\needspace{5\baselineskip}
\hypertarget{q210}{\textbf{\color{ink}210.}}  {\color{muted}\scriptsize Y1 \textbullet\ daraja 1 \textbullet\ DB dd0d5288}
\par
24-bit RGBda bir piksel uchun xom ma'lumot qancha?
\begin{enumerate}[label=\Alph*),leftmargin=9mm]
\item 3 bayt
\item 1 bayt
\item 8 bayt
\item 24 bayt
\end{enumerate}
\vspace{2mm}\hrule height 0.2pt \vspace{4mm}
\needspace{5\baselineskip}
\hypertarget{q211}{\textbf{\color{ink}211.}}  {\color{muted}\scriptsize Y1 \textbullet\ daraja 1 \textbullet\ DB fa16e573}
\par
Vector rasterization nimani anglatadi?
\begin{enumerate}[label=\Alph*),leftmargin=9mm]
\item Shaklni ko'chirish
\item Pikselni audio qilish
\item Vektor tavsifini piksel tasvirga aylantirish
\item Layerlarni yashirish
\end{enumerate}
\vspace{2mm}\hrule height 0.2pt \vspace{4mm}
\needspace{5\baselineskip}
\hypertarget{q212}{\textbf{\color{ink}212.}}  {\color{muted}\scriptsize Y1 \textbullet\ daraja 1 \textbullet\ DB fb5ccbd4}
\par
Horizontal flip nima?
\begin{enumerate}[label=\Alph*),leftmargin=9mm]
\item 90 daraja burish
\item Faylni eksport qilish
\item Chap-o'ng tomonni akslantirish
\item Qatlamni shifrlash
\end{enumerate}
\vspace{2mm}\hrule height 0.2pt \vspace{4mm}
\needspace{5\baselineskip}
\hypertarget{q213}{\textbf{\color{ink}213.}}  {\color{muted}\scriptsize Y1 \textbullet\ daraja 1 \textbullet\ DB fc18ea02}
\par
Vektorda group amali nima qiladi?
\begin{enumerate}[label=\Alph*),leftmargin=9mm]
\item Har bir obyektni JPEGga yozadi
\item Obyektlarni bir guruh sifatida boshqarishga birlashtiradi
\item Faylni parollaydi
\item Piksellarni siqadi
\end{enumerate}
\vspace{2mm}\hrule height 0.2pt \vspace{4mm}
\needspace{5\baselineskip}
\hypertarget{q214}{\textbf{\color{ink}214.}}  {\color{muted}\scriptsize Y1 \textbullet\ daraja 1 \textbullet\ DB fdfb5ae8}
\par
Resampling nima?
\begin{enumerate}[label=\Alph*),leftmargin=9mm]
\item HTML kodni formatlash
\item Piksellar sonini o'zgartirishda yangi namunalarni hisoblash
\item Vector group qilish
\item Fayl nomini almashtirish
\end{enumerate}
\vspace{2mm}\hrule height 0.2pt \vspace{4mm}
\needspace{5\baselineskip}
\hypertarget{q215}{\textbf{\color{ink}215.}}  {\color{muted}\scriptsize Y1 \textbullet\ daraja 2 \textbullet\ DB 17786e96}
\par
Palitrali 8-bitli rasm va RGB24 farqi nima?
\begin{enumerate}[label=\Alph*),leftmargin=9mm]
\item Birinchida 8-bit indeks, ikkinchida 24-bit rang qiymati
\item Palitra faqat 3D bo'ladi
\item RGB24 da 8 rang bor
\item Ikkalasi 24-bit
\end{enumerate}
\vspace{2mm}\hrule height 0.2pt \vspace{4mm}
\needspace{5\baselineskip}
\hypertarget{q216}{\textbf{\color{ink}216.}}  {\color{muted}\scriptsize Y1 \textbullet\ daraja 2 \textbullet\ DB 18b3ffba}
\par
Qaysi holat xom piksel formulasini aniq qo'llashga mos?
\begin{enumerate}[label=\Alph*),leftmargin=9mm]
\item Metadata hajmi hisoblanmagan BMPning aniq fayl hajmi
\item JPEG diskdagi hajmi so'ralgan
\item PSD faylining qatlamlari noma'lum
\item Siqilmagan, piksel soni va bit chuqurligi berilgan
\end{enumerate}
\vspace{2mm}\hrule height 0.2pt \vspace{4mm}
\needspace{5\baselineskip}
\hypertarget{q217}{\textbf{\color{ink}217.}}  {\color{muted}\scriptsize Y1 \textbullet\ daraja 2 \textbullet\ DB 2e21faf2}
\par
Vector union qanday amal?
\begin{enumerate}[label=\Alph*),leftmargin=9mm]
\item Faqat umumiy qismi
\item Qatlamni bekitish
\item Biridan boshqasini ayirish
\item Ikkita shaklning birlashmasi
\end{enumerate}
\vspace{2mm}\hrule height 0.2pt \vspace{4mm}
\needspace{5\baselineskip}
\hypertarget{q218}{\textbf{\color{ink}218.}}  {\color{muted}\scriptsize Y1 \textbullet\ daraja 2 \textbullet\ DB 2efee3c4}
\par
Vector intersection qanday amal?
\begin{enumerate}[label=\Alph*),leftmargin=9mm]
\item Faqat birinchi shakl
\item Ikkala shaklning birlashmasi
\item Shakllarning umumiy sohasi
\item JPEG siqish
\end{enumerate}
\vspace{2mm}\hrule height 0.2pt \vspace{4mm}
\needspace{5\baselineskip}
\hypertarget{q219}{\textbf{\color{ink}219.}}  {\color{muted}\scriptsize Y1 \textbullet\ daraja 2 \textbullet\ DB 63533b4c}
\par
Resample va PPI metama'lumotini o'zgartirish farqi nima?
\begin{enumerate}[label=\Alph*),leftmargin=9mm]
\item PPI faqat video uchun
\item Resample faqat CSSga tegishli
\item Ikkalasi har doim bir xil
\item Birinchisi piksel sonini o'zgartirishi mumkin, ikkinchisi shart emas
\end{enumerate}
\vspace{2mm}\hrule height 0.2pt \vspace{4mm}
\needspace{5\baselineskip}
\hypertarget{q220}{\textbf{\color{ink}220.}}  {\color{muted}\scriptsize Y1 \textbullet\ daraja 2 \textbullet\ DB 66a425c8}
\par
Bring to front qaysi natijani beradi?
\begin{enumerate}[label=\Alph*),leftmargin=9mm]
\item Shakl hajmini 2 baravar qiladi
\item Rang kanallarini almashtiradi
\item Obyektni ustki chizilish tartibiga chiqaradi
\item Har doim chapga suradi
\end{enumerate}
\vspace{2mm}\hrule height 0.2pt \vspace{4mm}
\needspace{5\baselineskip}
\hypertarget{q221}{\textbf{\color{ink}221.}}  {\color{muted}\scriptsize Y1 \textbullet\ daraja 2 \textbullet\ DB 7a45b6ab}
\par
Bilinear interpolyatsiya nimaga asoslanadi?
\begin{enumerate}[label=\Alph*),leftmargin=9mm]
\item Qo'shni 4 namuna qiymatini aralashtirishga
\item Faqat qatlam nomiga
\item Faqat bitta rang kanaliga
\item Faqat HTML matniga
\end{enumerate}
\vspace{2mm}\hrule height 0.2pt \vspace{4mm}
\needspace{5\baselineskip}
\hypertarget{q222}{\textbf{\color{ink}222.}}  {\color{muted}\scriptsize Y1 \textbullet\ daraja 2 \textbullet\ DB 7c84f2a6}
\par
Qaysi farq to'g'ri?
\begin{enumerate}[label=\Alph*),leftmargin=9mm]
\item Resize faqat vektor faylda ishlaydi
\item Ikkalasi doim bir xil
\item Crop faqat rangni o'zgartiradi, resize faylni o'chiradi
\item Crop tashqi qismlarni kesadi, resize o'lchamni qayta hisoblaydi
\end{enumerate}
\vspace{2mm}\hrule height 0.2pt \vspace{4mm}
\needspace{5\baselineskip}
\hypertarget{q223}{\textbf{\color{ink}223.}}  {\color{muted}\scriptsize Y1 \textbullet\ daraja 2 \textbullet\ DB 7d037d27}
\par
Qaysi usul piksel-artni kattalashtirishda qirralarni ko'proq saqlaydi?
\begin{enumerate}[label=\Alph*),leftmargin=9mm]
\item JPEG save
\item Nearest-neighbor
\item CMYK konversiya
\item Gaussian blur
\end{enumerate}
\vspace{2mm}\hrule height 0.2pt \vspace{4mm}
\needspace{5\baselineskip}
\hypertarget{q224}{\textbf{\color{ink}224.}}  {\color{muted}\scriptsize Y1 \textbullet\ daraja 2 \textbullet\ DB d6f11aaf}
\par
Hajmni qanday birlik bilan to'g'ri ifodalash mumkin?
\begin{enumerate}[label=\Alph*),leftmargin=9mm]
\item 1 bit = 8 bayt
\item 8 bit = 1 bayt
\item 1 KB har doim 1024 B
\item 1 piksel har doim 1 bayt
\end{enumerate}
\vspace{2mm}\hrule height 0.2pt \vspace{4mm}
\needspace{5\baselineskip}
\hypertarget{q225}{\textbf{\color{ink}225.}}  {\color{muted}\scriptsize Y1 \textbullet\ daraja 2 \textbullet\ DB ef30679c}
\par
Maskaning destruktiv o'chirishdan asosiy afzalligi nima?
\begin{enumerate}[label=\Alph*),leftmargin=9mm]
\item Fayl ovoz chiqaradi
\item Har doim 0 bayt joy egallaydi
\item JPEG rangini vektorga aylantiradi
\item Yashirilgan hududni ko'p hollarda qayta ochish mumkin
\end{enumerate}
\vspace{2mm}\hrule height 0.2pt \vspace{4mm}
\needspace{5\baselineskip}
\hypertarget{q226}{\textbf{\color{ink}226.}}  {\color{muted}\scriptsize Y1 \textbullet\ daraja 2 \textbullet\ DB f08fad52}
\par
Crop amali bilan Select amalini qanday farqlaysiz?
\begin{enumerate}[label=\Alph*),leftmargin=9mm]
\item Crop faqat matnga qo'llanadi
\item Ikkalasi bir xil saqlash formati
\item Select hududni belgilaydi, Crop yakuniy kadr chegarasini o'zgartiradi
\item Select JPEG formatini o'zgartiradi
\end{enumerate}
\vspace{2mm}\hrule height 0.2pt \vspace{4mm}
\needspace{5\baselineskip}
\hypertarget{q227}{\textbf{\color{ink}227.}}  {\color{muted}\scriptsize Y2 \textbullet\ daraja 3 \textbullet\ DB 0b403691}
\par
Tasvirga qarash va tasvirni o'zgartirish amallarini farqlang.
\textbf{\footnotesize Chap ustun:}
\begin{enumerate}[label=\arabic*),leftmargin=9mm]
\item Zoom
\item Flip horizontal
\item Rotate 90 deg 
\item Undo
\end{enumerate}
\textbf{\footnotesize O'ng ustun:}
\begin{enumerate}[label=\Alph*),leftmargin=9mm]
\item Tasvirni chap-o'ng oynaviy akslantirish
\item Tasvirni chorak aylana burish
\item Oxirgi bajarilgan tahrirni bekor qilish
\item Ekrandagi ko'rish masshtabini o'zgartirish
\end{enumerate}
\vspace{2mm}\hrule height 0.2pt \vspace{4mm}
\needspace{5\baselineskip}
\hypertarget{q228}{\textbf{\color{ink}228.}}  {\color{muted}\scriptsize Y1 \textbullet\ daraja 3 \textbullet\ DB 2b29d9b2}
\par
500x200 rasm har ikki tomondan 2 marta kattalashtirilsa yangi piksel o'lchamlari?
\begin{enumerate}[label=\Alph*),leftmargin=9mm]
\item 1000x200
\item 1000x400
\item 500x400
\item 2000x800
\end{enumerate}
\vspace{2mm}\hrule height 0.2pt \vspace{4mm}
\needspace{5\baselineskip}
\hypertarget{q229}{\textbf{\color{ink}229.}}  {\color{muted}\scriptsize Y1 \textbullet\ daraja 3 \textbullet\ DB 3bf11d7f}
\par
1200x900 pikselli tasvirning faqat markazdagi 600x300 maydoni kesib olindi. Yangi piksel soni nechta?
\begin{enumerate}[label=\Alph*),leftmargin=9mm]
\item 1080000
\item 90000
\item 180000
\item 540000
\end{enumerate}
\vspace{2mm}\hrule height 0.2pt \vspace{4mm}
\needspace{5\baselineskip}
\hypertarget{q230}{\textbf{\color{ink}230.}}  {\color{muted}\scriptsize Y2 \textbullet\ daraja 3 \textbullet\ DB 44049607}
\par
Qayta namunalash usulini uning xossasi bilan moslang.
\textbf{\footnotesize Chap ustun:}
\begin{enumerate}[label=\arabic*),leftmargin=9mm]
\item Nearest neighbor
\item Bilinear
\item Bicubic
\item Interpolation
\end{enumerate}
\textbf{\footnotesize O'ng ustun:}
\begin{enumerate}[label=\Alph*),leftmargin=9mm]
\item Eng yaqin manba piksel qiymatini olish
\item Yaqin to'rtta pikselni aralashtirish
\item Ko'proq qo'shni qiymatlar yordamida silliqlash
\item Oraliq nuqta qiymatini taxmin qilish jarayoni
\end{enumerate}
\vspace{2mm}\hrule height 0.2pt \vspace{4mm}
\needspace{5\baselineskip}
\hypertarget{q231}{\textbf{\color{ink}231.}}  {\color{muted}\scriptsize Y1 \textbullet\ daraja 3 \textbullet\ DB 448c3a7c}
\par
640x480 rasm har ikki tomondan 2 marta kattalashtirilsa yangi piksel o'lchamlari?
\begin{enumerate}[label=\Alph*),leftmargin=9mm]
\item 2560x1920
\item 1280x960
\item 640x960
\item 1280x480
\end{enumerate}
\vspace{2mm}\hrule height 0.2pt \vspace{4mm}
\needspace{5\baselineskip}
\hypertarget{q232}{\textbf{\color{ink}232.}}  {\color{muted}\scriptsize Y3 \textbullet\ daraja 3 \textbullet\ DB 48372d04}
\par
600x400 rastrni har ikki tomondan 150\% qilish hajmini hisoblash usulini tartiblang.
\begin{enumerate}[label=\Alph*),leftmargin=9mm]
\item Asl W=600 va H=400 piksellarini yozing
\item Yangi W=900 va H=600 ni ko'paytirib 540 000 pikselni toping
\item Asl 240 000 pikselga nisbat 2,25 ekanini tekshiring
\item Har ikkala o'lchamni 1,5 ga ko'paytiring
\end{enumerate}
\vspace{2mm}\hrule height 0.2pt \vspace{4mm}
\needspace{5\baselineskip}
\hypertarget{q233}{\textbf{\color{ink}233.}}  {\color{muted}\scriptsize Y1 \textbullet\ daraja 3 \textbullet\ DB 4f068431}
\par
300x150 rasm har ikki tomondan 3 marta kattalashtirilsa yangi piksel o'lchamlari?
\begin{enumerate}[label=\Alph*),leftmargin=9mm]
\item 900x150
\item 900x450
\item 300x450
\item 2700x1350
\end{enumerate}
\vspace{2mm}\hrule height 0.2pt \vspace{4mm}
\needspace{5\baselineskip}
\hypertarget{q234}{\textbf{\color{ink}234.}}  {\color{muted}\scriptsize Y1 \textbullet\ daraja 3 \textbullet\ DB 5716c754}
\par
120x80 rasm har ikki tomondan 4 marta kattalashtirilsa yangi piksel o'lchamlari?
\begin{enumerate}[label=\Alph*),leftmargin=9mm]
\item 480x80
\item 120x320
\item 1920x1280
\item 480x320
\end{enumerate}
\vspace{2mm}\hrule height 0.2pt \vspace{4mm}
\needspace{5\baselineskip}
\hypertarget{q235}{\textbf{\color{ink}235.}}  {\color{muted}\scriptsize Y1 \textbullet\ daraja 3 \textbullet\ DB 5999cc2e}
\par
200x100 rasm har ikki tomondan 2 marta kattalashtirilsa yangi piksel o'lchamlari?
\begin{enumerate}[label=\Alph*),leftmargin=9mm]
\item 200x200
\item 400x200
\item 400x100
\item 800x400
\end{enumerate}
\vspace{2mm}\hrule height 0.2pt \vspace{4mm}
\needspace{5\baselineskip}
\hypertarget{q236}{\textbf{\color{ink}236.}}  {\color{muted}\scriptsize Y3 \textbullet\ daraja 3 \textbullet\ DB 61640f1c}
\par
Rasmning markaziy qismini olib, keyin kichraytirish tartibini tiklang.
\begin{enumerate}[label=\Alph*),leftmargin=9mm]
\item Select orqali kerakli piksel maydonini belgilang
\item Resize bilan qolgan rasmning o'lchamini kamaytiring
\item Crop bilan kadr tashqarisini olib tashlang
\item Kerakli hudud koordinatalari va chegaralarini aniqlang
\end{enumerate}
\vspace{2mm}\hrule height 0.2pt \vspace{4mm}
\needspace{5\baselineskip}
\hypertarget{q237}{\textbf{\color{ink}237.}}  {\color{muted}\scriptsize Y1 \textbullet\ daraja 3 \textbullet\ DB 75a0e2d1}
\par
Rasm muharririda suratning faqat chap yuqori to'rtburchak qismi qoldirilishi, qolgan tashqi qismi esa olib tashlanishi kerak. Qaysi amalni bajarish lozim?
\begin{enumerate}[label=\Alph*),leftmargin=9mm]
\item Layer group
\item Hue o'zgartirish
\item Crop
\item Canvas kengaytirish
\end{enumerate}
\vspace{2mm}\hrule height 0.2pt \vspace{4mm}
\needspace{5\baselineskip}
\hypertarget{q238}{\textbf{\color{ink}238.}}  {\color{muted}\scriptsize Y2 \textbullet\ daraja 3 \textbullet\ DB 9bff3276}
\par
Rastr tahrirlash amallarini natijasi bilan moslang.
\textbf{\footnotesize Chap ustun:}
\begin{enumerate}[label=\arabic*),leftmargin=9mm]
\item Select
\item Crop
\item Resize
\item Canvas size
\end{enumerate}
\textbf{\footnotesize O'ng ustun:}
\begin{enumerate}[label=\Alph*),leftmargin=9mm]
\item Tasvirning piksel o'lchamlarini o'zgartirish
\item Ish maydoni chegaralarini o'zgartirish
\item Amal bajariladigan piksel sohasini belgilash
\item Tanlangan kadrdan tashqaridagi qismni kesish
\end{enumerate}
\vspace{2mm}\hrule height 0.2pt \vspace{4mm}
\needspace{5\baselineskip}
\hypertarget{q239}{\textbf{\color{ink}239.}}  {\color{muted}\scriptsize Y1 \textbullet\ daraja 3 \textbullet\ DB a23b4753}
\par
Rasmda 8x6 pikselli tasvirning bo'yalgan 4x3 to'rtburchagi crop qilindi. Natijadagi piksel soni?
\begin{enumerate}[label=\Alph*),leftmargin=9mm]
\item 16
\item 48
\item 24
\item 12
\end{enumerate}
\begin{quote}\footnotesize\color{muted}\textbf{Rasm:} 8x6 panjara ichida 4x3 belgilangan crop sohasi\\
\href{https://github.com/sarvar9417/attestatsiya/blob/main/public/graphics-web-v2/crop_grid.svg}{Rasmni onlayn ko'rish}
\end{quote}
\vspace{2mm}\hrule height 0.2pt \vspace{4mm}
\needspace{5\baselineskip}
\hypertarget{q240}{\textbf{\color{ink}240.}}  {\color{muted}\scriptsize Y2 \textbullet\ daraja 3 \textbullet\ DB ab27a424}
\par
Moslashtiring: 1) Crop; 2) Rotate; 3) Flip. Ta'riflar: a) Burish; b) Kadrni kesish; c) Akslantirish.
\textbf{\footnotesize Chap ustun:}
\begin{enumerate}[label=\arabic*),leftmargin=9mm]
\item Crop
\item Rotate
\item Flip
\end{enumerate}
\textbf{\footnotesize O'ng ustun:}
\begin{enumerate}[label=\Alph*),leftmargin=9mm]
\item Burish
\item Kadrni kesish
\item Akslantirish
\end{enumerate}
\vspace{2mm}\hrule height 0.2pt \vspace{4mm}
\needspace{5\baselineskip}
\hypertarget{q241}{\textbf{\color{ink}241.}}  {\color{muted}\scriptsize Y3 \textbullet\ daraja 3 \textbullet\ DB b1be4c34}
\par
Rastr tasvirning siqilmagan piksel hajmini topish tartibini belgilang.
\begin{enumerate}[label=\Alph*),leftmargin=9mm]
\item Bir pikselga to'g'ri keladigan bit chuqurligini aniqlang
\item Tasvir eni va bo'yini pikselda aniqlang
\item Natijani 8 ga bo'lib baytga aylantiring
\item WxHxb orqali jami bitni hisoblang
\end{enumerate}
\vspace{2mm}\hrule height 0.2pt \vspace{4mm}
\needspace{5\baselineskip}
\hypertarget{q242}{\textbf{\color{ink}242.}}  {\color{muted}\scriptsize Y1 \textbullet\ daraja 3 \textbullet\ DB bb575973}
\par
800x600 pikselli 24 bitli siqilmagan RGB tasvirda piksel ma'lumotlari necha bayt?
\begin{enumerate}[label=\Alph*),leftmargin=9mm]
\item 480000
\item 960000
\item 1440000
\item 1920000
\end{enumerate}
\vspace{2mm}\hrule height 0.2pt \vspace{4mm}
\needspace{5\baselineskip}
\hypertarget{q243}{\textbf{\color{ink}243.}}  {\color{muted}\scriptsize Y1 \textbullet\ daraja 3 \textbullet\ DB e8f59670}
\par
800x600 rasm har ikki tomondan 2 marta kattalashtirilsa yangi piksel o'lchamlari?
\begin{enumerate}[label=\Alph*),leftmargin=9mm]
\item 1600x1200
\item 800x1200
\item 1600x600
\item 3200x2400
\end{enumerate}
\vspace{2mm}\hrule height 0.2pt \vspace{4mm}
\needspace{5\baselineskip}
\hypertarget{q244}{\textbf{\color{ink}244.}}  {\color{muted}\scriptsize Y3 \textbullet\ daraja 3 \textbullet\ DB fb9b9063}
\par
Quyidagi amallarni to'g'ri tartibda bajaring: 1) Yangi o'lchamni tanlash; 2) Asl fayl nusxasini saqlash; 3) Resize/Resample amalini bajarish.
\begin{enumerate}[label=\Alph*),leftmargin=9mm]
\item Yangi o'lchamni tanlash
\item Asl fayl nusxasini saqlash
\item Resize/Resample amalini bajarish
\end{enumerate}
\vspace{2mm}\hrule height 0.2pt \vspace{4mm}
\needspace{5\baselineskip}
\hypertarget{q245}{\textbf{\color{ink}245.}}  {\color{muted}\scriptsize Y1 \textbullet\ daraja 4 \textbullet\ DB 0ae8e40d}
\par
Siqilmagan 720x480, 32 bit/piksel rastrning xom hajmi nechta bayt? (Header/padding yo'q.)
\begin{enumerate}[label=\Alph*),leftmargin=9mm]
\item 691200
\item 2764800
\item 345600
\item 1382400
\end{enumerate}
\vspace{2mm}\hrule height 0.2pt \vspace{4mm}
\needspace{5\baselineskip}
\hypertarget{q246}{\textbf{\color{ink}246.}}  {\color{muted}\scriptsize Y1 \textbullet\ daraja 4 \textbullet\ DB 0e537c48}
\par
Siqilmagan 300x200, 32 bit/piksel rastrning xom hajmi nechta bayt? (Header/padding yo'q.)
\begin{enumerate}[label=\Alph*),leftmargin=9mm]
\item 240000
\item 60000
\item 120000
\item 480000
\end{enumerate}
\vspace{2mm}\hrule height 0.2pt \vspace{4mm}
\needspace{5\baselineskip}
\hypertarget{q247}{\textbf{\color{ink}247.}}  {\color{muted}\scriptsize Y1 \textbullet\ daraja 4 \textbullet\ DB 22ec58ee}
\par
Siqilmagan 240x120, 24 bit/piksel rastrning xom hajmi nechta bayt? (Header/padding yo'q.)
\begin{enumerate}[label=\Alph*),leftmargin=9mm]
\item 28800
\item 86400
\item 172800
\item 43200
\end{enumerate}
\vspace{2mm}\hrule height 0.2pt \vspace{4mm}
\needspace{5\baselineskip}
\hypertarget{q248}{\textbf{\color{ink}248.}}  {\color{muted}\scriptsize Y3 \textbullet\ daraja 4 \textbullet\ DB 2ea617a0}
\par
Piksel-artni kattalashtirganda o'ziga xos kataklarni saqlash tartibini tiklang.
\begin{enumerate}[label=\Alph*),leftmargin=9mm]
\item Asl kichik pikselli rasmni tanlang
\item Resample uchun nearest-neighbor usulini tanlang
\item Chiqish o'lchamini belgilang
\item Kattalashgan katak qirralarini tekshiring
\end{enumerate}
\vspace{2mm}\hrule height 0.2pt \vspace{4mm}
\needspace{5\baselineskip}
\hypertarget{q249}{\textbf{\color{ink}249.}}  {\color{muted}\scriptsize Y2 \textbullet\ daraja 4 \textbullet\ DB 33f6d2af}
\par
Vektor konturlariga oid amallarni moslang.
\textbf{\footnotesize Chap ustun:}
\begin{enumerate}[label=\arabic*),leftmargin=9mm]
\item Union
\item Intersection
\item Difference
\item Group
\end{enumerate}
\textbf{\footnotesize O'ng ustun:}
\begin{enumerate}[label=\Alph*),leftmargin=9mm]
\item Ikki shaklning birlashgan hududini olish
\item Birinchi shakldan ikkinchisini ayirish
\item Bir nechta obyektni bir guruhda boshqarish
\item Faqat umumiy hududini qoldirish
\end{enumerate}
\vspace{2mm}\hrule height 0.2pt \vspace{4mm}
\needspace{5\baselineskip}
\hypertarget{q250}{\textbf{\color{ink}250.}}  {\color{muted}\scriptsize Y1 \textbullet\ daraja 4 \textbullet\ DB 3a992997}
\par
Siqilmagan 400x300, 24 bit/piksel rastrning xom hajmi nechta bayt? (Header/padding yo'q.)
\begin{enumerate}[label=\Alph*),leftmargin=9mm]
\item 360000
\item 180000
\item 720000
\item 120000
\end{enumerate}
\vspace{2mm}\hrule height 0.2pt \vspace{4mm}
\needspace{5\baselineskip}
\hypertarget{q251}{\textbf{\color{ink}251.}}  {\color{muted}\scriptsize Y2 \textbullet\ daraja 4 \textbullet\ DB 4a9be737}
\par
Rang chuqurligi bilan sig'imni aniq moslang; siqish va metama'lumot yo'q.
\textbf{\footnotesize Chap ustun:}
\begin{enumerate}[label=\arabic*),leftmargin=9mm]
\item 1 bit/piksel
\item 4 bit/piksel
\item 8 bit/piksel
\item 24 bit/piksel RGB
\end{enumerate}
\textbf{\footnotesize O'ng ustun:}
\begin{enumerate}[label=\Alph*),leftmargin=9mm]
\item Har piksel 16 indeksdan birini olishi mumkin
\item Har piksel 2 qiymatdan birini olishi mumkin
\item Har piksel xom ma'lumotda 3 bayt joy oladi
\item Har piksel 256 indeksdan birini olishi mumkin
\end{enumerate}
\vspace{2mm}\hrule height 0.2pt \vspace{4mm}
\needspace{5\baselineskip}
\hypertarget{q252}{\textbf{\color{ink}252.}}  {\color{muted}\scriptsize Y2 \textbullet\ daraja 4 \textbullet\ DB 4c210c9b}
\par
Moslashtiring: 1) Resampling; 2) Zoom; 3) Rasterize. Ta'riflar: a) Vektordan rastrga o'tish; b) Pikselni qayta hisoblash; c) Ko'rish masshtabi.
\textbf{\footnotesize Chap ustun:}
\begin{enumerate}[label=\arabic*),leftmargin=9mm]
\item Resampling
\item Zoom
\item Rasterize
\end{enumerate}
\textbf{\footnotesize O'ng ustun:}
\begin{enumerate}[label=\Alph*),leftmargin=9mm]
\item Vektordan rastrga o'tish
\item Pikselni qayta hisoblash
\item Ko'rish masshtabi
\end{enumerate}
\vspace{2mm}\hrule height 0.2pt \vspace{4mm}
\needspace{5\baselineskip}
\hypertarget{q253}{\textbf{\color{ink}253.}}  {\color{muted}\scriptsize Y1 \textbullet\ daraja 4 \textbullet\ DB 59b7aa95}
\par
Ikki doiraning faqat ustma-ust qismi kerak. Qaysi amal?
\begin{enumerate}[label=\Alph*),leftmargin=9mm]
\item Rasterization
\item Union
\item Difference
\item Intersection
\end{enumerate}
\vspace{2mm}\hrule height 0.2pt \vspace{4mm}
\needspace{5\baselineskip}
\hypertarget{q254}{\textbf{\color{ink}254.}}  {\color{muted}\scriptsize Y1 \textbullet\ daraja 4 \textbullet\ DB 704a80df}
\par
1200 piksel enli tasvir 300 PPI bilan chop etilsa nominal bosma eni qancha dyuym?
\begin{enumerate}[label=\Alph*),leftmargin=9mm]
\item 6
\item 3
\item 12
\item 4
\end{enumerate}
\vspace{2mm}\hrule height 0.2pt \vspace{4mm}
\needspace{5\baselineskip}
\hypertarget{q255}{\textbf{\color{ink}255.}}  {\color{muted}\scriptsize Y3 \textbullet\ daraja 4 \textbullet\ DB 81a5a3e7}
\par
Vektor figurani 90 deg  soat miliga teskari aylantirishda koordinatani toping.
\begin{enumerate}[label=\Alph*),leftmargin=9mm]
\item Aylanish markazi koordinata boshi ekanini tekshiring
\item 90 deg  CCW qoidasini qo'llang: (x,y)=\textgreater{}(-y,x)
\item Natija uchun (2,3)=\textgreater{}(-3,2) ekanini tekshiring
\item Boshlang'ich nuqtani (x,y) deb belgilang
\end{enumerate}
\vspace{2mm}\hrule height 0.2pt \vspace{4mm}
\needspace{5\baselineskip}
\hypertarget{q256}{\textbf{\color{ink}256.}}  {\color{muted}\scriptsize Y1 \textbullet\ daraja 4 \textbullet\ DB 83e5faf3}
\par
Siqilmagan 500x500, 8 bit/piksel rastrning xom hajmi nechta bayt? (Header/padding yo'q.)
\begin{enumerate}[label=\Alph*),leftmargin=9mm]
\item 750000
\item 500000
\item 125000
\item 250000
\end{enumerate}
\vspace{2mm}\hrule height 0.2pt \vspace{4mm}
\needspace{5\baselineskip}
\hypertarget{q257}{\textbf{\color{ink}257.}}  {\color{muted}\scriptsize Y1 \textbullet\ daraja 4 \textbullet\ DB a5618ade}
\par
1024x768 24-bitli rasmning xom piksel hajmini JPEG fayli hajmi bilan tenglashtirish mumkinmi?
\begin{enumerate}[label=\Alph*),leftmargin=9mm]
\item Ha, 24 bit xamisha 24 bayt
\item Ha, JPEG doim siqilmagan format
\item Yo'q, JPEG siqish va metama'lumotlar natijani o'zgartiradi
\item Yo'q, chunki JPEG vektor formati
\end{enumerate}
\vspace{2mm}\hrule height 0.2pt \vspace{4mm}
\needspace{5\baselineskip}
\hypertarget{q258}{\textbf{\color{ink}258.}}  {\color{muted}\scriptsize Y1 \textbullet\ daraja 4 \textbullet\ DB add66f90}
\par
Siqilmagan 1280x720, 24 bit/piksel rastrning xom hajmi nechta bayt? (Header/padding yo'q.)
\begin{enumerate}[label=\Alph*),leftmargin=9mm]
\item 5529600
\item 921600
\item 1382400
\item 2764800
\end{enumerate}
\vspace{2mm}\hrule height 0.2pt \vspace{4mm}
\needspace{5\baselineskip}
\hypertarget{q259}{\textbf{\color{ink}259.}}  {\color{muted}\scriptsize Y3 \textbullet\ daraja 4 \textbullet\ DB b203972c}
\par
Ikki vektor shaklning faqat kesishgan joyini olish tartibini belgilang.
\begin{enumerate}[label=\Alph*),leftmargin=9mm]
\item Kesishishi kerak bo'lgan ikkita yopiq shaklni yarating
\item Intersection amalini bajarib qolgan umumiy hududni tekshiring
\item Ikkala shaklni tanlang
\item Shakllarni kerakli o'zaro joylashuvga keltiring
\end{enumerate}
\vspace{2mm}\hrule height 0.2pt \vspace{4mm}
\needspace{5\baselineskip}
\hypertarget{q260}{\textbf{\color{ink}260.}}  {\color{muted}\scriptsize Y2 \textbullet\ daraja 4 \textbullet\ DB db86b682}
\par
Vektor va rastr obyekt amallarini vazifasiga bog'lang.
\textbf{\footnotesize Chap ustun:}
\begin{enumerate}[label=\arabic*),leftmargin=9mm]
\item Translation
\item Scale
\item Bring to front
\item Rasterization
\end{enumerate}
\textbf{\footnotesize O'ng ustun:}
\begin{enumerate}[label=\Alph*),leftmargin=9mm]
\item Obyektni ustma-ust tartibda yuqoriga chiqarish
\item Geometrik shakldan piksel tasvir hosil qilish
\item Koordinatalarni koeffitsiyentga ko'paytirish
\item Nuqtaga bir xil koordinata siljishini qo'shish
\end{enumerate}
\vspace{2mm}\hrule height 0.2pt \vspace{4mm}
\needspace{5\baselineskip}
\hypertarget{q261}{\textbf{\color{ink}261.}}  {\color{muted}\scriptsize Y3 \textbullet\ daraja 4 \textbullet\ DB e65bf1c4}
\par
Quyidagi amallarni to'g'ri tartibda bajaring: 1) Xom hajmni baytga hisoblash; 2) WxH piksel sonini aniqlash; 3) Bit chuqurligini o'qish.
\begin{enumerate}[label=\Alph*),leftmargin=9mm]
\item Xom hajmni baytga hisoblash
\item WxH piksel sonini aniqlash
\item Bit chuqurligini o'qish
\end{enumerate}
\vspace{2mm}\hrule height 0.2pt \vspace{4mm}
\needspace{5\baselineskip}
\hypertarget{q262}{\textbf{\color{ink}262.}}  {\color{muted}\scriptsize Y1 \textbullet\ daraja 4 \textbullet\ DB f6b164e2}
\par
Siqilmagan 640x200, 8 bit/piksel rastrning xom hajmi nechta bayt? (Header/padding yo'q.)
\begin{enumerate}[label=\Alph*),leftmargin=9mm]
\item 64000
\item 256000
\item 384000
\item 128000
\end{enumerate}
\vspace{2mm}\hrule height 0.2pt \vspace{4mm}
\needspace{5\baselineskip}
\hypertarget{q263}{\textbf{\color{ink}263.}}  {\color{muted}\scriptsize Y1 \textbullet\ daraja 5 \textbullet\ DB 0a7964e3}
\par
600x400 tasvirda crop natijasida 300x200 qismi qoldi. Piksel miqdori qanday o'zgaradi?
\begin{enumerate}[label=\Alph*),leftmargin=9mm]
\item 4 baravar kamayadi
\item 2 baravar kamayadi
\item 4 baravar oshadi
\item O'zgarmaydi
\end{enumerate}
\vspace{2mm}\hrule height 0.2pt \vspace{4mm}
\needspace{5\baselineskip}
\hypertarget{q264}{\textbf{\color{ink}264.}}  {\color{muted}\scriptsize Y1 \textbullet\ daraja 5 \textbullet\ DB 1a87cb16}
\par
Rastr tasvir dastlab 12x8 piksel edi. Crop amali bilan 6x4 pikselli qismi qoldirildi. Tasvirda qolgan piksel soni dastlabkiga nisbatan necha foizga kamaydi?
\begin{enumerate}[label=\Alph*),leftmargin=9mm]
\item 75\%
\item 50\%
\item 25\%
\item 150\%
\end{enumerate}
\vspace{2mm}\hrule height 0.2pt \vspace{4mm}
\needspace{5\baselineskip}
\hypertarget{q265}{\textbf{\color{ink}265.}}  {\color{muted}\scriptsize Y2 \textbullet\ daraja 5 \textbullet\ DB 3ff3ce49}
\par
Shaffof tahrirlash atamalarini moslang.
\textbf{\footnotesize Chap ustun:}
\begin{enumerate}[label=\arabic*),leftmargin=9mm]
\item Layer mask
\item Eraser
\item Clipping mask
\item Alpha channel
\end{enumerate}
\textbf{\footnotesize O'ng ustun:}
\begin{enumerate}[label=\Alph*),leftmargin=9mm]
\item Ko'rinishni boshqa obyekt shakli bilan cheklash
\item Pikselning o'zini o'chiruvchi tahrir
\item Manba piksellarni o'chirmay ko'rinishni boshqarish
\item Piksel yoki qatlam shaffofligi qiymatlari
\end{enumerate}
\vspace{2mm}\hrule height 0.2pt \vspace{4mm}
\needspace{5\baselineskip}
\hypertarget{q266}{\textbf{\color{ink}266.}}  {\color{muted}\scriptsize Y1 \textbullet\ daraja 5 \textbullet\ DB 4c1fb897}
\par
400x300 tasvirning ikkala piksel o'lchami 2 martadan oshirildi (24 bit/piksel). Xom massiv hajmi qanday o'zgaradi?
\begin{enumerate}[label=\Alph*),leftmargin=9mm]
\item 8 marta oshadi
\item 2 marta oshadi
\item O'zgarmaydi
\item 4 marta oshadi
\end{enumerate}
\vspace{2mm}\hrule height 0.2pt \vspace{4mm}
\needspace{5\baselineskip}
\hypertarget{q267}{\textbf{\color{ink}267.}}  {\color{muted}\scriptsize Y1 \textbullet\ daraja 5 \textbullet\ DB 60c0eb99}
\par
900 piksel enli rasm 150 PPIda bosilganda nominal eni nechta dyuym?
\begin{enumerate}[label=\Alph*),leftmargin=9mm]
\item 6
\item 12
\item 5
\item 7
\end{enumerate}
\vspace{2mm}\hrule height 0.2pt \vspace{4mm}
\needspace{5\baselineskip}
\hypertarget{q268}{\textbf{\color{ink}268.}}  {\color{muted}\scriptsize Y1 \textbullet\ daraja 5 \textbullet\ DB 7c2c2016}
\par
1600 piksel enli rasm 200 PPIda bosilganda nominal eni nechta dyuym?
\begin{enumerate}[label=\Alph*),leftmargin=9mm]
\item 8
\item 16
\item 7
\item 9
\end{enumerate}
\vspace{2mm}\hrule height 0.2pt \vspace{4mm}
\needspace{5\baselineskip}
\hypertarget{q269}{\textbf{\color{ink}269.}}  {\color{muted}\scriptsize Y3 \textbullet\ daraja 5 \textbullet\ DB 9b1e8953}
\par
Rastrdagi obyektni yo'qotmasdan vaqtincha berkitish tartibini tiklang.
\begin{enumerate}[label=\Alph*),leftmargin=9mm]
\item Maskadagi tegishli qismni yashiring
\item Piksellarni o'chirish o'rniga layer mask yarating
\item Berkitiladigan obyekt joylashgan qatlamni tanlang
\item Kerak bo'lganda maskani o'zgartirib obyektni qayta ko'rsating
\end{enumerate}
\vspace{2mm}\hrule height 0.2pt \vspace{4mm}
\needspace{5\baselineskip}
\hypertarget{q270}{\textbf{\color{ink}270.}}  {\color{muted}\scriptsize Y2 \textbullet\ daraja 5 \textbullet\ DB 9c357a44}
\par
Shartli rastr misollarini natijasi bilan moslang; xom ma'lumot, siqishsiz.
\textbf{\footnotesize Chap ustun:}
\begin{enumerate}[label=\arabic*),leftmargin=9mm]
\item 100x100, RGB24
\item 100x100, RGBA32
\item 200x100, RGB24
\item 200x100, RGBA32
\end{enumerate}
\textbf{\footnotesize O'ng ustun:}
\begin{enumerate}[label=\Alph*),leftmargin=9mm]
\item 40 000 bayt
\item 80 000 bayt
\item 30 000 bayt
\item 60 000 bayt
\end{enumerate}
\vspace{2mm}\hrule height 0.2pt \vspace{4mm}
\needspace{5\baselineskip}
\hypertarget{q271}{\textbf{\color{ink}271.}}  {\color{muted}\scriptsize Y1 \textbullet\ daraja 5 \textbullet\ DB a5609060}
\par
48x32 pikselli ikonka nearest-neighbor bilan 4 marta kattalashtirildi. Yangi piksel o'lchami?
\begin{enumerate}[label=\Alph*),leftmargin=9mm]
\item 192x128
\item 48x128
\item 52x36
\item 96x64
\end{enumerate}
\vspace{2mm}\hrule height 0.2pt \vspace{4mm}
\needspace{5\baselineskip}
\hypertarget{q272}{\textbf{\color{ink}272.}}  {\color{muted}\scriptsize Y1 \textbullet\ daraja 5 \textbullet\ DB a6a5f417}
\par
Tasvirda matnli konturlarni silliqlashtirish uchun nearest-neighbor va bilinear orasida qaysi farq muhim?
\begin{enumerate}[label=\Alph*),leftmargin=9mm]
\item Nearest-neighbor doim to'rt pikselni o'rtacha qiladi
\item Bilinear har doim vektor modelni qaytaradi
\item Bilinear qo'shni piksel ranglarini interpolyatsiya qiladi, nearest-neighbor bitta yaqin namunani oladi
\item Ikkisi JPEG siqish turi
\end{enumerate}
\vspace{2mm}\hrule height 0.2pt \vspace{4mm}
\needspace{5\baselineskip}
\hypertarget{q273}{\textbf{\color{ink}273.}}  {\color{muted}\scriptsize Y1 \textbullet\ daraja 5 \textbullet\ DB bf2d5514}
\par
Piksel soni o'zgarmagan holda faqat PPI 150 dan 300 ga o'zgarsa, asosiy ta'sir nima?
\begin{enumerate}[label=\Alph*),leftmargin=9mm]
\item Rasm 4 baravar ko'p rangli bo'ladi
\item Nominal chop o'lchami ikki baravar kichrayadi
\item Piksel soni ikki baravar ko'payadi
\item JPEG avtomatik SVGga aylanadi
\end{enumerate}
\vspace{2mm}\hrule height 0.2pt \vspace{4mm}
\needspace{5\baselineskip}
\hypertarget{q274}{\textbf{\color{ink}274.}}  {\color{muted}\scriptsize Y1 \textbullet\ daraja 5 \textbullet\ DB ca2cff0e}
\par
Rasmda 8 ustun va 6 qator bor. Har piksel 3 bayt RGB bo'lsa, xom ma'lumot jami necha bayt?
\begin{enumerate}[label=\Alph*),leftmargin=9mm]
\item 24
\item 48
\item 144
\item 384
\end{enumerate}
\begin{quote}\footnotesize\color{muted}\textbf{Rasm:} 8x6 pikselli rastr panjara\\
\href{https://github.com/sarvar9417/attestatsiya/blob/main/public/graphics-web-v2/raster_grid.svg}{Rasmni onlayn ko'rish}
\end{quote}
\vspace{2mm}\hrule height 0.2pt \vspace{4mm}
\needspace{5\baselineskip}
\hypertarget{q275}{\textbf{\color{ink}275.}}  {\color{muted}\scriptsize Y1 \textbullet\ daraja 5 \textbullet\ DB cdaf625a}
\par
1200 piksel enli rasm 300 PPIda bosilganda nominal eni nechta dyuym?
\begin{enumerate}[label=\Alph*),leftmargin=9mm]
\item 3
\item 8
\item 5
\item 4
\end{enumerate}
\vspace{2mm}\hrule height 0.2pt \vspace{4mm}
\needspace{5\baselineskip}
\hypertarget{q276}{\textbf{\color{ink}276.}}  {\color{muted}\scriptsize Y1 \textbullet\ daraja 5 \textbullet\ DB e884f61d}
\par
Vektor logotipni avval SVG loyiha sifatida, so'ng 200x200 PNGda saqladingiz. Qaysi fikr aniq?
\begin{enumerate}[label=\Alph*),leftmargin=9mm]
\item SVG faqat audio ma'lumotdir
\item SVGda geometrik obyektlar qolishi, PNGda esa piksellar saqlanishi mumkin
\item PNGda piksel o'lchami tushunchasi yo'q
\item PNG doim vektor manbadan ko'proq tugun saqlaydi
\end{enumerate}
\vspace{2mm}\hrule height 0.2pt \vspace{4mm}
\needspace{5\baselineskip}
\hypertarget{q277}{\textbf{\color{ink}277.}}  {\color{muted}\scriptsize Y1 \textbullet\ daraja 5 \textbullet\ DB f08749f5}
\par
Vektor grafikasida Union, Subtract va Intersect tushunchalarining umumiy vazifasi nima?
\begin{enumerate}[label=\Alph*),leftmargin=9mm]
\item Kontur shakllarini geometrik tarzda birlashtirish, ayirish yoki kesishish
\item Tasvirga faqat audio qo'shish
\item Bir xil fayl nomini qayta yozish
\item Faqat PDF sahifasini ochish
\end{enumerate}
\vspace{2mm}\hrule height 0.2pt \vspace{4mm}
\needspace{5\baselineskip}
\hypertarget{q278}{\textbf{\color{ink}278.}}  {\color{muted}\scriptsize Y3 \textbullet\ daraja 5 \textbullet\ DB f7bd458c}
\par
900 pikselli rasmni 300 PPI bilan chop etish enini aniqlashni tartiblang.
\begin{enumerate}[label=\Alph*),leftmargin=9mm]
\item Zarur bo'lsa 3x2,54=7,62 sm ga aylantiring
\item Tasvirning gorizontal piksel sonini yozing: 900
\item Bosma zichligini oling: 300 piksel/dyuym
\item 900 ni 300 ga bo'lib 3 dyuymni toping
\end{enumerate}
\vspace{2mm}\hrule height 0.2pt \vspace{4mm}
\subsubsection{MS Paint va Photoshop - nazariya V2}
\needspace{5\baselineskip}
\hypertarget{q279}{\textbf{\color{ink}279.}}  {\color{muted}\scriptsize Y1 \textbullet\ daraja 1 \textbullet\ DB 14082104}
\par
Microsoft Paintning Curved line uskunasi 5-sinf Cambridge darsligida qanday boshlanadi?
\begin{enumerate}[label=\Alph*),leftmargin=9mm]
\item Avval Photoshop PSD fayli eksport qilinadi
\item Avval to'g'ri segment chizib, so'ng uning egriligi sozlanadi
\item Faqat jadval katagi chiziladi
\item Faqat RGB kanali o'chiriladi
\end{enumerate}
\vspace{2mm}\hrule height 0.2pt \vspace{4mm}
\needspace{5\baselineskip}
\hypertarget{q280}{\textbf{\color{ink}280.}}  {\color{muted}\scriptsize Y1 \textbullet\ daraja 1 \textbullet\ DB 22f1df8c}
\par
Brush uskunasi nima qiladi?
\begin{enumerate}[label=\Alph*),leftmargin=9mm]
\item Tasvirni videoga aylantirish
\item CMYKni shifrlash
\item HTML kodni yozish
\item Cho'tka izi bilan bo'yash
\end{enumerate}
\vspace{2mm}\hrule height 0.2pt \vspace{4mm}
\needspace{5\baselineskip}
\hypertarget{q281}{\textbf{\color{ink}281.}}  {\color{muted}\scriptsize Y1 \textbullet\ daraja 1 \textbullet\ DB 3021b5fa}
\par
Marquee uskunasining asosiy vazifasi nima?
\begin{enumerate}[label=\Alph*),leftmargin=9mm]
\item Rasmni 3Dga o'tkazish
\item Qatlamlarni guruhlash
\item Loyihani eksport qilish
\item Geometrik shaklli tanlov yaratish
\end{enumerate}
\vspace{2mm}\hrule height 0.2pt \vspace{4mm}
\needspace{5\baselineskip}
\hypertarget{q282}{\textbf{\color{ink}282.}}  {\color{muted}\scriptsize Y1 \textbullet\ daraja 1 \textbullet\ DB 45ed22e7}
\par
Eyedropper uskunasi nima qiladi?
\begin{enumerate}[label=\Alph*),leftmargin=9mm]
\item Qatlamni qulfdan chiqaradi
\item Tasvirdan rang namunasini oladi
\item Har doim fonni o'chiradi
\item Jadvalni ranglaydi
\end{enumerate}
\vspace{2mm}\hrule height 0.2pt \vspace{4mm}
\needspace{5\baselineskip}
\hypertarget{q283}{\textbf{\color{ink}283.}}  {\color{muted}\scriptsize Y1 \textbullet\ daraja 1 \textbullet\ DB 7a3deec3}
\par
MS Paintning asosiy vazifasi nima?
\begin{enumerate}[label=\Alph*),leftmargin=9mm]
\item Ma'lumotlar bazasi boshqarish
\item SQL so'rov yozish
\item Audio miks yozish
\item Rastr rasmni sodda yaratish va tahrirlash
\end{enumerate}
\vspace{2mm}\hrule height 0.2pt \vspace{4mm}
\needspace{5\baselineskip}
\hypertarget{q284}{\textbf{\color{ink}284.}}  {\color{muted}\scriptsize Y1 \textbullet\ daraja 1 \textbullet\ DB 7f62a2ed}
\par
MS Paintdagi Add to Custom Colours buyrug'ining vazifasi nima?
\begin{enumerate}[label=\Alph*),leftmargin=9mm]
\item Qatlamni xotiradan butunlay o'chirish
\item HTML jadval yaratish
\item Photoshop filtrini yoqish
\item Foydalanuvchi yaratgan rangni maxsus palitraga qo'shish
\end{enumerate}
\vspace{2mm}\hrule height 0.2pt \vspace{4mm}
\needspace{5\baselineskip}
\hypertarget{q285}{\textbf{\color{ink}285.}}  {\color{muted}\scriptsize Y1 \textbullet\ daraja 1 \textbullet\ DB 9d3c3e18}
\par
Lasso uskunasi nima qiladi?
\begin{enumerate}[label=\Alph*),leftmargin=9mm]
\item Erkin konturli soha tanlaydi
\item Qatlam shaffofligini 0 qiladi
\item JPEG sifatini pasaytiradi
\item Faqat matn kiritadi
\end{enumerate}
\vspace{2mm}\hrule height 0.2pt \vspace{4mm}
\needspace{5\baselineskip}
\hypertarget{q286}{\textbf{\color{ink}286.}}  {\color{muted}\scriptsize Y1 \textbullet\ daraja 1 \textbullet\ DB b31e73e6}
\par
PSD formati nima uchun?
\begin{enumerate}[label=\Alph*),leftmargin=9mm]
\item Photoshop tahrir loyihasini qatlamlari bilan saqlash
\item Faqat audio yozish
\item Faqat kodni kompilyatsiya qilish
\item Faqat animatsion GIF saqlash
\end{enumerate}
\vspace{2mm}\hrule height 0.2pt \vspace{4mm}
\needspace{5\baselineskip}
\hypertarget{q287}{\textbf{\color{ink}287.}}  {\color{muted}\scriptsize Y1 \textbullet\ daraja 1 \textbullet\ DB bab91f3a}
\par
Photoshopda Layers paneli nima uchun?
\begin{enumerate}[label=\Alph*),leftmargin=9mm]
\item HTML formani yuborish
\item Kamerani ulash
\item Fayl kengaytmasini topish
\item Qatlamlarni boshqarish
\end{enumerate}
\vspace{2mm}\hrule height 0.2pt \vspace{4mm}
\needspace{5\baselineskip}
\hypertarget{q288}{\textbf{\color{ink}288.}}  {\color{muted}\scriptsize Y1 \textbullet\ daraja 1 \textbullet\ DB ee588489}
\par
Clone Stamp nima?
\begin{enumerate}[label=\Alph*),leftmargin=9mm]
\item HTML tegni nusxalaydi
\item Printerga ulanadi
\item Tasvirdan namuna olib boshqa hududga ko'chirib bo'yaydi
\item Qatlam nomini o'zgartiradi
\end{enumerate}
\vspace{2mm}\hrule height 0.2pt \vspace{4mm}
\needspace{5\baselineskip}
\hypertarget{q289}{\textbf{\color{ink}289.}}  {\color{muted}\scriptsize Y1 \textbullet\ daraja 1 \textbullet\ DB f6dd2e65}
\par
Opacity xossasi nima?
\begin{enumerate}[label=\Alph*),leftmargin=9mm]
\item Faylni aylantiradi
\item Qatlam umumiy shaffofligini nazorat qiladi
\item Tasvir kengligini aniqlaydi
\item RGB kanallari sonini oshiradi
\end{enumerate}
\vspace{2mm}\hrule height 0.2pt \vspace{4mm}
\needspace{5\baselineskip}
\hypertarget{q290}{\textbf{\color{ink}290.}}  {\color{muted}\scriptsize Y1 \textbullet\ daraja 1 \textbullet\ DB f75d3724}
\par
Layer mask nima?
\begin{enumerate}[label=\Alph*),leftmargin=9mm]
\item Sahifa sarlavhasi
\item Jadvaldagi ustun
\item Bitta JPEG filtri
\item Qatlamning ko'rinuvchi qismlarini boshqaruvchi niqob
\end{enumerate}
\vspace{2mm}\hrule height 0.2pt \vspace{4mm}
\needspace{5\baselineskip}
\hypertarget{q291}{\textbf{\color{ink}291.}}  {\color{muted}\scriptsize Y1 \textbullet\ daraja 2 \textbullet\ DB 0823609c}
\par
Photoshop'da suratning rang va yorug'ligini asl piksel qiymatlarini doimiy o'zgartirmasdan sozlash, keyinchalik esa tahrir parametrlari va niqobini qayta moslash zarur. Qaysi yondashuv eng mos?
\begin{enumerate}[label=\Alph*),leftmargin=9mm]
\item Rang va tonni faqat asl piksel qatlamida doimiy o'zgartirish
\item Rang/ton tahririni alohida qatlam orqali nazorat qilish
\item Tasvirning barcha qatlamlarini bitta qatlamga birlashtirib, shu yerda sozlash
\item Rangni bir marta sozlab, keyinchalik uni qayta tahrirlash imkonini yopish
\end{enumerate}
\vspace{2mm}\hrule height 0.2pt \vspace{4mm}
\needspace{5\baselineskip}
\hypertarget{q292}{\textbf{\color{ink}292.}}  {\color{muted}\scriptsize Y1 \textbullet\ daraja 2 \textbullet\ DB 320ad9e7}
\par
Photoshopda matn qo'shishga mos asbob qaysi?
\begin{enumerate}[label=\Alph*),leftmargin=9mm]
\item Type Tool
\item Crop Tool
\item Magic Wand
\item Clone Stamp
\end{enumerate}
\vspace{2mm}\hrule height 0.2pt \vspace{4mm}
\needspace{5\baselineskip}
\hypertarget{q293}{\textbf{\color{ink}293.}}  {\color{muted}\scriptsize Y1 \textbullet\ daraja 2 \textbullet\ DB 4134182f}
\par
JPEGga eksport bilan PSDni saqlashning muhim farqi?
\begin{enumerate}[label=\Alph*),leftmargin=9mm]
\item PSD hech qanday rastr saqlamaydi
\item JPEG tekislangan natija bo'lishi mumkin, PSD qatlamni saqlaydi
\item JPEG qatlamdan ko'proq imkoniyat beradi
\item JPEG faqat vektor
\end{enumerate}
\vspace{2mm}\hrule height 0.2pt \vspace{4mm}
\needspace{5\baselineskip}
\hypertarget{q294}{\textbf{\color{ink}294.}}  {\color{muted}\scriptsize Y1 \textbullet\ daraja 2 \textbullet\ DB 460c5246}
\par
Qaysi fikr Paint va Paint.NETni to'g'ri farqlaydi?
\begin{enumerate}[label=\Alph*),leftmargin=9mm]
\item Paint.NET MS Paintning shunchaki boshqa nomi
\item Paint faqat vektor dasturi
\item Paint.NET faqat Word qismi
\item Ikkalasi alohida muharrir va funksiyalari bir xil deb bo'lmaydi
\end{enumerate}
\vspace{2mm}\hrule height 0.2pt \vspace{4mm}
\needspace{5\baselineskip}
\hypertarget{q295}{\textbf{\color{ink}295.}}  {\color{muted}\scriptsize Y1 \textbullet\ daraja 2 \textbullet\ DB 8a88536c}
\par
Paintdagi Fill uskunasi qaysi vaziyatda xatoga olib kelishi mumkin?
\begin{enumerate}[label=\Alph*),leftmargin=9mm]
\item Rasm o'lchami juft bo'lganda
\item PNG tanlanganda
\item Fayl nomi kichik bo'lganda
\item Kontur yopilmaganda rang tashqariga oqib ketishi
\end{enumerate}
\vspace{2mm}\hrule height 0.2pt \vspace{4mm}
\needspace{5\baselineskip}
\hypertarget{q296}{\textbf{\color{ink}296.}}  {\color{muted}\scriptsize Y1 \textbullet\ daraja 2 \textbullet\ DB 9f7402ce}
\par
Crop bilan Move farqi?
\begin{enumerate}[label=\Alph*),leftmargin=9mm]
\item Ikkisi shriftdir
\item Crop kadrni kesadi, Move obyektni siljitadi
\item Crop faqat filtr
\item Move faqat fayl eksporti
\end{enumerate}
\vspace{2mm}\hrule height 0.2pt \vspace{4mm}
\needspace{5\baselineskip}
\hypertarget{q297}{\textbf{\color{ink}297.}}  {\color{muted}\scriptsize Y1 \textbullet\ daraja 2 \textbullet\ DB c3f10aa5}
\par
Paintda Ctrl+Z odatda nimaga xizmat qiladi?
\begin{enumerate}[label=\Alph*),leftmargin=9mm]
\item Oxirgi amalni bekor qiladi
\item Faylni ochadi
\item Tasvirni 90 daraja buradi
\item Rangni o'chiradi
\end{enumerate}
\vspace{2mm}\hrule height 0.2pt \vspace{4mm}
\needspace{5\baselineskip}
\hypertarget{q298}{\textbf{\color{ink}298.}}  {\color{muted}\scriptsize Y1 \textbullet\ daraja 2 \textbullet\ DB cec7a4dc}
\par
Layer mask va Eraser farqi nimada?
\begin{enumerate}[label=\Alph*),leftmargin=9mm]
\item Mask ko'rinishni yashirishi, Eraser piksellarga o'zgartirish kiritishi mumkin
\item Har doim bir xil
\item Mask faqat HTMLda
\item Eraser rang profilini o'zgartiradi
\end{enumerate}
\vspace{2mm}\hrule height 0.2pt \vspace{4mm}
\needspace{5\baselineskip}
\hypertarget{q299}{\textbf{\color{ink}299.}}  {\color{muted}\scriptsize Y1 \textbullet\ daraja 2 \textbullet\ DB d20b8459}
\par
Qatlamning eye belgisi o'chirildi. Nima yuz beradi?
\begin{enumerate}[label=\Alph*),leftmargin=9mm]
\item Pastdagi barcha qatlamlar yo'qoladi
\item Qatlam ko'rinmaydi, ma'lumoti saqlanishi mumkin
\item Fayl tarmoqdan o'chiriladi
\item Qatlam doim o'chiriladi
\end{enumerate}
\vspace{2mm}\hrule height 0.2pt \vspace{4mm}
\needspace{5\baselineskip}
\hypertarget{q300}{\textbf{\color{ink}300.}}  {\color{muted}\scriptsize Y1 \textbullet\ daraja 2 \textbullet\ DB df2bf2db}
\par
Qaysi format Photoshop qatlamlarini odatda saqlaydi?
\begin{enumerate}[label=\Alph*),leftmargin=9mm]
\item GIF
\item BMP
\item PSD
\item JPEG
\end{enumerate}
\vspace{2mm}\hrule height 0.2pt \vspace{4mm}
\needspace{5\baselineskip}
\hypertarget{q301}{\textbf{\color{ink}301.}}  {\color{muted}\scriptsize Y1 \textbullet\ daraja 2 \textbullet\ DB f24c34af}
\par
Paintdagi Rectangular select va Free-form selection qachon farq qiladi?
\begin{enumerate}[label=\Alph*),leftmargin=9mm]
\item Free-form faqat to'g'ri chiziq chizadi
\item Birinchisi to'rtburchak, ikkinchisi erkin shaklli sohaga mos
\item Birinchisi faqat audio, ikkinchisi video uchun
\item Ikkalasi ham faqat saqlash buyrug'i
\end{enumerate}
\vspace{2mm}\hrule height 0.2pt \vspace{4mm}
\needspace{5\baselineskip}
\hypertarget{q302}{\textbf{\color{ink}302.}}  {\color{muted}\scriptsize Y1 \textbullet\ daraja 2 \textbullet\ DB fe75f3a1}
\par
Qaysi amal fonni olib tashlashda qaytariluvchi yechim?
\begin{enumerate}[label=\Alph*),leftmargin=9mm]
\item Canvas size bilan bo'yash
\item Asl qatlamni butunlay o'chirib, loyihani yopish
\item Faqat JPEG export
\item Layer mask orqali yashirish
\end{enumerate}
\vspace{2mm}\hrule height 0.2pt \vspace{4mm}
\needspace{5\baselineskip}
\hypertarget{q303}{\textbf{\color{ink}303.}}  {\color{muted}\scriptsize Y1 \textbullet\ daraja 3 \textbullet\ DB 23114287}
\par
Aylana konturini yopib uning ichki qismini tez bo'yash kerak. Qaysi Paint uskuna?
\begin{enumerate}[label=\Alph*),leftmargin=9mm]
\item Crop
\item Fill with color
\item Eraser
\item Text
\end{enumerate}
\vspace{2mm}\hrule height 0.2pt \vspace{4mm}
\needspace{5\baselineskip}
\hypertarget{q304}{\textbf{\color{ink}304.}}  {\color{muted}\scriptsize Y1 \textbullet\ daraja 3 \textbullet\ DB 2c24c07c}
\par
Paintda bitta asalari belgisi bor. Aslini saqlagan holda jami uchta asalari qilish uchun nima qilinadi?
\begin{enumerate}[label=\Alph*),leftmargin=9mm]
\item Resize =\textgreater{} Zoom =\textgreater{} Export
\item Select =\textgreater{} Cut =\textgreater{} Save
\item Select =\textgreater{} Copy =\textgreater{} Paste =\textgreater{} Paste
\item Cut =\textgreater{} Delete =\textgreater{} Save
\end{enumerate}
\vspace{2mm}\hrule height 0.2pt \vspace{4mm}
\needspace{5\baselineskip}
\hypertarget{q305}{\textbf{\color{ink}305.}}  {\color{muted}\scriptsize Y1 \textbullet\ daraja 3 \textbullet\ DB 354d7bb8}
\par
Photoshopdagi qatlamlar yo'qolmasin. Qaysi fayl?
\begin{enumerate}[label=\Alph*),leftmargin=9mm]
\item PSD
\item JPG
\item MP3
\item CSV
\end{enumerate}
\vspace{2mm}\hrule height 0.2pt \vspace{4mm}
\needspace{5\baselineskip}
\hypertarget{q306}{\textbf{\color{ink}306.}}  {\color{muted}\scriptsize Y1 \textbullet\ daraja 3 \textbullet\ DB 391170ac}
\par
Paintdagi kichik tasvirdan to'rtburchak qismni qoldirish kerak. Qaysi juft amal?
\begin{enumerate}[label=\Alph*),leftmargin=9mm]
\item Select va Crop
\item Type va Save
\item Hue va Opacity
\item Group va Trace
\end{enumerate}
\vspace{2mm}\hrule height 0.2pt \vspace{4mm}
\needspace{5\baselineskip}
\hypertarget{q307}{\textbf{\color{ink}307.}}  {\color{muted}\scriptsize Y3 \textbullet\ daraja 3 \textbullet\ DB 43654fa4}
\par
Paintda bitta belgi nusxasidan ikkinchi belgi hosil qilish tartibini belgilang.
\begin{enumerate}[label=\Alph*),leftmargin=9mm]
\item Rectangular yoki Free-form Select bilan uni tanlang
\item Asl belgini chizing
\item Paste qilib nusxani yangi joyga ko'chiring
\item Copy amali orqali buferga nusxa oling
\end{enumerate}
\vspace{2mm}\hrule height 0.2pt \vspace{4mm}
\needspace{5\baselineskip}
\hypertarget{q308}{\textbf{\color{ink}308.}}  {\color{muted}\scriptsize Y3 \textbullet\ daraja 3 \textbullet\ DB 4c86d42d}
\par
Quyidagi amallarni to'g'ri tartibda bajaring: 1) Copy buyrug'ini bajarish; 2) Ikki nusxani Paste qilish; 3) Belgini Select bilan ajratish.
\begin{enumerate}[label=\Alph*),leftmargin=9mm]
\item Copy buyrug'ini bajarish
\item Ikki nusxani Paste qilish
\item Belgini Select bilan ajratish
\end{enumerate}
\vspace{2mm}\hrule height 0.2pt \vspace{4mm}
\needspace{5\baselineskip}
\hypertarget{q309}{\textbf{\color{ink}309.}}  {\color{muted}\scriptsize Y1 \textbullet\ daraja 3 \textbullet\ DB 77a1b399}
\par
Photoshopda qatlamning ko'rinishi o'chirildi (eye icon). Tahrirni tiklash nuqtai nazaridan qanday holat bo'lishi mumkin?
\begin{enumerate}[label=\Alph*),leftmargin=9mm]
\item Qatlam har doim butunlay o'chib ketadi
\item Ma'lumot saqlanib, qatlam vaqtincha yashirinadi
\item JPG avtomatik SVGga aylanadi
\item Pastdagi qatlamlar doim o'chadi
\end{enumerate}
\vspace{2mm}\hrule height 0.2pt \vspace{4mm}
\needspace{5\baselineskip}
\hypertarget{q310}{\textbf{\color{ink}310.}}  {\color{muted}\scriptsize Y1 \textbullet\ daraja 3 \textbullet\ DB 77f94908}
\par
Photoshopda ustki qatlam vaqtincha ko'rinmas bo'lsin. Qanday?
\begin{enumerate}[label=\Alph*),leftmargin=9mm]
\item Flatten image
\item Layers panelidagi eye belgisini o'chirish
\item Save as JPG
\item Delete layer
\end{enumerate}
\vspace{2mm}\hrule height 0.2pt \vspace{4mm}
\needspace{5\baselineskip}
\hypertarget{q311}{\textbf{\color{ink}311.}}  {\color{muted}\scriptsize Y1 \textbullet\ daraja 3 \textbullet\ DB 7b5e910e}
\par
Rasmda mavjud qizil rang aynan olinib yangi chiziqqa berilsin. Qaysi uskuna?
\begin{enumerate}[label=\Alph*),leftmargin=9mm]
\item Crop
\item Blur
\item Eyedropper / Color picker
\item Select all
\end{enumerate}
\vspace{2mm}\hrule height 0.2pt \vspace{4mm}
\needspace{5\baselineskip}
\hypertarget{q312}{\textbf{\color{ink}312.}}  {\color{muted}\scriptsize Y2 \textbullet\ daraja 3 \textbullet\ DB 826a0d30}
\par
Moslashtiring: 1) Pencil; 2) Curved line; 3) Save as. Ta'riflar: a) Avval chiziq, keyin egrilik; b) Boshqa nom bilan saqlash; c) Qo'lda chizish.
\textbf{\footnotesize Chap ustun:}
\begin{enumerate}[label=\arabic*),leftmargin=9mm]
\item Pencil
\item Curved line
\item Save as
\end{enumerate}
\textbf{\footnotesize O'ng ustun:}
\begin{enumerate}[label=\Alph*),leftmargin=9mm]
\item Avval chiziq, keyin egrilik
\item Boshqa nom bilan saqlash
\item Qo'lda chizish
\end{enumerate}
\vspace{2mm}\hrule height 0.2pt \vspace{4mm}
\needspace{5\baselineskip}
\hypertarget{q313}{\textbf{\color{ink}313.}}  {\color{muted}\scriptsize Y1 \textbullet\ daraja 3 \textbullet\ DB 98e1f818}
\par
Fotosuratdagi bitta dog'ni yaqin atrofdagi tekstura bilan ko'chirib tuzatish kerak. Qaysi uskuna?
\begin{enumerate}[label=\Alph*),leftmargin=9mm]
\item Crop
\item Text
\item Clone Stamp
\item Rectangle
\end{enumerate}
\vspace{2mm}\hrule height 0.2pt \vspace{4mm}
\needspace{5\baselineskip}
\hypertarget{q314}{\textbf{\color{ink}314.}}  {\color{muted}\scriptsize Y3 \textbullet\ daraja 3 \textbullet\ DB a9c45fc1}
\par
Paintda yopiq shaklni rang bilan to'ldirish jarayonini tartiblang.
\begin{enumerate}[label=\Alph*),leftmargin=9mm]
\item Yangi ish maydonini yarating
\item Fill with color va to'ldirish rangini tanlang
\item Shapes bilan yopiq kontur chizing
\item Kontur ichida bosib tutash hududni bo'yang
\end{enumerate}
\vspace{2mm}\hrule height 0.2pt \vspace{4mm}
\needspace{5\baselineskip}
\hypertarget{q315}{\textbf{\color{ink}315.}}  {\color{muted}\scriptsize Y2 \textbullet\ daraja 3 \textbullet\ DB ba7375c4}
\par
MS Paint uskunalarini to'g'ri vazifasi bilan moslashtiring.
\textbf{\footnotesize Chap ustun:}
\begin{enumerate}[label=\arabic*),leftmargin=9mm]
\item Pencil
\item Brush
\item Fill with color
\item Color picker
\end{enumerate}
\textbf{\footnotesize O'ng ustun:}
\begin{enumerate}[label=\Alph*),leftmargin=9mm]
\item Turli teksturali chiziq chizish
\item Tutash hududni rang bilan to'ldirish
\item Rasmdan mavjud rangni namuna olish
\item Erkin ingichka chiziq chizish
\end{enumerate}
\vspace{2mm}\hrule height 0.2pt \vspace{4mm}
\needspace{5\baselineskip}
\hypertarget{q316}{\textbf{\color{ink}316.}}  {\color{muted}\scriptsize Y1 \textbullet\ daraja 3 \textbullet\ DB c89e6db0}
\par
Rasmdagi uch qatlam tartibida yuqoridan pastga A, B, C. A va C ko'rinadigan, B yashirin. Qaysi qatlamlar natijaga hissa qo'shadi?
\begin{enumerate}[label=\Alph*),leftmargin=9mm]
\item B va C
\item Faqat B
\item Hech biri
\item A va C
\end{enumerate}
\begin{quote}\footnotesize\color{muted}\textbf{Rasm:} A ko'rinadigan, B yashirin, C ko'rinadigan qatlamlar paneli\\
\href{https://github.com/sarvar9417/attestatsiya/blob/main/public/graphics-web-v2/layer_stack.svg}{Rasmni onlayn ko'rish}
\end{quote}
\vspace{2mm}\hrule height 0.2pt \vspace{4mm}
\needspace{5\baselineskip}
\hypertarget{q317}{\textbf{\color{ink}317.}}  {\color{muted}\scriptsize Y3 \textbullet\ daraja 3 \textbullet\ DB dd37bf46}
\par
MS Paintda Curved line vositasida egri chiziq yasash tartibini tiklang.
\begin{enumerate}[label=\Alph*),leftmargin=9mm]
\item Chiziqni bir nuqtasidan tortib birinchi egrilikni bering
\item Zarur ikkinchi egishni sozlab chiziqni yakunlang
\item Curved line asbobini tanlang
\item Boshlanish va tugash nuqtalari bilan dastlab to'g'ri segment chizing
\end{enumerate}
\vspace{2mm}\hrule height 0.2pt \vspace{4mm}
\needspace{5\baselineskip}
\hypertarget{q318}{\textbf{\color{ink}318.}}  {\color{muted}\scriptsize Y1 \textbullet\ daraja 3 \textbullet\ DB f47b283a}
\par
Photoshopda erkin shaklli sohani ajratish talab qilindi. Qaysi uskuna?
\begin{enumerate}[label=\Alph*),leftmargin=9mm]
\item Move
\item Gradient
\item Type
\item Lasso
\end{enumerate}
\vspace{2mm}\hrule height 0.2pt \vspace{4mm}
\needspace{5\baselineskip}
\hypertarget{q319}{\textbf{\color{ink}319.}}  {\color{muted}\scriptsize Y2 \textbullet\ daraja 3 \textbullet\ DB fbe5504e}
\par
Paintda belgilash va bufer amallarini moslang.
\textbf{\footnotesize Chap ustun:}
\begin{enumerate}[label=\arabic*),leftmargin=9mm]
\item Rectangular select
\item Free-form select
\item Copy
\item Cut
\end{enumerate}
\textbf{\footnotesize O'ng ustun:}
\begin{enumerate}[label=\Alph*),leftmargin=9mm]
\item To'g'ri to'rtburchak hududni tanlash
\item Erkin kontur shaklida hudud tanlash
\item Tanlangan qism nusxasini clipboardga olish
\item Tanlangan qismini olib clipboardga joylashtirish
\end{enumerate}
\vspace{2mm}\hrule height 0.2pt \vspace{4mm}
\needspace{5\baselineskip}
\hypertarget{q320}{\textbf{\color{ink}320.}}  {\color{muted}\scriptsize Y1 \textbullet\ daraja 4 \textbullet\ DB 148214bc}
\par
Paintdagi to'liq bo'lmagan aylana chegarasiga Fill bilan rang berilganda nima mumkin?
\begin{enumerate}[label=\Alph*),leftmargin=9mm]
\item Canvas har doim yopiladi
\item Faqat aylana markazi bo'yaladi
\item Rasm avtomatik SVG bo'ladi
\item Rang tashqaridagi tutash sohaga ham tarqaladi
\end{enumerate}
\vspace{2mm}\hrule height 0.2pt \vspace{4mm}
\needspace{5\baselineskip}
\hypertarget{q321}{\textbf{\color{ink}321.}}  {\color{muted}\scriptsize Y1 \textbullet\ daraja 4 \textbullet\ DB 30a5a3f9}
\par
Ustki qatlam opacity 0\% bo'lsa, pastdagi qatlamga nima bo'ladi?
\begin{enumerate}[label=\Alph*),leftmargin=9mm]
\item Ustki qatlam to'liq xira bo'ladi
\item Canvas avtomatik o'zgaradi
\item Ustki qatlam ko'rinmaydi, pastdagi ko'rinishi mumkin
\item Ikki qatlam ham o'chadi
\end{enumerate}
\vspace{2mm}\hrule height 0.2pt \vspace{4mm}
\needspace{5\baselineskip}
\hypertarget{q322}{\textbf{\color{ink}322.}}  {\color{muted}\scriptsize Y1 \textbullet\ daraja 4 \textbullet\ DB 3d653c4c}
\par
Photoshopda pastdagi yozuvni tepadagi rasmdan yuqoriga chiqarish uchun nima qilinadi?
\begin{enumerate}[label=\Alph*),leftmargin=9mm]
\item Yozuv qatlamini Layers panelida yuqoriga o'tkazish
\item Export JPEG
\item Blur radiusini kamaytirish
\item Canvasni o'chirish
\end{enumerate}
\vspace{2mm}\hrule height 0.2pt \vspace{4mm}
\needspace{5\baselineskip}
\hypertarget{q323}{\textbf{\color{ink}323.}}  {\color{muted}\scriptsize Y1 \textbullet\ daraja 4 \textbullet\ DB 57ff3bc7}
\par
Ustki qatlam opacity 100\% va to'liq qizil bo'lsa, pastdagi ko'k qismga qanday ta'sir?
\begin{enumerate}[label=\Alph*),leftmargin=9mm]
\item Ikki qatlam RGBsiz bo'ladi
\item Ikkalasi avtomatik group
\item Pastdagi ko'k doim ustun
\item Ustki qizil yopib turgan hududda ko'k ko'rinmaydi
\end{enumerate}
\vspace{2mm}\hrule height 0.2pt \vspace{4mm}
\needspace{5\baselineskip}
\hypertarget{q324}{\textbf{\color{ink}324.}}  {\color{muted}\scriptsize Y1 \textbullet\ daraja 4 \textbullet\ DB 5ef3819b}
\par
Photoshopda maskaning qora qismi odatda qanday ma'noga ega?
\begin{enumerate}[label=\Alph*),leftmargin=9mm]
\item Shu qismni ko'rsatish
\item Tasvir rangini o'zgartirish
\item Yangi layer ochish
\item Shu qismni yashirish
\end{enumerate}
\vspace{2mm}\hrule height 0.2pt \vspace{4mm}
\needspace{5\baselineskip}
\hypertarget{q325}{\textbf{\color{ink}325.}}  {\color{muted}\scriptsize Y1 \textbullet\ daraja 4 \textbullet\ DB 88d1ab1e}
\par
Photoshopda maskaning oq qismi odatda qanday ma'noga ega?
\begin{enumerate}[label=\Alph*),leftmargin=9mm]
\item Shu qismni ko'rsatish
\item Layerni o'chirish
\item JPEGga o'tkazish
\item Shu qismni yashirish
\end{enumerate}
\vspace{2mm}\hrule height 0.2pt \vspace{4mm}
\needspace{5\baselineskip}
\hypertarget{q326}{\textbf{\color{ink}326.}}  {\color{muted}\scriptsize Y2 \textbullet\ daraja 4 \textbullet\ DB 8e467aa6}
\par
Photoshop tanlash asboblari nomi va vazifasini juftlang.
\textbf{\footnotesize Chap ustun:}
\begin{enumerate}[label=\arabic*),leftmargin=9mm]
\item Rectangular Marquee
\item Lasso
\item Magic Wand
\item Crop
\end{enumerate}
\textbf{\footnotesize O'ng ustun:}
\begin{enumerate}[label=\Alph*),leftmargin=9mm]
\item To'rtburchak sohani ajratish
\item Rang o'xshashligiga tayangan tanlash
\item Tasvir kadrini kesib qayta belgilash
\item Erkin chegarali sohani tanlash
\end{enumerate}
\vspace{2mm}\hrule height 0.2pt \vspace{4mm}
\needspace{5\baselineskip}
\hypertarget{q327}{\textbf{\color{ink}327.}}  {\color{muted}\scriptsize Y1 \textbullet\ daraja 4 \textbullet\ DB 93b49dbe}
\par
MS Paintda Fill asbobi ishlatilayotganda ochiq kontur nimaga olib kelishi mumkin?
\begin{enumerate}[label=\Alph*),leftmargin=9mm]
\item Rasm HTMLga aylanib qoladi
\item SVG avtomatik paydo bo'ladi
\item Bo'yoq kutilmagan qo'shni sohalarga tarqalishi mumkin
\item Fayl doim to'g'ri to'rtburchak bo'ladi
\end{enumerate}
\vspace{2mm}\hrule height 0.2pt \vspace{4mm}
\needspace{5\baselineskip}
\hypertarget{q328}{\textbf{\color{ink}328.}}  {\color{muted}\scriptsize Y2 \textbullet\ daraja 4 \textbullet\ DB abe17ce3}
\par
Moslashtiring: 1) Clone Stamp; 2) Move; 3) Type. Ta'riflar: a) Matn kiritish; b) Qatlamni joyidan siljitish; c) Namunadan nusxalash.
\textbf{\footnotesize Chap ustun:}
\begin{enumerate}[label=\arabic*),leftmargin=9mm]
\item Clone Stamp
\item Move
\item Type
\end{enumerate}
\textbf{\footnotesize O'ng ustun:}
\begin{enumerate}[label=\Alph*),leftmargin=9mm]
\item Matn kiritish
\item Qatlamni joyidan siljitish
\item Namunadan nusxalash
\end{enumerate}
\vspace{2mm}\hrule height 0.2pt \vspace{4mm}
\needspace{5\baselineskip}
\hypertarget{q329}{\textbf{\color{ink}329.}}  {\color{muted}\scriptsize Y1 \textbullet\ daraja 4 \textbullet\ DB bcecfb34}
\par
Photoshopda adjustment layer va asl fon alohida saqlansa, keyin nima qulay?
\begin{enumerate}[label=\Alph*),leftmargin=9mm]
\item JPEG har doim vektor bo'ladi
\item Asl piksel majburiy yo'qoladi
\item Tuzatish sozlamasini tahrirlash yoki o'chirish
\item Fayl doim kichiklashadi
\end{enumerate}
\vspace{2mm}\hrule height 0.2pt \vspace{4mm}
\needspace{5\baselineskip}
\hypertarget{q330}{\textbf{\color{ink}330.}}  {\color{muted}\scriptsize Y3 \textbullet\ daraja 4 \textbullet\ DB bd718390}
\par
Quyidagi amallarni to'g'ri tartibda bajaring: 1) Maska bilan ko'rinishni sozlash; 2) Layer Mask qo'shish; 3) Qatlamni tanlash.
\begin{enumerate}[label=\Alph*),leftmargin=9mm]
\item Maska bilan ko'rinishni sozlash
\item Layer Mask qo'shish
\item Qatlamni tanlash
\end{enumerate}
\vspace{2mm}\hrule height 0.2pt \vspace{4mm}
\needspace{5\baselineskip}
\hypertarget{q331}{\textbf{\color{ink}331.}}  {\color{muted}\scriptsize Y1 \textbullet\ daraja 4 \textbullet\ DB be1af769}
\par
Paintdagi Undo bir necha marta ishlatilsa, nima sodir bo'lishi mumkin?
\begin{enumerate}[label=\Alph*),leftmargin=9mm]
\item Har safar yangi hujjat ochiladi
\item Rasm faqat kattalashadi
\item Har safar bitta qatlam yaratiladi
\item Oxirgi bajarilgan amallar ketma-ket orqaga qaytariladi
\end{enumerate}
\vspace{2mm}\hrule height 0.2pt \vspace{4mm}
\needspace{5\baselineskip}
\hypertarget{q332}{\textbf{\color{ink}332.}}  {\color{muted}\scriptsize Y1 \textbullet\ daraja 4 \textbullet\ DB c39c142b}
\par
Fotosuratni xiralashtirish uchun qanday filtr turidan foydalaniladi?
\begin{enumerate}[label=\Alph*),leftmargin=9mm]
\item Text
\item Gaussian Blur
\item Levels
\item Sharpen
\end{enumerate}
\vspace{2mm}\hrule height 0.2pt \vspace{4mm}
\needspace{5\baselineskip}
\hypertarget{q333}{\textbf{\color{ink}333.}}  {\color{muted}\scriptsize Y1 \textbullet\ daraja 5 \textbullet\ DB 0cf89c97}
\par
Bir qatlam qulflangan. Move bilan nima kuzatilishi mumkin?
\begin{enumerate}[label=\Alph*),leftmargin=9mm]
\item RGB CMYK bo'lib qoladi
\item Ushbu qatlamni ko'chirish cheklanadi
\item Fayl avtomatik PDFga aylanadi
\item Barcha qatlam o'chadi
\end{enumerate}
\vspace{2mm}\hrule height 0.2pt \vspace{4mm}
\needspace{5\baselineskip}
\hypertarget{q334}{\textbf{\color{ink}334.}}  {\color{muted}\scriptsize Y1 \textbullet\ daraja 5 \textbullet\ DB 2e2d7f9c}
\par
Grafik loyiha ham tahrirlanadigan, ham vebda ko'rinadigan bo'lsin. Qaysi ish oqimi?
\begin{enumerate}[label=\Alph*),leftmargin=9mm]
\item Faqat MP3
\item Faqat JPEG manba
\item PSD manba + mos PNG/JPEG eksport
\item Faqat screenshot
\end{enumerate}
\vspace{2mm}\hrule height 0.2pt \vspace{4mm}
\needspace{5\baselineskip}
\hypertarget{q335}{\textbf{\color{ink}335.}}  {\color{muted}\scriptsize Y1 \textbullet\ daraja 5 \textbullet\ DB 3dbe6d54}
\par
Photoshopda Sharpen filtri nima uchun qo'llanishi mumkin?
\begin{enumerate}[label=\Alph*),leftmargin=9mm]
\item Ayrim konturlar tafsilotini kontrast orqali kuchaytirish
\item JPEGni tahrirlanadigan vektorga aylantirish
\item HTML jadvalni avtomatik chizish
\item Canvasni doim 5x kattalashtirish
\end{enumerate}
\vspace{2mm}\hrule height 0.2pt \vspace{4mm}
\needspace{5\baselineskip}
\hypertarget{q336}{\textbf{\color{ink}336.}}  {\color{muted}\scriptsize Y1 \textbullet\ daraja 5 \textbullet\ DB 560bb6f3}
\par
Paintda bitta obyekt nusxasi kerak, asl obyekt qolsin. Qaysi amal ketma-ketligi to'g'ri?
\begin{enumerate}[label=\Alph*),leftmargin=9mm]
\item Zoom =\textgreater{} Save =\textgreater{} Erase
\item Select =\textgreater{} Cut =\textgreater{} Delete
\item Select =\textgreater{} Copy =\textgreater{} Paste
\item Resize =\textgreater{} Undo =\textgreater{} Close
\end{enumerate}
\vspace{2mm}\hrule height 0.2pt \vspace{4mm}
\needspace{5\baselineskip}
\hypertarget{q337}{\textbf{\color{ink}337.}}  {\color{muted}\scriptsize Y2 \textbullet\ daraja 5 \textbullet\ DB 59dab3b0}
\par
Saqlash va eksport vazifalarini to'g'ri moslang.
\textbf{\footnotesize Chap ustun:}
\begin{enumerate}[label=\arabic*),leftmargin=9mm]
\item Save
\item Save As
\item PSD
\item PNG
\end{enumerate}
\textbf{\footnotesize O'ng ustun:}
\begin{enumerate}[label=\Alph*),leftmargin=9mm]
\item Mavjud ish faylidagi tahrirlarni saqlash
\item Shaffoflik zarur bo'lgan tarqatish uchun rastr nusxa
\item Qatlam va maskalari saqlanadigan Photoshop loyiha fayli
\item Boshqa nom yoki mos format bilan alohida saqlash
\end{enumerate}
\vspace{2mm}\hrule height 0.2pt \vspace{4mm}
\needspace{5\baselineskip}
\hypertarget{q338}{\textbf{\color{ink}338.}}  {\color{muted}\scriptsize Y1 \textbullet\ daraja 5 \textbullet\ DB 5a150897}
\par
Photoshopda maska qora bo'yalgan hudud nima bo'lishi mumkin?
\begin{enumerate}[label=\Alph*),leftmargin=9mm]
\item Rasm 180 deg  aylanadi
\item Qatlamning o'sha qismi yashirinadi
\item Qatlam nomi o'chadi
\item O'sha qismga doim yangi vektor tugun qo'shiladi
\end{enumerate}
\vspace{2mm}\hrule height 0.2pt \vspace{4mm}
\needspace{5\baselineskip}
\hypertarget{q339}{\textbf{\color{ink}339.}}  {\color{muted}\scriptsize Y1 \textbullet\ daraja 5 \textbullet\ DB 66b4dc83}
\par
Photoshopdagi Eraser va Layer Mask farqi haqida qaysi javob yaxshiroq?
\begin{enumerate}[label=\Alph*),leftmargin=9mm]
\item Mask pikselni o'chirmasdan ko'rinishni boshqarishi mumkin
\item Mask faqat audio qatlamida mavjud
\item Har doim ikki amal ham bir xil destruktiv natija beradi
\item Eraser faqat CMYK profilini tanlaydi
\end{enumerate}
\vspace{2mm}\hrule height 0.2pt \vspace{4mm}
\needspace{5\baselineskip}
\hypertarget{q340}{\textbf{\color{ink}340.}}  {\color{muted}\scriptsize Y1 \textbullet\ daraja 5 \textbullet\ DB 95d1fea6}
\par
Photoshopda bir obyekt konturi boshqa elementlar bilan chalkashib ketdi. Kesishdan oldin nima ma'qul?
\begin{enumerate}[label=\Alph*),leftmargin=9mm]
\item Fayl nomini o'zgartirish
\item JPEG sifatini oshirish
\item Tanlov va tegishli qatlamni aniq tekshirish
\item Darhol barcha qatlamni birlashtirish
\end{enumerate}
\vspace{2mm}\hrule height 0.2pt \vspace{4mm}
\needspace{5\baselineskip}
\hypertarget{q341}{\textbf{\color{ink}341.}}  {\color{muted}\scriptsize Y1 \textbullet\ daraja 5 \textbullet\ DB 9cf344bb}
\par
Photoshopda PSD qatlamli manba alohida saqlandi, JPEG eksport qilindi. Qaysi biri tahrir uchun asosiy?
\begin{enumerate}[label=\Alph*),leftmargin=9mm]
\item PSD
\item JPEG
\item Fayl kengaytmasisiz JPG
\item Har doim PNG
\end{enumerate}
\vspace{2mm}\hrule height 0.2pt \vspace{4mm}
\needspace{5\baselineskip}
\hypertarget{q342}{\textbf{\color{ink}342.}}  {\color{muted}\scriptsize Y1 \textbullet\ daraja 5 \textbullet\ DB c0d0b9c1}
\par
Microsoft Paintda Save as buyrug'i bilan qaysi amal mumkin?
\begin{enumerate}[label=\Alph*),leftmargin=9mm]
\item Faylni har doim vektor modelga aylantirish
\item Brauzerda formani jo'natish
\item Fayl nomini va qo'llab-quvvatlansa formatini tanlab saqlash
\item Qatlamlar sonini aniqlash
\end{enumerate}
\vspace{2mm}\hrule height 0.2pt \vspace{4mm}
\needspace{5\baselineskip}
\hypertarget{q343}{\textbf{\color{ink}343.}}  {\color{muted}\scriptsize Y1 \textbullet\ daraja 5 \textbullet\ DB c0f92bc2}
\par
Rasmdagi qatlamlar panelida qaysi amal B ni ko'rsatmasdan A va C ni saqlashga mos?
\begin{enumerate}[label=\Alph*),leftmargin=9mm]
\item B ni JPGga saqlab, qolganini o'chirish
\item B qatlamining RGB pikselini butkul o'chirish
\item B qatlamining ko'rinishini o'chirish
\item B bilan C ni majburan birlashtirish
\end{enumerate}
\begin{quote}\footnotesize\color{muted}\textbf{Rasm:} Photoshop qatlamlar ko'rinish holati\\
\href{https://github.com/sarvar9417/attestatsiya/blob/main/public/graphics-web-v2/layer_stack.svg}{Rasmni onlayn ko'rish}
\end{quote}
\vspace{2mm}\hrule height 0.2pt \vspace{4mm}
\needspace{5\baselineskip}
\hypertarget{q344}{\textbf{\color{ink}344.}}  {\color{muted}\scriptsize Y1 \textbullet\ daraja 5 \textbullet\ DB de3e3f17}
\par
MS Paintda birinchi marta egri chiziq chizayotganda Curved line uskunasining to'g'ri ish ketma-ketligi?
\begin{enumerate}[label=\Alph*),leftmargin=9mm]
\item Faqat palitrani o'chirish
\item Barcha fayllarni yopish =\textgreater{} chizish
\item Segment chizish =\textgreater{} egrilikni sozlash
\item Avval qatlam maskasi qo'shish =\textgreater{} PSDga o'tish
\end{enumerate}
\vspace{2mm}\hrule height 0.2pt \vspace{4mm}
\needspace{5\baselineskip}
\hypertarget{q345}{\textbf{\color{ink}345.}}  {\color{muted}\scriptsize Y1 \textbullet\ daraja 5 \textbullet\ DB e2dc6dac}
\par
Fotosuratdagi dog'ni boshqa hudud namunasi bilan berkitish uchun qaysi asbob qo'llanadi?
\begin{enumerate}[label=\Alph*),leftmargin=9mm]
\item Clone Stamp
\item Zoom Tool
\item Canvas Size
\item Type Tool
\end{enumerate}
\vspace{2mm}\hrule height 0.2pt \vspace{4mm}
\needspace{5\baselineskip}
\hypertarget{q346}{\textbf{\color{ink}346.}}  {\color{muted}\scriptsize Y3 \textbullet\ daraja 5 \textbullet\ DB e4235f4f}
\par
Photoshop qatlamli kompozitsiyasidagi yozuvni ko'rinadigan qilish tartibini belgilang.
\begin{enumerate}[label=\Alph*),leftmargin=9mm]
\item Type asbobi yordamida alohida matn qatlamini yarating
\item Matn qatlamini kerakli ustki tartibga olib, ko'rinishini tekshiring
\item Grafik obyekt qatlamini fon ustiga qo'ying
\item Fon qatlamini yarating
\end{enumerate}
\vspace{2mm}\hrule height 0.2pt \vspace{4mm}
\subsection{HTML va CSS}
\subsubsection{HTML: matn, rasm va ro'yxatlar}
\needspace{5\baselineskip}
\hypertarget{q347}{\textbf{\color{ink}347.}}  {\color{muted}\scriptsize Y1 \textbullet\ daraja 3 \textbullet\ DB 33cd4f8a}
\par
Ushbu HTML fragmentining veb brauzer sahifasida ko'rinishi qaysi javobda to'g'ri ko'rsatilgan? Kod \texttt{\textless{}body\textgreater{}} ichida joylashgan.

\par
\begin{lstlisting}
<ol start="2" type="a">
<b><i>Beginner and Fundamentals</i></b>
<li>Introduction to Computers</li>
<li>IT Support & Networking</li>
<b><u>Programming and Software</u></b>
<li>Python and Java:</li>
<li>Web Development</li>
</ol>
\end{lstlisting}

\begin{enumerate}[label=\Alph*),leftmargin=9mm]
\item Ro'yxat B, C, D, E bosh harflari bilan boshlanadi.
\item Ro'yxat A, B, C, D bosh harflari bilan boshlanadi.
\item Ro'yxat b, c, d, e kichik harflari bilan boshlanadi.
\item Ro'yxat a, b, c, d kichik harflari bilan boshlanadi.
\end{enumerate}
\vspace{2mm}\hrule height 0.2pt \vspace{4mm}
\subsubsection{HTML: jadval va forma elementlari}
\needspace{5\baselineskip}
\hypertarget{q348}{\textbf{\color{ink}348.}}  {\color{muted}\scriptsize Y1 \textbullet\ daraja 3 \textbullet\ DB 738fa38e}
\par
Quyidagi HTML ning vizual natijasini aniqlang. Jadval ustunlari \texttt{Title}, \texttt{Author}, \texttt{Pages}.

\par
\begin{lstlisting}
<table border="1" cellspacing="0" cellpadding="5">
<tr><th>Title</th><th>Author</th><th>Pages</th></tr>
<tr><td>Gone with the Wind</td><td>Margaret Mitchell</td><td rowspan="2">248</td></tr>
<tr><td>Harry Potter and the Chamber</td><td rowspan="2">Rowling</td></tr>
<tr><td>Harry Potter and the Prisoner of Azkaban</td><td>342</td></tr>
</table>
\end{lstlisting}


\par
Qaysi variantdagi \texttt{rowspan} ko'rinishi to'g'ri?
\begin{enumerate}[label=\Alph*),leftmargin=9mm]
\item \texttt{248} birinchi va ikkinchi qatorni egallaydi, \texttt{Rowling} esa ikkita alohida katakda takrorlanadi; sarlavha chapda
\item \texttt{248} birinchi va ikkinchi qatorni, \texttt{Rowling} esa ikkinchi va uchinchi qatorni bitta vertikal katakda egallaydi; sarlavha markazda
\item \texttt{248} va \texttt{Rowling} har qatorda takrorlanadi; quyi ikki qator qalin shrift bilan
\item Ikkita Harry Potter kitobi bitta Title katagiga birlashtiriladi, \texttt{Rowling} alohida ikki katakda
\end{enumerate}
\begin{quote}\footnotesize\color{muted}\textbf{Rasm:} HTML jadval kodi asl rasmdan\\
\href{https://github.com/sarvar9417/attestatsiya/blob/main/public/kuzgi-demo/demo_q27_html.svg}{Rasmni onlayn ko'rish}
\end{quote}
\vspace{2mm}\hrule height 0.2pt \vspace{4mm}
\subsubsection{CSS stillarini HTML elementlariga qo'llash}
\needspace{5\baselineskip}
\hypertarget{q349}{\textbf{\color{ink}349.}}  {\color{muted}\scriptsize Y1 \textbullet\ daraja 3 \textbullet\ DB 681d4cf7}
\par
Quyidagi HTML va CSS berilgan. Ekranda A va B matnlari qanday rangda ko'rinadi?

\par
\begin{lstlisting}
.box { color: red; }
.box p { color: blue; }
.box span { color: green; }
.box p + span { color: orange; }
\end{lstlisting}
\par
\begin{lstlisting}
<div class="box"><p>A</p><span>B</span></div>
\end{lstlisting}

\begin{enumerate}[label=\Alph*),leftmargin=9mm]
\item A - ko'k, B - yashil
\item A - qizil, B - yashil
\item A - ko'k, B - olovrang
\item A - qizil, B - olovrang
\end{enumerate}
\vspace{2mm}\hrule height 0.2pt \vspace{4mm}
\needspace{5\baselineskip}
\hypertarget{q350}{\textbf{\color{ink}350.}}  {\color{muted}\scriptsize Y1 \textbullet\ daraja 3 \textbullet\ DB fbe1ff56}
\par
Berilgan HTML sahifada 400px kenglikda \texttt{Spider-Man: Brand New Day} sarlavhasi och ko'k fonda, matn ko'k rangda, o'ng tomonda esa 95x140 px tasvir joylashgan. 400px lik och kulrang matn paneli berilgan. Qaysi CSS varianti ushbu ko'rinishga mos keladi? (Manba: 21-22-rasmlar.)
\begin{enumerate}[label=\Alph*),leftmargin=9mm]
\item body\{font-family:Times,san-serif;font-color:\#000000\} h2\{width:400px;background-color:\#e6e6ff;padding:10px;position:static\} p\{width:400px;height:auto;padding:2px 10px;background-color:\#f2f2f2\} img.title\{width:95px;height:140px\} img\{float:right;margin-left:0px\}
\item body\{font-family:italic,serif;color:\#000000\} h2\{width:400px;background-color:\#e6e6ff;margin:10px;position:static\} p\{width:400px;height:140px;padding:2px 10px;background-color:\#f2f2f2\} img.title\{width:95px;height:95px\} img\{float:right;margin-left:0px\}
\item body\{font-family:Georgia,Times,serif;color:\#000066\} h2\{width:400px;background-color:\#ffffff;padding:10px;position:static\} p\{width:400px;padding:50px 10px;background-color:\#f2f2f2\} img.title\{width:95px;height:140px\} img\{float:left;margin-left:0px\}
\item body\{font-family:Georgia,Times,serif;color:\#000066\} h2\{width:400px;background-color:\#e6e6ff;padding:10px;position:static\} p\{width:400px;height:140px;padding:2px 10px;background-color:\#f2f2f2\} img.title\{width:95px;height:140px\} img\{float:right;margin-left:0px\}
\end{enumerate}
\begin{quote}\footnotesize\color{muted}\textbf{Rasm:} CSS natijasidagi veb sahifa namunasi\\
\href{https://github.com/sarvar9417/attestatsiya/blob/main/public/kuzgi-demo/demo_q21_css.svg}{Rasmni onlayn ko'rish}
\end{quote}
\vspace{2mm}\hrule height 0.2pt \vspace{4mm}
\subsubsection{HTML: matn, rasm va ro'yxatlar - V2}
\needspace{5\baselineskip}
\hypertarget{q351}{\textbf{\color{ink}351.}}  {\color{muted}\scriptsize Y2 \textbullet\ daraja 3 \textbullet\ DB 1e07940d}
\par
Rasm va havola atributlarini asosiy vazifalari bilan moslang.
\textbf{\footnotesize Chap ustun:}
\begin{enumerate}[label=\arabic*),leftmargin=9mm]
\item \texttt{img src}
\item \texttt{img alt}
\item \texttt{a href}
\item \texttt{a target='\_blank'}
\end{enumerate}
\textbf{\footnotesize O'ng ustun:}
\begin{enumerate}[label=\Alph*),leftmargin=9mm]
\item Rasm ko'rinmasa muqobil mazmuni
\item Rasmning manba manzili
\item Havolani yangi ko'rish kontekstida ochishga ishora
\item Gipermurojaat manzili
\end{enumerate}
\vspace{2mm}\hrule height 0.2pt \vspace{4mm}
\needspace{5\baselineskip}
\hypertarget{q352}{\textbf{\color{ink}352.}}  {\color{muted}\scriptsize Y3 \textbullet\ daraja 3 \textbullet\ DB 39e9f9dd}
\par
HTML hujjatining mantiqiy tuzilishini bosqichma-bosqich tartiblang.
\begin{enumerate}[label=\Alph*),leftmargin=9mm]
\item \texttt{\textless{}html\textgreater{}} ichida \texttt{\textless{}head\textgreater{}} metama'lumotlari va \texttt{\textless{}title\textgreater{}} ni yozing
\item \texttt{\textless{}!doctype html\textgreater{}} bilan hujjat turini e'lon qiling
\item \texttt{\textless{}/body\textgreater{}} va \texttt{\textless{}/html\textgreater{}} bilan ochiq konteynerlarni yoping
\item \texttt{\textless{}body\textgreater{}} ichiga ko'rinadigan sarlavha va paragraflarni kiriting
\end{enumerate}
\vspace{2mm}\hrule height 0.2pt \vspace{4mm}
\needspace{5\baselineskip}
\hypertarget{q353}{\textbf{\color{ink}353.}}  {\color{muted}\scriptsize Y3 \textbullet\ daraja 3 \textbullet\ DB 67cb38bf}
\par
Ichma-ich HTML ro'yxatni to'g'ri qurish ketma-ketligini belgilang.
\begin{enumerate}[label=\Alph*),leftmargin=9mm]
\item Tashqi \texttt{\textless{}ol\textgreater{}} ro'yxatni oching
\item Shu \texttt{\textless{}li\textgreater{}} ichiga \texttt{\textless{}ul\textgreater{}} va ichki \texttt{\textless{}li\textgreater{}} bandlarni joylashtirib ichki ro'yxatni yoping
\item Ichida tashqi \texttt{\textless{}li\textgreater{}} bandni oching va matnini yozing
\item Tashqi \texttt{\textless{}/li\textgreater{}} va \texttt{\textless{}/ol\textgreater{}} teglarini yoping
\end{enumerate}
\vspace{2mm}\hrule height 0.2pt \vspace{4mm}
\needspace{5\baselineskip}
\hypertarget{q354}{\textbf{\color{ink}354.}}  {\color{muted}\scriptsize Y2 \textbullet\ daraja 3 \textbullet\ DB 6d01337a}
\par
HTML matn teglari va vazifalarini moslang.
\textbf{\footnotesize Chap ustun:}
\begin{enumerate}[label=\arabic*),leftmargin=9mm]
\item \texttt{\textless{}h2\textgreater{}}
\item \texttt{\textless{}p\textgreater{}}
\item \texttt{\textless{}strong\textgreater{}}
\item \texttt{\textless{}sub\textgreater{}}
\end{enumerate}
\textbf{\footnotesize O'ng ustun:}
\begin{enumerate}[label=\Alph*),leftmargin=9mm]
\item Ikkinchi darajali sarlavha
\item Abzats yoki paragraf
\item Mazmunan muhim matn
\item Pastki indeks
\end{enumerate}
\vspace{2mm}\hrule height 0.2pt \vspace{4mm}
\needspace{5\baselineskip}
\hypertarget{q355}{\textbf{\color{ink}355.}}  {\color{muted}\scriptsize Y2 \textbullet\ daraja 3 \textbullet\ DB 8e6c8b4e}
\par
Ro'yxat elementlarini ma'nosi bilan moslang.
\textbf{\footnotesize Chap ustun:}
\begin{enumerate}[label=\arabic*),leftmargin=9mm]
\item \texttt{\textless{}ul\textgreater{}}
\item \texttt{\textless{}ol\textgreater{}}
\item \texttt{\textless{}li\textgreater{}}
\item \texttt{ol start='4'}
\end{enumerate}
\textbf{\footnotesize O'ng ustun:}
\begin{enumerate}[label=\Alph*),leftmargin=9mm]
\item Tartibsiz bandlar konteyneri
\item Alohida ro'yxat bandi
\item Raqamlashni to'rtdan boshlash
\item Tartiblangan bandlar konteyneri
\end{enumerate}
\vspace{2mm}\hrule height 0.2pt \vspace{4mm}
\needspace{5\baselineskip}
\hypertarget{q356}{\textbf{\color{ink}356.}}  {\color{muted}\scriptsize Y3 \textbullet\ daraja 3 \textbullet\ DB ab41f271}
\par
Havola ichida bosiladigan rasm yaratish tartibini tuzing.
\begin{enumerate}[label=\Alph*),leftmargin=9mm]
\item \texttt{\textless{}a href='...'\textgreater{}} havola konteynerini oching
\item Konteyner ichiga src va alt bilan \texttt{\textless{}img\textgreater{}} ni joylang
\item Havolani \texttt{\textless{}/a\textgreater{}} bilan yoping va brauzerda tekshiring
\item Rasm uchun to'g'ri src manzilini va alt matnini aniqlang
\end{enumerate}
\vspace{2mm}\hrule height 0.2pt \vspace{4mm}
\needspace{5\baselineskip}
\hypertarget{q357}{\textbf{\color{ink}357.}}  {\color{muted}\scriptsize Y2 \textbullet\ daraja 4 \textbullet\ DB 2531ff1d}
\par
HTML hujjatidagi asosiy qismlarni vazifasi bilan moslang.
\textbf{\footnotesize Chap ustun:}
\begin{enumerate}[label=\arabic*),leftmargin=9mm]
\item \texttt{\textless{}head\textgreater{}}
\item \texttt{\textless{}title\textgreater{}}
\item \texttt{\textless{}body\textgreater{}}
\item \texttt{meta charset='utf-8'}
\end{enumerate}
\textbf{\footnotesize O'ng ustun:}
\begin{enumerate}[label=\Alph*),leftmargin=9mm]
\item Ko'rinadigan asosiy sahifa tarkibi
\item Matn kodlashini ko'rsatish
\item Brauzer yorlig'idagi sahifa nomi
\item Metama'lumotlar va resurs ulanishlari
\end{enumerate}
\vspace{2mm}\hrule height 0.2pt \vspace{4mm}
\needspace{5\baselineskip}
\hypertarget{q358}{\textbf{\color{ink}358.}}  {\color{muted}\scriptsize Y3 \textbullet\ daraja 4 \textbullet\ DB 82c57795}
\par
\texttt{\textless{}ol start='3'\textgreater{}} ichida ikkita band natijasini tekshirish tartibini tuzing.
\begin{enumerate}[label=\Alph*),leftmargin=9mm]
\item Tartiblangan \texttt{\textless{}ol\textgreater{}} va start=3 atributini aniqlang
\item Ikkinchi \texttt{\textless{}li\textgreater{}} ni birinchidan keyin joylang
\item Raqamlar 3 va 4 bo'lib ko'rinishini tekshiring
\item Birinchi \texttt{\textless{}li\textgreater{}} matnini ro'yxatga qo'shing
\end{enumerate}
\vspace{2mm}\hrule height 0.2pt \vspace{4mm}
\subsubsection{HTML: jadvallar va formalar - V2}
\needspace{5\baselineskip}
\hypertarget{q359}{\textbf{\color{ink}359.}}  {\color{muted}\scriptsize Y1 \textbullet\ daraja 1 \textbullet\ DB 0a6f17e7}
\par
HTMLda brauzer yorlig'i nomi qaysi tegda?
\begin{enumerate}[label=\Alph*),leftmargin=9mm]
\item h1
\item title
\item caption
\item p
\end{enumerate}
\vspace{2mm}\hrule height 0.2pt \vspace{4mm}
\needspace{5\baselineskip}
\hypertarget{q360}{\textbf{\color{ink}360.}}  {\color{muted}\scriptsize Y1 \textbullet\ daraja 1 \textbullet\ DB 16328bef}
\par
HTMLda paragrafni qaysi teg yaratadi?
\begin{enumerate}[label=\Alph*),leftmargin=9mm]
\item ul
\item tr
\item p
\item img
\end{enumerate}
\vspace{2mm}\hrule height 0.2pt \vspace{4mm}
\needspace{5\baselineskip}
\hypertarget{q361}{\textbf{\color{ink}361.}}  {\color{muted}\scriptsize Y1 \textbullet\ daraja 1 \textbullet\ DB 19c052f6}
\par
HTMLda havola manzili uchun atribut?
\begin{enumerate}[label=\Alph*),leftmargin=9mm]
\item src
\item href
\item rowspan
\item alt
\end{enumerate}
\vspace{2mm}\hrule height 0.2pt \vspace{4mm}
\needspace{5\baselineskip}
\hypertarget{q362}{\textbf{\color{ink}362.}}  {\color{muted}\scriptsize Y1 \textbullet\ daraja 1 \textbullet\ DB 39a76761}
\par
HTMLda jadval qatorini qaysi teg belgilaydi?
\begin{enumerate}[label=\Alph*),leftmargin=9mm]
\item tr
\item td
\item th
\item caption
\end{enumerate}
\vspace{2mm}\hrule height 0.2pt \vspace{4mm}
\needspace{5\baselineskip}
\hypertarget{q363}{\textbf{\color{ink}363.}}  {\color{muted}\scriptsize Y1 \textbullet\ daraja 1 \textbullet\ DB 4ac38aa4}
\par
Tartiblangan ro'yxat tegi qaysi?
\begin{enumerate}[label=\Alph*),leftmargin=9mm]
\item ol
\item dl
\item li
\item ul
\end{enumerate}
\vspace{2mm}\hrule height 0.2pt \vspace{4mm}
\needspace{5\baselineskip}
\hypertarget{q364}{\textbf{\color{ink}364.}}  {\color{muted}\scriptsize Y1 \textbullet\ daraja 1 \textbullet\ DB 5561a093}
\par
HTMLda \textless{}input type="reset"\textgreater{} elementi odatda nima qiladi?
\begin{enumerate}[label=\Alph*),leftmargin=9mm]
\item Forma ma'lumotlarini majburan POST qiladi
\item Serverdagi bazadan barcha yozuvlarni o'chiradi
\item Forma elementlarini boshlang'ich qiymatlariga qaytaradi
\item HTML jadvalning qatorini qo'shadi
\end{enumerate}
\vspace{2mm}\hrule height 0.2pt \vspace{4mm}
\needspace{5\baselineskip}
\hypertarget{q365}{\textbf{\color{ink}365.}}  {\color{muted}\scriptsize Y1 \textbullet\ daraja 1 \textbullet\ DB 6e8b6eaa}
\par
HTMLda asosiy ko'rinadigan mazmun qaysi element ichida?
\begin{enumerate}[label=\Alph*),leftmargin=9mm]
\item meta
\item head
\item body
\item title
\end{enumerate}
\vspace{2mm}\hrule height 0.2pt \vspace{4mm}
\needspace{5\baselineskip}
\hypertarget{q366}{\textbf{\color{ink}366.}}  {\color{muted}\scriptsize Y1 \textbullet\ daraja 1 \textbullet\ DB 8c0000a4}
\par
HTMLdagi \textless{}input type="hidden"\textgreater{} qiymatiga nisbatan qaysi fikr to'g'ri?
\begin{enumerate}[label=\Alph*),leftmargin=9mm]
\item Har doim serverda parolni o'zgartiradi
\item Ko'rinmas bo'lsa ham sahifa manbasidan o'qilishi mumkin
\item Foydalanuvchining yuborishiga hech ta'sir qilmaydi
\item Mutlaqo shifrlangan sir bo'ladi
\end{enumerate}
\vspace{2mm}\hrule height 0.2pt \vspace{4mm}
\needspace{5\baselineskip}
\hypertarget{q367}{\textbf{\color{ink}367.}}  {\color{muted}\scriptsize Y1 \textbullet\ daraja 1 \textbullet\ DB 95831dd7}
\par
HTMLda rasm joylash tegi qaysi?
\begin{enumerate}[label=\Alph*),leftmargin=9mm]
\item pictureurl
\item image
\item img
\item src
\end{enumerate}
\vspace{2mm}\hrule height 0.2pt \vspace{4mm}
\needspace{5\baselineskip}
\hypertarget{q368}{\textbf{\color{ink}368.}}  {\color{muted}\scriptsize Y1 \textbullet\ daraja 1 \textbullet\ DB a2d1e556}
\par
Tartibsiz ro'yxat tegi qaysi?
\begin{enumerate}[label=\Alph*),leftmargin=9mm]
\item ol
\item tr
\item th
\item ul
\end{enumerate}
\vspace{2mm}\hrule height 0.2pt \vspace{4mm}
\needspace{5\baselineskip}
\hypertarget{q369}{\textbf{\color{ink}369.}}  {\color{muted}\scriptsize Y1 \textbullet\ daraja 1 \textbullet\ DB c11d8b7b}
\par
HTML nimaning qisqartmasi?
\begin{enumerate}[label=\Alph*),leftmargin=9mm]
\item Home Text Module Language
\item HyperTool Machine Language
\item HighText Media List
\item HyperText Markup Language
\end{enumerate}
\vspace{2mm}\hrule height 0.2pt \vspace{4mm}
\needspace{5\baselineskip}
\hypertarget{q370}{\textbf{\color{ink}370.}}  {\color{muted}\scriptsize Y1 \textbullet\ daraja 1 \textbullet\ DB f0f9e659}
\par
HTMLda rasming URL manzili uchun atribut?
\begin{enumerate}[label=\Alph*),leftmargin=9mm]
\item src
\item colspan
\item target
\item href
\end{enumerate}
\vspace{2mm}\hrule height 0.2pt \vspace{4mm}
\needspace{5\baselineskip}
\hypertarget{q371}{\textbf{\color{ink}371.}}  {\color{muted}\scriptsize Y1 \textbullet\ daraja 2 \textbullet\ DB 25670633}
\par
strong va b farqi nimada?
\begin{enumerate}[label=\Alph*),leftmargin=9mm]
\item b faqat rasm
\item strong semantik muhimlikni, b ko'proq vizual ajratishni bildiradi
\item Ikkisi HTML formaga tegishli
\item strong faqat jadval
\end{enumerate}
\vspace{2mm}\hrule height 0.2pt \vspace{4mm}
\needspace{5\baselineskip}
\hypertarget{q372}{\textbf{\color{ink}372.}}  {\color{muted}\scriptsize Y1 \textbullet\ daraja 2 \textbullet\ DB 263beaf4}
\par
Checkbox va radio farqi qaysi?
\begin{enumerate}[label=\Alph*),leftmargin=9mm]
\item Radio parol kiritadi
\item Ikkalasi faqat rasm
\item Checkbox mustaqil ko'p tanlov, radio bir name guruhi ichida bittasi
\item Radio doim bir nechta, checkbox bittasi
\end{enumerate}
\vspace{2mm}\hrule height 0.2pt \vspace{4mm}
\needspace{5\baselineskip}
\hypertarget{q373}{\textbf{\color{ink}373.}}  {\color{muted}\scriptsize Y1 \textbullet\ daraja 2 \textbullet\ DB 30b00548}
\par
rowspan va colspan farqi nima?
\begin{enumerate}[label=\Alph*),leftmargin=9mm]
\item Ikkalasi faqat rang
\item Ikkalasi faqat HTML sarlavha
\item Birinchisi qator, ikkinchisi ustun qamrovi
\item Birinchisi ustun, ikkinchisi qator
\end{enumerate}
\vspace{2mm}\hrule height 0.2pt \vspace{4mm}
\needspace{5\baselineskip}
\hypertarget{q374}{\textbf{\color{ink}374.}}  {\color{muted}\scriptsize Y1 \textbullet\ daraja 2 \textbullet\ DB 4db4482a}
\par
Qaysi juftlik to'g'ri?
\begin{enumerate}[label=\Alph*),leftmargin=9mm]
\item ul:method va form:src
\item a:href va img:src
\item a:src va img:href
\item p:colspan va tr:alt
\end{enumerate}
\vspace{2mm}\hrule height 0.2pt \vspace{4mm}
\needspace{5\baselineskip}
\hypertarget{q375}{\textbf{\color{ink}375.}}  {\color{muted}\scriptsize Y1 \textbullet\ daraja 2 \textbullet\ DB 5fe20d0f}
\par
required va disabled haqida qaysi fikr to'g'ri?
\begin{enumerate}[label=\Alph*),leftmargin=9mm]
\item disabled serverga har doim yuboriladi
\item required maydonni doim yashiradi
\item Ikkisi faqat matnni qalin qiladi
\item required kiritishni talab qiladi, disabled odatda yuborilmaydi
\end{enumerate}
\vspace{2mm}\hrule height 0.2pt \vspace{4mm}
\needspace{5\baselineskip}
\hypertarget{q376}{\textbf{\color{ink}376.}}  {\color{muted}\scriptsize Y1 \textbullet\ daraja 2 \textbullet\ DB 75e3e661}
\par
GET va POST haqida qaysi fikr to'g'ri?
\begin{enumerate}[label=\Alph*),leftmargin=9mm]
\item GET odatda URL query parametrlariga, POST so'rov tanasiga yozadi
\item POST faqat CSSga tegishli
\item GET doim serverga hech nima yubormaydi
\item POST o'zi HTTPS bilan teng
\end{enumerate}
\vspace{2mm}\hrule height 0.2pt \vspace{4mm}
\needspace{5\baselineskip}
\hypertarget{q377}{\textbf{\color{ink}377.}}  {\color{muted}\scriptsize Y1 \textbullet\ daraja 2 \textbullet\ DB 870ba2dc}
\par
HTMLdagi font tegi haqida zamonaviy fikr?
\begin{enumerate}[label=\Alph*),leftmargin=9mm]
\item Faqat jadval uchun
\item Tarixiy/uslubiy, ko'rinish uchun CSS afzal
\item Brauzerda hech qachon ishlamaydi
\item HTML5da yagona majburiy usul
\end{enumerate}
\vspace{2mm}\hrule height 0.2pt \vspace{4mm}
\needspace{5\baselineskip}
\hypertarget{q378}{\textbf{\color{ink}378.}}  {\color{muted}\scriptsize Y1 \textbullet\ daraja 2 \textbullet\ DB 8fc76c5c}
\par
Rasm uchun alt atributi nima beradi?
\begin{enumerate}[label=\Alph*),leftmargin=9mm]
\item Qatorlar sonini
\item JPEG siqish darajasini
\item Rasmning muqobil tavsifini
\item CSS rangini
\end{enumerate}
\vspace{2mm}\hrule height 0.2pt \vspace{4mm}
\needspace{5\baselineskip}
\hypertarget{q379}{\textbf{\color{ink}379.}}  {\color{muted}\scriptsize Y1 \textbullet\ daraja 2 \textbullet\ DB 99888599}
\par
th va td farqi nima?
\begin{enumerate}[label=\Alph*),leftmargin=9mm]
\item th faqat rasm, td faqat forma
\item th sarlavha, td oddiy katak
\item Ikkalasi har doim bir xil
\item td sarlavha, th ro'yxat
\end{enumerate}
\vspace{2mm}\hrule height 0.2pt \vspace{4mm}
\needspace{5\baselineskip}
\hypertarget{q380}{\textbf{\color{ink}380.}}  {\color{muted}\scriptsize Y1 \textbullet\ daraja 2 \textbullet\ DB ac659173}
\par
Qaysi elementning o'zi odatda yopuvchi teg talab qilmaydi?
\begin{enumerate}[label=\Alph*),leftmargin=9mm]
\item p
\item img
\item table
\item ol
\end{enumerate}
\vspace{2mm}\hrule height 0.2pt \vspace{4mm}
\needspace{5\baselineskip}
\hypertarget{q381}{\textbf{\color{ink}381.}}  {\color{muted}\scriptsize Y1 \textbullet\ daraja 2 \textbullet\ DB e9386b8b}
\par
file input orqali haqiqiy fayl yuboriladigan formaga odatda qaysi enctype mos?
\begin{enumerate}[label=\Alph*),leftmargin=9mm]
\item text/css
\item image/svg+xml
\item application/javascript
\item multipart/form-data
\end{enumerate}
\vspace{2mm}\hrule height 0.2pt \vspace{4mm}
\needspace{5\baselineskip}
\hypertarget{q382}{\textbf{\color{ink}382.}}  {\color{muted}\scriptsize Y1 \textbullet\ daraja 2 \textbullet\ DB f7f3f1d2}
\par
Formadagi id va name farqi nima?
\begin{enumerate}[label=\Alph*),leftmargin=9mm]
\item Ikkalasi faqat rang
\item id URL, name faqat CSS
\item name hech qachon yuborilmaydi
\item id elementni bog'lashda, name yuboriladigan maydon kalitida ishlaydi
\end{enumerate}
\vspace{2mm}\hrule height 0.2pt \vspace{4mm}
\needspace{5\baselineskip}
\hypertarget{q383}{\textbf{\color{ink}383.}}  {\color{muted}\scriptsize Y1 \textbullet\ daraja 3 \textbullet\ DB 03aef049}
\par
HTML sahifasida \texttt{photo.png} rasmini ko'rsatish va rasm yuklanmasa \texttt{Foto} degan muqobil matnni berish uchun qaysi yozuv to'g'ri?
\begin{enumerate}[label=\Alph*),leftmargin=9mm]
\item \textless{}img href="photo.png"\textgreater{}
\item \textless{}a src="photo.png"\textgreater{}
\item \textless{}img src="photo.png" alt="Foto"\textgreater{}
\item \textless{}img file="photo.png"\textgreater{}
\end{enumerate}
\vspace{2mm}\hrule height 0.2pt \vspace{4mm}
\needspace{5\baselineskip}
\hypertarget{q384}{\textbf{\color{ink}384.}}  {\color{muted}\scriptsize Y1 \textbullet\ daraja 3 \textbullet\ DB 140ceba5}
\par
Bir veb-hujjat ichida aloqa bo'limiga o'tish uchun?
\begin{enumerate}[label=\Alph*),leftmargin=9mm]
\item \textless{}a src="aloqa"\textgreater{}Aloqa\textless{}/a\textgreater{}
\item \textless{}p target="aloqa"\textgreater{}
\item \textless{}img href="aloqa"\textgreater{}
\item \textless{}a href="\#aloqa"\textgreater{}Aloqa\textless{}/a\textgreater{}
\end{enumerate}
\vspace{2mm}\hrule height 0.2pt \vspace{4mm}
\needspace{5\baselineskip}
\hypertarget{q385}{\textbf{\color{ink}385.}}  {\color{muted}\scriptsize Y1 \textbullet\ daraja 3 \textbullet\ DB 15c6a570}
\par
Formani odatda URL query orqali yuborish?
\begin{enumerate}[label=\Alph*),leftmargin=9mm]
\item method="post"
\item method="get"
\item method="put"
\item method="file"
\end{enumerate}
\vspace{2mm}\hrule height 0.2pt \vspace{4mm}
\needspace{5\baselineskip}
\hypertarget{q386}{\textbf{\color{ink}386.}}  {\color{muted}\scriptsize Y1 \textbullet\ daraja 3 \textbullet\ DB 378f1faa}
\par
Rasmda ko'rsatilgan 2 darajali HTML ro'yxat uchun ichki \textless{}ul\textgreater{} eng to'g'ri qayerda joylashadi?
\begin{enumerate}[label=\Alph*),leftmargin=9mm]
\item \textless{}img\textgreater{} tegi ichida
\item \textless{}meta charset\textgreater{} atributida
\item Tashqi \textless{}li\textgreater{} elementi ichida
\item \textless{}title\textgreater{} ichida
\end{enumerate}
\begin{quote}\footnotesize\color{muted}\textbf{Rasm:} Grafika ichida rastr-vektor va tashqi Veb bandi ko'rsatilgan ro'yxat\\
\href{https://github.com/sarvar9417/attestatsiya/blob/main/public/graphics-web-v2/nested_list.svg}{Rasmni onlayn ko'rish}
\end{quote}
\vspace{2mm}\hrule height 0.2pt \vspace{4mm}
\needspace{5\baselineskip}
\hypertarget{q387}{\textbf{\color{ink}387.}}  {\color{muted}\scriptsize Y1 \textbullet\ daraja 3 \textbullet\ DB 5b5bded9}
\par
\textless{}ol start="3"\textgreater{} ichiga 5 ta \textless{}li\textgreater{} yozildi. Oxirgi band standart raqamlanishda nechchi bo'ladi?
\begin{enumerate}[label=\Alph*),leftmargin=9mm]
\item 7
\item 8
\item 21
\item 3
\end{enumerate}
\vspace{2mm}\hrule height 0.2pt \vspace{4mm}
\needspace{5\baselineskip}
\hypertarget{q388}{\textbf{\color{ink}388.}}  {\color{muted}\scriptsize Y1 \textbullet\ daraja 3 \textbullet\ DB 6ebbf71b}
\par
\textless{}ol start="3"\textgreater{}\textless{}li\textgreater{}A\textless{}/li\textgreater{}\textless{}li\textgreater{}B\textless{}/li\textgreater{}\textless{}li\textgreater{}C\textless{}/li\textgreater{}\textless{}/ol\textgreater{} ning oxirgi band raqami nechchi?
\begin{enumerate}[label=\Alph*),leftmargin=9mm]
\item 6
\item 4
\item 5
\item 3
\end{enumerate}
\vspace{2mm}\hrule height 0.2pt \vspace{4mm}
\needspace{5\baselineskip}
\hypertarget{q389}{\textbf{\color{ink}389.}}  {\color{muted}\scriptsize Y1 \textbullet\ daraja 3 \textbullet\ DB 9bd6867c}
\par
\textless{}ol start="7"\textgreater{} ichiga 2 ta \textless{}li\textgreater{} yozildi. Oxirgi band standart raqamlanishda nechchi bo'ladi?
\begin{enumerate}[label=\Alph*),leftmargin=9mm]
\item 7
\item 8
\item 9
\item 24
\end{enumerate}
\vspace{2mm}\hrule height 0.2pt \vspace{4mm}
\needspace{5\baselineskip}
\hypertarget{q390}{\textbf{\color{ink}390.}}  {\color{muted}\scriptsize Y3 \textbullet\ daraja 3 \textbullet\ DB 9fc3c2a5}
\par
Quyidagi amallarni to'g'ri tartibda bajaring: 1) \textless{}form\textgreater{}ni ochish; 2) Maydonlarni joylashtirish; 3) Submit tugmasini qo'yish.
\begin{enumerate}[label=\Alph*),leftmargin=9mm]
\item \textless{}form\textgreater{}ni ochish
\item Maydonlarni joylashtirish
\item Submit tugmasini qo'yish
\end{enumerate}
\vspace{2mm}\hrule height 0.2pt \vspace{4mm}
\needspace{5\baselineskip}
\hypertarget{q391}{\textbf{\color{ink}391.}}  {\color{muted}\scriptsize Y1 \textbullet\ daraja 3 \textbullet\ DB da044b25}
\par
Bir xil name="g" bo'lgan ikkita radio va alohida name="h" bo'lgan radio mavjud. Qaysi fikr to'g'ri?
\begin{enumerate}[label=\Alph*),leftmargin=9mm]
\item g guruhidan odatda bittasi, h esa alohida tanlanadi
\item Radio uchun name ishlamaydi
\item Checkbox kabi hammasi bir vaqtda g guruhida tanlanadi
\item Uchalasidan har doim bittasi
\end{enumerate}
\vspace{2mm}\hrule height 0.2pt \vspace{4mm}
\needspace{5\baselineskip}
\hypertarget{q392}{\textbf{\color{ink}392.}}  {\color{muted}\scriptsize Y2 \textbullet\ daraja 3 \textbullet\ DB dea86e68}
\par
Moslashtiring: 1) reset; 2) hidden; 3) file. Ta'riflar: a) Boshlang'ich qiymatga qaytarish; b) Ko'rinmas maydon; c) Fayl tanlash.
\textbf{\footnotesize Chap ustun:}
\begin{enumerate}[label=\arabic*),leftmargin=9mm]
\item reset
\item hidden
\item file
\end{enumerate}
\textbf{\footnotesize O'ng ustun:}
\begin{enumerate}[label=\Alph*),leftmargin=9mm]
\item Boshlang'ich qiymatga qaytarish
\item Ko'rinmas maydon
\item Fayl tanlash
\end{enumerate}
\vspace{2mm}\hrule height 0.2pt \vspace{4mm}
\needspace{5\baselineskip}
\hypertarget{q393}{\textbf{\color{ink}393.}}  {\color{muted}\scriptsize Y1 \textbullet\ daraja 3 \textbullet\ DB ed93fe78}
\par
\textless{}ol start="5"\textgreater{} ichiga 4 ta \textless{}li\textgreater{} yozildi. Oxirgi band standart raqamlanishda nechchi bo'ladi?
\begin{enumerate}[label=\Alph*),leftmargin=9mm]
\item 24
\item 5
\item 9
\item 8
\end{enumerate}
\vspace{2mm}\hrule height 0.2pt \vspace{4mm}
\needspace{5\baselineskip}
\hypertarget{q394}{\textbf{\color{ink}394.}}  {\color{muted}\scriptsize Y1 \textbullet\ daraja 3 \textbullet\ DB f5843da9}
\par
\textless{}ol start="1"\textgreater{} ichiga 3 ta \textless{}li\textgreater{} yozildi. Oxirgi band standart raqamlanishda nechchi bo'ladi?
\begin{enumerate}[label=\Alph*),leftmargin=9mm]
\item 3
\item 4
\item 1
\item 9
\end{enumerate}
\vspace{2mm}\hrule height 0.2pt \vspace{4mm}
\needspace{5\baselineskip}
\hypertarget{q395}{\textbf{\color{ink}395.}}  {\color{muted}\scriptsize Y1 \textbullet\ daraja 4 \textbullet\ DB 0302c80c}
\par
\textless{}tr\textgreater{} ichida 2 ta oddiy \textless{}td\textgreater{} va yana bitta colspan="4" katak bor. Jami mantiqiy ustunlar nechta? (Boshqa qamrov yo'q.)
\begin{enumerate}[label=\Alph*),leftmargin=9mm]
\item 3
\item 7
\item 5
\item 6
\end{enumerate}
\vspace{2mm}\hrule height 0.2pt \vspace{4mm}
\needspace{5\baselineskip}
\hypertarget{q396}{\textbf{\color{ink}396.}}  {\color{muted}\scriptsize Y1 \textbullet\ daraja 4 \textbullet\ DB 1764a435}
\par
\textless{}tr\textgreater{} ichida 3 ta oddiy \textless{}td\textgreater{} va yana bitta colspan="2" katak bor. Jami mantiqiy ustunlar nechta? (Boshqa qamrov yo'q.)
\begin{enumerate}[label=\Alph*),leftmargin=9mm]
\item 15
\item 5
\item 4
\item 6
\end{enumerate}
\vspace{2mm}\hrule height 0.2pt \vspace{4mm}
\needspace{5\baselineskip}
\hypertarget{q397}{\textbf{\color{ink}397.}}  {\color{muted}\scriptsize Y1 \textbullet\ daraja 4 \textbullet\ DB 2c979995}
\par
HTML: \textless{}ol\textgreater{}\textless{}li\textgreater{}A\textless{}/li\textgreater{}\textless{}li\textgreater{}B\textless{}ul\textgreater{}\textless{}li\textgreater{}C\textless{}/li\textgreater{}\textless{}/ul\textgreater{}\textless{}/li\textgreater{}\textless{}/ol\textgreater{} . Yuqori darajali raqamli bandlar soni?
\begin{enumerate}[label=\Alph*),leftmargin=9mm]
\item 2
\item 3
\item 4
\item 1
\end{enumerate}
\vspace{2mm}\hrule height 0.2pt \vspace{4mm}
\needspace{5\baselineskip}
\hypertarget{q398}{\textbf{\color{ink}398.}}  {\color{muted}\scriptsize Y1 \textbullet\ daraja 4 \textbullet\ DB 2fc8f7ad}
\par
\textless{}tr\textgreater{} ichida 5 ta oddiy \textless{}td\textgreater{} va yana bitta colspan="2" katak bor. Jami mantiqiy ustunlar nechta? (Boshqa qamrov yo'q.)
\begin{enumerate}[label=\Alph*),leftmargin=9mm]
\item 21
\item 6
\item 7
\item 8
\end{enumerate}
\vspace{2mm}\hrule height 0.2pt \vspace{4mm}
\needspace{5\baselineskip}
\hypertarget{q399}{\textbf{\color{ink}399.}}  {\color{muted}\scriptsize Y1 \textbullet\ daraja 4 \textbullet\ DB 4b14ce6c}
\par
\textless{}tr\textgreater{} ichida 2 ta oddiy \textless{}td\textgreater{} va yana bitta colspan="3" katak bor. Jami mantiqiy ustunlar nechta? (Boshqa qamrov yo'q.)
\begin{enumerate}[label=\Alph*),leftmargin=9mm]
\item 5
\item 6
\item 3
\item 4
\end{enumerate}
\vspace{2mm}\hrule height 0.2pt \vspace{4mm}
\needspace{5\baselineskip}
\hypertarget{q400}{\textbf{\color{ink}400.}}  {\color{muted}\scriptsize Y1 \textbullet\ daraja 4 \textbullet\ DB 59d74479}
\par
Qaysi tugma forma qiymatlarini jo'natishga, qaysisi dastlabki qiymatga qaytarishga xizmat qiladi?
\begin{enumerate}[label=\Alph*),leftmargin=9mm]
\item ikkalasi faylni siqadi
\item submit - tiklaydi; reset - jo'natadi
\item ikkalasi CSS selektoridir
\item submit - jo'natadi; reset - tiklaydi
\end{enumerate}
\vspace{2mm}\hrule height 0.2pt \vspace{4mm}
\needspace{5\baselineskip}
\hypertarget{q401}{\textbf{\color{ink}401.}}  {\color{muted}\scriptsize Y3 \textbullet\ daraja 4 \textbullet\ DB 5d0f75cf}
\par
Uch mantiqiy ustunli jadval sarlavhasini tuzish ketma-ketligini belgilang.
\begin{enumerate}[label=\Alph*),leftmargin=9mm]
\item Birinchi \texttt{\textless{}tr\textgreater{}} sarlavha qatorini yarating
\item Sarlavha qatorini yoping va keyingi ma'lumot qatorlarini yozing
\item \texttt{\textless{}table\textgreater{}} konteynerini oching
\item Qator ichiga oddiy \texttt{\textless{}th\textgreater{}} va \texttt{colspan='2'} li ikkinchi \texttt{\textless{}th\textgreater{}} ni kiriting
\end{enumerate}
\vspace{2mm}\hrule height 0.2pt \vspace{4mm}
\needspace{5\baselineskip}
\hypertarget{q402}{\textbf{\color{ink}402.}}  {\color{muted}\scriptsize Y1 \textbullet\ daraja 4 \textbullet\ DB 6461363f}
\par
Formada id="email" bor, lekin name atributi yo'q. Oddiy form submitda shu maydon kaliti qanday?
\begin{enumerate}[label=\Alph*),leftmargin=9mm]
\item Kalit har doim type bo'ladi
\item HTML xato tufayli sahifa ochilmaydi
\item Maydon odatda yuborilmaydi
\item Kalit avtomatik id bo'ladi
\end{enumerate}
\vspace{2mm}\hrule height 0.2pt \vspace{4mm}
\needspace{5\baselineskip}
\hypertarget{q403}{\textbf{\color{ink}403.}}  {\color{muted}\scriptsize Y3 \textbullet\ daraja 4 \textbullet\ DB 65f3647b}
\par
Fayl yuklash formasini mazmunli tartibda tuzing.
\begin{enumerate}[label=\Alph*),leftmargin=9mm]
\item \texttt{\textless{}form\textgreater{}} uchun POST, action va \texttt{enctype='multipart/form-data'} ni belgilang
\item Fayl tanlash uchun name bilan \texttt{\textless{}input type='file'\textgreater{}} yozing
\item Kerakli ko'rinmaydigan ma'lumot uchun \texttt{\textless{}input type='hidden'\textgreater{}} qo'shing
\item Submit tugmasini qo'shib formani yoping va tekshiring
\end{enumerate}
\vspace{2mm}\hrule height 0.2pt \vspace{4mm}
\needspace{5\baselineskip}
\hypertarget{q404}{\textbf{\color{ink}404.}}  {\color{muted}\scriptsize Y1 \textbullet\ daraja 4 \textbullet\ DB 69154c84}
\par
\textless{}tr\textgreater{} ichida 4 ta oddiy \textless{}td\textgreater{} va yana bitta colspan="2" katak bor. Jami mantiqiy ustunlar nechta? (Boshqa qamrov yo'q.)
\begin{enumerate}[label=\Alph*),leftmargin=9mm]
\item 7
\item 6
\item 5
\item 18
\end{enumerate}
\vspace{2mm}\hrule height 0.2pt \vspace{4mm}
\needspace{5\baselineskip}
\hypertarget{q405}{\textbf{\color{ink}405.}}  {\color{muted}\scriptsize Y3 \textbullet\ daraja 4 \textbullet\ DB 89bde414}
\par
Quyidagi amallarni to'g'ri tartibda bajaring: 1) tr ichidagi katak teglarni aniqlash; 2) Jadvalning mantiqiy ustunlarini hisoblash; 3) colspan/rowspan qiymatini o'qish.
\begin{enumerate}[label=\Alph*),leftmargin=9mm]
\item tr ichidagi katak teglarni aniqlash
\item Jadvalning mantiqiy ustunlarini hisoblash
\item colspan/rowspan qiymatini o'qish
\end{enumerate}
\vspace{2mm}\hrule height 0.2pt \vspace{4mm}
\needspace{5\baselineskip}
\hypertarget{q406}{\textbf{\color{ink}406.}}  {\color{muted}\scriptsize Y2 \textbullet\ daraja 4 \textbullet\ DB 9cf9960d}
\par
Jadval teglarini ularning vazifasi bilan moslang.
\textbf{\footnotesize Chap ustun:}
\begin{enumerate}[label=\arabic*),leftmargin=9mm]
\item \texttt{\textless{}table\textgreater{}}
\item \texttt{\textless{}tr\textgreater{}}
\item \texttt{\textless{}th\textgreater{}}
\item \texttt{\textless{}td\textgreater{}}
\end{enumerate}
\textbf{\footnotesize O'ng ustun:}
\begin{enumerate}[label=\Alph*),leftmargin=9mm]
\item Sarlavha katagi
\item Jadval qatori
\item Jadval konteyneri
\item Oddiy ma'lumot katagi
\end{enumerate}
\vspace{2mm}\hrule height 0.2pt \vspace{4mm}
\needspace{5\baselineskip}
\hypertarget{q407}{\textbf{\color{ink}407.}}  {\color{muted}\scriptsize Y1 \textbullet\ daraja 4 \textbullet\ DB c1c677c6}
\par
HTML \textless{}input\textgreater{} uchun zamonaviy sintaksis qaysi?
\begin{enumerate}[label=\Alph*),leftmargin=9mm]
\item \textless{}input\textgreater{}\textless{}p\textgreater{}\textless{}/input\textgreater{} shart
\item \textless{}input name="a"\textgreater{}\textless{}/input\textgreater{} doim majburiy
\item \textless{}input src="a"\textgreater{} har doim rasm
\item \textless{}input name="a"\textgreater{} (yopuvchi teg shart emas)
\end{enumerate}
\vspace{2mm}\hrule height 0.2pt \vspace{4mm}
\needspace{5\baselineskip}
\hypertarget{q408}{\textbf{\color{ink}408.}}  {\color{muted}\scriptsize Y2 \textbullet\ daraja 4 \textbullet\ DB c3f03dd4}
\par
Moslashtiring: 1) ol; 2) td; 3) textarea. Ta'riflar: a) Tartiblangan ro'yxat; b) Ko'p qatorli matn maydoni; c) Oddiy jadval katagi.
\textbf{\footnotesize Chap ustun:}
\begin{enumerate}[label=\arabic*),leftmargin=9mm]
\item ol
\item td
\item textarea
\end{enumerate}
\textbf{\footnotesize O'ng ustun:}
\begin{enumerate}[label=\Alph*),leftmargin=9mm]
\item Tartiblangan ro'yxat
\item Ko'p qatorli matn maydoni
\item Oddiy jadval katagi
\end{enumerate}
\vspace{2mm}\hrule height 0.2pt \vspace{4mm}
\needspace{5\baselineskip}
\hypertarget{q409}{\textbf{\color{ink}409.}}  {\color{muted}\scriptsize Y2 \textbullet\ daraja 4 \textbullet\ DB cba700e1}
\par
Jadval atributlari va teglarini moslang.
\textbf{\footnotesize Chap ustun:}
\begin{enumerate}[label=\arabic*),leftmargin=9mm]
\item \texttt{colspan='3'}
\item \texttt{rowspan='2'}
\item \texttt{\textless{}caption\textgreater{}}
\item \texttt{\textless{}th\textgreater{}}
\end{enumerate}
\textbf{\footnotesize O'ng ustun:}
\begin{enumerate}[label=\Alph*),leftmargin=9mm]
\item Katak ikki qatorni qamrashi
\item Sarlavha katagi yaratish
\item Katak uchta mantiqiy ustunni qamrashi
\item Jadvalning nomi/sarlavhasi
\end{enumerate}
\vspace{2mm}\hrule height 0.2pt \vspace{4mm}
\needspace{5\baselineskip}
\hypertarget{q410}{\textbf{\color{ink}410.}}  {\color{muted}\scriptsize Y1 \textbullet\ daraja 4 \textbullet\ DB d25684fa}
\par
\textless{}tr\textgreater{} ichida 3 ta oddiy \textless{}td\textgreater{} va yana bitta colspan="3" katak bor. Jami mantiqiy ustunlar nechta? (Boshqa qamrov yo'q.)
\begin{enumerate}[label=\Alph*),leftmargin=9mm]
\item 5
\item 6
\item 7
\item 4
\end{enumerate}
\vspace{2mm}\hrule height 0.2pt \vspace{4mm}
\needspace{5\baselineskip}
\hypertarget{q411}{\textbf{\color{ink}411.}}  {\color{muted}\scriptsize Y3 \textbullet\ daraja 5 \textbullet\ DB 058370d8}
\par
Forma yorlig'ini input maydoniga ulash tartibini tiklang.
\begin{enumerate}[label=\Alph*),leftmargin=9mm]
\item \texttt{\textless{}input id='login' name='username'\textgreater{}} ni yarating
\item Label bosilganda tegishli input faollashishini sinang
\item \texttt{\textless{}label for='login'\textgreater{}} bilan foydalanuvchiga matn ko'rsating
\item Maydon uchun noyob \texttt{id='login'} ni tanlang
\end{enumerate}
\vspace{2mm}\hrule height 0.2pt \vspace{4mm}
\needspace{5\baselineskip}
\hypertarget{q412}{\textbf{\color{ink}412.}}  {\color{muted}\scriptsize Y2 \textbullet\ daraja 5 \textbullet\ DB 0c735ddb}
\par
Forma input turlarini ularning foydalanuvchi amali bilan moslang.
\textbf{\footnotesize Chap ustun:}
\begin{enumerate}[label=\arabic*),leftmargin=9mm]
\item \texttt{type='file'}
\item \texttt{type='hidden'}
\item \texttt{type='reset'}
\item \texttt{type='submit'}
\end{enumerate}
\textbf{\footnotesize O'ng ustun:}
\begin{enumerate}[label=\Alph*),leftmargin=9mm]
\item Ko'rinmaydigan qiymat maydoni
\item Formani jo'natishni boshlash
\item Dastlabki forma qiymatlarini tiklash
\item Fayl tanlash
\end{enumerate}
\vspace{2mm}\hrule height 0.2pt \vspace{4mm}
\needspace{5\baselineskip}
\hypertarget{q413}{\textbf{\color{ink}413.}}  {\color{muted}\scriptsize Y1 \textbullet\ daraja 5 \textbullet\ DB 0fd0aa63}
\par
\textless{}form method="get"\textgreater{}\textless{}input name="x" value="7"\textgreater{}\textless{}input id="y" value="9"\textgreater{}\textless{}/form\textgreater{} yuborildi. Oddiy holatda qaysi kalit mavjud?
\begin{enumerate}[label=\Alph*),leftmargin=9mm]
\item x=7
\item x=7 va y=9 albatta
\item name bo'lmagani uchun hech biri
\item y=9
\end{enumerate}
\vspace{2mm}\hrule height 0.2pt \vspace{4mm}
\needspace{5\baselineskip}
\hypertarget{q414}{\textbf{\color{ink}414.}}  {\color{muted}\scriptsize Y1 \textbullet\ daraja 5 \textbullet\ DB 21f7e6b1}
\par
HTML jadvalda bitta \textless{}th colspan="2"\textgreater{} va yana bitta oddiy \textless{}td\textgreater{} bir qator ichida bo'lsa mantiqiy ustunlar soni nechta?
\begin{enumerate}[label=\Alph*),leftmargin=9mm]
\item 3
\item 1
\item 2
\item 4
\end{enumerate}
\vspace{2mm}\hrule height 0.2pt \vspace{4mm}
\needspace{5\baselineskip}
\hypertarget{q415}{\textbf{\color{ink}415.}}  {\color{muted}\scriptsize Y1 \textbullet\ daraja 5 \textbullet\ DB 2faf923a}
\par
\textless{}input name="x" value="1" disabled\textgreater{} va \textless{}input name="y" value="2" readonly\textgreater{} oddiy submit qilinsa qaysi biri ko'pincha yuboriladi?
\begin{enumerate}[label=\Alph*),leftmargin=9mm]
\item Faqat y=2
\item Ikkalasi ham albatta
\item Har doim ikkalasi ham yuborilmaydi
\item Faqat x=1
\end{enumerate}
\vspace{2mm}\hrule height 0.2pt \vspace{4mm}
\needspace{5\baselineskip}
\hypertarget{q416}{\textbf{\color{ink}416.}}  {\color{muted}\scriptsize Y1 \textbullet\ daraja 5 \textbullet\ DB 45730eae}
\par
\textless{}input type="checkbox" name="ok" value="1"\textgreater{} belgilanmagan bo'lsa, odatda nima yuboriladi?
\begin{enumerate}[label=\Alph*),leftmargin=9mm]
\item radio guruhiga aylanadi
\item ok=1 doim yuboriladi
\item ok=0 doim yuboriladi
\item ok maydoni umuman yuborilmaydi
\end{enumerate}
\vspace{2mm}\hrule height 0.2pt \vspace{4mm}
\needspace{5\baselineskip}
\hypertarget{q417}{\textbf{\color{ink}417.}}  {\color{muted}\scriptsize Y1 \textbullet\ daraja 5 \textbullet\ DB 562ff37d}
\par
Rasmdagi jadvalning yuqori chap katagi 2 qatorni egallaydi, uning yonidagi sarlavha 2 ustunni egallaydi. Mantiqiy ustunlar soni nechta?
\begin{enumerate}[label=\Alph*),leftmargin=9mm]
\item 3
\item 4
\item 5
\item 2
\end{enumerate}
\begin{quote}\footnotesize\color{muted}\textbf{Rasm:} Bir katagi ikki qatorga, sarlavhasi ikki ustunga birlashgan jadval\\
\href{https://github.com/sarvar9417/attestatsiya/blob/main/public/graphics-web-v2/html_table.svg}{Rasmni onlayn ko'rish}
\end{quote}
\vspace{2mm}\hrule height 0.2pt \vspace{4mm}
\needspace{5\baselineskip}
\hypertarget{q418}{\textbf{\color{ink}418.}}  {\color{muted}\scriptsize Y1 \textbullet\ daraja 5 \textbullet\ DB 66f7552d}
\par
HTMLda \textless{}td colspan="3"\textgreater{} 1 ta teg bo'lsa, u mantiqiy nechta ustunni egallaydi?
\begin{enumerate}[label=\Alph*),leftmargin=9mm]
\item 4
\item 3
\item 2
\item 1
\end{enumerate}
\vspace{2mm}\hrule height 0.2pt \vspace{4mm}
\needspace{5\baselineskip}
\hypertarget{q419}{\textbf{\color{ink}419.}}  {\color{muted}\scriptsize Y1 \textbullet\ daraja 5 \textbullet\ DB 6ca0a1bb}
\par
Jadvalda birinchi qatorda \textless{}th colspan="2"\textgreater{} bor. Qaysi xulosa aniq?
\begin{enumerate}[label=\Alph*),leftmargin=9mm]
\item Ikki alohida th tegi hosil bo'ladi
\item Ikki tr hosil bo'ladi
\item Jadvalga ikkita CSS fayl ulanadi
\item Bitta th tegi ikki mantiqiy ustunni qamraydi
\end{enumerate}
\vspace{2mm}\hrule height 0.2pt \vspace{4mm}
\needspace{5\baselineskip}
\hypertarget{q420}{\textbf{\color{ink}420.}}  {\color{muted}\scriptsize Y3 \textbullet\ daraja 5 \textbullet\ DB 71d508fa}
\par
Radio tugmalardan faqat bittasini tanlash guruhini tuzing.
\begin{enumerate}[label=\Alph*),leftmargin=9mm]
\item Ikkinchi radio inputga ham shu name, lekin boshqa value qo'shing
\item Birini tanlaganda boshqa radio belgisi bekor bo'lishini sinang
\item Guruh uchun yagona \texttt{name='level'} ni belgilang
\item Birinchi radio inputga shu name va o'z value qiymatini qo'shing
\end{enumerate}
\vspace{2mm}\hrule height 0.2pt \vspace{4mm}
\needspace{5\baselineskip}
\hypertarget{q421}{\textbf{\color{ink}421.}}  {\color{muted}\scriptsize Y1 \textbullet\ daraja 5 \textbullet\ DB 722ada40}
\par
\textless{}ol start="5"\textgreater{}\textless{}li\textgreater{}A\textless{}/li\textgreater{}\textless{}li\textgreater{}B\textless{}ul\textgreater{}\textless{}li\textgreater{}X\textless{}/li\textgreater{}\textless{}/ul\textgreater{}\textless{}/li\textgreater{}\textless{}/ol\textgreater{}da tashqi ro'yxatning ikkinchi band raqami nechchi?
\begin{enumerate}[label=\Alph*),leftmargin=9mm]
\item 1
\item 2
\item 6
\item 5
\end{enumerate}
\vspace{2mm}\hrule height 0.2pt \vspace{4mm}
\needspace{5\baselineskip}
\hypertarget{q422}{\textbf{\color{ink}422.}}  {\color{muted}\scriptsize Y1 \textbullet\ daraja 5 \textbullet\ DB 72f538c2}
\par
Bir sahifada \textless{}input type="radio" name="a"\textgreater{} va \textless{}input type="radio" name="b"\textgreater{} bor. Tanlov bog'lanishi?
\begin{enumerate}[label=\Alph*),leftmargin=9mm]
\item Faqat bittasini butun sahifada tanlash mumkin
\item Barcha radio majburiy tanlanadi
\item Ular ikki mustaqil guruh; har biri tanlanishi mumkin
\item HTMLni brauzer ochmaydi
\end{enumerate}
\vspace{2mm}\hrule height 0.2pt \vspace{4mm}
\needspace{5\baselineskip}
\hypertarget{q423}{\textbf{\color{ink}423.}}  {\color{muted}\scriptsize Y1 \textbullet\ daraja 5 \textbullet\ DB 94bb698a}
\par
POST formasi HTTP orqali yuborildi. Maxfiylik haqida to'g'ri fikr?
\begin{enumerate}[label=\Alph*),leftmargin=9mm]
\item POST har doim shifrlaydi
\item HTTP avtomatik HTTPS
\item GET POST bilan bir xil xavfsiz
\item POST o'zi shifrlamaydi; HTTPS kerak
\end{enumerate}
\vspace{2mm}\hrule height 0.2pt \vspace{4mm}
\needspace{5\baselineskip}
\hypertarget{q424}{\textbf{\color{ink}424.}}  {\color{muted}\scriptsize Y2 \textbullet\ daraja 5 \textbullet\ DB a893bd44}
\par
Forma atributlarini ularning vazifasi bilan moslang.
\textbf{\footnotesize Chap ustun:}
\begin{enumerate}[label=\arabic*),leftmargin=9mm]
\item \texttt{form action}
\item \texttt{form method}
\item \texttt{input name}
\item \texttt{input required}
\end{enumerate}
\textbf{\footnotesize O'ng ustun:}
\begin{enumerate}[label=\Alph*),leftmargin=9mm]
\item Ma'lumot jo'natiladigan manzil
\item Bo'sh maydon yuborilishini brauzerda cheklash
\item Yuborilgan ma'lumotdagi kalit nomi
\item GET yoki POST yuborish usuli
\end{enumerate}
\vspace{2mm}\hrule height 0.2pt \vspace{4mm}
\needspace{5\baselineskip}
\hypertarget{q425}{\textbf{\color{ink}425.}}  {\color{muted}\scriptsize Y1 \textbullet\ daraja 5 \textbullet\ DB ee1d69b1}
\par
\textless{}a href="docs/a.html"\textgreater{}\textless{}img src="logo.png" alt="Logo"\textgreater{}\textless{}/a\textgreater{} nimani bildiradi?
\begin{enumerate}[label=\Alph*),leftmargin=9mm]
\item src qiymati havolaning manzilini ham o'zgartiradi
\item Rasmni brauzer umuman ko'rsatmaydi
\item alt CSS sinfi hisoblanadi
\item Logo rasmi docs/a.htmlga olib boruvchi havola sifatida ishlatiladi
\end{enumerate}
\vspace{2mm}\hrule height 0.2pt \vspace{4mm}
\needspace{5\baselineskip}
\hypertarget{q426}{\textbf{\color{ink}426.}}  {\color{muted}\scriptsize Y1 \textbullet\ daraja 5 \textbullet\ DB f05dc296}
\par
HTMLda \textless{}a href="../docs/a.html"\textgreater{} havolasi nimani bildiradi?
\begin{enumerate}[label=\Alph*),leftmargin=9mm]
\item Faqat tashqi FTP serverni
\item Joriy katalogdan bir yuqori katalogga chiqib, docs/a.html ga o'tishni
\item Faqat sayt ildizidagi docsni
\item Faqat shu katalogdagi docs papkasini
\end{enumerate}
\vspace{2mm}\hrule height 0.2pt \vspace{4mm}
\subsubsection{CSS stillari - nazariya V2}
\needspace{5\baselineskip}
\hypertarget{q427}{\textbf{\color{ink}427.}}  {\color{muted}\scriptsize Y1 \textbullet\ daraja 1 \textbullet\ DB 14550edd}
\par
CSSdagi \#title selektori nimani tanlaydi?
\begin{enumerate}[label=\Alph*),leftmargin=9mm]
\item Har qanday h1 ni
\item class="title" bo'lgan barcha elementni
\item id="title" bo'lgan elementni
\item Faqat title tegini
\end{enumerate}
\vspace{2mm}\hrule height 0.2pt \vspace{4mm}
\needspace{5\baselineskip}
\hypertarget{q428}{\textbf{\color{ink}428.}}  {\color{muted}\scriptsize Y1 \textbullet\ daraja 1 \textbullet\ DB 57e39605}
\par
CSSda elementni tanlash qismi nima deyiladi?
\begin{enumerate}[label=\Alph*),leftmargin=9mm]
\item Attribyut qiymati
\item Qator indeksi
\item Kodek
\item Selektor
\end{enumerate}
\vspace{2mm}\hrule height 0.2pt \vspace{4mm}
\needspace{5\baselineskip}
\hypertarget{q429}{\textbf{\color{ink}429.}}  {\color{muted}\scriptsize Y1 \textbullet\ daraja 1 \textbullet\ DB 68876dcc}
\par
CSSda fon rangi uchun xossa qaysi?
\begin{enumerate}[label=\Alph*),leftmargin=9mm]
\item text-align
\item src
\item background-color
\item line-height
\end{enumerate}
\vspace{2mm}\hrule height 0.2pt \vspace{4mm}
\needspace{5\baselineskip}
\hypertarget{q430}{\textbf{\color{ink}430.}}  {\color{muted}\scriptsize Y1 \textbullet\ daraja 1 \textbullet\ DB 9a586ab7}
\par
CSSda border va content oralig'idagi ichki masofa nima?
\begin{enumerate}[label=\Alph*),leftmargin=9mm]
\item opacity
\item margin
\item padding
\item gap
\end{enumerate}
\vspace{2mm}\hrule height 0.2pt \vspace{4mm}
\needspace{5\baselineskip}
\hypertarget{q431}{\textbf{\color{ink}431.}}  {\color{muted}\scriptsize Y1 \textbullet\ daraja 1 \textbullet\ DB 9fe84b30}
\par
Class selektori qanday belgidan boshlanadi?
\begin{enumerate}[label=\Alph*),leftmargin=9mm]
\item Nuqta (.)
\item Panjara (\#)
\item Yulduz (*)
\item Ikki nuqta (:)
\end{enumerate}
\vspace{2mm}\hrule height 0.2pt \vspace{4mm}
\needspace{5\baselineskip}
\hypertarget{q432}{\textbf{\color{ink}432.}}  {\color{muted}\scriptsize Y1 \textbullet\ daraja 1 \textbullet\ DB b51098db}
\par
CSSda matnni markazga tekislash qoidasi qaysi?
\begin{enumerate}[label=\Alph*),leftmargin=9mm]
\item text-align: center
\item padding: center
\item font-weight: center
\item color: center
\end{enumerate}
\vspace{2mm}\hrule height 0.2pt \vspace{4mm}
\needspace{5\baselineskip}
\hypertarget{q433}{\textbf{\color{ink}433.}}  {\color{muted}\scriptsize Y1 \textbullet\ daraja 1 \textbullet\ DB b969d515}
\par
ID selektori qanday belgidan boshlanadi?
\begin{enumerate}[label=\Alph*),leftmargin=9mm]
\item Panjara (\#)
\item Vergul (,)
\item Slash (/)
\item Nuqta (.)
\end{enumerate}
\vspace{2mm}\hrule height 0.2pt \vspace{4mm}
\needspace{5\baselineskip}
\hypertarget{q434}{\textbf{\color{ink}434.}}  {\color{muted}\scriptsize Y1 \textbullet\ daraja 1 \textbullet\ DB c3bca99a}
\par
CSSdagi .note selektori nimani tanlaydi?
\begin{enumerate}[label=\Alph*),leftmargin=9mm]
\item Faqat tashqi faylni
\item id="note" ni albatta
\item Faqat \textless{}note\textgreater{} tegini
\item class="note" bilan belgilangan elementlarni
\end{enumerate}
\vspace{2mm}\hrule height 0.2pt \vspace{4mm}
\needspace{5\baselineskip}
\hypertarget{q435}{\textbf{\color{ink}435.}}  {\color{muted}\scriptsize Y1 \textbullet\ daraja 1 \textbullet\ DB d6b00db3}
\par
CSSda shrift o'lchami qanday belgilanadi?
\begin{enumerate}[label=\Alph*),leftmargin=9mm]
\item type
\item text-height
\item font-size
\item font-color
\end{enumerate}
\vspace{2mm}\hrule height 0.2pt \vspace{4mm}
\needspace{5\baselineskip}
\hypertarget{q436}{\textbf{\color{ink}436.}}  {\color{muted}\scriptsize Y1 \textbullet\ daraja 1 \textbullet\ DB d71e3233}
\par
Jadval chiziqlarini birlashtirishga mos xossa?
\begin{enumerate}[label=\Alph*),leftmargin=9mm]
\item line-height
\item border-style-color
\item object-fit
\item border-collapse
\end{enumerate}
\vspace{2mm}\hrule height 0.2pt \vspace{4mm}
\needspace{5\baselineskip}
\hypertarget{q437}{\textbf{\color{ink}437.}}  {\color{muted}\scriptsize Y1 \textbullet\ daraja 1 \textbullet\ DB da30acd7}
\par
CSS nimaning qisqartmasi?
\begin{enumerate}[label=\Alph*),leftmargin=9mm]
\item Cascading Style Sheets
\item Coded Style Search
\item Computer Script Source
\item Creative System Syntax
\end{enumerate}
\vspace{2mm}\hrule height 0.2pt \vspace{4mm}
\needspace{5\baselineskip}
\hypertarget{q438}{\textbf{\color{ink}438.}}  {\color{muted}\scriptsize Y1 \textbullet\ daraja 1 \textbullet\ DB ddb56b39}
\par
CSSda matn rangi uchun xossa qaysi?
\begin{enumerate}[label=\Alph*),leftmargin=9mm]
\item font-weight
\item height
\item color
\item background-color
\end{enumerate}
\vspace{2mm}\hrule height 0.2pt \vspace{4mm}
\needspace{5\baselineskip}
\hypertarget{q439}{\textbf{\color{ink}439.}}  {\color{muted}\scriptsize Y1 \textbullet\ daraja 2 \textbullet\ DB 1c4ef4e4}
\par
Tashqi CSS qanday ulanadi?
\begin{enumerate}[label=\Alph*),leftmargin=9mm]
\item form method bilan
\item link rel="stylesheet" va href bilan
\item img va src bilan
\item table colspan bilan
\end{enumerate}
\vspace{2mm}\hrule height 0.2pt \vspace{4mm}
\needspace{5\baselineskip}
\hypertarget{q440}{\textbf{\color{ink}440.}}  {\color{muted}\scriptsize Y1 \textbullet\ daraja 2 \textbullet\ DB 385bb797}
\par
CSS'da \texttt{\textless{}img\textgreater{}} elementining belgilangan qutisi tasvirdan boshqa tomonlar nisbatiga ega. Tasvir butunlay ko'rinsin, qirqilmasin va dastlabki tomonlar nisbati saqlansin. Qaysi \texttt{object-fit} qiymati mos?
\begin{enumerate}[label=\Alph*),leftmargin=9mm]
\item contain
\item fill
\item none
\item cover
\end{enumerate}
\vspace{2mm}\hrule height 0.2pt \vspace{4mm}
\needspace{5\baselineskip}
\hypertarget{q441}{\textbf{\color{ink}441.}}  {\color{muted}\scriptsize Y1 \textbullet\ daraja 2 \textbullet\ DB 436c3195}
\par
content-box va border-box farqi nima?
\begin{enumerate}[label=\Alph*),leftmargin=9mm]
\item Ikkalasi bir xil
\item content-box CSS selektoridir
\item Birinchi width faqat content, ikkinchi padding va borderni ham o'z ichiga oladi
\item border-box marginni width ichiga oladi
\end{enumerate}
\vspace{2mm}\hrule height 0.2pt \vspace{4mm}
\needspace{5\baselineskip}
\hypertarget{q442}{\textbf{\color{ink}442.}}  {\color{muted}\scriptsize Y1 \textbullet\ daraja 2 \textbullet\ DB 46b3e813}
\par
CSSda bir xil spetsifiklikdagi keyingi qoida qanday?
\begin{enumerate}[label=\Alph*),leftmargin=9mm]
\item Doim birinchisi yutadi
\item Qoidalar avtomatik o'chiriladi
\item Faqat HTML tegi hal qiladi
\item Odatda keyingisi kuchga kiradi
\end{enumerate}
\vspace{2mm}\hrule height 0.2pt \vspace{4mm}
\needspace{5\baselineskip}
\hypertarget{q443}{\textbf{\color{ink}443.}}  {\color{muted}\scriptsize Y1 \textbullet\ daraja 2 \textbullet\ DB 50b7718a}
\par
Jadval katak ichki masofasi uchun qaysi xossa?
\begin{enumerate}[label=\Alph*),leftmargin=9mm]
\item padding
\item list-style
\item margin
\item opacity
\end{enumerate}
\vspace{2mm}\hrule height 0.2pt \vspace{4mm}
\needspace{5\baselineskip}
\hypertarget{q444}{\textbf{\color{ink}444.}}  {\color{muted}\scriptsize Y1 \textbullet\ daraja 2 \textbullet\ DB 52435cb0}
\par
HTMLda inline CSS qanday beriladi?
\begin{enumerate}[label=\Alph*),leftmargin=9mm]
\item Faqat \textless{}link\textgreater{} orqali
\item Elementning style atributida
\item Faqat \textless{}meta\textgreater{} orqali
\item Faqat alohida CSS faylda
\end{enumerate}
\vspace{2mm}\hrule height 0.2pt \vspace{4mm}
\needspace{5\baselineskip}
\hypertarget{q445}{\textbf{\color{ink}445.}}  {\color{muted}\scriptsize Y1 \textbullet\ daraja 2 \textbullet\ DB 5ce9f3ca}
\par
color va background-color farqi nima?
\begin{enumerate}[label=\Alph*),leftmargin=9mm]
\item Ikkisi faqat shrift turi
\item Birinchisi matn, ikkinchisi fon rangi
\item background-color rasm manzili
\item color faqat chiziq qalinligi
\end{enumerate}
\vspace{2mm}\hrule height 0.2pt \vspace{4mm}
\needspace{5\baselineskip}
\hypertarget{q446}{\textbf{\color{ink}446.}}  {\color{muted}\scriptsize Y1 \textbullet\ daraja 2 \textbullet\ DB 83fba835}
\par
CSSda \#box nimani tanlaydi?
\begin{enumerate}[label=\Alph*),leftmargin=9mm]
\item Barcha box teglarni
\item Barcha qatorlarni
\item class="box" elementini
\item id="box" elementini
\end{enumerate}
\vspace{2mm}\hrule height 0.2pt \vspace{4mm}
\needspace{5\baselineskip}
\hypertarget{q447}{\textbf{\color{ink}447.}}  {\color{muted}\scriptsize Y1 \textbullet\ daraja 2 \textbullet\ DB 8c437350}
\par
.card p va .card \textgreater{} p farqi nima?
\begin{enumerate}[label=\Alph*),leftmargin=9mm]
\item Birinchisi barcha avlod p, ikkinchisi bevosita farzand p
\item Ikkalasi classni o'chiradi
\item Birinchisi idga tegishli
\item Ikkinchisi hamma p larni tanlaydi
\end{enumerate}
\vspace{2mm}\hrule height 0.2pt \vspace{4mm}
\needspace{5\baselineskip}
\hypertarget{q448}{\textbf{\color{ink}448.}}  {\color{muted}\scriptsize Y1 \textbullet\ daraja 2 \textbullet\ DB e33d91d4}
\par
style atributi, \textless{}style\textgreater{} bloki va tashqi CSS fayl orasidagi farq qanday?
\begin{enumerate}[label=\Alph*),leftmargin=9mm]
\item External doim faqat imgga ishlaydi
\item Inline bitta elementda, internal hujjatda, external alohida .css faylda
\item Uchalasi ham bir xil HTML teg
\item Inline faqat brauzer keshini boshqaradi
\end{enumerate}
\vspace{2mm}\hrule height 0.2pt \vspace{4mm}
\needspace{5\baselineskip}
\hypertarget{q449}{\textbf{\color{ink}449.}}  {\color{muted}\scriptsize Y1 \textbullet\ daraja 2 \textbullet\ DB eaf8a0f0}
\par
CSSda display:none va visibility:hidden qanday farq qiladi?
\begin{enumerate}[label=\Alph*),leftmargin=9mm]
\item visibility doim yangi satr yaratadi
\item display faqat rang o'zgartiradi
\item Ikkalasi DOMdan o'chiradi
\item Birinchisi layoutdan olib tashlaydi, ikkinchisi joyini saqlashi mumkin
\end{enumerate}
\vspace{2mm}\hrule height 0.2pt \vspace{4mm}
\needspace{5\baselineskip}
\hypertarget{q450}{\textbf{\color{ink}450.}}  {\color{muted}\scriptsize Y1 \textbullet\ daraja 2 \textbullet\ DB faaa9e19}
\par
CSSda p bilan .p farqi nimada?
\begin{enumerate}[label=\Alph*),leftmargin=9mm]
\item p tegni, .p class="p"ni tanlaydi
\item .p idni tanlaydi
\item Ikkalasi aynan bir xil
\item p faqat rasmni tanlaydi
\end{enumerate}
\vspace{2mm}\hrule height 0.2pt \vspace{4mm}
\needspace{5\baselineskip}
\hypertarget{q451}{\textbf{\color{ink}451.}}  {\color{muted}\scriptsize Y2 \textbullet\ daraja 3 \textbullet\ DB 0c8128e1}
\par
CSSni ulash usulini koddagi o'rni bilan moslang.
\textbf{\footnotesize Chap ustun:}
\begin{enumerate}[label=\arabic*),leftmargin=9mm]
\item Inline
\item Internal
\item External
\item Media query
\end{enumerate}
\textbf{\footnotesize O'ng ustun:}
\begin{enumerate}[label=\Alph*),leftmargin=9mm]
\item Alohida CSS faylni link orqali ulash
\item Muayyan ekran shartida qoida qo'llash
\item Hujjatdagi style elementi
\item Elementning style atributi
\end{enumerate}
\vspace{2mm}\hrule height 0.2pt \vspace{4mm}
\needspace{5\baselineskip}
\hypertarget{q452}{\textbf{\color{ink}452.}}  {\color{muted}\scriptsize Y1 \textbullet\ daraja 3 \textbullet\ DB 0eef6725}
\par
Oddiy \textless{}h2 id="header"\textgreater{} uchun qaysi selektor ishlaydi?
\begin{enumerate}[label=\Alph*),leftmargin=9mm]
\item header\#note
\item h1.header
\item \#header
\item .header
\end{enumerate}
\vspace{2mm}\hrule height 0.2pt \vspace{4mm}
\needspace{5\baselineskip}
\hypertarget{q453}{\textbf{\color{ink}453.}}  {\color{muted}\scriptsize Y3 \textbullet\ daraja 3 \textbullet\ DB 13f54550}
\par
Quyidagi amallarni to'g'ri tartibda bajaring: 1) Selektorni yozish; 2) HTML class/idlarini aniqlash; 3) Deklaratsiya blokini yozish.
\begin{enumerate}[label=\Alph*),leftmargin=9mm]
\item Selektorni yozish
\item HTML class/idlarini aniqlash
\item Deklaratsiya blokini yozish
\end{enumerate}
\vspace{2mm}\hrule height 0.2pt \vspace{4mm}
\needspace{5\baselineskip}
\hypertarget{q454}{\textbf{\color{ink}454.}}  {\color{muted}\scriptsize Y1 \textbullet\ daraja 3 \textbullet\ DB 160296fa}
\par
CSS kodi: .x \{text-align:center;\} . Qaysi natija?
\begin{enumerate}[label=\Alph*),leftmargin=9mm]
\item Barcha fon qora
\item Faqat jadval yonadi
\item Rasm o'chadi
\item Mos class matni markazda
\end{enumerate}
\vspace{2mm}\hrule height 0.2pt \vspace{4mm}
\needspace{5\baselineskip}
\hypertarget{q455}{\textbf{\color{ink}455.}}  {\color{muted}\scriptsize Y1 \textbullet\ daraja 3 \textbullet\ DB 2bbc1cd5}
\par
CSS kodi: ul \{list-style-type:square;\} . Qaysi natija?
\begin{enumerate}[label=\Alph*),leftmargin=9mm]
\item Markerlar rom raqamiga o'tadi
\item Matn yonadi
\item Markerlar kvadrat bo'ladi
\item Jadval yo'qoladi
\end{enumerate}
\vspace{2mm}\hrule height 0.2pt \vspace{4mm}
\needspace{5\baselineskip}
\hypertarget{q456}{\textbf{\color{ink}456.}}  {\color{muted}\scriptsize Y3 \textbullet\ daraja 3 \textbullet\ DB 34fbfe11}
\par
Bir nechta CSS qoidasi rangga ta'sir etsa, natijani tekshirish tartibini belgilang.
\begin{enumerate}[label=\Alph*),leftmargin=9mm]
\item Elementga mos barcha selektor va deklaratsiyalarni ajrating
\item Soddalashtirilgan holatda inline/id/class/element spetsifikligini solishtiring
\item Bir xil spetsifiklikda keyingi qoidani tanlab natijani tekshiring
\item !important, kelib chiqish yoki layer sharti bor-yo'qligini tekshiring
\end{enumerate}
\vspace{2mm}\hrule height 0.2pt \vspace{4mm}
\needspace{5\baselineskip}
\hypertarget{q457}{\textbf{\color{ink}457.}}  {\color{muted}\scriptsize Y1 \textbullet\ daraja 3 \textbullet\ DB 3694240a}
\par
\textless{}div class="a"\textgreater{}\textless{}p\textgreater{}Bir\textless{}/p\textgreater{}\textless{}section\textgreater{}\textless{}p\textgreater{}Ikki\textless{}/p\textgreater{}\textless{}/section\textgreater{}\textless{}/div\textgreater{} uchun .a \textgreater{} p selektori nechta p ga ta'sir qiladi?
\begin{enumerate}[label=\Alph*),leftmargin=9mm]
\item 3
\item 1
\item 2
\item 0
\end{enumerate}
\vspace{2mm}\hrule height 0.2pt \vspace{4mm}
\needspace{5\baselineskip}
\hypertarget{q458}{\textbf{\color{ink}458.}}  {\color{muted}\scriptsize Y2 \textbullet\ daraja 3 \textbullet\ DB 991e481d}
\par
CSS selektori bilan tanlanadigan HTML elementini moslang.
\textbf{\footnotesize Chap ustun:}
\begin{enumerate}[label=\arabic*),leftmargin=9mm]
\item \texttt{p}
\item \texttt{.card}
\item \texttt{\#main}
\item \texttt{p.note}
\end{enumerate}
\textbf{\footnotesize O'ng ustun:}
\begin{enumerate}[label=\Alph*),leftmargin=9mm]
\item class='card' bo'lgan elementlar
\item Barcha p elementlari
\item class='note' bo'lgan p elementlari
\item id='main' bo'lgan element
\end{enumerate}
\vspace{2mm}\hrule height 0.2pt \vspace{4mm}
\needspace{5\baselineskip}
\hypertarget{q459}{\textbf{\color{ink}459.}}  {\color{muted}\scriptsize Y2 \textbullet\ daraja 3 \textbullet\ DB a5156c3e}
\par
Moslashtiring: 1) color; 2) padding; 3) margin. Ta'riflar: a) Ichki masofa; b) Matn rangi; c) Tashqi masofa.
\textbf{\footnotesize Chap ustun:}
\begin{enumerate}[label=\arabic*),leftmargin=9mm]
\item color
\item padding
\item margin
\end{enumerate}
\textbf{\footnotesize O'ng ustun:}
\begin{enumerate}[label=\Alph*),leftmargin=9mm]
\item Ichki masofa
\item Matn rangi
\item Tashqi masofa
\end{enumerate}
\vspace{2mm}\hrule height 0.2pt \vspace{4mm}
\needspace{5\baselineskip}
\hypertarget{q460}{\textbf{\color{ink}460.}}  {\color{muted}\scriptsize Y1 \textbullet\ daraja 3 \textbullet\ DB b6447c89}
\par
\textless{}p id="a" class="x" style="color:green"\textgreater{} uchun p\{color:blue\}, .x\{color:red\}, \#a\{color:orange\} bo'lsa, important yo'q. Rang nima?
\begin{enumerate}[label=\Alph*),leftmargin=9mm]
\item Qizil
\item To'q sariq
\item Yashil
\item Ko'k
\end{enumerate}
\vspace{2mm}\hrule height 0.2pt \vspace{4mm}
\needspace{5\baselineskip}
\hypertarget{q461}{\textbf{\color{ink}461.}}  {\color{muted}\scriptsize Y1 \textbullet\ daraja 3 \textbullet\ DB b6cc553c}
\par
CSS kodi: p \{color:red;\} . Qaysi natija?
\begin{enumerate}[label=\Alph*),leftmargin=9mm]
\item Fon qizil
\item Matn ko'k rangda
\item Matn markazda
\item Matn qizil rangda
\end{enumerate}
\vspace{2mm}\hrule height 0.2pt \vspace{4mm}
\needspace{5\baselineskip}
\hypertarget{q462}{\textbf{\color{ink}462.}}  {\color{muted}\scriptsize Y3 \textbullet\ daraja 3 \textbullet\ DB c60954c5}
\par
Yangi CSS qoidasi tuzib natijasini tekshirish tartibini belgilang.
\begin{enumerate}[label=\Alph*),leftmargin=9mm]
\item \texttt{property: value;} deklaratsiyalarini qo'shib blokni yoping
\item Mos selektorni yozib deklaratsiya blokini oching
\item Uslub beriladigan HTML element yoki classni aniqlang
\item Brauzerda tanlangan element ko'rinishini tekshiring
\end{enumerate}
\vspace{2mm}\hrule height 0.2pt \vspace{4mm}
\needspace{5\baselineskip}
\hypertarget{q463}{\textbf{\color{ink}463.}}  {\color{muted}\scriptsize Y1 \textbullet\ daraja 3 \textbullet\ DB d4bcd114}
\par
Faqat id="title" element rangini qizil qilish?
\begin{enumerate}[label=\Alph*),leftmargin=9mm]
\item .title \{color:red;\}
\item title \{background:red;\}
\item \#title \{src:red;\}
\item \#title \{color:red;\}
\end{enumerate}
\vspace{2mm}\hrule height 0.2pt \vspace{4mm}
\needspace{5\baselineskip}
\hypertarget{q464}{\textbf{\color{ink}464.}}  {\color{muted}\scriptsize Y1 \textbullet\ daraja 3 \textbullet\ DB e5d1537f}
\par
CSS kodi: th \{background-color:yellow;\} . Qaysi natija?
\begin{enumerate}[label=\Alph*),leftmargin=9mm]
\item th kataklarning foni sariq
\item Jadval ustunlari birlashadi
\item Barcha matn sariq
\item Har bir rasm sariq
\end{enumerate}
\vspace{2mm}\hrule height 0.2pt \vspace{4mm}
\needspace{5\baselineskip}
\hypertarget{q465}{\textbf{\color{ink}465.}}  {\color{muted}\scriptsize Y1 \textbullet\ daraja 3 \textbullet\ DB edd7d493}
\par
Barcha p matnini 20px qilish?
\begin{enumerate}[label=\Alph*),leftmargin=9mm]
\item \textless{}p size="20px"\textgreater{}
\item p \{size:20px;\}
\item p \{font-size:20px;\}
\item p \{font:20;\}
\end{enumerate}
\vspace{2mm}\hrule height 0.2pt \vspace{4mm}
\needspace{5\baselineskip}
\hypertarget{q466}{\textbf{\color{ink}466.}}  {\color{muted}\scriptsize Y1 \textbullet\ daraja 3 \textbullet\ DB eec29775}
\par
Rasmda DOM daraxtida .card \textgreater{} p selektori qaysi belgilangan elementga mos?
\begin{enumerate}[label=\Alph*),leftmargin=9mm]
\item B
\item A va B
\item Hech biri
\item A
\end{enumerate}
\begin{quote}\footnotesize\color{muted}\textbf{Rasm:} Bevosita p(A) va ichki section p(B) daraxti\\
\href{https://github.com/sarvar9417/attestatsiya/blob/main/public/graphics-web-v2/dom_tree.svg}{Rasmni onlayn ko'rish}
\end{quote}
\vspace{2mm}\hrule height 0.2pt \vspace{4mm}
\needspace{5\baselineskip}
\hypertarget{q467}{\textbf{\color{ink}467.}}  {\color{muted}\scriptsize Y3 \textbullet\ daraja 3 \textbullet\ DB f3557a83}
\par
Bitta CSS faylni bir nechta HTML sahifaga ulash tartibini tuzing.
\begin{enumerate}[label=\Alph*),leftmargin=9mm]
\item Har bir sahifaning head qismiga \texttt{link rel='stylesheet'} yozing
\item Sahifalarni ochib umumiy uslub ishlashini tekshiring
\item linkning href qiymatini \texttt{style.css} manziliga moslang
\item Sayt uchun umumiy uslublarni \texttt{style.css} faylida yozing
\end{enumerate}
\vspace{2mm}\hrule height 0.2pt \vspace{4mm}
\needspace{5\baselineskip}
\hypertarget{q468}{\textbf{\color{ink}468.}}  {\color{muted}\scriptsize Y2 \textbullet\ daraja 3 \textbullet\ DB f7ed1fe3}
\par
Kombinatsiyalangan CSS selektorlarining ma'nolarini moslang.
\textbf{\footnotesize Chap ustun:}
\begin{enumerate}[label=\arabic*),leftmargin=9mm]
\item \texttt{.box p}
\item \texttt{.box \textgreater{} p}
\item \texttt{h1, h2}
\item \texttt{tr:nth-child(even)}
\end{enumerate}
\textbf{\footnotesize O'ng ustun:}
\begin{enumerate}[label=\Alph*),leftmargin=9mm]
\item box ichidagi barcha p avlodlar
\item H1 va H2 elementlarining ikkalasi
\item O'z ota elementidagi juft tartibli trlar
\item boxning bevosita p farzandlari
\end{enumerate}
\vspace{2mm}\hrule height 0.2pt \vspace{4mm}
\needspace{5\baselineskip}
\hypertarget{q469}{\textbf{\color{ink}469.}}  {\color{muted}\scriptsize Y1 \textbullet\ daraja 4 \textbullet\ DB 27c76e41}
\par
CSSda \textless{}p style="color:red"\textgreater{} va p\{color:blue\} bo'lsa, odatdagi kaskadda natija?
\begin{enumerate}[label=\Alph*),leftmargin=9mm]
\item Ko'k
\item Qizil
\item Qaytadan hisoblab bo'lmaydi
\item Qora
\end{enumerate}
\vspace{2mm}\hrule height 0.2pt \vspace{4mm}
\needspace{5\baselineskip}
\hypertarget{q470}{\textbf{\color{ink}470.}}  {\color{muted}\scriptsize Y3 \textbullet\ daraja 4 \textbullet\ DB 29bddd75}
\par
200px content, 10px padding, 2px border, 5px marginli qutining tashqi enini topish tartibini tuzing.
\begin{enumerate}[label=\Alph*),leftmargin=9mm]
\item Content kengligini 200px deb oling
\item Ikki yon marginni qo'shib jami 234px aniqlang
\item box-sizing \texttt{content-box} ekanini tasdiqlang
\item Ikki yon padding va borderni qo'shib 224px toping
\end{enumerate}
\vspace{2mm}\hrule height 0.2pt \vspace{4mm}
\needspace{5\baselineskip}
\hypertarget{q471}{\textbf{\color{ink}471.}}  {\color{muted}\scriptsize Y3 \textbullet\ daraja 4 \textbullet\ DB 2faf9478}
\par
Quyidagi amallarni to'g'ri tartibda bajaring: 1) Ustuvor deklaratsiyani tanlash; 2) Mos keluvchi selektorlarni belgilash; 3) Spetsifiklikni solishtirish.
\begin{enumerate}[label=\Alph*),leftmargin=9mm]
\item Ustuvor deklaratsiyani tanlash
\item Mos keluvchi selektorlarni belgilash
\item Spetsifiklikni solishtirish
\end{enumerate}
\vspace{2mm}\hrule height 0.2pt \vspace{4mm}
\needspace{5\baselineskip}
\hypertarget{q472}{\textbf{\color{ink}472.}}  {\color{muted}\scriptsize Y1 \textbullet\ daraja 4 \textbullet\ DB 346437d0}
\par
CSSda border-collapse:collapse qaysi obyektga ko'proq ta'sir qiladi?
\begin{enumerate}[label=\Alph*),leftmargin=9mm]
\item Jadval chegaralarini tutashtirishga
\item Fayl siqish tezligiga
\item Video subtitriga
\item HTML input typega
\end{enumerate}
\vspace{2mm}\hrule height 0.2pt \vspace{4mm}
\needspace{5\baselineskip}
\hypertarget{q473}{\textbf{\color{ink}473.}}  {\color{muted}\scriptsize Y1 \textbullet\ daraja 4 \textbullet\ DB 6152b0fe}
\par
content-box: width=300px, padding chap/o'ng 15px, border chap/o'ng 3px, margin chap/o'ng 10px. Oddiy modelda jami gorizontal egallash?
\begin{enumerate}[label=\Alph*),leftmargin=9mm]
\item 362px
\item 330px
\item 356px
\item 300px
\end{enumerate}
\vspace{2mm}\hrule height 0.2pt \vspace{4mm}
\needspace{5\baselineskip}
\hypertarget{q474}{\textbf{\color{ink}474.}}  {\color{muted}\scriptsize Y2 \textbullet\ daraja 4 \textbullet\ DB 7b933718}
\par
Tasvirga beriladigan CSS deklaratsiyalarini natijasi bilan moslang.
\textbf{\footnotesize Chap ustun:}
\begin{enumerate}[label=\arabic*),leftmargin=9mm]
\item \texttt{max-width:100\%}
\item \texttt{height:auto}
\item \texttt{object-fit:cover}
\item \texttt{object-fit:contain}
\end{enumerate}
\textbf{\footnotesize O'ng ustun:}
\begin{enumerate}[label=\Alph*),leftmargin=9mm]
\item Rasmga ota kengligidan oshmaslik cheklovi
\item To'liq tasvir sig'ishi uchun bo'sh joy qoldirishi mumkin
\item Qutini to'ldirish uchun tasvir chetini kesishi mumkin
\item Avtomatik nisbatga mos balandlik
\end{enumerate}
\vspace{2mm}\hrule height 0.2pt \vspace{4mm}
\needspace{5\baselineskip}
\hypertarget{q475}{\textbf{\color{ink}475.}}  {\color{muted}\scriptsize Y1 \textbullet\ daraja 4 \textbullet\ DB 850a679a}
\par
content-box: width=240px, padding chap/o'ng 12px, border chap/o'ng 4px, margin chap/o'ng 6px. Oddiy modelda jami gorizontal egallash?
\begin{enumerate}[label=\Alph*),leftmargin=9mm]
\item 264px
\item 292px
\item 284px
\item 240px
\end{enumerate}
\vspace{2mm}\hrule height 0.2pt \vspace{4mm}
\needspace{5\baselineskip}
\hypertarget{q476}{\textbf{\color{ink}476.}}  {\color{muted}\scriptsize Y1 \textbullet\ daraja 4 \textbullet\ DB b1f5a21b}
\par
CSS: \#a\{color:red\} .b\{color:green\} va \textless{}p id="a" class="b"\textgreater{}. Odatdagi natija?
\begin{enumerate}[label=\Alph*),leftmargin=9mm]
\item Yashil
\item Ko'k
\item Qizil
\item Matn ko'rinmaydi
\end{enumerate}
\vspace{2mm}\hrule height 0.2pt \vspace{4mm}
\needspace{5\baselineskip}
\hypertarget{q477}{\textbf{\color{ink}477.}}  {\color{muted}\scriptsize Y1 \textbullet\ daraja 4 \textbullet\ DB b437c539}
\par
content-box: width=100px, padding chap/o'ng 20px, border chap/o'ng 5px, margin chap/o'ng 0px. Oddiy modelda jami gorizontal egallash?
\begin{enumerate}[label=\Alph*),leftmargin=9mm]
\item 140px
\item 160px
\item 100px
\item 150px
\end{enumerate}
\vspace{2mm}\hrule height 0.2pt \vspace{4mm}
\needspace{5\baselineskip}
\hypertarget{q478}{\textbf{\color{ink}478.}}  {\color{muted}\scriptsize Y2 \textbullet\ daraja 4 \textbullet\ DB b4b1eaba}
\par
Standart box model qismlarini ma'nolari bilan moslang.
\textbf{\footnotesize Chap ustun:}
\begin{enumerate}[label=\arabic*),leftmargin=9mm]
\item content
\item padding
\item border
\item margin
\end{enumerate}
\textbf{\footnotesize O'ng ustun:}
\begin{enumerate}[label=\Alph*),leftmargin=9mm]
\item Mazmun bilan border orasidagi masofa
\item Border tashqarisidagi bo'sh masofa
\item Matn yoki rasm joylashgan ichki mazmun sohasi
\item Qutining chegara chizig'i
\end{enumerate}
\vspace{2mm}\hrule height 0.2pt \vspace{4mm}
\needspace{5\baselineskip}
\hypertarget{q479}{\textbf{\color{ink}479.}}  {\color{muted}\scriptsize Y2 \textbullet\ daraja 4 \textbullet\ DB bdd38405}
\par
CSS matn xossalari va natijalarini moslang.
\textbf{\footnotesize Chap ustun:}
\begin{enumerate}[label=\arabic*),leftmargin=9mm]
\item \texttt{color}
\item \texttt{background-color}
\item \texttt{font-weight}
\item \texttt{text-align}
\end{enumerate}
\textbf{\footnotesize O'ng ustun:}
\begin{enumerate}[label=\Alph*),leftmargin=9mm]
\item Shrift qalinligi
\item Element fon rangi
\item Matn rangi
\item Matnning gorizontal tekislanishi
\end{enumerate}
\vspace{2mm}\hrule height 0.2pt \vspace{4mm}
\needspace{5\baselineskip}
\hypertarget{q480}{\textbf{\color{ink}480.}}  {\color{muted}\scriptsize Y3 \textbullet\ daraja 4 \textbullet\ DB bf31db47}
\par
Rasmni nisbatini saqlagan holda sig'imli ko'rsatish tartibini tiklang.
\begin{enumerate}[label=\Alph*),leftmargin=9mm]
\item HTMLdagi img elementini aniqlang
\item CSSda unga \texttt{max-width:100\%} bering
\item \texttt{height:auto} xossasini qo'shing
\item Kichik va katta konteynerda rasm nisbatini tekshiring
\end{enumerate}
\vspace{2mm}\hrule height 0.2pt \vspace{4mm}
\needspace{5\baselineskip}
\hypertarget{q481}{\textbf{\color{ink}481.}}  {\color{muted}\scriptsize Y1 \textbullet\ daraja 4 \textbullet\ DB c532f988}
\par
content-box: width=200px, padding chap/o'ng 10px, border chap/o'ng 2px, margin chap/o'ng 5px. Oddiy modelda jami gorizontal egallash?
\begin{enumerate}[label=\Alph*),leftmargin=9mm]
\item 200px
\item 238px
\item 234px
\item 220px
\end{enumerate}
\vspace{2mm}\hrule height 0.2pt \vspace{4mm}
\needspace{5\baselineskip}
\hypertarget{q482}{\textbf{\color{ink}482.}}  {\color{muted}\scriptsize Y1 \textbullet\ daraja 4 \textbullet\ DB cd682e28}
\par
content-box: width=120px, padding chap/o'ng 8px, border chap/o'ng 1px, margin chap/o'ng 0px. Oddiy modelda jami gorizontal egallash?
\begin{enumerate}[label=\Alph*),leftmargin=9mm]
\item 140px
\item 138px
\item 136px
\item 120px
\end{enumerate}
\vspace{2mm}\hrule height 0.2pt \vspace{4mm}
\needspace{5\baselineskip}
\hypertarget{q483}{\textbf{\color{ink}483.}}  {\color{muted}\scriptsize Y3 \textbullet\ daraja 4 \textbullet\ DB d39ca95f}
\par
CSS orqali HTML jadvalini bezashning ketma-ketligini belgilang.
\begin{enumerate}[label=\Alph*),leftmargin=9mm]
\item \texttt{th,td} uchun border va padding yozing
\item Brauzerda chegaralar va katak masofasini tekshiring
\item \texttt{table\{border-collapse:collapse\}} qoidasini belgilang
\item HTML table, tr, th va td tuzilmasini to'g'ri tayyorlang
\end{enumerate}
\vspace{2mm}\hrule height 0.2pt \vspace{4mm}
\needspace{5\baselineskip}
\hypertarget{q484}{\textbf{\color{ink}484.}}  {\color{muted}\scriptsize Y1 \textbullet\ daraja 4 \textbullet\ DB db89444a}
\par
content-box: width=80px, padding chap/o'ng 6px, border chap/o'ng 2px, margin chap/o'ng 2px. Oddiy modelda jami gorizontal egallash?
\begin{enumerate}[label=\Alph*),leftmargin=9mm]
\item 104px
\item 100px
\item 80px
\item 92px
\end{enumerate}
\vspace{2mm}\hrule height 0.2pt \vspace{4mm}
\needspace{5\baselineskip}
\hypertarget{q485}{\textbf{\color{ink}485.}}  {\color{muted}\scriptsize Y1 \textbullet\ daraja 4 \textbullet\ DB e59ca7c2}
\par
content-box: width=150px, padding chap/o'ng 5px, border chap/o'ng 2px, margin chap/o'ng 8px. Oddiy modelda jami gorizontal egallash?
\begin{enumerate}[label=\Alph*),leftmargin=9mm]
\item 184px
\item 150px
\item 180px
\item 160px
\end{enumerate}
\vspace{2mm}\hrule height 0.2pt \vspace{4mm}
\needspace{5\baselineskip}
\hypertarget{q486}{\textbf{\color{ink}486.}}  {\color{muted}\scriptsize Y2 \textbullet\ daraja 4 \textbullet\ DB eabbb227}
\par
Moslashtiring: 1) font-weight; 2) border-collapse; 3) list-style-type. Ta'riflar: a) Ro'yxat markeri; b) Shrift qalinligi; c) Jadval chegaralari.
\textbf{\footnotesize Chap ustun:}
\begin{enumerate}[label=\arabic*),leftmargin=9mm]
\item font-weight
\item border-collapse
\item list-style-type
\end{enumerate}
\textbf{\footnotesize O'ng ustun:}
\begin{enumerate}[label=\Alph*),leftmargin=9mm]
\item Ro'yxat markeri
\item Shrift qalinligi
\item Jadval chegaralari
\end{enumerate}
\vspace{2mm}\hrule height 0.2pt \vspace{4mm}
\needspace{5\baselineskip}
\hypertarget{q487}{\textbf{\color{ink}487.}}  {\color{muted}\scriptsize Y3 \textbullet\ daraja 5 \textbullet\ DB 235dc2ec}
\par
600px va undan tor ko'rish maydoni uchun kartalarni ustun qilish tartibini tuzing.
\begin{enumerate}[label=\Alph*),leftmargin=9mm]
\item Kartalar konteynerini \texttt{.cards} sinfi bilan belgilang
\item Keng va tor viewportda qator/ustun natijasini tekshiring
\item \texttt{@media(max-width:600px)} ichida \texttt{flex-direction:column} bering
\item Asosiy qoidada \texttt{display:flex} va kerakli gapni yozing
\end{enumerate}
\vspace{2mm}\hrule height 0.2pt \vspace{4mm}
\needspace{5\baselineskip}
\hypertarget{q488}{\textbf{\color{ink}488.}}  {\color{muted}\scriptsize Y1 \textbullet\ daraja 5 \textbullet\ DB 2e3b387c}
\par
CSS: table, th \{background-color:lightblue\} tr \{background-color:red\} td \{color:green\}. \textless{}td\textgreater{}A\textless{}/td\textgreater{} ichidagi A matni qanday rangda bo'ladi? (Boshqa qoidalar yo'q.)
\begin{enumerate}[label=\Alph*),leftmargin=9mm]
\item och ko'k
\item yashil
\item qizil
\item qora
\end{enumerate}
\vspace{2mm}\hrule height 0.2pt \vspace{4mm}
\needspace{5\baselineskip}
\hypertarget{q489}{\textbf{\color{ink}489.}}  {\color{muted}\scriptsize Y1 \textbullet\ daraja 5 \textbullet\ DB 463454f6}
\par
\textless{}p class="note"\textgreater{} uchun p\{color:red\}, .note\{color:blue\} yozilgan. Important yo'q. Rang?
\begin{enumerate}[label=\Alph*),leftmargin=9mm]
\item Qizil
\item Sariq
\item Oq
\item Ko'k
\end{enumerate}
\vspace{2mm}\hrule height 0.2pt \vspace{4mm}
\needspace{5\baselineskip}
\hypertarget{q490}{\textbf{\color{ink}490.}}  {\color{muted}\scriptsize Y2 \textbullet\ daraja 5 \textbullet\ DB 52e866f8}
\par
Bir xil CSS origin/importance holatida ustuvorlikni misol bilan moslang.
\textbf{\footnotesize Chap ustun:}
\begin{enumerate}[label=\arabic*),leftmargin=9mm]
\item \texttt{p\{color:blue\}}
\item \texttt{.note\{color:green\}}
\item \texttt{\#msg\{color:red\}}
\item \texttt{style='color:purple'}
\end{enumerate}
\textbf{\footnotesize O'ng ustun:}
\begin{enumerate}[label=\Alph*),leftmargin=9mm]
\item Oddiy element selektori
\item Classdan kuchliroq id selektori
\item Elementdan kuchliroq class selektori
\item Oddiy muallif qoidalaridan ko'pincha ustun inline stil
\end{enumerate}
\vspace{2mm}\hrule height 0.2pt \vspace{4mm}
\needspace{5\baselineskip}
\hypertarget{q491}{\textbf{\color{ink}491.}}  {\color{muted}\scriptsize Y1 \textbullet\ daraja 5 \textbullet\ DB 540dee5a}
\par
Rasmdagi quti modeli: width=200px (content-box), padding=10px, border=2px, margin=5px (ikki yonda). Umumiy gorizontal joy qancha?
\begin{enumerate}[label=\Alph*),leftmargin=9mm]
\item 210px
\item 200px
\item 224px
\item 234px
\end{enumerate}
\begin{quote}\footnotesize\color{muted}\textbf{Rasm:} CSS content-box padding border margin o'lcham sxemasi\\
\href{https://github.com/sarvar9417/attestatsiya/blob/main/public/graphics-web-v2/css_box.svg}{Rasmni onlayn ko'rish}
\end{quote}
\vspace{2mm}\hrule height 0.2pt \vspace{4mm}
\needspace{5\baselineskip}
\hypertarget{q492}{\textbf{\color{ink}492.}}  {\color{muted}\scriptsize Y1 \textbullet\ daraja 5 \textbullet\ DB 55f938fb}
\par
p\{color:red!important\} va \textless{}p style="color:blue"\textgreater{} uchun muhimlik darajasi bir xil origin, layer yo'q. Natija?
\begin{enumerate}[label=\Alph*),leftmargin=9mm]
\item Qizil
\item Sariq
\item Ko'k
\item Yashil
\end{enumerate}
\vspace{2mm}\hrule height 0.2pt \vspace{4mm}
\needspace{5\baselineskip}
\hypertarget{q493}{\textbf{\color{ink}493.}}  {\color{muted}\scriptsize Y1 \textbullet\ daraja 5 \textbullet\ DB 666999cd}
\par
width:100px; padding:10px; border:2px solid; margin:5px; box-sizing:content-box. Gorizontal umumiy joy?
\begin{enumerate}[label=\Alph*),leftmargin=9mm]
\item 134px
\item 110px
\item 100px
\item 124px
\end{enumerate}
\vspace{2mm}\hrule height 0.2pt \vspace{4mm}
\needspace{5\baselineskip}
\hypertarget{q494}{\textbf{\color{ink}494.}}  {\color{muted}\scriptsize Y1 \textbullet\ daraja 5 \textbullet\ DB 945a062d}
\par
CSSda width:100px;height:100px;object-fit:contain bo'lsa, 3:1 rastr rasm qanday joylashadi?
\begin{enumerate}[label=\Alph*),leftmargin=9mm]
\item Rasm nisbatini saqlagan holda quti ichiga sig'ishi mumkin
\item Butun quti doim rasmsiz qoladi
\item Tasvirning yarmi har doim kesiladi
\item Nisbat majburan 1:1 bo'lib cho'ziladi
\end{enumerate}
\vspace{2mm}\hrule height 0.2pt \vspace{4mm}
\needspace{5\baselineskip}
\hypertarget{q495}{\textbf{\color{ink}495.}}  {\color{muted}\scriptsize Y1 \textbullet\ daraja 5 \textbullet\ DB a1c23c96}
\par
CSS: \#x\{color:red\} p\{color:blue\} va \textless{}p id="x" style="color:green"\textgreater{}. Oddiy, important yo'q. Natija?
\begin{enumerate}[label=\Alph*),leftmargin=9mm]
\item Qizil
\item Yashil
\item Rangsiz
\item Ko'k
\end{enumerate}
\vspace{2mm}\hrule height 0.2pt \vspace{4mm}
\needspace{5\baselineskip}
\hypertarget{q496}{\textbf{\color{ink}496.}}  {\color{muted}\scriptsize Y3 \textbullet\ daraja 5 \textbullet\ DB a2c834ca}
\par
\texttt{.card \textgreater{} p} selektori qaysi paragrafga tegishini tahlil qilish tartibini tuzing.
\begin{enumerate}[label=\Alph*),leftmargin=9mm]
\item Faqat bevosita pga qoida tatbiq bo'lishini tasdiqlang
\item HTMLdagi class='card' konteynerini aniqlang
\item Uning ichidagi p elementlarning ota-ona bog'lanishini tekshiring
\item Bevosita farzand pni ichkaridagi section pdan ajrating
\end{enumerate}
\vspace{2mm}\hrule height 0.2pt \vspace{4mm}
\needspace{5\baselineskip}
\hypertarget{q497}{\textbf{\color{ink}497.}}  {\color{muted}\scriptsize Y1 \textbullet\ daraja 5 \textbullet\ DB a612c3fa}
\par
img\{width:100px;height:100px;object-fit:cover\} va 3:1 nisbatdagi manba. Qanday natija mumkin?
\begin{enumerate}[label=\Alph*),leftmargin=9mm]
\item Kvadrat quti to'ladi, rasmning ayrim qismi qirqiladi
\item Har doim tasvir butunligicha ko'rinadi
\item Manba fayl avtomatik 100x100ga yoziladi
\item Rasm doim yo'qoladi
\end{enumerate}
\vspace{2mm}\hrule height 0.2pt \vspace{4mm}
\needspace{5\baselineskip}
\hypertarget{q498}{\textbf{\color{ink}498.}}  {\color{muted}\scriptsize Y1 \textbullet\ daraja 5 \textbullet\ DB c5a0678e}
\par
width:200px; padding:10px; border:2px solid; margin:5px; box-sizing:border-box holatida jami gorizontal joy?
\begin{enumerate}[label=\Alph*),leftmargin=9mm]
\item 234px
\item 224px
\item 200px
\item 210px
\end{enumerate}
\vspace{2mm}\hrule height 0.2pt \vspace{4mm}
\needspace{5\baselineskip}
\hypertarget{q499}{\textbf{\color{ink}499.}}  {\color{muted}\scriptsize Y2 \textbullet\ daraja 5 \textbullet\ DB dd58dcee}
\par
Ro'yxat va jadvalning CSS xossalarini moslang.
\textbf{\footnotesize Chap ustun:}
\begin{enumerate}[label=\arabic*),leftmargin=9mm]
\item \texttt{list-style-type:square}
\item \texttt{border-collapse:collapse}
\item \texttt{td\{padding:8px\}}
\item \texttt{list-style-position:inside}
\end{enumerate}
\textbf{\footnotesize O'ng ustun:}
\begin{enumerate}[label=\Alph*),leftmargin=9mm]
\item Jadvaldagi qo'shni chegaralarni birlashtirish
\item Markerning band ichida joylashishi
\item Katak ichida sakkiz piksel bo'sh joy
\item Tartibsiz ro'yxatda kvadrat marker
\end{enumerate}
\vspace{2mm}\hrule height 0.2pt \vspace{4mm}
\needspace{5\baselineskip}
\hypertarget{q500}{\textbf{\color{ink}500.}}  {\color{muted}\scriptsize Y1 \textbullet\ daraja 5 \textbullet\ DB eff1aa35}
\par
\textless{}p class="a" id="x"\textgreater{} ga p\{color:red\}, .a\{color:green\}, \#x\{color:blue\} berildi. Inline va important yo'q. Natija?
\begin{enumerate}[label=\Alph*),leftmargin=9mm]
\item Qizil
\item Ko'k
\item Yashil
\item Brauzer fon rangi
\end{enumerate}
\vspace{2mm}\hrule height 0.2pt \vspace{4mm}
\needspace{5\baselineskip}
\hypertarget{q501}{\textbf{\color{ink}501.}}  {\color{muted}\scriptsize Y1 \textbullet\ daraja 5 \textbullet\ DB f8caa350}
\par
CSS: .box p\{color:blue\} va .box \textgreater{} p\{color:red\}. Birinchi bevosita p, ikkinchisi ichki sectiondagi p. Qanday ranglar?
\begin{enumerate}[label=\Alph*),leftmargin=9mm]
\item Birinchi ko'k, ikkinchisi qizil
\item Ikkalasi qizil
\item Ikkalasi ko'k
\item Birinchi qizil, ikkinchi ko'k
\end{enumerate}
\vspace{2mm}\hrule height 0.2pt \vspace{4mm}
\needspace{5\baselineskip}
\hypertarget{q502}{\textbf{\color{ink}502.}}  {\color{muted}\scriptsize Y1 \textbullet\ daraja 5 \textbullet\ DB fbd7c752}
\par
CSSda .box p \{color:blue\} va .box \textgreater{} p \{color:red\} berilgan; bevosita p va ichki sectiondagi p qanday bo'ladi?
\begin{enumerate}[label=\Alph*),leftmargin=9mm]
\item Bevosita ko'k, ichki qizil
\item Bevosita qizil, ichki ko'k
\item Ikkisi qizil
\item Ikkisi ko'k
\end{enumerate}
\vspace{2mm}\hrule height 0.2pt \vspace{4mm}
\clearpage\section{Javoblar va izohlar}
\textit{Javoblar bazadagi question\_keys yozuvlaridan olingan.}\\
Har bir raqam ustiga bossangiz, savolning o'ziga qaytasiz.
\subsection{Nashr qilingan savollar}
\needspace{3\baselineskip}
\noindent\hyperlink{q1}{\textbf{1.}} \textbf{A}.  Rastr grafika - tasvirni piksellar to'ri ko'rinishida saqlaydigan grafika turi.
\par\vspace{3mm}
\needspace{3\baselineskip}
\noindent\hyperlink{q2}{\textbf{2.}} \textbf{B}.  Ta'rif Vektor grafika tushunchasiga tegishli.
\par\vspace{3mm}
\needspace{3\baselineskip}
\noindent\hyperlink{q3}{\textbf{3.}} \textbf{B}.  Rastr va vektor farqi - rastr piksel, vektor geometrik obyektlarga asoslanishini ajratish.
\par\vspace{3mm}
\needspace{3\baselineskip}
\noindent\hyperlink{q4}{\textbf{4.}} \textbf{A}.  Hex rang yozuvi \texttt{\#RRGGBB} shaklida qizil, yashil va ko'k - RGB komponentlarining qiymatlarini beradi.
\par\vspace{3mm}
\needspace{3\baselineskip}
\noindent\hyperlink{q5}{\textbf{5.}} \textbf{A}.  Ta'rif Rastr grafika tushunchasiga tegishli.
\par\vspace{3mm}
\needspace{3\baselineskip}
\noindent\hyperlink{q6}{\textbf{6.}} \textbf{B}.  Vektor grafika - tasvirni geometrik obyektlar va formulalar orqali ifodalaydigan grafika turi.
\par\vspace{3mm}
\needspace{3\baselineskip}
\noindent\hyperlink{q7}{\textbf{7.}} \textbf{B}.  Vaziyat Vektor grafikaga misol: tasvirni geometrik obyektlar va formulalar orqali ifodalaydigan grafika turi.
\par\vspace{3mm}
\needspace{3\baselineskip}
\noindent\hyperlink{q8}{\textbf{8.}} \textbf{A}.  Vaziyat Rastr grafikaga misol: tasvirni piksellar to'ri ko'rinishida saqlaydigan grafika turi.
\par\vspace{3mm}
\needspace{3\baselineskip}
\noindent\hyperlink{q9}{\textbf{9.}} \textbf{D}.  Ta'rif Rastr va vektor farqi tushunchasiga tegishli.
\par\vspace{3mm}
\needspace{3\baselineskip}
\noindent\hyperlink{q10}{\textbf{10.}} \textbf{B}.  Vaziyat Rastr va vektor farqiga misol: rastr piksel, vektor geometrik obyektlarga asoslanishini ajratish.
\par\vspace{3mm}
\needspace{3\baselineskip}
\noindent\hyperlink{q11}{\textbf{11.}} \textbf{A}.  Fotosuratlar odatda piksellardan tashkil topgan rastr tasvir sifatida tahrirlanadi.
\par\vspace{3mm}
\needspace{3\baselineskip}
\noindent\hyperlink{q12}{\textbf{12.}} \textbf{A}.  Chizish uchun varaqsimon ish maydoni Canvas (kanvas), standart egri chiziq asbobi Curve, palitrada yo'q rangni tanlash Edit Colors oynasi bilan bog'liq. Qolgan variantlar vazifalarni boshqa turdagi menyu yoki asbob bilan almashtiradi.
\par\vspace{3mm}
\needspace{3\baselineskip}
\noindent\hyperlink{q13}{\textbf{13.}} \textbf{B}.  Grafik muharrir rasmni yaratish va qayta ishlash uchun xizmat qiladi. [2025 kuzgi demo YECHIMLARI, sahifa 1.]
\par\vspace{3mm}
\needspace{3\baselineskip}
\noindent\hyperlink{q14}{\textbf{14.}} \textbf{A}.  Vektor grafika geometrik obyektlar bilan ifodalanadi va o'lcham o'zgarganda sifatni yaxshi saqlaydi.
\par\vspace{3mm}
\needspace{3\baselineskip}
\noindent\hyperlink{q15}{\textbf{15.}} \textbf{A}.  SVG geometrik yo'l/obyektlar bilan ishlaydi; konturlarni o'lchamdan mustaqil qayta hisoblash mumkin. JPEG/BMP/GIF esa asosan rastr modeliga tayangan.
\par\vspace{3mm}
\needspace{3\baselineskip}
\noindent\hyperlink{q16}{\textbf{16.}} \textbf{1-C, 2-A, 3-D, 4-B}.  SVG geometrik manba; PSD qatlamli loyiha; PNG alpha va lossless; JPEG fotolar uchun keng qo'llanadigan yo'qotishli eksport. Bular bir vazifani bajarmaydi.
\par\vspace{3mm}
\needspace{3\baselineskip}
\noindent\hyperlink{q17}{\textbf{17.}} \textbf{D}.  Ekranda RGB yorug'likka asoslanadi, bosmada esa ko'pincha CMYK bo'yoq modeli ishlatiladi. Bosma profil/gamut farqi sabab ko'rinish o'zgarishi mumkin.
\par\vspace{3mm}
\needspace{3\baselineskip}
\noindent\hyperlink{q18}{\textbf{18.}} \textbf{C - A - B - D}.  Obyektlar va shaffoflik talabi aniqlanmasdan eksport nazorati qilinmaydi; manba va chiqish nusxalarini ajratish qayta tahrir imkonini saqlaydi.
\par\vspace{3mm}
\needspace{3\baselineskip}
\noindent\hyperlink{q19}{\textbf{19.}} \textbf{C}.  120x80=9 600 piksel. 24 bit=3 bayt. 9 600x3=28 800 bayt. Faylning haqiqiy disk hajmi format/metama'lumotga bog'liq.
\par\vspace{3mm}
\needspace{3\baselineskip}
\noindent\hyperlink{q20}{\textbf{20.}} \textbf{B}.  Eni 960, bo'yi 540 bo'ladi. Jami 960x540=518 400 piksel. Ikki o'lchamning har biri yarmiga tushsa umumiy piksel soni to'rt baravar kamayadi.
\par\vspace{3mm}
\needspace{3\baselineskip}
\noindent\hyperlink{q21}{\textbf{21.}} \textbf{D}.  Alfa kanal - pikselning shaffoflik darajasini saqlaydigan qo'shimcha kanal.
\par\vspace{3mm}
\needspace{3\baselineskip}
\noindent\hyperlink{q22}{\textbf{22.}} \textbf{D}.  PPI - rasm yoki ekranning bir dyuymiga to'g'ri keladigan piksel soni.
\par\vspace{3mm}
\needspace{3\baselineskip}
\noindent\hyperlink{q23}{\textbf{23.}} \textbf{A}.  Obyektlar proporsiyasi saqlangan, lekin yon qismlar olib tashlanib rasm toraygan. Bu o'lchamni cho'zish emas, portret nisbatida crop qilishdir.
\par\vspace{3mm}
\needspace{3\baselineskip}
\noindent\hyperlink{q24}{\textbf{24.}} \textbf{C}.  Kenglik:balandlik munosabati aspect ratio, ya'ni tomonlar nisbati yoki mutanosiblik deyiladi.
\par\vspace{3mm}
\needspace{3\baselineskip}
\noindent\hyperlink{q25}{\textbf{25.}} \textbf{C}.  Piksel - rastr tasvirning rang qiymatiga ega eng kichik elementi.
\par\vspace{3mm}
\needspace{3\baselineskip}
\noindent\hyperlink{q26}{\textbf{26.}} \textbf{A}.  Rang chuqurligi - bitta piksel rangini kodlash uchun ajratilgan bitlar soni.
\par\vspace{3mm}
\needspace{3\baselineskip}
\noindent\hyperlink{q27}{\textbf{27.}} \textbf{C}.  CMYK - sian, magenta, sariq va qora bo'yoqlarni ayrish tamoyiliga asoslangan model.
\par\vspace{3mm}
\needspace{3\baselineskip}
\noindent\hyperlink{q28}{\textbf{28.}} \textbf{B}.  RGB - qizil, yashil va ko'k yorug'likni qo'shishga asoslangan rang modeli.
\par\vspace{3mm}
\needspace{3\baselineskip}
\noindent\hyperlink{q29}{\textbf{29.}} \textbf{A}.  Bo'laklar 5+2+4+5+3=19 kadr; 4 ta bo'sh kadr bilan 23 kadr. 23/15≈1,53 s.
\par\vspace{3mm}
\needspace{3\baselineskip}
\noindent\hyperlink{q30}{\textbf{30.}} \textbf{C}.  Ta'rif CMYK tushunchasiga tegishli.
\par\vspace{3mm}
\needspace{3\baselineskip}
\noindent\hyperlink{q31}{\textbf{31.}} \textbf{D}.  Ta'rif PPI tushunchasiga tegishli.
\par\vspace{3mm}
\needspace{3\baselineskip}
\noindent\hyperlink{q32}{\textbf{32.}} \textbf{C}.  Ta'rif Piksel tushunchasiga tegishli.
\par\vspace{3mm}
\needspace{3\baselineskip}
\noindent\hyperlink{q33}{\textbf{33.}} \textbf{D}.  Ta'rif Alfa kanal tushunchasiga tegishli.
\par\vspace{3mm}
\needspace{3\baselineskip}
\noindent\hyperlink{q34}{\textbf{34.}} \textbf{B}.  Ta'rif RGB tushunchasiga tegishli.
\par\vspace{3mm}
\needspace{3\baselineskip}
\noindent\hyperlink{q35}{\textbf{35.}} \textbf{A}.  Elementlar jami 4+5+5+5+6=25; ular 25x2=50 kadr. To'rtta oraliq uchun 4x3=12 bo'sh kadr: jami 62 kadr. 62/10=6,2 soniya.
\par\vspace{3mm}
\needspace{3\baselineskip}
\noindent\hyperlink{q36}{\textbf{36.}} \textbf{D}.  Vaziyat PPIga misol: rasm yoki ekranning bir dyuymiga to'g'ri keladigan piksel soni.
\par\vspace{3mm}
\needspace{3\baselineskip}
\noindent\hyperlink{q37}{\textbf{37.}} \textbf{A}.  Ta'rif Rang chuqurligi tushunchasiga tegishli.
\par\vspace{3mm}
\needspace{3\baselineskip}
\noindent\hyperlink{q38}{\textbf{38.}} \textbf{C}.  Vaziyat Pikselga misol: rastr tasvirning rang qiymatiga ega eng kichik elementi.
\par\vspace{3mm}
\needspace{3\baselineskip}
\noindent\hyperlink{q39}{\textbf{39.}} \textbf{D}.  Vaziyat RGBga misol: qizil, yashil va ko'k yorug'likni qo'shishga asoslangan rang modeli.
\par\vspace{3mm}
\needspace{3\baselineskip}
\noindent\hyperlink{q40}{\textbf{40.}} \textbf{C}.  Vaziyat Rang chuqurligiga misol: bitta piksel rangini kodlash uchun ajratilgan bitlar soni.
\par\vspace{3mm}
\needspace{3\baselineskip}
\noindent\hyperlink{q41}{\textbf{41.}} \textbf{1-D, 2-B, 3-C, 4-A}.  Zoom ko'rishni, Crop kadrni, Resize piksel o'lchamini, Canvas size ish maydoni chegarasini boshqaradi.
\par\vspace{3mm}
\needspace{3\baselineskip}
\noindent\hyperlink{q42}{\textbf{42.}} \textbf{A}.  Rastr tasvir piksellardan tuzilgani uchun kuchli kattalashtirishda piksellanish ko'rinadi.
\par\vspace{3mm}
\needspace{3\baselineskip}
\noindent\hyperlink{q43}{\textbf{43.}} \textbf{A}.  Zoom faqat dasturdagi ko'rish masshtabini o'zgartiradi. Asl rasm 600x400=240 000 piksel bo'lib qoladi.
\par\vspace{3mm}
\needspace{3\baselineskip}
\noindent\hyperlink{q44}{\textbf{44.}} \textbf{A}.  Crop tasvirning keraksiz tashqi qismlarini kesib tashlaydi.
\par\vspace{3mm}
\needspace{3\baselineskip}
\noindent\hyperlink{q45}{\textbf{45.}} \textbf{C}.  Qo'shiluvchi RGB modelida uch kanal intensivligi nol bo'lsa qora hosil bo'ladi.
\par\vspace{3mm}
\needspace{3\baselineskip}
\noindent\hyperlink{q46}{\textbf{46.}} \textbf{D}.  Vaziyat CMYKga misol: sian, magenta, sariq va qora bo'yoqlarni ayrish tamoyiliga asoslangan model.
\par\vspace{3mm}
\needspace{3\baselineskip}
\noindent\hyperlink{q47}{\textbf{47.}} \textbf{D}.  Bosmadagi eni=piksel eni/PPI=900/300=3 dyuym (taxminan 7,62 sm). Faqat metadata PPI ni o'zgartirish yangi tafsilot yaratmaydi.
\par\vspace{3mm}
\needspace{3\baselineskip}
\noindent\hyperlink{q48}{\textbf{48.}} \textbf{B - D - C - A}.  Asl manbani saqlash, kadrlash, Crop va eksport tartibi ma'lumot yo'qotilishini boshqaradi. Rastrni oldin siqish shart emas.
\par\vspace{3mm}
\needspace{3\baselineskip}
\noindent\hyperlink{q49}{\textbf{49.}} \textbf{B}.  SVG vektor obyektlarni saqlaydi va shaffoflikni qo'llashi mumkin.
\par\vspace{3mm}
\needspace{3\baselineskip}
\noindent\hyperlink{q50}{\textbf{50.}} \textbf{A}.  Vaziyat Alfa kanalga misol: pikselning shaffoflik darajasini saqlaydigan qo'shimcha kanal.
\par\vspace{3mm}
\needspace{3\baselineskip}
\noindent\hyperlink{q51}{\textbf{51.}} \textbf{C}.  Crop berilgan 600x500 maydon tashqarisini tashlaydi; qolgan rastr 600x500=300 000 piksel. Bu resize emas.
\par\vspace{3mm}
\needspace{3\baselineskip}
\noindent\hyperlink{q52}{\textbf{52.}} \textbf{B}.  Intersection to'plamlar kesishmasidir. Union birlashma, Difference ayirma, Group esa geometrik kesim emas.
\par\vspace{3mm}
\needspace{3\baselineskip}
\noindent\hyperlink{q53}{\textbf{53.}} \textbf{B}.  Vertikal va gorizontal akslantirishni ketma-ket bajarish 180 deg  aylantirishga teng. Qolgan ketma-ketliklar bitta akslantirish yoki 90 deg  bilan bog'liq boshqa natija beradi.
\par\vspace{3mm}
\needspace{3\baselineskip}
\noindent\hyperlink{q54}{\textbf{54.}} \textbf{B}.  Asl qatlamdan nusxa olib nusxa ustida ishlash qaytish imkonini saqlaydi va asl tasvirni himoya qiladi.
\par\vspace{3mm}
\needspace{3\baselineskip}
\noindent\hyperlink{q55}{\textbf{55.}} \textbf{B}.  Qora rang layer mask sohasini yashiradi; tasvirning asl piksellari saqlanadi. Oq rang bilan qayta ochish mumkin. https://helpx.adobe.com/photoshop/using/editing-layer-masks.html
\par\vspace{3mm}
\needspace{3\baselineskip}
\noindent\hyperlink{q56}{\textbf{56.}} \textbf{1-B, 2-A, 3-D, 4-C}.  Asbobni tanlash bilan uning Options bar parametrini sozlash bir xil amal emas.
\par\vspace{3mm}
\needspace{3\baselineskip}
\noindent\hyperlink{q57}{\textbf{57.}} \textbf{A}.  Ko'rinish belgisi qatlamni o'chirmasdan yashirish imkonini beradi.
\par\vspace{3mm}
\needspace{3\baselineskip}
\noindent\hyperlink{q58}{\textbf{58.}} \textbf{A}.  Qatlamlar obyektlarni mustaqil tahrirlash, yashirish va joylashtirish imkonini beradi.
\par\vspace{3mm}
\needspace{3\baselineskip}
\noindent\hyperlink{q59}{\textbf{59.}} \textbf{A}.  Maskada oq ko'rsatadi, qora yashiradi. Bu bevosita piksellarni o'chirib tashlashdan farqli non-destructive usul.
\par\vspace{3mm}
\needspace{3\baselineskip}
\noindent\hyperlink{q60}{\textbf{60.}} \textbf{1-C, 2-A, 3-B, 4-D}.  Color Picker rangni tanlaydi, Lasso chegarani, Layer Mask ko'rinishni, Clone Stamp manba piksellar namunasi bilan ishlaydi.
\par\vspace{3mm}
\needspace{3\baselineskip}
\noindent\hyperlink{q61}{\textbf{61.}} \textbf{A}.  Adobe ko'rsatmasiga ko'ra, Windows'da Move Tool bilan Alt bosilgan holda tanlangan hududni sudrash nusxa ko'chirishni amalga oshiradi. https://helpx.adobe.com/photoshop/desktop/make-selections/refine-modify-selections/create-multiple-copies-of-a-selection-within-an-image.html
\par\vspace{3mm}
\needspace{3\baselineskip}
\noindent\hyperlink{q62}{\textbf{62.}} \textbf{D - A - B - C}.  Save As asl nomni saqlab qolgan holda alohida yangi fayl chiqarish uchun qulay.
\par\vspace{3mm}
\needspace{3\baselineskip}
\noindent\hyperlink{q63}{\textbf{63.}} \textbf{C}.  Delete Cropped Pixels belgilanmasa, kesilgan hududlar yashiriladi va keyinchalik qaytarilishi mumkin. Bu Background emas, oddiy qatlam bilan ishlashga taalluqli. https://helpx.adobe.com/photoshop/desktop/create-manage-layers/create-layer-compositions/non-destructive-editing-techniques.html
\par\vspace{3mm}
\needspace{3\baselineskip}
\noindent\hyperlink{q64}{\textbf{64.}} \textbf{1-D, 2-C, 3-A, 4-B}.  Ko'rinmas qatlam manba tasvirni o'chirib yuborgan degani emas.
\par\vspace{3mm}
\needspace{3\baselineskip}
\noindent\hyperlink{q65}{\textbf{65.}} \textbf{B - C - D - A}.  Maskdagi qora yashiradi, oq ko'rsatadi; asl piksel saqlanadi.
\par\vspace{3mm}
\needspace{3\baselineskip}
\noindent\hyperlink{q66}{\textbf{66.}} \textbf{1-C, 2-B, 3-A, 4-D}.  Nusxalar, maskalar va adjustment layer asl kontentga zarar yetkazmasdan ishlashga yordam beradi.
\par\vspace{3mm}
\needspace{3\baselineskip}
\noindent\hyperlink{q67}{\textbf{67.}} \textbf{A - D - C - B}.  Avval obyekt belgilanadi, so'ng nusxasi alohida qatlamga olinadi.
\par\vspace{3mm}
\needspace{3\baselineskip}
\noindent\hyperlink{q68}{\textbf{68.}} \textbf{D}.  Ko'rinadigan qatlamlarning stacking order tartibi qaysi obyekt oldinda bo'lishini belgilaydi. Zoom manba qatlam tartibini o'zgartirmaydi.
\par\vspace{3mm}
\needspace{3\baselineskip}
\noindent\hyperlink{q69}{\textbf{69.}} \textbf{B - D - A - C}.  Tanlov faqat maqsad sohani ajratadi; Layer Mask uni qaytariladigan tarzda yashiradi. PSD manbani saqlashdan oldin flatten qilish kerak emas.
\par\vspace{3mm}
\needspace{3\baselineskip}
\noindent\hyperlink{q70}{\textbf{70.}} \textbf{A - B - C - D}.  Adjustment layer boshqa qatlamdagi manba piksellarni bevosita qayta yozmasdan tuzatadi.
\par\vspace{3mm}
\needspace{3\baselineskip}
\noindent\hyperlink{q71}{\textbf{71.}} \textbf{C}.  Cambridge+ 5-sinf Paint mavzusidagi Curved line asbobi avval segmentni, keyin uning egriligini sozlashga asoslanadi.
\par\vspace{3mm}
\needspace{3\baselineskip}
\noindent\hyperlink{q72}{\textbf{72.}} \textbf{B}.  PSD qatlamli ishni tahrirlashga mo'ljallangan; PNG/JPEG esa yetkaziladigan ko'rish/rastr nusxasi. Bir maqsadni ikkinchisidan ajratish kerak.
\par\vspace{3mm}
\needspace{3\baselineskip}
\noindent\hyperlink{q73}{\textbf{73.}} \textbf{1-C, 2-A, 3-D, 4-B}.  Sharpen Blur bilan bir xil emas; Brightness va Contrast ham farqli.
\par\vspace{3mm}
\needspace{3\baselineskip}
\noindent\hyperlink{q74}{\textbf{74.}} \textbf{C}.  Distort filtrlari tasvir geometriyasini egish, burish va deformatsiya qilish uchun xizmat qiladi.
\par\vspace{3mm}
\needspace{3\baselineskip}
\noindent\hyperlink{q75}{\textbf{75.}} \textbf{C}.  180 deg  aylantirish gorizontal va vertikal akslantirishlarning kompozitsiyasidir. Undan keyingi vertikal akslantirish vertikal o'zgarishni bekor qilib, faqat chap-o'ng akslantirishni qoldiradi.
\par\vspace{3mm}
\needspace{3\baselineskip}
\noindent\hyperlink{q76}{\textbf{76.}} \textbf{C}.  180 deg  aylantirish gorizontal va vertikal akslantirishlar kompozitsiyasiga teng. Keyingi vertikal akslantirish vertikal o'zgarishni bekor qilib, faqat chap-o'ng akslantirishni qoldiradi.
\par\vspace{3mm}
\needspace{3\baselineskip}
\noindent\hyperlink{q77}{\textbf{77.}} \textbf{C}.  \#EE00EE rangida qizil va ko'k komponentlar yuqori, yashil esa 0; bu magenta rangidir.
\par\vspace{3mm}
\needspace{3\baselineskip}
\noindent\hyperlink{q78}{\textbf{78.}} \textbf{B}.  Sahifa sarlavhasini belgilovchi \texttt{\textless{}title\textgreater{}} elementi HTML hujjatining \texttt{\textless{}head\textgreater{}} qismida joylashadi.
\par\vspace{3mm}
\needspace{3\baselineskip}
\noindent\hyperlink{q79}{\textbf{79.}} \textbf{A}.  \textless{}h1\textgreater{} eng yuqori darajadagi sarlavha tegi.
\par\vspace{3mm}
\needspace{3\baselineskip}
\noindent\hyperlink{q80}{\textbf{80.}} \textbf{A}.  Tartibsiz ro'yxat uchun \textless{}ul\textgreater{}, elementlar uchun \textless{}li\textgreater{} ishlatiladi.
\par\vspace{3mm}
\needspace{3\baselineskip}
\noindent\hyperlink{q81}{\textbf{81.}} \textbf{A}.  Giperhavola \textless{}a\textgreater{} tegi va href atributi bilan yaratiladi.
\par\vspace{3mm}
\needspace{3\baselineskip}
\noindent\hyperlink{q82}{\textbf{82.}} \textbf{B - D - A - C}.  Document root =\textgreater{} head =\textgreater{} body =\textgreater{} close/test.
\par\vspace{3mm}
\needspace{3\baselineskip}
\noindent\hyperlink{q83}{\textbf{83.}} \textbf{A}.  Tashqi \texttt{\textless{}ul\textgreater{}} ichida Aralov bor. Uning ichki \texttt{\textless{}ul\textgreater{}} ro'yxatida Frontend va Backend turadi. Frontend ichidagi \texttt{\textless{}ol\textgreater{}} raqamli, Backend ichidagi \texttt{\textless{}ul\textgreater{}} esa markerli ro'yxat hosil qiladi.
\par\vspace{3mm}
\needspace{3\baselineskip}
\noindent\hyperlink{q84}{\textbf{84.}} \textbf{A}.  Rasm uchun img tegi va manba uchun src atributi ishlatiladi.
\par\vspace{3mm}
\needspace{3\baselineskip}
\noindent\hyperlink{q85}{\textbf{85.}} \textbf{1-B, 2-D, 3-A, 4-C}.  h1 - heading, a - anchor, img - image, ul - unordered list.
\par\vspace{3mm}
\needspace{3\baselineskip}
\noindent\hyperlink{q86}{\textbf{86.}} \textbf{A}.  \textless{}br\textgreater{} matnda qator uzilishini hosil qiladi.
\par\vspace{3mm}
\needspace{3\baselineskip}
\noindent\hyperlink{q87}{\textbf{87.}} \textbf{C - A - D - B}.  ol ichki ro'yxat tashqi li ichida joylashadi. li ni ichki ol dan oldin yopish nesting ma'nosini o'zgartiradi.
\par\vspace{3mm}
\needspace{3\baselineskip}
\noindent\hyperlink{q88}{\textbf{88.}} \textbf{C}.  \#008000 yashil rangni bildiradi; text-align:right o'ng tomonga tekislaydi; font-family:serif esa umumiy serif oilasini tanlaydi. Shuning uchun chap tekislash haqidagi mulohaza noto'g'ri.
\par\vspace{3mm}
\needspace{3\baselineskip}
\noindent\hyperlink{q89}{\textbf{89.}} \textbf{B}.  ol start="3" tashqi bandlarni 3 va 4 deb belgilaydi. Veb li ichida ul bo'lgani uchun HTML/CSS uning ichki ro'yxatidir.
\par\vspace{3mm}
\needspace{3\baselineskip}
\noindent\hyperlink{q90}{\textbf{90.}} \textbf{B}.  \texttt{rowspan} katakni bir necha qator bo'ylab, ya'ni vertikal birlashtiradi.
\par\vspace{3mm}
\needspace{3\baselineskip}
\noindent\hyperlink{q91}{\textbf{91.}} \textbf{B}.  \texttt{rowspan} katakni vertikal yo'nalishda bir nechta qator bo'ylab birlashtiradi; \texttt{colspan} esa gorizontal yo'nalishda ustunlarni qamrab oladi.
\par\vspace{3mm}
\needspace{3\baselineskip}
\noindent\hyperlink{q92}{\textbf{92.}} \textbf{A}.  \texttt{h1} sarlavha beradi; \texttt{input type="text"} matn maydonlarini, \texttt{value} ularning boshlang'ich qiymatini, \texttt{button} esa ikkita tugmani ko'rsatadi.
\par\vspace{3mm}
\needspace{3\baselineskip}
\noindent\hyperlink{q93}{\textbf{93.}} \textbf{B}.  Birlashtirish atributi yo'q kataklar: 160, 80, -160 va 85; yig'indi 165. Kodda jami 6 ta \texttt{\textless{}td\textgreater{}} elementi bor. 165·6=990.
\par\vspace{3mm}
\needspace{3\baselineskip}
\noindent\hyperlink{q94}{\textbf{94.}} \textbf{B}.  \texttt{th} odatda qalin; rowspan ikki satrni, colspan ikki ustunni birlashtiradi. \texttt{cellspacing=0} kataklar orasidagi masofani nol qiladi.
\par\vspace{3mm}
\needspace{3\baselineskip}
\noindent\hyperlink{q95}{\textbf{95.}} \textbf{B}.  \texttt{colspan=2} qo'llangan va \texttt{\textless{}ol\textgreater{}} ichida joylashgan nom - Facebook. Facebook asoschisi sifatida berilgan javob Mark Zuckerberg.
\par\vspace{3mm}
\needspace{3\baselineskip}
\noindent\hyperlink{q96}{\textbf{96.}} \textbf{C}.  Uchinchi ustun 400x25\%=100px; Lenovo=100px. Qolgan ikki ustunni birlashtirgan Asus=300px. a=400px, eng qisqa katak b=100px; a+b=500px. Manbadagi o'zaro zid width atributlari bir ma'noli jadval sxemasiga keltirildi.
\par\vspace{3mm}
\needspace{3\baselineskip}
\noindent\hyperlink{q97}{\textbf{97.}} \textbf{A}.  rowspan=2 bir td katagini joriy va keyingi qator bo'ylab yoyadi. https://developer.mozilla.org/en-US/docs/Web/HTML/Reference/Elements/td
\par\vspace{3mm}
\needspace{3\baselineskip}
\noindent\hyperlink{q98}{\textbf{98.}} \textbf{A}.  Bir qatorli matn kiritish uchun input type=text ishlatiladi.
\par\vspace{3mm}
\needspace{3\baselineskip}
\noindent\hyperlink{q99}{\textbf{99.}} \textbf{A}.  \textless{}tr\textgreater{} jadval qatorini, \textless{}td\textgreater{}/\textless{}th\textgreater{} esa uning kataklarini belgilaydi.
\par\vspace{3mm}
\needspace{3\baselineskip}
\noindent\hyperlink{q100}{\textbf{100.}} \textbf{A}.  rowspan="2" A katagini ikkita qator bo'ylab birlashtiradi. Birinchi qatorning ikkinchi katagi B, ikkinchi qatorning ikkinchi katagi C bo'ladi.
\par\vspace{3mm}
\needspace{3\baselineskip}
\noindent\hyperlink{q101}{\textbf{101.}} \textbf{1-B, 2-D, 3-C, 4-A}.  h2 sarlavha, start raqamlashni, colspan ustunlarni, img src/alt esa rasm manbasi va muqobil matnni belgilaydi.
\par\vspace{3mm}
\needspace{3\baselineskip}
\noindent\hyperlink{q102}{\textbf{102.}} \textbf{A}.  Bir guruhdan bitta variant tanlash uchun radio inputlar mos.
\par\vspace{3mm}
\needspace{3\baselineskip}
\noindent\hyperlink{q103}{\textbf{103.}} \textbf{B}.  method='post' POST usulini, action='/qabul' yuborish manzilini, input name va value esa yuboriladigan ism=Ali juftligini belgilaydi. https://developer.mozilla.org/en-US/docs/Web/API/HTMLFormElement/method
\par\vspace{3mm}
\needspace{3\baselineskip}
\noindent\hyperlink{q104}{\textbf{104.}} \textbf{D}.  Hidden maydon vizual ko'rsatilmaydi, ammo name/value bilan yuboriladi; foydalanuvchi uni devtools orqali o'zgartirishi mumkin. Client hidden qiymatni sir tutmaydi.
\par\vspace{3mm}
\needspace{3\baselineskip}
\noindent\hyperlink{q105}{\textbf{105.}} \textbf{A}.  \texttt{h2} o'ngga tekislanadi va qizil. Birinchi \texttt{p} bodydan qora rangni meros qilib oladi, 16px bo'ladi. Ikkinchi \texttt{p.t1} esa yashil va qalin. Har bir \texttt{p} 5px margin oladi; oddiy oqimda qo'shni vertikal marginlar collapse bo'lishi mumkin.
\par\vspace{3mm}
\needspace{3\baselineskip}
\noindent\hyperlink{q106}{\textbf{106.}} \textbf{A}.  \texttt{h2} inline uslubi sarlavhani qizil, o'ngga tekislangan va 40px qiladi. \texttt{float:right} rasmni tegishli konteynerning o'ng qismiga joylaydi. Keyingi oddiy \texttt{p} sarlavha uslubini meros qilib olmaydi. Rasm resursi topilmasa muqobil matn ko'rinishi mumkin.
\par\vspace{3mm}
\needspace{3\baselineskip}
\noindent\hyperlink{q107}{\textbf{107.}} \textbf{A}.  Qo'shni bloklarning vertikal marginlari odatda qo'shilmaydi, balki collapse bo'lib kattasi olinadi. Ikkalasi ham 15px, shuning uchun masofa 15px.
\par\vspace{3mm}
\needspace{3\baselineskip}
\noindent\hyperlink{q108}{\textbf{108.}} \textbf{A}.  Ikki qiymatli margin yozuvida yuqori/pastki margin 10px. Qo'shni vertikal marginlar qo'shilmay, 10px ga collapse bo'ladi.
\par\vspace{3mm}
\needspace{3\baselineskip}
\noindent\hyperlink{q109}{\textbf{109.}} \textbf{A}.  \texttt{td \{ color: green; \}} katak ichidagi matn rangini yashil qiladi. \texttt{background-color} qoidalari faqat fon rangiga ta'sir qiladi.
\par\vspace{3mm}
\needspace{3\baselineskip}
\noindent\hyperlink{q110}{\textbf{110.}} \textbf{A}.  Universal \texttt{*} selektori barcha elementlarga rangni bevosita beradi, shuning uchun \texttt{h2} bodydan orange rangni meros olmaydi: sarlavha purple va qalin. \texttt{li} uchun qalinlik normal; \texttt{.p2} birinchi bandni black, \texttt{\#id3} ikkinchi bandni darkred qiladi. \texttt{ol} markerlar yuqori Rim raqamlari. II marker to'q qizil, ichki \texttt{p} esa teal, italic va universal qoida sabab bold. III li brown va normal.
\par\vspace{3mm}
\needspace{3\baselineskip}
\noindent\hyperlink{q111}{\textbf{111.}} \textbf{B}.  Normal block formatting contextda qo'shni vertikal marginlar collapse bo'ladi, ya'ni 25px+25px emas, kattasiga teng 25px vertikal oralig' hosil bo'ladi. Border, padding, flex yoki grid holatlari natijani o'zgartirishi mumkin; shart bularni chiqarib tashladi.
\par\vspace{3mm}
\needspace{3\baselineskip}
\noindent\hyperlink{q112}{\textbf{112.}} \textbf{B - D - A - C}.  Target =\textgreater{} selector =\textgreater{} property/value =\textgreater{} verify.
\par\vspace{3mm}
\needspace{3\baselineskip}
\noindent\hyperlink{q113}{\textbf{113.}} \textbf{A}.  ID selektori \# belgisi bilan yoziladi.
\par\vspace{3mm}
\needspace{3\baselineskip}
\noindent\hyperlink{q114}{\textbf{114.}} \textbf{A}.  Class selektori nuqta bilan boshlanadi.
\par\vspace{3mm}
\needspace{3\baselineskip}
\noindent\hyperlink{q115}{\textbf{115.}} \textbf{A}.  width xossasi element kengligini belgilaydi.
\par\vspace{3mm}
\needspace{3\baselineskip}
\noindent\hyperlink{q116}{\textbf{116.}} \textbf{A}.  CSSda selektor p va color xossasi ishlatiladi.
\par\vspace{3mm}
\needspace{3\baselineskip}
\noindent\hyperlink{q117}{\textbf{117.}} \textbf{A}.  background-color element fon rangini belgilaydi.
\par\vspace{3mm}
\needspace{3\baselineskip}
\noindent\hyperlink{q118}{\textbf{118.}} \textbf{A}.  \texttt{\textgreater{}} bevosita farzand kombinatoridir. \texttt{\#one} p div.boxning bevosita farzandi, \texttt{\#two} esa section ichida avlod bo'lib turibdi.
\par\vspace{3mm}
\needspace{3\baselineskip}
\noindent\hyperlink{q119}{\textbf{119.}} \textbf{1-B, 2-D, 3-A, 4-C}.  Bo'sh joy avlod kombinatori, \texttt{\textgreater{}} farzand kombinatori, vergul guruhlash, \texttt{p.note} esa element va class kesishmasi.
\par\vspace{3mm}
\needspace{3\baselineskip}
\noindent\hyperlink{q120}{\textbf{120.}} \textbf{1-B, 2-D, 3-A, 4-C}.  color - text; background-color - fon; width - kenglik; text-align - tekislash.
\par\vspace{3mm}
\needspace{3\baselineskip}
\noindent\hyperlink{q121}{\textbf{121.}} \textbf{D}.  Tashqi kenglik=content 200 + padding 20 + border 6 + margin 10=236px. Width content-box rejimida faqat contentga tegishli.
\par\vspace{3mm}
\needspace{3\baselineskip}
\noindent\hyperlink{q122}{\textbf{122.}} \textbf{D - B - C - A}.  CSS kaskadi faqat manbadagi oxirgi qoidaga bog'liq emas: avval origin/importance/layer, keyin specificity, so'ng manba tartibi hal qiladi.
\par\vspace{3mm}
\needspace{3\baselineskip}
\noindent\hyperlink{q123}{\textbf{123.}} \textbf{C}.  Uchala selektor elementga mos keladi, lekin id \#msg (1,0,0) spetsifikligi class .note (0,1,0) va element p (0,0,1)dan yuqori. Shu bois maroon ishlaydi.
\par\vspace{3mm}
\needspace{3\baselineskip}
\noindent\hyperlink{q124}{\textbf{124.}} \textbf{B}.  cover tasvirning ko'rinadigan qismini quti to'lsin deb masshtablaydi va chekkalarini kesishi mumkin; contain esa butun tasvirni sig'dirib bo'sh joy qoldirishi mumkin.
\par\vspace{3mm}
\needspace{3\baselineskip}
\noindent\hyperlink{q125}{\textbf{125.}} \textbf{A}.  Samsung \texttt{360·25\%=90 px}, Xiaomi \texttt{72 px}, Huawei \texttt{100 px}, Honor \texttt{360-90-72-100=98 px}. Nokia ikki ustunni egallab \texttt{90+98=188 px}. Eng keng - Nokia, eng tor - Xiaomi.
\par\vspace{3mm}
\subsection{Tekshiruvdagi savollar}
\needspace{3\baselineskip}
\noindent\hyperlink{q126}{\textbf{126.}} \textbf{A}.  Rang berishda Paint Bucket, qatlam chegarasini chizishda Stroke, qatlamni nusxalab surishda Alt+Move ishlatiladi.
\par\vspace{3mm}
\needspace{3\baselineskip}
\noindent\hyperlink{q127}{\textbf{127.}} \textbf{D}.  Vektor manbada konturlar matematik tavsiflanadi.
\par\vspace{3mm}
\needspace{3\baselineskip}
\noindent\hyperlink{q128}{\textbf{128.}} \textbf{B}.  PSD qatlam/maskalarni saqlashi mumkin.
\par\vspace{3mm}
\needspace{3\baselineskip}
\noindent\hyperlink{q129}{\textbf{129.}} \textbf{A}.  CMYK bo'yoqqa asoslangan ayirish modelidir.
\par\vspace{3mm}
\needspace{3\baselineskip}
\noindent\hyperlink{q130}{\textbf{130.}} \textbf{A}.  3D koordinatalar x,y,z.
\par\vspace{3mm}
\needspace{3\baselineskip}
\noindent\hyperlink{q131}{\textbf{131.}} \textbf{B}.  Rastr tasvir piksellar matritsasidan iborat.
\par\vspace{3mm}
\needspace{3\baselineskip}
\noindent\hyperlink{q132}{\textbf{132.}} \textbf{A}.  ICT7 PDF 40-42: manba modeli bilan eksport faylining turi farqlanadi.
\par\vspace{3mm}
\needspace{3\baselineskip}
\noindent\hyperlink{q133}{\textbf{133.}} \textbf{B}.  PNG lossless rejimda ishlaydi va alpha saqlashi mumkin.
\par\vspace{3mm}
\needspace{3\baselineskip}
\noindent\hyperlink{q134}{\textbf{134.}} \textbf{B}.  SVG geometriya/elementlar bilan ishlaydi.
\par\vspace{3mm}
\needspace{3\baselineskip}
\noindent\hyperlink{q135}{\textbf{135.}} \textbf{B}.  RGB: Red, Green, Blue.
\par\vspace{3mm}
\needspace{3\baselineskip}
\noindent\hyperlink{q136}{\textbf{136.}} \textbf{C}.  JPEG foto uchun rastrli siqilgan format.
\par\vspace{3mm}
\needspace{3\baselineskip}
\noindent\hyperlink{q137}{\textbf{137.}} \textbf{B}.  ICT7 PDF 40-42: fraktal tasvir tenglamalar yordamida yaratiladi.
\par\vspace{3mm}
\needspace{3\baselineskip}
\noindent\hyperlink{q138}{\textbf{138.}} \textbf{C}.  Alpha kompozitsiyadagi shaffoflikni nazorat qiladi.
\par\vspace{3mm}
\needspace{3\baselineskip}
\noindent\hyperlink{q139}{\textbf{139.}} \textbf{A}.  ICT7 PDF 44: Hue - rang toni.
\par\vspace{3mm}
\needspace{3\baselineskip}
\noindent\hyperlink{q140}{\textbf{140.}} \textbf{B}.  RGB additiv, CMYK bosma/subtraktiv.
\par\vspace{3mm}
\needspace{3\baselineskip}
\noindent\hyperlink{q141}{\textbf{141.}} \textbf{D}.  JPEG manba geometriyasini vektor sifatida saqlamaydi.
\par\vspace{3mm}
\needspace{3\baselineskip}
\noindent\hyperlink{q142}{\textbf{142.}} \textbf{D}.  PSD tahrirlash ma'lumotini saqlaydi.
\par\vspace{3mm}
\needspace{3\baselineskip}
\noindent\hyperlink{q143}{\textbf{143.}} \textbf{D}.  Vektor kontur qayta masshtablanadi.
\par\vspace{3mm}
\needspace{3\baselineskip}
\noindent\hyperlink{q144}{\textbf{144.}} \textbf{D}.  Rastrda manba tafsilot chegaralangan.
\par\vspace{3mm}
\needspace{3\baselineskip}
\noindent\hyperlink{q145}{\textbf{145.}} \textbf{D}.  GIFda oddiy kadrli animatsiya va palitra mavjud.
\par\vspace{3mm}
\needspace{3\baselineskip}
\noindent\hyperlink{q146}{\textbf{146.}} \textbf{C}.  Uchala yorug'lik kanali maksimal.
\par\vspace{3mm}
\needspace{3\baselineskip}
\noindent\hyperlink{q147}{\textbf{147.}} \textbf{C}.  SVG vektor, PNG rastr formatdir.
\par\vspace{3mm}
\needspace{3\baselineskip}
\noindent\hyperlink{q148}{\textbf{148.}} \textbf{D}.  PNG alpha kanali qo'llashi mumkin.
\par\vspace{3mm}
\needspace{3\baselineskip}
\noindent\hyperlink{q149}{\textbf{149.}} \textbf{B}.  Nurlar yo'qligi qora.
\par\vspace{3mm}
\needspace{3\baselineskip}
\noindent\hyperlink{q150}{\textbf{150.}} \textbf{A}.  Stroke kontur/chegara xossasi.
\par\vspace{3mm}
\needspace{3\baselineskip}
\noindent\hyperlink{q151}{\textbf{151.}} \textbf{B}.  Chizmada ko'k, yashil, to'q sariq va kulrang - jami 4 xil rang bor.
\par\vspace{3mm}
\needspace{3\baselineskip}
\noindent\hyperlink{q152}{\textbf{152.}} \textbf{C}.  Har kanal qiymatini additiv rang modeli bo'yicha birlashtiramiz.
\par\vspace{3mm}
\needspace{3\baselineskip}
\noindent\hyperlink{q153}{\textbf{153.}} \textbf{A - B - C}.  To'g'ri ketma-ketlik: 1 =\textgreater{} 2 =\textgreater{} 3.
\par\vspace{3mm}
\needspace{3\baselineskip}
\noindent\hyperlink{q154}{\textbf{154.}} \textbf{C}.  ICT7 PDF 43-44: RGB qo'shiluvchi rang modeli.
\par\vspace{3mm}
\needspace{3\baselineskip}
\noindent\hyperlink{q155}{\textbf{155.}} \textbf{A - D - B - C}.  Asl nusxani saqlash yo'qotishli siqishdan avval amalga oshiriladi.
\par\vspace{3mm}
\needspace{3\baselineskip}
\noindent\hyperlink{q156}{\textbf{156.}} \textbf{C - B - A - D}.  Vektor manba saqlanib, talabga mos eksport undan keyin olinadi.
\par\vspace{3mm}
\needspace{3\baselineskip}
\noindent\hyperlink{q157}{\textbf{157.}} \textbf{1-D, 2-A, 3-B, 4-C}.  RGB yorug'lik, CMYK bo'yoq; HSB va LAB boshqa koordinata modellari.
\par\vspace{3mm}
\needspace{3\baselineskip}
\noindent\hyperlink{q158}{\textbf{158.}} \textbf{1-A, 2-D, 3-B, 4-C}.  SVG vektorni, PSD tahrirlanadigan loyihani, JPEG foto siqishini, PNG esa alfa va lossless rejimni qo'llaydi.
\par\vspace{3mm}
\needspace{3\baselineskip}
\noindent\hyperlink{q159}{\textbf{159.}} \textbf{D - A - C - B}.  Shaffoflik tekshirilmasdan eksport qilishda fon keraksiz rangga aylanishi mumkin.
\par\vspace{3mm}
\needspace{3\baselineskip}
\noindent\hyperlink{q160}{\textbf{160.}} \textbf{A}.  Har kanal qiymatini additiv rang modeli bo'yicha birlashtiramiz.
\par\vspace{3mm}
\needspace{3\baselineskip}
\noindent\hyperlink{q161}{\textbf{161.}} \textbf{1-A, 2-B, 3-C}.  To'g'ri moslik: 1-a; 2-b; 3-c.
\par\vspace{3mm}
\needspace{3\baselineskip}
\noindent\hyperlink{q162}{\textbf{162.}} \textbf{B}.  Har kanal qiymatini additiv rang modeli bo'yicha birlashtiramiz.
\par\vspace{3mm}
\needspace{3\baselineskip}
\noindent\hyperlink{q163}{\textbf{163.}} \textbf{A}.  Har kanal qiymatini additiv rang modeli bo'yicha birlashtiramiz.
\par\vspace{3mm}
\needspace{3\baselineskip}
\noindent\hyperlink{q164}{\textbf{164.}} \textbf{C}.  ICT7 PDF 41: fraktal turlar matematik qonuniyatga tayanadi.
\par\vspace{3mm}
\needspace{3\baselineskip}
\noindent\hyperlink{q165}{\textbf{165.}} \textbf{B}.  Vazifada formatning maqsadga mos xususiyatini tanlash kerak.
\par\vspace{3mm}
\needspace{3\baselineskip}
\noindent\hyperlink{q166}{\textbf{166.}} \textbf{D}.  \$2\textasciicircum{}\{2\}=4\$ indeks.
\par\vspace{3mm}
\needspace{3\baselineskip}
\noindent\hyperlink{q167}{\textbf{167.}} \textbf{1-B, 2-C, 3-D, 4-A}.  Ko'rinish piksel bo'lsa ham manba modeli turlicha bo'lishi mumkin.
\par\vspace{3mm}
\needspace{3\baselineskip}
\noindent\hyperlink{q168}{\textbf{168.}} \textbf{D}.  Vazifada formatning maqsadga mos xususiyatini tanlash kerak.
\par\vspace{3mm}
\needspace{3\baselineskip}
\noindent\hyperlink{q169}{\textbf{169.}} \textbf{B}.  Piksel soni 1280x720=921600.
\par\vspace{3mm}
\needspace{3\baselineskip}
\noindent\hyperlink{q170}{\textbf{170.}} \textbf{D}.  Piksel soni 300x200=60000.
\par\vspace{3mm}
\needspace{3\baselineskip}
\noindent\hyperlink{q171}{\textbf{171.}} \textbf{C - A - D - B}.  Fazoviy model oldin, sahna parametrlaridan keyin rendering keladi.
\par\vspace{3mm}
\needspace{3\baselineskip}
\noindent\hyperlink{q172}{\textbf{172.}} \textbf{D}.  ICT7 PDF 44: Saturation - to'yinganlik; nol bo'lsa rangsiz/kulrang.
\par\vspace{3mm}
\needspace{3\baselineskip}
\noindent\hyperlink{q173}{\textbf{173.}} \textbf{D - A - B - C}.  Tahrir manbasi va taqdimot nusxasi alohida saqlanadi.
\par\vspace{3mm}
\needspace{3\baselineskip}
\noindent\hyperlink{q174}{\textbf{174.}} \textbf{C}.  ICT7 PDF 45-47: TIFF tasvirlarni saqlash variantlariga ega.
\par\vspace{3mm}
\needspace{3\baselineskip}
\noindent\hyperlink{q175}{\textbf{175.}} \textbf{D}.  Piksel soni 1024x768=786432.
\par\vspace{3mm}
\needspace{3\baselineskip}
\noindent\hyperlink{q176}{\textbf{176.}} \textbf{A - B - C - D}.  PPI = piksel kengligi / bosma kenglik; metama'lumotni shunchaki almashtirish detallarni ko'paytirmaydi.
\par\vspace{3mm}
\needspace{3\baselineskip}
\noindent\hyperlink{q177}{\textbf{177.}} \textbf{A}.  Tahrir manbasi va chiqish formatini farqlang.
\par\vspace{3mm}
\needspace{3\baselineskip}
\noindent\hyperlink{q178}{\textbf{178.}} \textbf{C}.  Maksimum \$2\textasciicircum{}\{8\}=256\$ indeks.
\par\vspace{3mm}
\needspace{3\baselineskip}
\noindent\hyperlink{q179}{\textbf{179.}} \textbf{A}.  Piksel soni 640x480=307200.
\par\vspace{3mm}
\needspace{3\baselineskip}
\noindent\hyperlink{q180}{\textbf{180.}} \textbf{1-A, 2-D, 3-C, 4-B}.  8 bitli RGB kanal 0-255 oralig'ida; alpha=0 shaffoflikni bildiradi.
\par\vspace{3mm}
\needspace{3\baselineskip}
\noindent\hyperlink{q181}{\textbf{181.}} \textbf{1-B, 2-A, 3-D, 4-C}.  PPI, pixel count va bit depth turli o'lchovlardir.
\par\vspace{3mm}
\needspace{3\baselineskip}
\noindent\hyperlink{q182}{\textbf{182.}} \textbf{1-A, 2-C, 3-B}.  To'g'ri moslik: 1-a; 2-c; 3-b.
\par\vspace{3mm}
\needspace{3\baselineskip}
\noindent\hyperlink{q183}{\textbf{183.}} \textbf{C}.  Piksel soni 1920x1080=2073600.
\par\vspace{3mm}
\needspace{3\baselineskip}
\noindent\hyperlink{q184}{\textbf{184.}} \textbf{A - C - B}.  To'g'ri ketma-ketlik: 1 =\textgreater{} 3 =\textgreater{} 2.
\par\vspace{3mm}
\needspace{3\baselineskip}
\noindent\hyperlink{q185}{\textbf{185.}} \textbf{1-C, 2-B, 3-A, 4-D}.  Formatni faqat kengaytma emas, ichki model va ish vazifasi aniqlaydi.
\par\vspace{3mm}
\needspace{3\baselineskip}
\noindent\hyperlink{q186}{\textbf{186.}} \textbf{C}.  Piksel soni 1200x800=960000.
\par\vspace{3mm}
\needspace{3\baselineskip}
\noindent\hyperlink{q187}{\textbf{187.}} \textbf{A}.  ICT7 PDF 40-42: vektor obyekti kattalashtirishda qayta hisoblanadi.
\par\vspace{3mm}
\needspace{3\baselineskip}
\noindent\hyperlink{q188}{\textbf{188.}} \textbf{A}.  Manba modeli va natijaviy eksportni alohida tekshiring.
\par\vspace{3mm}
\needspace{3\baselineskip}
\noindent\hyperlink{q189}{\textbf{189.}} \textbf{A - D - C - B}.  Kengaytma nomi yoki ekrandagi tasvirning o'zi yetarli dalil emas.
\par\vspace{3mm}
\needspace{3\baselineskip}
\noindent\hyperlink{q190}{\textbf{190.}} \textbf{1-A, 2-C, 3-B, 4-D}.  Rastr va vektor vazifasini, loyiha bilan eksport nusxasini adashtirmang.
\par\vspace{3mm}
\needspace{3\baselineskip}
\noindent\hyperlink{q191}{\textbf{191.}} \textbf{A}.  640x480x4=1 228 800 bayt.
\par\vspace{3mm}
\needspace{3\baselineskip}
\noindent\hyperlink{q192}{\textbf{192.}} \textbf{B}.  ICT7 PDF 45-47: grafik formatlar cheklovlarini farqlash kerak.
\par\vspace{3mm}
\needspace{3\baselineskip}
\noindent\hyperlink{q193}{\textbf{193.}} \textbf{C}.  ICT7 PDF 41,45-47: grafik tur va format farqlanadi.
\par\vspace{3mm}
\needspace{3\baselineskip}
\noindent\hyperlink{q194}{\textbf{194.}} \textbf{B}.  JPEG fotosuratli rastr siqish bo'lib, alfa kanalli oddiy PNGdan farq qiladi.
\par\vspace{3mm}
\needspace{3\baselineskip}
\noindent\hyperlink{q195}{\textbf{195.}} \textbf{B - C - D - A}.  Fraktalning matematik manbasi yakuniy ekranda rastr sifatida ko'rinishi mumkin.
\par\vspace{3mm}
\needspace{3\baselineskip}
\noindent\hyperlink{q196}{\textbf{196.}} \textbf{C}.  ICT7 PDF 44: HSB Hue burchaklarining darslik modeli.
\par\vspace{3mm}
\needspace{3\baselineskip}
\noindent\hyperlink{q197}{\textbf{197.}} \textbf{1-D, 2-C, 3-A, 4-B}.  2D tasvirga rendering 3D modelning fayl ko'rinishini avtomatik saqlab qolmaydi.
\par\vspace{3mm}
\needspace{3\baselineskip}
\noindent\hyperlink{q198}{\textbf{198.}} \textbf{C}.  ICT7 PDF 44: LAB modeli yorqinlik va rang o'qlariga asoslanadi.
\par\vspace{3mm}
\needspace{3\baselineskip}
\noindent\hyperlink{q199}{\textbf{199.}} \textbf{A}.  Tasvir modeli, format va maqsad uyg'un bo'lishi zarur.
\par\vspace{3mm}
\needspace{3\baselineskip}
\noindent\hyperlink{q200}{\textbf{200.}} \textbf{A}.  3D loyiha bilan undan olingan 2D render bir xil ma'lumot modeliga ega emas.
\par\vspace{3mm}
\needspace{3\baselineskip}
\noindent\hyperlink{q201}{\textbf{201.}} \textbf{A}.  HSV Saturation rang to'yinganligidir.
\par\vspace{3mm}
\needspace{3\baselineskip}
\noindent\hyperlink{q202}{\textbf{202.}} \textbf{D}.  Tasvir modeli, format va maqsad uyg'un bo'lishi zarur.
\par\vspace{3mm}
\needspace{3\baselineskip}
\noindent\hyperlink{q203}{\textbf{203.}} \textbf{A}.  ICT7 PDF 29-34: indekslangan rang va rang chuqurligini farqlash.
\par\vspace{3mm}
\needspace{3\baselineskip}
\noindent\hyperlink{q204}{\textbf{204.}} \textbf{B}.  Canvas size hujjat kadr/ish maydonini o'zgartiradi.
\par\vspace{3mm}
\needspace{3\baselineskip}
\noindent\hyperlink{q205}{\textbf{205.}} \textbf{C}.  Crop kadr chegarasini o'zgartiradi.
\par\vspace{3mm}
\needspace{3\baselineskip}
\noindent\hyperlink{q206}{\textbf{206.}} \textbf{B}.  Resize o'lchamlarni o'zgartiradi.
\par\vspace{3mm}
\needspace{3\baselineskip}
\noindent\hyperlink{q207}{\textbf{207.}} \textbf{A}.  32/8=4 bayt.
\par\vspace{3mm}
\needspace{3\baselineskip}
\noindent\hyperlink{q208}{\textbf{208.}} \textbf{C}.  Clipping mask yuqoridagi tasvirni tegishli kontur orqali cheklaydi.
\par\vspace{3mm}
\needspace{3\baselineskip}
\noindent\hyperlink{q209}{\textbf{209.}} \textbf{B}.  Amaliy tahrir amallari: ko'rinish va resampling alohida.
\par\vspace{3mm}
\needspace{3\baselineskip}
\noindent\hyperlink{q210}{\textbf{210.}} \textbf{A}.  24/8=3 bayt.
\par\vspace{3mm}
\needspace{3\baselineskip}
\noindent\hyperlink{q211}{\textbf{211.}} \textbf{C}.  Rastrlashtirish natijasi piksellar bo'ladi.
\par\vspace{3mm}
\needspace{3\baselineskip}
\noindent\hyperlink{q212}{\textbf{212.}} \textbf{C}.  Horizontal flip ko'zguli akslantirishdir.
\par\vspace{3mm}
\needspace{3\baselineskip}
\noindent\hyperlink{q213}{\textbf{213.}} \textbf{B}.  Grouping obyektlarning guruhlanishidir.
\par\vspace{3mm}
\needspace{3\baselineskip}
\noindent\hyperlink{q214}{\textbf{214.}} \textbf{B}.  Resampling qayta namunalashdir.
\par\vspace{3mm}
\needspace{3\baselineskip}
\noindent\hyperlink{q215}{\textbf{215.}} \textbf{A}.  Index palitradagi rangga murojaat qiladi.
\par\vspace{3mm}
\needspace{3\baselineskip}
\noindent\hyperlink{q216}{\textbf{216.}} \textbf{D}.  W*H*b/8 xom modeldir.
\par\vspace{3mm}
\needspace{3\baselineskip}
\noindent\hyperlink{q217}{\textbf{217.}} \textbf{D}.  Union birlashmadir.
\par\vspace{3mm}
\needspace{3\baselineskip}
\noindent\hyperlink{q218}{\textbf{218.}} \textbf{C}.  Intersection kesishma.
\par\vspace{3mm}
\needspace{3\baselineskip}
\noindent\hyperlink{q219}{\textbf{219.}} \textbf{D}.  PPI ta'rifni, resample namunalarni o'zgartirishi mumkin.
\par\vspace{3mm}
\needspace{3\baselineskip}
\noindent\hyperlink{q220}{\textbf{220.}} \textbf{C}.  Stak tartibi o'zgaradi.
\par\vspace{3mm}
\needspace{3\baselineskip}
\noindent\hyperlink{q221}{\textbf{221.}} \textbf{A}.  Ikki yo'nalish bo'yicha chiziqli interpolyatsiya.
\par\vspace{3mm}
\needspace{3\baselineskip}
\noindent\hyperlink{q222}{\textbf{222.}} \textbf{D}.  Crop va resize vazifalari boshqa.
\par\vspace{3mm}
\needspace{3\baselineskip}
\noindent\hyperlink{q223}{\textbf{223.}} \textbf{B}.  Nearest yaqin manbani tanlaydi.
\par\vspace{3mm}
\needspace{3\baselineskip}
\noindent\hyperlink{q224}{\textbf{224.}} \textbf{B}.  Bit va baytni aralashtirmaslik zarur.
\par\vspace{3mm}
\needspace{3\baselineskip}
\noindent\hyperlink{q225}{\textbf{225.}} \textbf{D}.  Mask nondestruktiv yondashuvdir.
\par\vspace{3mm}
\needspace{3\baselineskip}
\noindent\hyperlink{q226}{\textbf{226.}} \textbf{C}.  ICT5 PDF 76-83, Cambridge5 PDF 38-41: belgilash va kesish farqlanadi.
\par\vspace{3mm}
\needspace{3\baselineskip}
\noindent\hyperlink{q227}{\textbf{227.}} \textbf{1-D, 2-A, 3-B, 4-C}.  Zoom manba piksel sonini o'zgartirmaydi, Flip esa 90 deg  burish emas.
\par\vspace{3mm}
\needspace{3\baselineskip}
\noindent\hyperlink{q228}{\textbf{228.}} \textbf{B}.  Har ikki o'lchamni 2ga ko'paytiramiz.
\par\vspace{3mm}
\needspace{3\baselineskip}
\noindent\hyperlink{q229}{\textbf{229.}} \textbf{C}.  Kesib olingan soha 600x300=180000 piksel.
\par\vspace{3mm}
\needspace{3\baselineskip}
\noindent\hyperlink{q230}{\textbf{230.}} \textbf{1-A, 2-B, 3-C, 4-D}.  Interpolation yo'q tafsilotni kafolatli tiklamaydi.
\par\vspace{3mm}
\needspace{3\baselineskip}
\noindent\hyperlink{q231}{\textbf{231.}} \textbf{B}.  Har ikki o'lchamni 2ga ko'paytiramiz.
\par\vspace{3mm}
\needspace{3\baselineskip}
\noindent\hyperlink{q232}{\textbf{232.}} \textbf{A - D - B - C}.  Ikki chiziqli o'lcham 1,5 baravar bo'lsa yuza 1,5\textasciicircum{}2=2,25 baravar bo'ladi.
\par\vspace{3mm}
\needspace{3\baselineskip}
\noindent\hyperlink{q233}{\textbf{233.}} \textbf{B}.  Har ikki o'lchamni 3ga ko'paytiramiz.
\par\vspace{3mm}
\needspace{3\baselineskip}
\noindent\hyperlink{q234}{\textbf{234.}} \textbf{D}.  Har ikki o'lchamni 4ga ko'paytiramiz.
\par\vspace{3mm}
\needspace{3\baselineskip}
\noindent\hyperlink{q235}{\textbf{235.}} \textbf{B}.  Har ikki o'lchamni 2ga ko'paytiramiz.
\par\vspace{3mm}
\needspace{3\baselineskip}
\noindent\hyperlink{q236}{\textbf{236.}} \textbf{D - A - C - B}.  Oldin Crop, keyin Resize bajarilganda faqat kerakli kadr qayta namunalashdan o'tadi.
\par\vspace{3mm}
\needspace{3\baselineskip}
\noindent\hyperlink{q237}{\textbf{237.}} \textbf{C}.  Crop (qirqish) tasvirning kerakli chegaralar ichidagi qismini qoldirib, tashqaridagi hududni ko'rinishdan chiqaradi. Hue rangni, canvas kengaytirish esa ish maydonini o'zgartiradi; Layer group qatlamlarni tashkil etadi. Manba: https://helpx.adobe.com/photoshop/using/nondestructive-editing.html .
\par\vspace{3mm}
\needspace{3\baselineskip}
\noindent\hyperlink{q238}{\textbf{238.}} \textbf{1-C, 2-D, 3-A, 4-B}.  Select maydon tanlaydi; Crop kadrni qisqartiradi; Resize piksel o'lchamini, canvas esa ish maydoni chegarasini boshqaradi.
\par\vspace{3mm}
\needspace{3\baselineskip}
\noindent\hyperlink{q239}{\textbf{239.}} \textbf{D}.  Diagrammadagi bo'yalgan to'rtburchak 4x3=12 piksel.
\par\vspace{3mm}
\needspace{3\baselineskip}
\noindent\hyperlink{q240}{\textbf{240.}} \textbf{1-B, 2-A, 3-C}.  To'g'ri moslik: 1-b; 2-a; 3-c.
\par\vspace{3mm}
\needspace{3\baselineskip}
\noindent\hyperlink{q241}{\textbf{241.}} \textbf{B - A - D - C}.  WxHxb/8 faqat xom piksel massivi uchun qo'llanadi.
\par\vspace{3mm}
\needspace{3\baselineskip}
\noindent\hyperlink{q242}{\textbf{242.}} \textbf{C}.  800x600x24/8 = 1 440 000 bayt; header/siqishsiz model.
\par\vspace{3mm}
\needspace{3\baselineskip}
\noindent\hyperlink{q243}{\textbf{243.}} \textbf{A}.  Har ikki o'lchamni 2ga ko'paytiramiz.
\par\vspace{3mm}
\needspace{3\baselineskip}
\noindent\hyperlink{q244}{\textbf{244.}} \textbf{B - A - C}.  To'g'ri ketma-ketlik: 2 =\textgreater{} 1 =\textgreater{} 3.
\par\vspace{3mm}
\needspace{3\baselineskip}
\noindent\hyperlink{q245}{\textbf{245.}} \textbf{D}.  \$W\textbackslash{}times H\textbackslash{}times b/8=1382400\$ B.
\par\vspace{3mm}
\needspace{3\baselineskip}
\noindent\hyperlink{q246}{\textbf{246.}} \textbf{A}.  \$W\textbackslash{}times H\textbackslash{}times b/8=240000\$ B.
\par\vspace{3mm}
\needspace{3\baselineskip}
\noindent\hyperlink{q247}{\textbf{247.}} \textbf{B}.  \$W\textbackslash{}times H\textbackslash{}times b/8=86400\$ B.
\par\vspace{3mm}
\needspace{3\baselineskip}
\noindent\hyperlink{q248}{\textbf{248.}} \textbf{A - C - B - D}.  Nearest-neighbor yaqin manba pikselini tanlaydi va kataklar o'rtasini aralashtirmaydi.
\par\vspace{3mm}
\needspace{3\baselineskip}
\noindent\hyperlink{q249}{\textbf{249.}} \textbf{1-A, 2-D, 3-B, 4-C}.  Kontur operatsiyalari geometrik hududni o'zgartiradi; Group esa tahrir boshqaruvini birlashtiradi.
\par\vspace{3mm}
\needspace{3\baselineskip}
\noindent\hyperlink{q250}{\textbf{250.}} \textbf{A}.  \$W\textbackslash{}times H\textbackslash{}times b/8=360000\$ B.
\par\vspace{3mm}
\needspace{3\baselineskip}
\noindent\hyperlink{q251}{\textbf{251.}} \textbf{1-B, 2-A, 3-D, 4-C}.  b bit qiymat uchun 2\textasciicircum{}b imkoniyat; 24 bit = 3 bayt/piksel.
\par\vspace{3mm}
\needspace{3\baselineskip}
\noindent\hyperlink{q252}{\textbf{252.}} \textbf{1-B, 2-C, 3-A}.  To'g'ri moslik: 1-b; 2-c; 3-a.
\par\vspace{3mm}
\needspace{3\baselineskip}
\noindent\hyperlink{q253}{\textbf{253.}} \textbf{D}.  Vektor mantiqiy amallari va z-tartib farqlanadi.
\par\vspace{3mm}
\needspace{3\baselineskip}
\noindent\hyperlink{q254}{\textbf{254.}} \textbf{D}.  1200/300=4 dyuym; PPI bu yerda pikselning fizik chiqarish zichligi.
\par\vspace{3mm}
\needspace{3\baselineskip}
\noindent\hyperlink{q255}{\textbf{255.}} \textbf{D - A - B - C}.  90 deg  CCW aylanish qoidasida yangi x=-y, yangi y=x.
\par\vspace{3mm}
\needspace{3\baselineskip}
\noindent\hyperlink{q256}{\textbf{256.}} \textbf{D}.  \$W\textbackslash{}times H\textbackslash{}times b/8=250000\$ B.
\par\vspace{3mm}
\needspace{3\baselineskip}
\noindent\hyperlink{q257}{\textbf{257.}} \textbf{C}.  ICT7 PDF 29-34,45-47: kodlash hajmi va fayl siqishi farq qiladi.
\par\vspace{3mm}
\needspace{3\baselineskip}
\noindent\hyperlink{q258}{\textbf{258.}} \textbf{D}.  \$W\textbackslash{}times H\textbackslash{}times b/8=2764800\$ B.
\par\vspace{3mm}
\needspace{3\baselineskip}
\noindent\hyperlink{q259}{\textbf{259.}} \textbf{A - D - C - B}.  Intersection faqat shakllar tutash hududini qoldiradi.
\par\vspace{3mm}
\needspace{3\baselineskip}
\noindent\hyperlink{q260}{\textbf{260.}} \textbf{1-D, 2-C, 3-A, 4-B}.  Ko'chirish, masshtab va qatlam tartibi turli jarayonlar; rasterlash modellararo o'tish.
\par\vspace{3mm}
\needspace{3\baselineskip}
\noindent\hyperlink{q261}{\textbf{261.}} \textbf{B - C - A}.  To'g'ri ketma-ketlik: 2 =\textgreater{} 3 =\textgreater{} 1.
\par\vspace{3mm}
\needspace{3\baselineskip}
\noindent\hyperlink{q262}{\textbf{262.}} \textbf{D}.  \$W\textbackslash{}times H\textbackslash{}times b/8=128000\$ B.
\par\vspace{3mm}
\needspace{3\baselineskip}
\noindent\hyperlink{q263}{\textbf{263.}} \textbf{A}.  600x400=240000; 300x200=60000, nisbati 4:1.
\par\vspace{3mm}
\needspace{3\baselineskip}
\noindent\hyperlink{q264}{\textbf{264.}} \textbf{A}.  Dastlab 12x8=96 ta piksel, qirqishdan so'ng 6x4=24 ta piksel. Kamaygan qism 96-24=72; 72/96x100\%=75\%. Rasm alohida berilmagan, chunki hisoblash uchun barcha o'lchamlar savol matnida mavjud.
\par\vspace{3mm}
\needspace{3\baselineskip}
\noindent\hyperlink{q265}{\textbf{265.}} \textbf{1-C, 2-B, 3-A, 4-D}.  Layer mask va clipping maskni bevosita o'chirg'ich bilan tenglashtirmang.
\par\vspace{3mm}
\needspace{3\baselineskip}
\noindent\hyperlink{q266}{\textbf{266.}} \textbf{D}.  W va H ikkalasi x2 bo'lsa, maydon x4 bo'ladi.
\par\vspace{3mm}
\needspace{3\baselineskip}
\noindent\hyperlink{q267}{\textbf{267.}} \textbf{A}.  900/150=6 dyuym. PPI o'zi piksel o'lchamini o'zgartirmaydi.
\par\vspace{3mm}
\needspace{3\baselineskip}
\noindent\hyperlink{q268}{\textbf{268.}} \textbf{A}.  1600/200=8 dyuym. PPI o'zi piksel o'lchamini o'zgartirmaydi.
\par\vspace{3mm}
\needspace{3\baselineskip}
\noindent\hyperlink{q269}{\textbf{269.}} \textbf{C - B - A - D}.  Mask manba pikselni majburiy o'chirmasdan ko'rinishni boshqaradi.
\par\vspace{3mm}
\needspace{3\baselineskip}
\noindent\hyperlink{q270}{\textbf{270.}} \textbf{1-C, 2-A, 3-D, 4-B}.  WxHxb/8 bilan hisoblang; bu faylning siqilgan disk hajmi emas.
\par\vspace{3mm}
\needspace{3\baselineskip}
\noindent\hyperlink{q271}{\textbf{271.}} \textbf{A}.  48x4=192; 32x4=128. Nearest-neighbor ko'paytirish usulini bildiradi.
\par\vspace{3mm}
\needspace{3\baselineskip}
\noindent\hyperlink{q272}{\textbf{272.}} \textbf{C}.  Resampling usullari tasvirning silliq yoki qatlamli ko'rinishiga ta'sir qiladi.
\par\vspace{3mm}
\needspace{3\baselineskip}
\noindent\hyperlink{q273}{\textbf{273.}} \textbf{B}.  Berilgan piksel eni PPIga teskari proporsional; resample yo'q.
\par\vspace{3mm}
\needspace{3\baselineskip}
\noindent\hyperlink{q274}{\textbf{274.}} \textbf{C}.  8x6x3=144 bayt. Diagramma o'lchami 8x6.
\par\vspace{3mm}
\needspace{3\baselineskip}
\noindent\hyperlink{q275}{\textbf{275.}} \textbf{D}.  1200/300=4 dyuym. PPI o'zi piksel o'lchamini o'zgartirmaydi.
\par\vspace{3mm}
\needspace{3\baselineskip}
\noindent\hyperlink{q276}{\textbf{276.}} \textbf{B}.  Vektor manba va rastr eksportning tahrirlash modeli farqli.
\par\vspace{3mm}
\needspace{3\baselineskip}
\noindent\hyperlink{q277}{\textbf{277.}} \textbf{A}.  Vektor obyektlar bilan shakl algebra amallari; qo'shimcha texnik izoh.
\par\vspace{3mm}
\needspace{3\baselineskip}
\noindent\hyperlink{q278}{\textbf{278.}} \textbf{B - C - D - A}.  900 piksel / 300 PPI = 3 dyuym.
\par\vspace{3mm}
\needspace{3\baselineskip}
\noindent\hyperlink{q279}{\textbf{279.}} \textbf{B}.  Cambridge5 PDF 30-31: chiziq avval to'g'ri olinadi, keyin egiladi.
\par\vspace{3mm}
\needspace{3\baselineskip}
\noindent\hyperlink{q280}{\textbf{280.}} \textbf{D}.  Brush bo'yash uchun.
\par\vspace{3mm}
\needspace{3\baselineskip}
\noindent\hyperlink{q281}{\textbf{281.}} \textbf{D}.  Marquee tanlov asbobi.
\par\vspace{3mm}
\needspace{3\baselineskip}
\noindent\hyperlink{q282}{\textbf{282.}} \textbf{B}.  Eyedropper rang olish asbobi.
\par\vspace{3mm}
\needspace{3\baselineskip}
\noindent\hyperlink{q283}{\textbf{283.}} \textbf{D}.  MS Paint rastr grafik muharriridir.
\par\vspace{3mm}
\needspace{3\baselineskip}
\noindent\hyperlink{q284}{\textbf{284.}} \textbf{D}.  Cambridge5 PDF 35-36: maxsus rang palitrasi.
\par\vspace{3mm}
\needspace{3\baselineskip}
\noindent\hyperlink{q285}{\textbf{285.}} \textbf{A}.  Lasso erkin tanlov.
\par\vspace{3mm}
\needspace{3\baselineskip}
\noindent\hyperlink{q286}{\textbf{286.}} \textbf{A}.  PSD Photoshop loyihasi.
\par\vspace{3mm}
\needspace{3\baselineskip}
\noindent\hyperlink{q287}{\textbf{287.}} \textbf{D}.  Layers qatlam tartibi va xossalarini boshqaradi.
\par\vspace{3mm}
\needspace{3\baselineskip}
\noindent\hyperlink{q288}{\textbf{288.}} \textbf{C}.  Clone Stamp namunadan nusxa olish uskunasi.
\par\vspace{3mm}
\needspace{3\baselineskip}
\noindent\hyperlink{q289}{\textbf{289.}} \textbf{B}.  Opacity 0-100\% oralig'ida ishlatilishi mumkin.
\par\vspace{3mm}
\needspace{3\baselineskip}
\noindent\hyperlink{q290}{\textbf{290.}} \textbf{D}.  Mask ko'rinishni boshqaradi.
\par\vspace{3mm}
\needspace{3\baselineskip}
\noindent\hyperlink{q291}{\textbf{291.}} \textbf{B}.  Adjustment layer rang va tonal tuzatishlarni asl piksel qiymatlarini qayta yozmasdan alohida qatlamda qo'llaydi. Uning parametrlari va maskasi keyinchalik ham o'zgartirilishi mumkin. Boshqa variantlar tahrirning qayta tiklanishini rad etadi yoki asl piksellarni bevosita o'zgartirishni talab qiladi. Adobe rasmiy manbasi: https://helpx.adobe.com/photoshop/using/nondestructive-editing.html ; https://helpx.adobe.com/photoshop/desktop/create-manage-layers/color-adjustment-fill-layers/change-adjustment-and-fill-layer-options.html . Savol \texttt{review}, inson muharrir tasdig'i talab etiladi.
\par\vspace{3mm}
\needspace{3\baselineskip}
\noindent\hyperlink{q292}{\textbf{292.}} \textbf{A}.  Type matn yaratadi.
\par\vspace{3mm}
\needspace{3\baselineskip}
\noindent\hyperlink{q293}{\textbf{293.}} \textbf{B}.  Muharrir loyihasi va eksport farqlanadi.
\par\vspace{3mm}
\needspace{3\baselineskip}
\noindent\hyperlink{q294}{\textbf{294.}} \textbf{D}.  Darslikdagi Paint.NET bo'limi MS Paint versiyasini isbotlamaydi.
\par\vspace{3mm}
\needspace{3\baselineskip}
\noindent\hyperlink{q295}{\textbf{295.}} \textbf{D}.  Bo'yash tutash hududga bog'liq.
\par\vspace{3mm}
\needspace{3\baselineskip}
\noindent\hyperlink{q296}{\textbf{296.}} \textbf{B}.  Vazifalar mutlaqo boshqacha.
\par\vspace{3mm}
\needspace{3\baselineskip}
\noindent\hyperlink{q297}{\textbf{297.}} \textbf{A}.  Undo -- oxirgi amalni qaytarish.
\par\vspace{3mm}
\needspace{3\baselineskip}
\noindent\hyperlink{q298}{\textbf{298.}} \textbf{A}.  Nondestruktiv niqob farqlanadi.
\par\vspace{3mm}
\needspace{3\baselineskip}
\noindent\hyperlink{q299}{\textbf{299.}} \textbf{B}.  Visibility ma'lumotni o'chirishdan farq qiladi.
\par\vspace{3mm}
\needspace{3\baselineskip}
\noindent\hyperlink{q300}{\textbf{300.}} \textbf{C}.  PSD tahrirlash formati.
\par\vspace{3mm}
\needspace{3\baselineskip}
\noindent\hyperlink{q301}{\textbf{301.}} \textbf{B}.  Cambridge5 PDF 38-39: ikkita tanlov shakli.
\par\vspace{3mm}
\needspace{3\baselineskip}
\noindent\hyperlink{q302}{\textbf{302.}} \textbf{D}.  Mask ko'rinishni yashiradi.
\par\vspace{3mm}
\needspace{3\baselineskip}
\noindent\hyperlink{q303}{\textbf{303.}} \textbf{B}.  Bo'yoq to'kish tutash hududga ta'sir qiladi.
\par\vspace{3mm}
\needspace{3\baselineskip}
\noindent\hyperlink{q304}{\textbf{304.}} \textbf{C}.  Cambridge5 PDF 41: tanlash va ikki marta joylashtirish.
\par\vspace{3mm}
\needspace{3\baselineskip}
\noindent\hyperlink{q305}{\textbf{305.}} \textbf{A}.  PSD loyihani saqlaydi.
\par\vspace{3mm}
\needspace{3\baselineskip}
\noindent\hyperlink{q306}{\textbf{306.}} \textbf{A}.  Tanlang va kadrni kesing.
\par\vspace{3mm}
\needspace{3\baselineskip}
\noindent\hyperlink{q307}{\textbf{307.}} \textbf{B - A - D - C}.  Copy + Paste manba obyektni o'rnida qoldirgan holda nusxa hosil qiladi.
\par\vspace{3mm}
\needspace{3\baselineskip}
\noindent\hyperlink{q308}{\textbf{308.}} \textbf{C - A - B}.  To'g'ri ketma-ketlik: 3 =\textgreater{} 1 =\textgreater{} 2.
\par\vspace{3mm}
\needspace{3\baselineskip}
\noindent\hyperlink{q309}{\textbf{309.}} \textbf{B}.  ICT7 PDF 62-67: Layers paneldagi visibility va qatlam o'chirish farqi.
\par\vspace{3mm}
\needspace{3\baselineskip}
\noindent\hyperlink{q310}{\textbf{310.}} \textbf{B}.  Ko'rinish (visibility) o'chiriladi.
\par\vspace{3mm}
\needspace{3\baselineskip}
\noindent\hyperlink{q311}{\textbf{311.}} \textbf{C}.  Rang namunasi olish uchun Eyedropper.
\par\vspace{3mm}
\needspace{3\baselineskip}
\noindent\hyperlink{q312}{\textbf{312.}} \textbf{1-C, 2-A, 3-B}.  To'g'ri moslik: 1-c; 2-a; 3-b.
\par\vspace{3mm}
\needspace{3\baselineskip}
\noindent\hyperlink{q313}{\textbf{313.}} \textbf{C}.  Clone Stamp bir joydan namuna oladi.
\par\vspace{3mm}
\needspace{3\baselineskip}
\noindent\hyperlink{q314}{\textbf{314.}} \textbf{A - C - B - D}.  Kontur ochiq qolsa rang kutilmagan sohalarga tarqalishi mumkin.
\par\vspace{3mm}
\needspace{3\baselineskip}
\noindent\hyperlink{q315}{\textbf{315.}} \textbf{1-D, 2-A, 3-B, 4-C}.  Qalam va cho'tka chizadi, Fill tutash sohani bo'yaydi, Color picker mavjud rangni tanlaydi.
\par\vspace{3mm}
\needspace{3\baselineskip}
\noindent\hyperlink{q316}{\textbf{316.}} \textbf{D}.  Yashirilgan B qatlam compositingda ko'rinmaydi; A va C ko'rinadi.
\par\vspace{3mm}
\needspace{3\baselineskip}
\noindent\hyperlink{q317}{\textbf{317.}} \textbf{C - D - A - B}.  Cambridge 5-sinf Paintdagi Curved line avval to'g'ri segment, so'ng egilish bosqichlari bilan ishlaydi.
\par\vspace{3mm}
\needspace{3\baselineskip}
\noindent\hyperlink{q318}{\textbf{318.}} \textbf{D}.  Lasso erkin tanlov yaratadi.
\par\vspace{3mm}
\needspace{3\baselineskip}
\noindent\hyperlink{q319}{\textbf{319.}} \textbf{1-A, 2-B, 3-C, 4-D}.  Nusxalashda asl soha saqlanadi; kesishda manbadan olinishi mumkin.
\par\vspace{3mm}
\needspace{3\baselineskip}
\noindent\hyperlink{q320}{\textbf{320.}} \textbf{D}.  Yopilmagan kontur to'ldirishga ta'sir qiladi.
\par\vspace{3mm}
\needspace{3\baselineskip}
\noindent\hyperlink{q321}{\textbf{321.}} \textbf{C}.  Opacity 0\% to'liq shaffoflik.
\par\vspace{3mm}
\needspace{3\baselineskip}
\noindent\hyperlink{q322}{\textbf{322.}} \textbf{A}.  Layer stack tartibga ta'sir qiladi.
\par\vspace{3mm}
\needspace{3\baselineskip}
\noindent\hyperlink{q323}{\textbf{323.}} \textbf{D}.  To'liq xira ustki piksel pastdagini yopadi.
\par\vspace{3mm}
\needspace{3\baselineskip}
\noindent\hyperlink{q324}{\textbf{324.}} \textbf{D}.  Black conceals.
\par\vspace{3mm}
\needspace{3\baselineskip}
\noindent\hyperlink{q325}{\textbf{325.}} \textbf{A}.  White reveals.
\par\vspace{3mm}
\needspace{3\baselineskip}
\noindent\hyperlink{q326}{\textbf{326.}} \textbf{1-A, 2-D, 3-B, 4-C}.  Crop hudud tanlab qo'yishdan ko'ra kadr chegarasini o'zgartiradi.
\par\vspace{3mm}
\needspace{3\baselineskip}
\noindent\hyperlink{q327}{\textbf{327.}} \textbf{C}.  Cambridge5 PDF 37-38: bo'yash tutash soha chegarasiga bog'liq.
\par\vspace{3mm}
\needspace{3\baselineskip}
\noindent\hyperlink{q328}{\textbf{328.}} \textbf{1-C, 2-B, 3-A}.  To'g'ri moslik: 1-c; 2-b; 3-a.
\par\vspace{3mm}
\needspace{3\baselineskip}
\noindent\hyperlink{q329}{\textbf{329.}} \textbf{C}.  Nondestruktiv ish oqimi.
\par\vspace{3mm}
\needspace{3\baselineskip}
\noindent\hyperlink{q330}{\textbf{330.}} \textbf{C - B - A}.  To'g'ri ketma-ketlik: 3 =\textgreater{} 2 =\textgreater{} 1.
\par\vspace{3mm}
\needspace{3\baselineskip}
\noindent\hyperlink{q331}{\textbf{331.}} \textbf{D}.  Cambridge5 PDF 29: Undo o'tmishdagi amallarni qaytaradi.
\par\vspace{3mm}
\needspace{3\baselineskip}
\noindent\hyperlink{q332}{\textbf{332.}} \textbf{B}.  Blur tasvirni yumshatadi.
\par\vspace{3mm}
\needspace{3\baselineskip}
\noindent\hyperlink{q333}{\textbf{333.}} \textbf{B}.  Lock ayrim amallarni bloklaydi.
\par\vspace{3mm}
\needspace{3\baselineskip}
\noindent\hyperlink{q334}{\textbf{334.}} \textbf{C}.  Ishchi loyiha va tarqatish formati ajraladi.
\par\vspace{3mm}
\needspace{3\baselineskip}
\noindent\hyperlink{q335}{\textbf{335.}} \textbf{A}.  ICT7 PDF 67-79: tasvirga effekt berish va filtrlash.
\par\vspace{3mm}
\needspace{3\baselineskip}
\noindent\hyperlink{q336}{\textbf{336.}} \textbf{C}.  Cambridge5 PDF 38-41: nusxa olish va joylashtirish.
\par\vspace{3mm}
\needspace{3\baselineskip}
\noindent\hyperlink{q337}{\textbf{337.}} \textbf{1-A, 2-D, 3-C, 4-B}.  PSD bilan eksportning maqsadi farq qiladi; JPEG odatda qatlamlarni saqlamaydi.
\par\vspace{3mm}
\needspace{3\baselineskip}
\noindent\hyperlink{q338}{\textbf{338.}} \textbf{B}.  Oddiy raster layer maskda qora - yashirish, oq - ko'rsatish.
\par\vspace{3mm}
\needspace{3\baselineskip}
\noindent\hyperlink{q339}{\textbf{339.}} \textbf{A}.  ICT7 PDF 62-67: qatlamlarni tahrirlash va qaytariluvchan ishlash.
\par\vspace{3mm}
\needspace{3\baselineskip}
\noindent\hyperlink{q340}{\textbf{340.}} \textbf{C}.  Selection va layer scope muhim.
\par\vspace{3mm}
\needspace{3\baselineskip}
\noindent\hyperlink{q341}{\textbf{341.}} \textbf{A}.  Qatlamli tahrir ma'lumotlari PSD loyihada saqlanishi mumkin.
\par\vspace{3mm}
\needspace{3\baselineskip}
\noindent\hyperlink{q342}{\textbf{342.}} \textbf{C}.  Cambridge5 PDF 42: Save as nom berib saqlash.
\par\vspace{3mm}
\needspace{3\baselineskip}
\noindent\hyperlink{q343}{\textbf{343.}} \textbf{C}.  Diagramma qatlamlarning ko'rinish holatini ko'rsatadi.
\par\vspace{3mm}
\needspace{3\baselineskip}
\noindent\hyperlink{q344}{\textbf{344.}} \textbf{C}.  Cambridge5 PDF 30-31: avval segment, so'ng egri chiziq.
\par\vspace{3mm}
\needspace{3\baselineskip}
\noindent\hyperlink{q345}{\textbf{345.}} \textbf{A}.  ICT7 PDF 50-62: grafika muharirlarida asbob vazifalari; qo'shimcha amaliy izoh.
\par\vspace{3mm}
\needspace{3\baselineskip}
\noindent\hyperlink{q346}{\textbf{346.}} \textbf{D - C - A - B}.  Matn ustki qatlamda va ko'rinadigan bo'lsa fon va grafikdan yashirinib qolmaydi.
\par\vspace{3mm}
\needspace{3\baselineskip}
\noindent\hyperlink{q347}{\textbf{347.}} \textbf{C}.  type='a' ro'yxatni kichik lotin harflarida raqamlaydi; start='2' sabab birinchi element \texttt{b} bo'ladi. Qalin/kursiv sarlavhalar ro'yxat elementi emas.
\par\vspace{3mm}
\needspace{3\baselineskip}
\noindent\hyperlink{q348}{\textbf{348.}} \textbf{B}.  \texttt{rowspan='2'} birinchi holda 248 ni 1-2 qatorlarga, ikkinchi holda Rowling ni 2-3 qatorlarga yoyadi; \texttt{\textless{}th\textgreater{}} sarlavhalarni markazlashtiradi.
\par\vspace{3mm}
\needspace{3\baselineskip}
\noindent\hyperlink{q349}{\textbf{349.}} \textbf{C}.  A \texttt{.box p} selektori bo'yicha ko'k; B uchun \texttt{.box p + span} \texttt{.box span} dan xosroq, olovrang.
\par\vspace{3mm}
\needspace{3\baselineskip}
\noindent\hyperlink{q350}{\textbf{350.}} \textbf{D}.  To'g'ri tanlovda ko'k \#000066 matn, \#e6e6ff sarlavha, \#f2f2f2 matn paneli, 95x140 px surat va o'ngga float mavjud.
\par\vspace{3mm}
\needspace{3\baselineskip}
\noindent\hyperlink{q351}{\textbf{351.}} \textbf{1-B, 2-A, 3-D, 4-C}.  href va src manzilga xizmat qiladi, alt muqobil tavsif; target ko'rish kontekstini tanlaydi.
\par\vspace{3mm}
\needspace{3\baselineskip}
\noindent\hyperlink{q352}{\textbf{352.}} \textbf{B - A - D - C}.  HTMLda head metama'lumoti body tarkibidan oldin keltiriladi.
\par\vspace{3mm}
\needspace{3\baselineskip}
\noindent\hyperlink{q353}{\textbf{353.}} \textbf{A - C - B - D}.  Ichki ul tashqi li ning ichida bo'lishi kerak; li oldin yopilsa tuzilma buziladi.
\par\vspace{3mm}
\needspace{3\baselineskip}
\noindent\hyperlink{q354}{\textbf{354.}} \textbf{1-A, 2-B, 3-C, 4-D}.  Sarlavha h1-h6, paragraf p, strong semantik urg'u, sub pastki indeks beradi.
\par\vspace{3mm}
\needspace{3\baselineskip}
\noindent\hyperlink{q355}{\textbf{355.}} \textbf{1-A, 2-D, 3-B, 4-C}.  Ichki ro'yxat tegini tashqi li ichida joylashtirish kerak.
\par\vspace{3mm}
\needspace{3\baselineskip}
\noindent\hyperlink{q356}{\textbf{356.}} \textbf{D - A - B - C}.  img havolaning farzandi bo'lganda rasm bosiladigan havolaga aylanadi.
\par\vspace{3mm}
\needspace{3\baselineskip}
\noindent\hyperlink{q357}{\textbf{357.}} \textbf{1-D, 2-C, 3-A, 4-B}.  Title brauzer tabida, h1 esa body ichida sahifada ko'rinadi.
\par\vspace{3mm}
\needspace{3\baselineskip}
\noindent\hyperlink{q358}{\textbf{358.}} \textbf{A - D - B - C}.  Start raqamlashning boshlang'ich qiymatini belgilaydi; ikkita band 3 va 4 bo'ladi.
\par\vspace{3mm}
\needspace{3\baselineskip}
\noindent\hyperlink{q359}{\textbf{359.}} \textbf{B}.  Title odatda head qismida.
\par\vspace{3mm}
\needspace{3\baselineskip}
\noindent\hyperlink{q360}{\textbf{360.}} \textbf{C}.  p paragraf.
\par\vspace{3mm}
\needspace{3\baselineskip}
\noindent\hyperlink{q361}{\textbf{361.}} \textbf{B}.  a href havolani bildiradi.
\par\vspace{3mm}
\needspace{3\baselineskip}
\noindent\hyperlink{q362}{\textbf{362.}} \textbf{A}.  tr table row.
\par\vspace{3mm}
\needspace{3\baselineskip}
\noindent\hyperlink{q363}{\textbf{363.}} \textbf{A}.  ol ordered list.
\par\vspace{3mm}
\needspace{3\baselineskip}
\noindent\hyperlink{q364}{\textbf{364.}} \textbf{C}.  ICT7 PDF 138-140: reset forma qiymatlarini tiklaydi.
\par\vspace{3mm}
\needspace{3\baselineskip}
\noindent\hyperlink{q365}{\textbf{365.}} \textbf{C}.  Body sahifaning ko'rinadigan tarkibi.
\par\vspace{3mm}
\needspace{3\baselineskip}
\noindent\hyperlink{q366}{\textbf{366.}} \textbf{B}.  ICT7 PDF 138-140: hidden maxfiylik kafolati emas.
\par\vspace{3mm}
\needspace{3\baselineskip}
\noindent\hyperlink{q367}{\textbf{367.}} \textbf{C}.  img element rasmni ko'rsatadi.
\par\vspace{3mm}
\needspace{3\baselineskip}
\noindent\hyperlink{q368}{\textbf{368.}} \textbf{D}.  ul unordered list.
\par\vspace{3mm}
\needspace{3\baselineskip}
\noindent\hyperlink{q369}{\textbf{369.}} \textbf{D}.  HTML gipermatnli belgilash tili.
\par\vspace{3mm}
\needspace{3\baselineskip}
\noindent\hyperlink{q370}{\textbf{370.}} \textbf{A}.  img src tasvir manbasidir.
\par\vspace{3mm}
\needspace{3\baselineskip}
\noindent\hyperlink{q371}{\textbf{371.}} \textbf{B}.  Semantik ta'kidlash farqi.
\par\vspace{3mm}
\needspace{3\baselineskip}
\noindent\hyperlink{q372}{\textbf{372.}} \textbf{C}.  Tanlov mexanizmi farqlanadi.
\par\vspace{3mm}
\needspace{3\baselineskip}
\noindent\hyperlink{q373}{\textbf{373.}} \textbf{C}.  Rowspan qator, colspan ustun.
\par\vspace{3mm}
\needspace{3\baselineskip}
\noindent\hyperlink{q374}{\textbf{374.}} \textbf{B}.  Havola href, rasm src.
\par\vspace{3mm}
\needspace{3\baselineskip}
\noindent\hyperlink{q375}{\textbf{375.}} \textbf{D}.  Holat va submit farqi.
\par\vspace{3mm}
\needspace{3\baselineskip}
\noindent\hyperlink{q376}{\textbf{376.}} \textbf{A}.  Yuborish usuli va HTTPS alohida.
\par\vspace{3mm}
\needspace{3\baselineskip}
\noindent\hyperlink{q377}{\textbf{377.}} \textbf{B}.  Eskirgan ko'rinish usullari CSSga ko'chiriladi.
\par\vspace{3mm}
\needspace{3\baselineskip}
\noindent\hyperlink{q378}{\textbf{378.}} \textbf{C}.  Alt muqobil matn.
\par\vspace{3mm}
\needspace{3\baselineskip}
\noindent\hyperlink{q379}{\textbf{379.}} \textbf{B}.  Semantic th va td.
\par\vspace{3mm}
\needspace{3\baselineskip}
\noindent\hyperlink{q380}{\textbf{380.}} \textbf{B}.  img void element.
\par\vspace{3mm}
\needspace{3\baselineskip}
\noindent\hyperlink{q381}{\textbf{381.}} \textbf{D}.  Form fayllarini yuborishda multipart/form-data odatiy format; qo'shimcha texnik izoh.
\par\vspace{3mm}
\needspace{3\baselineskip}
\noindent\hyperlink{q382}{\textbf{382.}} \textbf{D}.  Submit uchun name muhim.
\par\vspace{3mm}
\needspace{3\baselineskip}
\noindent\hyperlink{q383}{\textbf{383.}} \textbf{C}.  Rasmni HTML hujjatiga joylash uchun \texttt{\textless{}img\textgreater{}} tegi ishlatiladi. \texttt{src="photo.png"} manba yo'lini, \texttt{alt="Foto"} esa rasm ko'rinmaganda yoki ekran o'quvchisida ishlatiladigan muqobil matnni belgilaydi. \texttt{href} havolaga, \texttt{file} esa bu vazifada kerakli \texttt{\textless{}img\textgreater{}} atributiga to'g'ri kelmaydi. MDN: https://developer.mozilla.org/en-US/docs/Web/HTML/Reference/Elements/img .
\par\vspace{3mm}
\needspace{3\baselineskip}
\noindent\hyperlink{q384}{\textbf{384.}} \textbf{D}.  Teglar ichma-ich tuzilishi va atributni tekshiring.
\par\vspace{3mm}
\needspace{3\baselineskip}
\noindent\hyperlink{q385}{\textbf{385.}} \textbf{B}.  Formalar va rasm atributlari funksiyasini farqlang.
\par\vspace{3mm}
\needspace{3\baselineskip}
\noindent\hyperlink{q386}{\textbf{386.}} \textbf{C}.  ICT7 PDF 120-123: ichma-ich ro'yxatli li tuzilishi.
\par\vspace{3mm}
\needspace{3\baselineskip}
\noindent\hyperlink{q387}{\textbf{387.}} \textbf{A}.  Boshlanish 3, 5 ta band: 3+5-1=7.
\par\vspace{3mm}
\needspace{3\baselineskip}
\noindent\hyperlink{q388}{\textbf{388.}} \textbf{C}.  Boshlanish 3, uchta band: 3,4,5; ICT7 PDF 120-123.
\par\vspace{3mm}
\needspace{3\baselineskip}
\noindent\hyperlink{q389}{\textbf{389.}} \textbf{B}.  Boshlanish 7, 2 ta band: 7+2-1=8.
\par\vspace{3mm}
\needspace{3\baselineskip}
\noindent\hyperlink{q390}{\textbf{390.}} \textbf{A - B - C}.  To'g'ri ketma-ketlik: 1 =\textgreater{} 2 =\textgreater{} 3.
\par\vspace{3mm}
\needspace{3\baselineskip}
\noindent\hyperlink{q391}{\textbf{391.}} \textbf{A}.  ICT7 PDF 139-141: radio name guruhini belgilaydi.
\par\vspace{3mm}
\needspace{3\baselineskip}
\noindent\hyperlink{q392}{\textbf{392.}} \textbf{1-A, 2-B, 3-C}.  To'g'ri moslik: 1-a; 2-b; 3-c.
\par\vspace{3mm}
\needspace{3\baselineskip}
\noindent\hyperlink{q393}{\textbf{393.}} \textbf{D}.  Boshlanish 5, 4 ta band: 5+4-1=8.
\par\vspace{3mm}
\needspace{3\baselineskip}
\noindent\hyperlink{q394}{\textbf{394.}} \textbf{A}.  Boshlanish 1, 3 ta band: 1+3-1=3.
\par\vspace{3mm}
\needspace{3\baselineskip}
\noindent\hyperlink{q395}{\textbf{395.}} \textbf{D}.  Har oddiy td 1, birlashtirilgan katak 4; jami 2+4=6.
\par\vspace{3mm}
\needspace{3\baselineskip}
\noindent\hyperlink{q396}{\textbf{396.}} \textbf{B}.  Har oddiy td 1, birlashtirilgan katak 2; jami 3+2=5.
\par\vspace{3mm}
\needspace{3\baselineskip}
\noindent\hyperlink{q397}{\textbf{397.}} \textbf{A}.  DOM tuzilishi va brauzer elementlarini sanang.
\par\vspace{3mm}
\needspace{3\baselineskip}
\noindent\hyperlink{q398}{\textbf{398.}} \textbf{C}.  Har oddiy td 1, birlashtirilgan katak 2; jami 5+2=7.
\par\vspace{3mm}
\needspace{3\baselineskip}
\noindent\hyperlink{q399}{\textbf{399.}} \textbf{A}.  Har oddiy td 1, birlashtirilgan katak 3; jami 2+3=5.
\par\vspace{3mm}
\needspace{3\baselineskip}
\noindent\hyperlink{q400}{\textbf{400.}} \textbf{D}.  ICT7 PDF 138-140: submit va reset vazifasi.
\par\vspace{3mm}
\needspace{3\baselineskip}
\noindent\hyperlink{q401}{\textbf{401.}} \textbf{C - A - D - B}.  1+2 ustun qamrovi mantiqiy 3 ustun beradi.
\par\vspace{3mm}
\needspace{3\baselineskip}
\noindent\hyperlink{q402}{\textbf{402.}} \textbf{C}.  HTML formadagi yuboriladigan qiymat mexanizmini tekshiring.
\par\vspace{3mm}
\needspace{3\baselineskip}
\noindent\hyperlink{q403}{\textbf{403.}} \textbf{A - B - C - D}.  Fayl kiritish odatda multipart/form-data bilan birga ishlatiladi.
\par\vspace{3mm}
\needspace{3\baselineskip}
\noindent\hyperlink{q404}{\textbf{404.}} \textbf{B}.  Har oddiy td 1, birlashtirilgan katak 2; jami 4+2=6.
\par\vspace{3mm}
\needspace{3\baselineskip}
\noindent\hyperlink{q405}{\textbf{405.}} \textbf{A - C - B}.  To'g'ri ketma-ketlik: 1 =\textgreater{} 3 =\textgreater{} 2.
\par\vspace{3mm}
\needspace{3\baselineskip}
\noindent\hyperlink{q406}{\textbf{406.}} \textbf{1-C, 2-B, 3-A, 4-D}.  tr qator, th sarlavha, td odatiy katak vazifasini bajaradi.
\par\vspace{3mm}
\needspace{3\baselineskip}
\noindent\hyperlink{q407}{\textbf{407.}} \textbf{D}.  HTMLda input void element; ICT7 PDF 137 dagi tarixiy namunadagi shaklga izoh.
\par\vspace{3mm}
\needspace{3\baselineskip}
\noindent\hyperlink{q408}{\textbf{408.}} \textbf{1-A, 2-C, 3-B}.  To'g'ri moslik: 1-a; 2-c; 3-b.
\par\vspace{3mm}
\needspace{3\baselineskip}
\noindent\hyperlink{q409}{\textbf{409.}} \textbf{1-C, 2-A, 3-D, 4-B}.  Colspan qator bo'yicha, rowspan ustun bo'yicha kengaytiradi; katak soni va mantiqiy ustun sonini farqlang.
\par\vspace{3mm}
\needspace{3\baselineskip}
\noindent\hyperlink{q410}{\textbf{410.}} \textbf{B}.  Har oddiy td 1, birlashtirilgan katak 3; jami 3+3=6.
\par\vspace{3mm}
\needspace{3\baselineskip}
\noindent\hyperlink{q411}{\textbf{411.}} \textbf{D - C - A - B}.  label for va input id bir xil bo'lsa bog'lanadi; name yuborish kalitidir.
\par\vspace{3mm}
\needspace{3\baselineskip}
\noindent\hyperlink{q412}{\textbf{412.}} \textbf{1-D, 2-A, 3-C, 4-B}.  Hidden qiymat maxfiy emas, reset esa server bazasini o'chirmaydi.
\par\vspace{3mm}
\needspace{3\baselineskip}
\noindent\hyperlink{q413}{\textbf{413.}} \textbf{A}.  Brauzer yuboradigan successful control uchun name muhim.
\par\vspace{3mm}
\needspace{3\baselineskip}
\noindent\hyperlink{q414}{\textbf{414.}} \textbf{A}.  colspan=2 ikki ustunni, oddiy td bir ustunni egallaydi.
\par\vspace{3mm}
\needspace{3\baselineskip}
\noindent\hyperlink{q415}{\textbf{415.}} \textbf{A}.  Disabled odatda yuborilmaydi; readonly maydon nomi bo'lsa yuborilishi mumkin.
\par\vspace{3mm}
\needspace{3\baselineskip}
\noindent\hyperlink{q416}{\textbf{416.}} \textbf{D}.  ICT7 PDF 139-141: checkbox holatiga qarab yuboriladi.
\par\vspace{3mm}
\needspace{3\baselineskip}
\noindent\hyperlink{q417}{\textbf{417.}} \textbf{A}.  Birinchi katak bitta ustun, keyingi katak 2 ustun; jami 3.
\par\vspace{3mm}
\needspace{3\baselineskip}
\noindent\hyperlink{q418}{\textbf{418.}} \textbf{B}.  ICT7 PDF 128-135: colspan ustun qamrovi.
\par\vspace{3mm}
\needspace{3\baselineskip}
\noindent\hyperlink{q419}{\textbf{419.}} \textbf{D}.  Teglar atributlari va tashqi kontekst o'rtasidagi munosabatni tekshiring.
\par\vspace{3mm}
\needspace{3\baselineskip}
\noindent\hyperlink{q420}{\textbf{420.}} \textbf{C - D - A - B}.  Bir xil name bo'lgan radio tugmalar bitta tanlash guruhida ishlaydi.
\par\vspace{3mm}
\needspace{3\baselineskip}
\noindent\hyperlink{q421}{\textbf{421.}} \textbf{C}.  ol start=5; tashqi li lar 5 va 6. Ichki ul tashqi raqamlanishga qo'shilmaydi.
\par\vspace{3mm}
\needspace{3\baselineskip}
\noindent\hyperlink{q422}{\textbf{422.}} \textbf{C}.  Teglar atributlari va tashqi kontekst o'rtasidagi munosabatni tekshiring.
\par\vspace{3mm}
\needspace{3\baselineskip}
\noindent\hyperlink{q423}{\textbf{423.}} \textbf{D}.  Teglar atributlari va tashqi kontekst o'rtasidagi munosabatni tekshiring.
\par\vspace{3mm}
\needspace{3\baselineskip}
\noindent\hyperlink{q424}{\textbf{424.}} \textbf{1-A, 2-D, 3-C, 4-B}.  Required serverdagi mustaqil validatsiyaning o'rnini bosmaydi.
\par\vspace{3mm}
\needspace{3\baselineskip}
\noindent\hyperlink{q425}{\textbf{425.}} \textbf{D}.  a elementi ichida img bo'lishi rasmni bosiladigan havolaga aylantiradi.
\par\vspace{3mm}
\needspace{3\baselineskip}
\noindent\hyperlink{q426}{\textbf{426.}} \textbf{B}.  Cambridge9 PDF 291-294: ../ nisbiy katalog manzili.
\par\vspace{3mm}
\needspace{3\baselineskip}
\noindent\hyperlink{q427}{\textbf{427.}} \textbf{C}.  Cambridge9 PDF 296: id selektori \# bilan yoziladi.
\par\vspace{3mm}
\needspace{3\baselineskip}
\noindent\hyperlink{q428}{\textbf{428.}} \textbf{D}.  Selektor elementni tanlaydi.
\par\vspace{3mm}
\needspace{3\baselineskip}
\noindent\hyperlink{q429}{\textbf{429.}} \textbf{C}.  background-color fon rangi.
\par\vspace{3mm}
\needspace{3\baselineskip}
\noindent\hyperlink{q430}{\textbf{430.}} \textbf{C}.  Padding ichki bo'shliq.
\par\vspace{3mm}
\needspace{3\baselineskip}
\noindent\hyperlink{q431}{\textbf{431.}} \textbf{A}.  .class class selektori.
\par\vspace{3mm}
\needspace{3\baselineskip}
\noindent\hyperlink{q432}{\textbf{432.}} \textbf{A}.  text-align center.
\par\vspace{3mm}
\needspace{3\baselineskip}
\noindent\hyperlink{q433}{\textbf{433.}} \textbf{A}.  \#id identifikator selektori.
\par\vspace{3mm}
\needspace{3\baselineskip}
\noindent\hyperlink{q434}{\textbf{434.}} \textbf{D}.  Cambridge9 PDF 296: class selektori . belgisi bilan boshlanadi.
\par\vspace{3mm}
\needspace{3\baselineskip}
\noindent\hyperlink{q435}{\textbf{435.}} \textbf{C}.  font-size shrift o'lchamini belgilaydi.
\par\vspace{3mm}
\needspace{3\baselineskip}
\noindent\hyperlink{q436}{\textbf{436.}} \textbf{D}.  border-collapse collapse.
\par\vspace{3mm}
\needspace{3\baselineskip}
\noindent\hyperlink{q437}{\textbf{437.}} \textbf{A}.  CSS kaskadli uslub jadvallari.
\par\vspace{3mm}
\needspace{3\baselineskip}
\noindent\hyperlink{q438}{\textbf{438.}} \textbf{C}.  color matnning oldingi fon rangini emas rangini belgilaydi.
\par\vspace{3mm}
\needspace{3\baselineskip}
\noindent\hyperlink{q439}{\textbf{439.}} \textbf{B}.  External CSS head link bilan.
\par\vspace{3mm}
\needspace{3\baselineskip}
\noindent\hyperlink{q440}{\textbf{440.}} \textbf{A}.  \texttt{object-fit: contain} tasvirni nisbatni saqlagan holda butun qutiga sig'diradi; mos kelmagan tomonlarda bo'sh joy qolishi mumkin. \texttt{cover} qutini to'ldirish uchun qirqishi, \texttt{fill} tasvirni cho'zib nisbatni buzishi, \texttt{none} esa o'lchamni o'zgartirmasligi mumkin. MDN: https://developer.mozilla.org/en-US/docs/Web/CSS/Reference/Properties/object-fit .
\par\vspace{3mm}
\needspace{3\baselineskip}
\noindent\hyperlink{q441}{\textbf{441.}} \textbf{C}.  Box sizing formula farqlanadi.
\par\vspace{3mm}
\needspace{3\baselineskip}
\noindent\hyperlink{q442}{\textbf{442.}} \textbf{D}.  Tie break source order.
\par\vspace{3mm}
\needspace{3\baselineskip}
\noindent\hyperlink{q443}{\textbf{443.}} \textbf{A}.  Padding cell ichki masofa.
\par\vspace{3mm}
\needspace{3\baselineskip}
\noindent\hyperlink{q444}{\textbf{444.}} \textbf{B}.  Inline style atributi.
\par\vspace{3mm}
\needspace{3\baselineskip}
\noindent\hyperlink{q445}{\textbf{445.}} \textbf{B}.  Oldingi/fon rang.
\par\vspace{3mm}
\needspace{3\baselineskip}
\noindent\hyperlink{q446}{\textbf{446.}} \textbf{D}.  \# ID selektori.
\par\vspace{3mm}
\needspace{3\baselineskip}
\noindent\hyperlink{q447}{\textbf{447.}} \textbf{A}.  Bo'shliq descendent, \textgreater{} child.
\par\vspace{3mm}
\needspace{3\baselineskip}
\noindent\hyperlink{q448}{\textbf{448.}} \textbf{B}.  Cambridge9 PDF 295-296: CSS ulashning 3 yo'li.
\par\vspace{3mm}
\needspace{3\baselineskip}
\noindent\hyperlink{q449}{\textbf{449.}} \textbf{D}.  Visibility joy saqlashi mumkin.
\par\vspace{3mm}
\needspace{3\baselineskip}
\noindent\hyperlink{q450}{\textbf{450.}} \textbf{A}.  Element selektori va class boshqacha.
\par\vspace{3mm}
\needspace{3\baselineskip}
\noindent\hyperlink{q451}{\textbf{451.}} \textbf{1-D, 2-C, 3-A, 4-B}.  Tashqi uslub ko'p sahifalar uchun qayta foydalanishga qulay.
\par\vspace{3mm}
\needspace{3\baselineskip}
\noindent\hyperlink{q452}{\textbf{452.}} \textbf{C}.  CSS selektor/sintaksis funksiyasi muhim.
\par\vspace{3mm}
\needspace{3\baselineskip}
\noindent\hyperlink{q453}{\textbf{453.}} \textbf{B - A - C}.  To'g'ri ketma-ketlik: 2 =\textgreater{} 1 =\textgreater{} 3.
\par\vspace{3mm}
\needspace{3\baselineskip}
\noindent\hyperlink{q454}{\textbf{454.}} \textbf{D}.  Selektor mos elementga berilgan xossani tatbiq etadi.
\par\vspace{3mm}
\needspace{3\baselineskip}
\noindent\hyperlink{q455}{\textbf{455.}} \textbf{C}.  Selektor mos elementga berilgan xossani tatbiq etadi.
\par\vspace{3mm}
\needspace{3\baselineskip}
\noindent\hyperlink{q456}{\textbf{456.}} \textbf{A - D - B - C}.  Faqat so'nggi qoida yutadi degan umumiy xulosa noto'g'ri.
\par\vspace{3mm}
\needspace{3\baselineskip}
\noindent\hyperlink{q457}{\textbf{457.}} \textbf{B}.  Bevosita farzand faqat Bir p; Ikki section ichida.
\par\vspace{3mm}
\needspace{3\baselineskip}
\noindent\hyperlink{q458}{\textbf{458.}} \textbf{1-B, 2-A, 3-D, 4-C}.  Element, class va id selektorlari turlicha tanlaydi; p.note faqat p.note kesishmasidir.
\par\vspace{3mm}
\needspace{3\baselineskip}
\noindent\hyperlink{q459}{\textbf{459.}} \textbf{1-B, 2-A, 3-C}.  To'g'ri moslik: 1-b; 2-a; 3-c.
\par\vspace{3mm}
\needspace{3\baselineskip}
\noindent\hyperlink{q460}{\textbf{460.}} \textbf{C}.  Inline style oddiy author qoidalaridan ustuvor bo'lishi mumkin.
\par\vspace{3mm}
\needspace{3\baselineskip}
\noindent\hyperlink{q461}{\textbf{461.}} \textbf{D}.  Selektor mos elementga berilgan xossani tatbiq etadi.
\par\vspace{3mm}
\needspace{3\baselineskip}
\noindent\hyperlink{q462}{\textbf{462.}} \textbf{C - B - A - D}.  Selector maqsadni belgilaydi; deklaratsiyalar uning xossalarini tayinlaydi.
\par\vspace{3mm}
\needspace{3\baselineskip}
\noindent\hyperlink{q463}{\textbf{463.}} \textbf{D}.  To'g'ri sintaksis va xossani tanlang.
\par\vspace{3mm}
\needspace{3\baselineskip}
\noindent\hyperlink{q464}{\textbf{464.}} \textbf{A}.  Selektor mos elementga berilgan xossani tatbiq etadi.
\par\vspace{3mm}
\needspace{3\baselineskip}
\noindent\hyperlink{q465}{\textbf{465.}} \textbf{C}.  To'g'ri sintaksis va xossani tanlang.
\par\vspace{3mm}
\needspace{3\baselineskip}
\noindent\hyperlink{q466}{\textbf{466.}} \textbf{D}.  DOM daraxtida A bevosita farzand p, B esa section ichidagi p.
\par\vspace{3mm}
\needspace{3\baselineskip}
\noindent\hyperlink{q467}{\textbf{467.}} \textbf{D - A - C - B}.  External CSS alohida faylni HTML head orqali ulaydi.
\par\vspace{3mm}
\needspace{3\baselineskip}
\noindent\hyperlink{q468}{\textbf{468.}} \textbf{1-A, 2-D, 3-B, 4-C}.  Bo'sh joy va \textgreater{} kombinatorlari bir xil emas; vergul guruhlaydi.
\par\vspace{3mm}
\needspace{3\baselineskip}
\noindent\hyperlink{q469}{\textbf{469.}} \textbf{B}.  Cambridge9 PDF 295-296 inline va internal farqi; important yo'q deb olinadi.
\par\vspace{3mm}
\needspace{3\baselineskip}
\noindent\hyperlink{q470}{\textbf{470.}} \textbf{C - A - D - B}.  200+2x10+2x2+2x5=234px, soddalashtirilgan gorizontal box modelda.
\par\vspace{3mm}
\needspace{3\baselineskip}
\noindent\hyperlink{q471}{\textbf{471.}} \textbf{B - C - A}.  To'g'ri ketma-ketlik: 2 =\textgreater{} 3 =\textgreater{} 1.
\par\vspace{3mm}
\needspace{3\baselineskip}
\noindent\hyperlink{q472}{\textbf{472.}} \textbf{A}.  Cambridge9 PDF 296-299: veb ko'rinishini CSS bilan boshqarish.
\par\vspace{3mm}
\needspace{3\baselineskip}
\noindent\hyperlink{q473}{\textbf{473.}} \textbf{C}.  300+2x15+2x3+2x10=356px.
\par\vspace{3mm}
\needspace{3\baselineskip}
\noindent\hyperlink{q474}{\textbf{474.}} \textbf{1-A, 2-D, 3-C, 4-B}.  cover va contain bir xil emas; ko'rinish manba fayl piksellarini eksport qilmaydi.
\par\vspace{3mm}
\needspace{3\baselineskip}
\noindent\hyperlink{q475}{\textbf{475.}} \textbf{C}.  240+2x12+2x4+2x6=284px.
\par\vspace{3mm}
\needspace{3\baselineskip}
\noindent\hyperlink{q476}{\textbf{476.}} \textbf{C}.  Oddiy kaskad/selektor va ro'yxat xossalari qo'llanadi.
\par\vspace{3mm}
\needspace{3\baselineskip}
\noindent\hyperlink{q477}{\textbf{477.}} \textbf{D}.  100+2x20+2x5+2x0=150px.
\par\vspace{3mm}
\needspace{3\baselineskip}
\noindent\hyperlink{q478}{\textbf{478.}} \textbf{1-C, 2-A, 3-D, 4-B}.  Standart content-box width faqat content kengligini bildiradi.
\par\vspace{3mm}
\needspace{3\baselineskip}
\noindent\hyperlink{q479}{\textbf{479.}} \textbf{1-C, 2-B, 3-A, 4-D}.  Rang, fon, qalinlik va tekislash bir-biridan farq qiladi.
\par\vspace{3mm}
\needspace{3\baselineskip}
\noindent\hyperlink{q480}{\textbf{480.}} \textbf{A - B - C - D}.  Max-width kontent chegarasini, height auto esa nisbatni saqlashni qo'llaydi.
\par\vspace{3mm}
\needspace{3\baselineskip}
\noindent\hyperlink{q481}{\textbf{481.}} \textbf{C}.  200+2x10+2x2+2x5=234px.
\par\vspace{3mm}
\needspace{3\baselineskip}
\noindent\hyperlink{q482}{\textbf{482.}} \textbf{B}.  120+2x8+2x1+2x0=138px.
\par\vspace{3mm}
\needspace{3\baselineskip}
\noindent\hyperlink{q483}{\textbf{483.}} \textbf{D - C - A - B}.  CSS jadvaldagi kataklar sonini emas, tashqi ko'rinishini boshqaradi.
\par\vspace{3mm}
\needspace{3\baselineskip}
\noindent\hyperlink{q484}{\textbf{484.}} \textbf{B}.  80+2x6+2x2+2x2=100px.
\par\vspace{3mm}
\needspace{3\baselineskip}
\noindent\hyperlink{q485}{\textbf{485.}} \textbf{C}.  150+2x5+2x2+2x8=180px.
\par\vspace{3mm}
\needspace{3\baselineskip}
\noindent\hyperlink{q486}{\textbf{486.}} \textbf{1-B, 2-C, 3-A}.  To'g'ri moslik: 1-b; 2-c; 3-a.
\par\vspace{3mm}
\needspace{3\baselineskip}
\noindent\hyperlink{q487}{\textbf{487.}} \textbf{A - D - C - B}.  Media query shart bajarilganda flex yo'nalishini o'zgartiradi.
\par\vspace{3mm}
\needspace{3\baselineskip}
\noindent\hyperlink{q488}{\textbf{488.}} \textbf{B}.  td selektori color:green beradi; table/th background-color va tr background-color matn rangini o'zgartirmaydi. Bu 13.04.2026 namunasidagi CSS-kod tahlili uslubidan yangi tuzilgan savol.
\par\vspace{3mm}
\needspace{3\baselineskip}
\noindent\hyperlink{q489}{\textbf{489.}} \textbf{D}.  Class selektori element selektoridan spetsifikroq.
\par\vspace{3mm}
\needspace{3\baselineskip}
\noindent\hyperlink{q490}{\textbf{490.}} \textbf{1-A, 2-C, 3-B, 4-D}.  Barcha misollar oddiy muallif deklaratsiyalari, !important va cascade layerlarsiz qabul qilinadi.
\par\vspace{3mm}
\needspace{3\baselineskip}
\noindent\hyperlink{q491}{\textbf{491.}} \textbf{D}.  200 + 2x(10+2+5) = 234 px.
\par\vspace{3mm}
\needspace{3\baselineskip}
\noindent\hyperlink{q492}{\textbf{492.}} \textbf{A}.  Oddiy inline ko'kdan !important qizil qoidasi ustuvor.
\par\vspace{3mm}
\needspace{3\baselineskip}
\noindent\hyperlink{q493}{\textbf{493.}} \textbf{A}.  100+2x(10+2+5)=134px.
\par\vspace{3mm}
\needspace{3\baselineskip}
\noindent\hyperlink{q494}{\textbf{494.}} \textbf{A}.  contain rasmning to'liq sig'ishini afzal ko'radi.
\par\vspace{3mm}
\needspace{3\baselineskip}
\noindent\hyperlink{q495}{\textbf{495.}} \textbf{B}.  Selektor, inline va ko'rinish xossalarini birgalikda tahlil qiling.
\par\vspace{3mm}
\needspace{3\baselineskip}
\noindent\hyperlink{q496}{\textbf{496.}} \textbf{B - C - D - A}.  \texttt{\textgreater{}} bevosita farzand kombinatori; \texttt{.card p} esa barcha avlod pga mos bo'ladi.
\par\vspace{3mm}
\needspace{3\baselineskip}
\noindent\hyperlink{q497}{\textbf{497.}} \textbf{A}.  Selektor, inline va ko'rinish xossalarini birgalikda tahlil qiling.
\par\vspace{3mm}
\needspace{3\baselineskip}
\noindent\hyperlink{q498}{\textbf{498.}} \textbf{D}.  border-box width 200 px ga padding/border kiradi, margin ikki tomonda +10.
\par\vspace{3mm}
\needspace{3\baselineskip}
\noindent\hyperlink{q499}{\textbf{499.}} \textbf{1-D, 2-A, 3-C, 4-B}.  Ro'yxat markerlari va jadval kataklarini bezash HTMLdagi mazmunni o'zgartirmaydi.
\par\vspace{3mm}
\needspace{3\baselineskip}
\noindent\hyperlink{q500}{\textbf{500.}} \textbf{B}.  ID selektori yuqori spetsifiklikka ega.
\par\vspace{3mm}
\needspace{3\baselineskip}
\noindent\hyperlink{q501}{\textbf{501.}} \textbf{D}.  Selektor, inline va ko'rinish xossalarini birgalikda tahlil qiling.
\par\vspace{3mm}
\needspace{3\baselineskip}
\noindent\hyperlink{q502}{\textbf{502.}} \textbf{B}.  Bevosita p ikkala selektorga tushadi; teng spetsifiklikda keyingi qoida red yutadi.
\par\vspace{3mm}
\clearpage\section{Nashr tarkibi va nazorat izohi}
502 ta faol yozuv (125 published + 377 review) kitobga kiritilgan; 452 archived versiya eskirgan/arxiv yozuvlari sifatida chiqarilgan.
Manbadagi to'g'ri javoblar va izohlar alohida yangi variantlar bilan almashtirilmadi.
Bu to'plam mavjud baza holatining nusxasi; review savollarining faktik ekspert baholashdan o'tganini bildirmaydi.
Savollarning ko'rinishi va havolalarini chopdan oldin tekshirish tavsiya qilinadi.
\end{document}